%!TEX program = xelatex
% Theoretische Physik C — deutsche Fassung (Leitfassung)
%
% Die englische Fassung TheoreticalPhysicsC.tex spiegelt dieses Dokument
% Vorlesung für Vorlesung.  scripts/check_parity.py prüft, dass beide
% strukturell übereinstimmen; ./ci.sh führt die Prüfung aus.

\documentclass[12pt,oneside]{amsart}

\usepackage[german]{tpc-style}
\usepackage{tpc-macros}
\usepackage{tpc-tikz-templates}

\title[Theoretische Physik C]{Theoretische Physik C}
\author{Tobias J.\ Osborne}
\date{\today}

\begin{document}

\maketitle
\tableofcontents

%!TEX root = TheoretischePhysikC.tex
% Skript v0.3, Kapitel 1: Kern von C0 und C1; Quellen in docs/plan/script-v0.3/sync-ch1.md.

\section{Vorlesung: Die Geburt der Quantenmechanik}
Wir befinden uns in einem Quanten-Universum. Bis 1900 haben wir die Quantenaspekte des Universums nicht bemerkt, weil uns die klassische Physik eine sehr gute Annäherung an die Physik des Alltags geliefert hat. Erklärungsbedürftig ist also nicht die Quantenmechanik, sondern die Güte der klassischen Näherung. Das Stichwort dazu heißt \emph{Dekohärenz}; sie gehört in die fortgeschrittene Quantenmechanik.

Eine Bemerkung zur Methode, die für den ganzen Kurs gilt: Man kann erst die falschen, klassischen oder halbklassischen Modelle durchrechnen, um dann zu zeigen, dass sie nicht funktionieren. Das ist nicht mein Stil. Wir betrachten von Anfang an alles richtig; das Bohrsche Atommodell rechnen wir nirgends durch. Dieses Skript enthält nur den Kern; Herleitungen, Vertiefungen und interaktive Abbildungen stehen im Webkurs.

\subsection{Warum eine neue Theorie nötig war}
Die Vorgeschichte der Quantenmechanik ist eine Reihe von Beobachtungen, die die These, unser Universum werde durch die klassische Physik beschrieben, allmählich in Frage stellten.
\begin{enumerate}
	\item \emph{Schwarzkörperstrahlung} (Planck, 1900). Ein Schwarzkörper reflektiert keine einfallende Strahlung, strahlt bei einer Temperatur $T$ aber dennoch. Das Spektrum dieser Strahlung ist mit klassischer Physik nicht erklärbar. Das richtige liefert die Annahme, dass Licht der Frequenz $\nu$ aus \emph{Lichtquanten} der Energie $E = h\nu$ besteht. Planck hat die Konstante $h$ ursprünglich als ,,Hilfsgröße`` postuliert, um das Strahlungsgesetz herzuleiten:
	\begin{equation*}
		u(\nu, T)\,\dd\nu = \frac{8\pi h \nu^3}{c^3}\,\frac{\dd\nu}{e^{h\nu/(\kB T)} - 1}
	\end{equation*}
	ist die Energie pro Volumen im Frequenzintervall $\dd\nu$, mit der Lichtgeschwindigkeit $c$ und der Boltzmann-Konstante $\kB$. Dieses Gesetz leiten wir im statistischen Teil selbst her. In der theoretischen Physik verwendet man vorwiegend die \emph{reduzierte Planck-Konstante}
	\begin{equation*}
		\eqbox{\hbar = \frac{h}{2\pi}}
	\end{equation*}
	\item \emph{Photoeffekt} (Lenard; Einstein ab 1905). Licht kann Elektronen auslösen. Bei festem Material hängt ihre \emph{maximale} kinetische Energie von der Lichtfrequenz ab, nicht von der Intensität. Klassisch ist das merkwürdig: Sonst bedeutet mehr Intensität mehr Energie. Gibt ein Lichtquant seine Energie an ein einzelnes Elektron ab, so lautet die Bilanz
	\begin{equation*}
		E_{\mathrm{kin,max}} = \tfrac12 m_{\mathrm{e}} v_{\max}^2 = h\nu - W
	\end{equation*}
	mit der Bindungsenergie $W$ des Materials (nichtrelativistisch geschrieben). Unterhalb der Frequenz $W/h$ löst in diesem Modell ein einzelnes Lichtquant kein Elektron aus.
	\item \emph{Stabilität der Atome, Linienspektren, Franck-Hertz-Versuch}. Nach den Liénard-Wiechert-Potentialen strahlt jede beschleunigte Punktladung. Ein Elektron auf einer Bahn um den Kern müsste also ständig Energie verlieren und in den Kern stürzen, was es nicht tut. Linienspektren und der Franck-Hertz-Versuch (1914) deuten darauf hin, dass in Atomen nur bestimmte, diskrete Energieniveaus erlaubt sind.
	\item \emph{Compton-Effekt}. Wie sich bei der Streuung von Licht an Elektronen Energie und Impuls verteilen, beschreibt man richtig mit Lichtquanten der Energie $h\nu$ und des Impulses $p = h/\lambda$ und mit Energie- und Impulserhaltung.
	\item \emph{Stern-Gerlach-Versuch} (1922). Ein Strahl von Silberatomen läuft durch ein inhomogenes Magnetfeld. Klassisch sollte er sich kontinuierlich auffächern; beobachtet werden zwei getrennte Strahlen. Bei fester Drehimpulsquantenzahl nimmt die Komponente des Drehimpulses in Feldrichtung nur endlich viele Werte an. Aus solchen quantisierten Drehimpulsen kann man heute Qubits bauen.
	\item \emph{Interferenz bei Elektronen}. Die Beugung von Elektronen an Kristallen (Davisson und Germer, 1927) war ein Hinweis auf Welleneigenschaften des Elektrons, die de Broglie 1924 postuliert hatte.
\end{enumerate}
Ich fasse die Befunde in zwei Sätzen zusammen.

\begin{merksatz}
Materie und Strahlung besitzen sowohl Wellen- als auch Teilcheneigenschaften.
\end{merksatz}

\begin{merksatz}[Die wichtigste Aussage über die Quantenmechanik]
Die Quantenmechanik ist eine \emph{probabilistische Theorie}. Sie liefert uns somit nur Vorhersagen in Form von \emph{Wahrscheinlichkeiten}.
\end{merksatz}

Diese zweite Aussage ist der ersten übergeordnet: ,,Wellen- und Teilcheneigenschaften`` bleibt ein wenig schwammig, die probabilistische Aussage dagegen ist präzise und liegt allen Quantenphänomenen zugrunde. Klassisch lauten die Vorhersagen definitiv: Das Teilchen ist zu dieser Zeit an dieser Stelle. Wahrscheinlichkeiten treten dort allein wegen fehlender Information auf; fundamentale Wahrscheinlichkeiten gibt es in der klassischen Physik nicht. Die statistische Deutung hat Born 1926 geliefert, und sie ist für mich der wichtigste Aspekt der Theorie. Nehmen Sie den Satz bitte mit: In jeder mündlichen Prüfung frage ich nach der wichtigsten Aussage über die Quantenmechanik.

Die Interpretation der Quantenmechanik ist trotz vieler Ideen bis heute umstritten; einen Widerspruch zu ihr gibt es bis heute nicht. Unsere Annahme ist, dass die Quantenmechanik die korrekte Beschreibung des Universums ist.

\subsection{Was eine physikalische Theorie liefert}
Aus dem zweiten Merksatz mache ich einen Grundsatz für jede physikalische Theorie: Physikalische Theorien sind probabilistische Theorien, ihre Vorhersagen haben die Form von Wahrscheinlichkeiten. Deterministische Theorien fasse ich als Teilmenge der probabilistischen auf. Eine physikalische Theorie liefert Daten für Präparationen, Dynamik und Messung.

\begin{greybox}[Jedes Experiment in drei Schritten]
Für den Theoretiker sehen alle Experimente der Physik gleich aus: Etwas wird \emph{präpariert}, etwa ein Würfel oder im Labor ein Lichtstrahl; das Experiment wird \emph{durchgeführt}; am Ende \emph{messen} wir und erhalten eines der möglichen Ergebnisse.
% fig_vl01_experiment.tex — Jedes Experiment in drei Stadien: Präparation,
% Durchführung, Messung; die Klammer fasst die ersten beiden zu einer
% komplizierteren Präparation zusammen.
%
% Adaptiert aus content/c00-theorie/figures/experiment.tex (c00.fig.experiment),
% das wiederum den grauen Kasten ,,Jedes Experiment in drei Schritten``
% des Skripts v0.2 (vl01.tex:118-126) zeichnet.  Inhalt unverändert; nur die
% Stile sind die des Skripts (tpc-style.sty, tpc-macros.sty) statt der
% Design-Token des Kurses (tpc state, tpc dynamics, tpc effect), damit die
% Abbildung mit scripts/check_lecture.sh allein baut.
%
% Geometrie: ein Schema, keine Kurve.  Drei gleich große Kästen der Breite
% w = 2.9 und Höhe 1.3, getrennt durch Pfeile der Länge g = 0.9; der i-te
% Kasten (i = 0, 1, 2) beginnt bei x_i = i (w + g), also bei 0, 3.8, 7.6.
% Die Klammer reicht von x_0 = 0 bis x_1 + w = 6.7.
% Rollen wie im Kurs: Präparation durchgezogen, Durchführung mit dickerem
% Strich, Messung gestrichelt (graustufentauglich, Farbe nur zusätzlich).

\begin{center}
\begin{tikzpicture}[x=1cm, y=1cm]
  \def\w{2.9}\def\g{0.9}\def\h{1.3}
  % 1 Präparation: durchgezogen
  \draw[cgblue, line width=0.9pt] (0,0) rectangle (\w,\h);
  \node[lbl, align=center] at (0.5*\w,0.5*\h)
    {1\\[1pt]\tpcifenglish{preparation}{Präparation}};
  \draw[axisstyle] (\w,0.5*\h) -- (\w+\g,0.5*\h);
  % 2 Durchführung: dickerer Strich
  \draw[spacecadet, line width=1.6pt] (\w+\g,0) rectangle (2*\w+\g,\h);
  \node[lbl, align=center] at (1.5*\w+\g,0.5*\h)
    {2\\[1pt]\tpcifenglish{run}{Durchführung}};
  \draw[axisstyle] (2*\w+\g,0.5*\h) -- (2*\w+2*\g,0.5*\h);
  % 3 Messung: gestrichelt
  \draw[birrenorange, line width=0.9pt, dashed]
    (2*\w+2*\g,0) rectangle (3*\w+2*\g,\h);
  \node[lbl, align=center] at (2.5*\w+2*\g,0.5*\h)
    {3\\[1pt]\tpcifenglish{measurement}{Messung}};
  \draw[axisstyle] (3*\w+2*\g,0.5*\h) -- (3*\w+3*\g,0.5*\h);
  \node[lbl, anchor=west] at (3*\w+3*\g+0.05,0.5*\h)
    {\tpcifenglish{outcome}{Ergebnis}};
  % Klammer: Stadien 1 und 2 sind zusammen wieder eine Präparation
  \draw[birrengray, line width=0.6pt]
    (0,-0.25) -- (0,-0.45) -- (2*\w+\g,-0.45) -- (2*\w+\g,-0.25);
  \node[lbl, anchor=north] at (\w+0.5*\g,-0.5)
    {\tpcifenglish{together: one more complicated preparation}{zusammen: eine kompliziertere Präparation}};
\end{tikzpicture}
\end{center}

Die ersten beiden Schritte lassen sich zu \emph{einer} komplizierteren Präparation zusammenfassen: Ein gezinkter Würfel ist nichts anderes als ein normaler Würfel, an dem man vorher etwas gemacht hat. Wo man die Grenze zieht, ist Geschmackssache; oft kommt man mit Präparation und Messung aus.
\end{greybox}

Auch die Reihenfolge des Stoffs ist festgelegt: Wir gehen vom Mikroskopischen zum Makroskopischen, von der Quantenmechanik einzelner Systeme über viele Teilchen zur statistischen Mechanik und Thermodynamik. Auf jeder Stufe lernt man, dass jede Theorie eine effektive Theorie der darunterliegenden ist, eine Näherung an die grundlegende Physik darunter.

\subsection{Wahrscheinlichkeit für Physiker}
Die Zahlen, welche die Quantenmechanik als Ergebnis liefert, sind Wahrscheinlichkeiten. Erwarten Sie nichts konzeptionell Neues: Ein intuitives Verständnis setze ich voraus; es geht um die \emph{Notation}, in der wir später die Quantenmechanik formulieren. Wir sprechen dabei die Sprache der Physik: Was die Mathematik eine Wahrscheinlichkeitsverteilung nennt, heißt hier \emph{Präparation} (Zustand), und eine Zufallsvariable heißt \emph{Messung}.

\subsubsection{Ergebnisse und Präparationen}
Um ein Experiment zu modellieren, brauchen wir drei Dinge: die Menge der möglichen Ergebnisse, die Präparation und die Messung. Jedem System wird eine Menge $X$ zugeordnet, die \emph{Ergebnismenge} (Stichprobenraum); man erhält sie, indem man alle möglichen Ergebnisse auflistet. Teilmengen $S \subset X$ heißen \emph{Ereignisse}. Zunächst sei $X$ endlich. Eine \emph{Präparation} wird durch eine Wahrscheinlichkeit $p(x)$ für jedes Ergebnis $x$ charakterisiert, ein \emph{Gewicht} mit
\begin{equation}
	p(x) \geq 0 \quad \text{und} \quad \sum_{x \in X} p(x) = 1.
\end{equation}
Die \emph{Wahrscheinlichkeit} eines Ereignisses ist $\mathbb{P}[S] = \sum_{x \in S} p(x)$.

\begin{beispiel}
	Bei einem Würfel ist $X = \{1, 2, 3, 4, 5, 6\}$, und ,,die Augenzahl ist gerade`` ist das Ereignis $S = \{2, 4, 6\}$. Für einen fairen Würfel ist $p(x) = 1/6$ für alle $x$, also $\mathbb{P}[S] = 1/2$.
\end{beispiel}

\begin{übung}
	Beschreiben Sie die Ergebnismenge für \emph{zwei} Würfel. Wie viele Elemente hat sie?
\end{übung}

Für paarweise disjunkte Ereignisse $S_1, \dots, S_k$ gilt
\begin{equation}
	\mathbb{P}[S_1 \cup \dots \cup S_k] = \mathbb{P}[S_1] + \dots + \mathbb{P}[S_k].
\end{equation}
Die Disjunktheit ist wesentlich: Überlappen sich die Ereignisse, so zählt die rechte Seite Ergebnisse doppelt. Eine solche Mengenfunktion heißt in der Mathematik ein \emph{Maß}, mit $\mathbb{P}[X] = 1$ ein \emph{Wahrscheinlichkeitsmaß}. Mehr Maßtheorie brauchen wir nicht.

\subsubsection{Messungen und Erwartungswert}
Eine \emph{Messung} (Observable, Zufallsvariable) ist eine Funktion $f\colon X \to \mathbb{R}$. Der Wert $f(x)$ ist der Messwert, wenn sicher das Ergebnis $x$ vorliegt. Beachten Sie, wie großzügig das ist: An eine Messung stellen wir keine Bedingung, an die Präparation zwei. Präparation und Messung zusammen ergeben den \emph{Erwartungswert}
\begin{equation}
	\eqbox{\mathbb{E}[f] = \sum_{x \in X} p(x) f(x)}
\end{equation}
Er ist linear, $\mathbb{E}[\alpha f + \gamma g] = \alpha\,\mathbb{E}[f] + \gamma\,\mathbb{E}[g]$, und für eine Konstante ist $\mathbb{E}[c] = c$.

Woher kommt die Formel? Man wiederholt das Experiment $N$-mal mit derselben Präparation und erhält die Ergebnisse $x_1, \dots, x_N$. Kommt $x$ dabei $N(x)$-mal vor, so ist $h_N(x) = N(x)/N$ die \emph{relative Häufigkeit}; sie ist selbst ein Gewicht, und für den Mittelwert der Messreihe gilt exakt
\begin{equation}
	\frac{1}{N} \sum_{i=1}^{N} f(x_i) = \sum_{x \in X} h_N(x)\, f(x).
\end{equation}
Bei vielen unabhängigen Wiederholungen liegt $h_N(x)$ mit hoher Wahrscheinlichkeit nahe bei $p(x)$, und der Mittelwert der Messreihe nahe beim Erwartungswert. Diese Lesart ist eine Deutung: In einer einzelnen Messreihe kann $h_N$ von $p$ abweichen.

\subsubsection{Ja-Nein-Messungen}
Die wichtigsten Messungen nehmen nur die Werte $0$ und $1$ an. Sie sind das Modell für Detektoren, etwa einen Photonendetektor: Klick oder kein Klick. An diesen zweiwertigen Messungen hängt die Quantenmechanik. Eine solche Messung ist durch die Menge $S = \{x \in X \colon f(x) = 1\}$ bestimmt; man schreibt
\begin{equation}
	\chi_S(x) = \begin{cases} 1, & x \in S, \\ 0, & \text{sonst,} \end{cases}
\end{equation}
und nennt $\chi_S$ die \emph{Indikatorfunktion} von $S$. Halten Sie fest: Indikatorfunktionen \emph{sind} Messungen, kein eigenes Objekt daneben. Ihr Erwartungswert ist die Wahrscheinlichkeit ihres Ereignisses:
\begin{equation}\label{eq:indicatorfunction}
	\mathbb{E}[\chi_S] = \sum_{x \in X} p(x)\, \chi_S(x) = \sum_{x \in S} p(x) = \mathbb{P}[S].
\end{equation}
Wahrscheinlichkeiten sind also besondere Erwartungswerte. Umgekehrt legt $\mathbb{P}$ alle Erwartungswerte fest, denn $p(x) = \mathbb{P}[\{x\}]$.

\begin{übung}
	Sei $X = \{1, 2, 3\}$ und $p(x) = 1/3$ für alle $x$. Berechnen Sie $\mathbb{P}[S]$ für $S = \{1, 2\}$ und den Erwartungswert der Messung $f(x) = x^2$.
\end{übung}

\subsubsection{Varianz und Determinismus}
Der Erwartungswert sagt, wo die Messwerte im Mittel liegen; die \emph{Varianz} sagt, wie weit sie streuen:
\begin{equation}
	\operatorname{Var}[f] = \mathbb{E}\bigl[(f - \mathbb{E}[f])^2\bigr] = \mathbb{E}[f^2] - \mathbb{E}[f]^2 .
\end{equation}
Die zweite Form folgt aus der Linearität. Ausgeschrieben ist $\operatorname{Var}[f] = \sum_x p(x)\,(f(x) - \mathbb{E}[f])^2$, eine Summe von Zahlen $\geq 0$. Also ist $\operatorname{Var}[f] \geq 0$, und $\operatorname{Var}[f] = 0$ gilt genau dann, wenn $f(x) = \mathbb{E}[f]$ für alle $x$ mit $p(x) > 0$ ist. In der Quantenmechanik spielt die Varianz eine große Rolle; bei der Unschärferelation begegnet sie uns wieder.

\begin{beispiel}
	Für den fairen Würfel und die Augenzahl $f(x) = x$ ist $\mathbb{E}[f] = 21/6 = 7/2$; der Erwartungswert muss kein möglicher Messwert sein. Mit $\mathbb{E}[f^2] = 91/6$ ist $\operatorname{Var}[f] = 91/6 - 49/4 = 35/12$.
\end{beispiel}

Für eine Ja-Nein-Messung ist $\chi_S^2 = \chi_S$, also
\begin{equation}
	\operatorname{Var}[\chi_S] = \mathbb{P}[S]\,\bigl(1 - \mathbb{P}[S]\bigr).
\end{equation}
Sie verschwindet genau dann, wenn die Antwort feststeht. Ist $p(a) = 1$ für ein Ergebnis $a$, so ist $\mathbb{E}[f] = f(a)$ für jede Messung.

\begin{merksatz}[Determinismus ist ein Sonderfall]
Für ein Gewicht $p$ auf einer endlichen Ergebnismenge sind gleichwertig: $p(a) = 1$ für ein Ergebnis $a$; jedes Ereignis hat die Wahrscheinlichkeit $0$ oder $1$; jede Messung hat die Varianz $0$. Solche Präparationen nennen wir deterministisch: Das Ergebnis jeder Messung steht fest. So passen die definitiven Vorhersagen der klassischen Physik in eine probabilistische Theorie.
\end{merksatz}

\subsubsection{Kontinuierliche Ergebnisse}
Für die Quantenmechanik ist $X = \mathbb{R}$ entscheidend: Ein Teilchen in einer Dimension hat eine Position, und die möglichen Positionen sind die reellen Zahlen. Wir betrachten nur Präparationen mit einer \emph{Wahrscheinlichkeitsdichte}, einer Funktion $p$ auf $\mathbb{R}$ mit
\begin{equation}\label{eq:ctsprob}
	\begin{aligned}
		p(x) &\geq 0, & \int_{-\infty}^{\infty} p(x)\,\dd x &= 1, \\
		\mathbb{P}[[a,b]] &= \int_a^b p(x)\,\dd x, & \mathbb{E}[f] &= \int_{-\infty}^{\infty} p(x) f(x)\,\dd x .
	\end{aligned}
\end{equation}
Dabei ist $a \leq b$, und den Erwartungswert verwenden wir nur, sofern $\int p(x) \abs{f(x)}\,\dd x < \infty$ ist. Gegenüber dem endlichen Fall ist also nur die Summe durch ein Integral ersetzt. In der Literatur steht dafür auch $\int \mathbb{P}(\dd x)\, f(x)$; ich mag diese Schreibweise nicht besonders, Sie sollen sie nur erkennen. Wichtig: Der Wert $p(x)$ einer Dichte ist keine Wahrscheinlichkeit. Wahrscheinlichkeiten sind Flächen unter der Kurve; ein einzelner Punkt hat die Wahrscheinlichkeit $0$, und $p(x)$ darf größer als $1$ sein.

\begin{beispiel}
	Die Dichte
	\begin{equation}
		p(x) = \frac{1}{\sqrt{2\pi}}\, e^{-\frac12 x^2}
	\end{equation}
	heißt \emph{Normal-} oder \emph{Gauß-Verteilung}; sie ist normiert. Ist das Teilchen links oder rechts vom Ursprung? Weil $p(-x) = p(x)$ ist, sind die Flächen links und rechts vom Nullpunkt gleich groß, zusammen $1$; also ist $\int_0^\infty p(x)\,\dd x = 1/2$. Vorsicht: Das Argument funktioniert \emph{nur} für symmetrische Dichten.
\end{beispiel}

Die Erwartungswerte der Koordinatenfunktion und ihrer Potenzen heißen \emph{Momente}, sofern die Integrale $\int \abs{x}^n p(x)\,\dd x$ endlich sind:
\begin{equation}
	M_n = \int_{-\infty}^{\infty} x^n\, p(x)\,\dd x .
\end{equation}
Nur $M_0 = 1$, die Normierung, existiert immer; $M_1$ heißt, sofern es existiert, \emph{Mittelwert}. Existieren $M_1$ und $M_2$, so ist die Varianz der Koordinate wie im endlichen Fall
\begin{equation}
	\operatorname{Var}[\mathbb{P}] = \int_{-\infty}^{\infty} (x - M_1)^2\, p(x)\,\dd x = M_2 - M_1^2 .
\end{equation}
Mittelwert und Varianz tragen dann dieselbe Information wie $M_1$ und $M_2$. Für die hier betrachteten Modelle steht damit die Sprache bereit: Präparationen sind Gewichte oder Dichten, Messungen sind Funktionen, und beide zusammen ergeben bei endlichen Integralen Erwartungswerte.

%!TEX root = TheoretischePhysikC.tex

\section{Vorlesung: Hilbertraum}
Um die Zustände der Quantenmechanik zu beschreiben wird der \emph{Hilbertraum} eine größere Rolle spielen. Um diesen mathematischen Raum zu definieren, benötigen wir erst \emph{komplexe Vektorräume}. Solche Räume werden analog mit dem reellen Fall definiert. 

Warum hier \emph{komplexe} und nicht reelle Vektoren auftreten, wissen wir nicht: die Natur ist nun einmal so gebaut, und eine bessere Alternative wurde bisher nicht gefunden.

Ein \emph{komplexer Vektor} $\mbf{v}\in \mathbb{C}^n$ ist eine Spalte von $n$ komplexen Zahlen
\begin{equation}
	\mbf{v} = 
	\begin{pmatrix}
		v_1 \\ v_2 \\ \vdots \\ v_n
	\end{pmatrix}.
\end{equation}
Es gelten die folgenden Rechenoperationen
\begin{enumerate}
	\item Multiplikation mit Skalar $a\in\mathbb{C}$
	\begin{equation}
		a\mbf{v} = 
	\begin{pmatrix}
		av_1 \\ av_2 \\ \vdots \\ av_n
	\end{pmatrix}.
	\end{equation}
	\item Addition
	\begin{equation}
		\mbf{v}+\mbf{w} = \begin{pmatrix}
		v_1 \\ v_2 \\ \vdots \\ v_n
	\end{pmatrix} + \begin{pmatrix}
		w_1 \\ w_2 \\ \vdots \\ w_n
	\end{pmatrix} =
	\begin{pmatrix}
		v_1+w_1 \\ v_2+w_2 \\ \vdots \\ v_n+w_n
	\end{pmatrix}.
	\end{equation}
	\item Kreuz-Operation
	\begin{equation}
		v^\dag = \begin{pmatrix}
		v_1 \\ v_2 \\ \vdots \\ v_n
	\end{pmatrix}^\dag \equiv \begin{pmatrix}
		v_1^* & v_2^* & \cdots & v_n^*
	\end{pmatrix}
	\end{equation}
	\item Skalarprodukt
	\begin{equation}
		\mbf{v}^\dag\mbf{w} = \begin{pmatrix}
		v_1^* & v_2^* & \cdots & v_n^*
	\end{pmatrix} \begin{pmatrix}
		w_1 \\ w_2 \\ \vdots \\ w_n
	\end{pmatrix} =
	\sum_{j=1}^n v_j^*w_j.
	\end{equation}
\end{enumerate}
Der Raum $V$ aller solcher Vektoren mit festem $n$ wird als \emph{komplexer Vektorraum} bezeichnet. 

Für den Vektorraum selbst genügen die ersten beiden Regeln; Kreuz-Operation und Skalarprodukt sind Zusatzstruktur, die ihn erst zum Hilbertraum macht.

Hierbei wird $n$ als \emph{Dimension} von $V$ definiert. Oft verwendet wird die Notation $\mathbb{C}^n$ für den oben definierten Vektorraum $V$.

\begin{beispiel}
	Der Vektorraum $\mathbb{C}^2$ besteht aus allen zweidimensionalen Vektoren
	\begin{equation}
		\mbf{v} = 
		\begin{pmatrix}
			v_1 \\ v_2
		\end{pmatrix}.
	\end{equation}
	Seien
	\begin{equation}
		\mbf{v} = \begin{pmatrix}
			1+i \\ 1-i
		\end{pmatrix} \quad \text{und} \quad \mbf{w} = \begin{pmatrix}
			1 \\ -1
		\end{pmatrix}
	\end{equation}
	zwei Vektoren. Deren Skalarprodukt lautet
	\begin{equation}
		\mbf{v}^\dag \mbf{w} = \begin{pmatrix}
			1-i & 1+i
		\end{pmatrix} \begin{pmatrix}
			1 \\ -1
		\end{pmatrix} = -2i.
	\end{equation}
\end{beispiel}

Unendlich dimensionale Vektorräume sind ebenfalls möglich und spielen eine sehr wichtige Rolle in der Beschreibung quantenmechanischer Systeme. Der Grund ist physikalisch: unsere besten Modelle sind Felder, und Felder haben unendlich viele Freiheitsgrade. Wie sind solche Vektorräume zu verstehen? Man kann sich unendlich dimensionale Vektoren als Spalten mit unendlich vielen Einträge vorstellen, z.B.:
\begin{equation}
	\mbf{v} = \begin{pmatrix}
		v_1 \\ v_2 \\ \vdots \\
	\end{pmatrix}
\end{equation}
Alle Rechenregeln für Vektorräume gelten auch in den unendlichen dimensionalen Fällen. Das Skalarprodukt $\mbf{v}^\dag \mbf{w}$ wird wie folgt berechnet
\begin{equation}
	\mbf{v}^\dag\mbf{w} =
	\sum_{j=1}^\infty v_j^*w_j.
\end{equation}
Das heißt, in der Regel, dass man in allen Gleichungen $n=\infty$ setzt. Es gibt mehrere subtile Aspekte, die man in Betracht ziehen müsste, um rigorös mit unendlich dimensionalen Vektorräumen zu arbeiten. Allerdings werden wir in dieser Vorlesung diese Aspekte nur nebenbei erwähnen.

Einer davon tritt sofort auf. Der Vektor mit lauter Einsen,
\begin{equation}
	\mbf{v} = \begin{pmatrix}
		1 \\ 1 \\ 1 \\ \vdots
	\end{pmatrix},
\end{equation}
ist eine legale unendliche Liste, aber sein Skalarprodukt mit sich selbst divergiert:
\begin{equation}
	\mbf{v}^\dag\mbf{v} = \sum_{j=1}^\infty 1 = \infty.
\end{equation}
Physikalisch ist die Antwort darauf einfach --- wir lassen solche Vektoren nicht zu. Erlaubt ist ein Vektor genau dann, wenn
\begin{equation}
	\eqbox{\mbf{v}^\dag\mbf{v} = \sum_{j=1}^\infty |v_j|^2 < \infty}
\end{equation}
gilt. Diese Menge ist wieder ein komplexer Vektorraum; man bezeichnet ihn mit $\ell^2$, wobei die Zwei vom Quadrat in der Summe kommt. Damit kann bei der Berechnung von Wahrscheinlichkeiten nichts unendlich groß werden.

Für uns ist das wichtigste Beispiel einen unendlich dimensionalen Vektorraumes der Raum der \emph{quadratintegrierbaren Funktionen} (der als $L^{2}(\mathbb{R})$ bezeichnet wird) welcher $\mathbb{R}$ auf $\mathbb{C}$ abbildet. Hierbei bedeutet quadratintegrierbar, dass
\begin{equation}
	\int_{-\infty}^\infty |f(x)|^2 \, dx < \infty.
\end{equation}
Quadratintegrierbare Funktionen können addiert, subtrahiert, und mit Skalargrößen multipliziert werden. Bezüglich dieser Operationen bleiben die Ergebnisse immer quadratintegrierbar. Daher bildet die Menge aller quadratintegrierbaren Funktionen einen abstrakten Vektorraum. Man betrachtet eine quadratintegrierbare Funktion $f(x)$ als einen Vektor. Auch hier gibt es ein Skalarprodukt; an die Stelle der Summe tritt ein Integral:
\begin{equation}
	\langle f, g\rangle \equiv \int_{-\infty}^\infty f^*(x)\, g(x)\, dx.
\end{equation}

\begin{beispiel}
	Ein \emph{Nicht}-Beispiel von einer quadratintegrierbaren Funktion wird von der konstanten Funktion $f(x) = 1$ geliefert:
	\begin{equation}
		\int_{-\infty}^\infty |f(x)|^2\, dx = \int_{-\infty}^\infty 1\, dx = \infty.
	\end{equation}
	Das ist das Gegenstück zum Vektor mit lauter Einsen von vorhin.
\end{beispiel}

% fig_vl02_raeume.tex — Der Zulassungstest für Hilbertraum-Vektoren:
% endliche Norm.  Links die Summe (ℓ²), rechts das Integral (L²(R)) —
% dieselbe Frage, einmal über einen diskreten und einmal über einen
% kontinuierlichen Index.
%
% Was die Abbildung zeigt, das der Text nicht sagt: es genügt nicht, dass
% die Einträge (bzw. die Funktionswerte) verschwinden — es kommt auf die
% Rate an.  1/j und 1/sqrt(j) gehen beide gegen null; nur die erste Folge
% ist in ℓ².  Das ist ein Skalierungsverhalten, und ein Skalierungs-
% verhalten sieht man, statt es zu erklären.
%
% Geometrie — jede Kurve ist die geschlossene Form der jeweiligen
% Partialsumme bzw. des Partialintegrals, nicht freihand:
%
%   (a)  S_N = sum_{j<=N} |v_j|^2
%          v_j = 1/j      : S_N -> pi^2/6 = 1.6449...   (Basler Reihe)
%          v_j = 1/sqrt(j): S_N = H_N ~ ln N + gamma    (log-divergent)
%          v_j = 1        : S_N = N                     (linear divergent)
%        Die Punkte sind die exakt aufsummierten Werte (N = 1..25); die
%        Asymptote liegt bei pi^2/6, nicht bei einem gerundeten Strich.
%
%   (b)  F(X) = int_{-X}^{X} |f|^2 dx
%          f = e^{-x^2/2}    : F = sqrt(pi) erf(X) -> sqrt(pi) = 1.7725...
%          f = 1/sqrt(1+|x|) : F = 2 ln(1+X)             (log-divergent)
%          f = 1             : F = 2X                    (linear divergent)
%
% Beide Panels teilen die y-Skala (0 bis 4.6), damit die beiden endlichen
% Grenzwerte pi^2/6 und sqrt(pi) direkt vergleichbar nebeneinander liegen
% und die beiden log-Divergenzen erkennbar dieselbe Form haben.
%
% Entwurfsentscheidungen:
%   * Farbe trägt genau eine Bedeutung, nämlich das Ergebnis des Tests:
%     kalt (cgblue) = zugelassen, warm = durchgefallen.  Innerhalb von
%     „warm“ unterscheidet der Ton den subtilen vom groben Fall
%     (birrenorange: Einträge gehen gegen null und trotzdem divergent;
%     birrenred: das Beispiel der Vorlesung, v_j = 1 bzw. f = 1).
%     Zusätzlich unterscheiden Strichmuster die drei Kurven, damit die
%     Aussage den Graustufendruck überlebt.
%   * Die Grenzwerte pi^2/6 und sqrt(pi) sind neutral graue, gestrichelte
%     Referenzlinien (levellinegrey) — sie sind keine Daten.
%   * Panel (a) trägt Marker auf jedem ganzzahligen N, Panel (b) nicht:
%     die Textur der Kurve zeigt selbst, dass links über einen diskreten
%     und rechts über einen kontinuierlichen Index summiert wird.
%   * Die divergenten Kurven laufen absichtlich oben aus dem Rahmen; sie
%     sind dort beschriftet, wo sie ihn verlassen.  Die log-divergente
%     Kurve tut das nicht — sie biegt nur ab und koennte als saettigend
%     missverstanden werden.  Deshalb steht ihr Wachstumsgesetz
%     (~ ln N bzw. ~ 2 ln X) als dritte Zeile an ihrer Beschriftung.
%   * ymax = 5 statt 4.6: die Vorlesung setzt 12pt, dort werden
%     \footnotesize-Knoten spuerbar groesser als in der 10pt-Vorschau.
%     Die Beschriftungen wurden gegen eine 12pt-Vorschau gesetzt, nicht
%     gegen scripts/tikz-preview.sh allein — bei 10pt sah eine Kollision
%     zwischen der blauen Kurvenbeschriftung und der C^n-Randnotiz noch
%     harmlos aus, im Satz ueberlappten die beiden.
%   * C^n erscheint als graue Randnotiz statt als vierte Kurve: dort hat
%     die Summe n Terme und ist immer endlich — es gibt nichts zu prüfen.
%     Das ist der ganze Unterschied zwischen C^2 und den beiden anderen
%     Räumen der Vorlesung.

\begin{center}
\begin{tikzpicture}

% ─────────────────────────────────────────────────────────────
%  (a)  Partialsummen: der Test in ℓ²
% ─────────────────────────────────────────────────────────────
\begin{axis}[
    name=seq,
    tpcplot,
    width=6.6cm, height=5.9cm,
    xlabel={$N$},
    ylabel={$\sum_{j=1}^{N}|v_j|^{2}$},
    every axis y label/.style={at={(axis description cs:-0.14,0.5)},
                               rotate=90, anchor=south, font=\small},
    xmin=0, xmax=26, ymin=0, ymax=5.0,
    xtick={0,5,10,15,20,25}, ytick={0,1,2,3,4},
    grid=major,
    title={(a) \tpcifenglish{Partial sums in $\ell^2$}%
                            {Partialsummen in $\ell^2$}},
]

  % Grenzwert der Basler Reihe: neutrale Referenzlinie, keine Daten
  \addplot[levellinegrey, domain=0:26, samples=2] {1.6449341};

  % v_j = 1/j : S_N = sum 1/j^2 -> pi^2/6
  \addplot[cgblue, line width=1.0pt, mark=*, mark size=1.0pt,
           mark options={draw=cgblue, fill=cgblue}] coordinates {
    (1,1.0000)(2,1.2500)(3,1.3611)(4,1.4236)(5,1.4636)(6,1.4914)
    (7,1.5118)(8,1.5274)(9,1.5398)(10,1.5498)(11,1.5580)(12,1.5650)
    (13,1.5709)(14,1.5760)(15,1.5804)(16,1.5843)(17,1.5878)(18,1.5909)
    (19,1.5937)(20,1.5962)(21,1.5984)(22,1.6005)(23,1.6024)(24,1.6041)
    (25,1.6057)};

  % v_j = 1/sqrt(j) : S_N = H_N, logarithmisch divergent
  \addplot[birrenorange, line width=1.0pt, densely dashed,
           mark=*, mark size=1.0pt,
           mark options={solid, draw=birrenorange, fill=birrenorange}]
  coordinates {
    (1,1.0000)(2,1.5000)(3,1.8333)(4,2.0833)(5,2.2833)(6,2.4500)
    (7,2.5929)(8,2.7179)(9,2.8290)(10,2.9290)(11,3.0199)(12,3.1032)
    (13,3.1801)(14,3.2516)(15,3.3182)(16,3.3807)(17,3.4396)(18,3.4951)
    (19,3.5477)(20,3.5977)(21,3.6454)(22,3.6908)(23,3.7343)(24,3.7760)
    (25,3.8160)};

  % v_j = 1 : S_N = N, das Nicht-Beispiel der Vorlesung
  \addplot[birrenred, line width=1.0pt, dash pattern=on 4pt off 2pt on 1pt off 2pt,
           mark=*, mark size=1.0pt,
           mark options={solid, draw=birrenred, fill=birrenred}]
  coordinates {(1,1)(2,2)(3,3)(4,4)(5.3,5.3)};

  % Beschriftung direkt an den Kurven
  \node[lbl, text=cgblue, align=left, anchor=west]
    at (axis cs:5.4,0.80) {$v_j=\frac{1}{j}$\\[1pt]$\in\ell^2$};
  \node[lbl, text=birrenorange, align=left, anchor=west]
    at (axis cs:14.6,2.45) {$v_j=\frac{1}{\sqrt{j}}$\\[1pt]$\notin\ell^2$\\[1pt]
                          {\scriptsize$\sim\ln N$}};
  \node[lbl, text=birrenred, align=left, anchor=west]
    at (axis cs:5.4,4.35) {$v_j=1$\\[1pt]$\notin\ell^2$};

  % Grenzwert der Basler Reihe, direkt über seiner Referenzlinie
  \node[lbl, text=birrengray, font=\scriptsize, anchor=west]
    at (axis cs:6.0,1.80) {$\pi^{2}/6$};

  % C^n: der endlichdimensionale Fall, in dem nichts zu prüfen ist
  \node[lbl, text=birrengray, font=\scriptsize, align=right,
        anchor=south east]
    at (axis cs:25.6,0.10)
    {$\C^{n}$: $n$ \tpcifenglish{terms}{Terme}\\[1pt]
     \tpcifenglish{always finite}{stets endlich}};
\end{axis}

% ─────────────────────────────────────────────────────────────
%  (b)  Partialintegrale: derselbe Test in L^2(R)
% ─────────────────────────────────────────────────────────────
\begin{axis}[
    tpcplot,
    at={(seq.south east)}, anchor=south west, xshift=11mm,
    width=6.6cm, height=5.9cm,
    xlabel={$X$},
    ylabel={$\int_{-X}^{X}|f|^{2}\,dx$},
    every axis y label/.style={at={(axis description cs:-0.14,0.5)},
                               rotate=90, anchor=south, font=\small},
    xmin=0, xmax=4.3, ymin=0, ymax=5.0,
    xtick={0,1,2,3,4}, ytick={0,1,2,3,4},
    grid=major,
    title={(b) \tpcifenglish{Partial integrals in $L^2(\R)$}%
                            {Partialintegrale in $L^2(\R)$}},
]

  % Grenzwert des Gauss-Integrals: neutrale Referenzlinie
  \addplot[levellinegrey, domain=0:4.3, samples=2] {1.7724539};

  % f = exp(-x^2/2) : F(X) = sqrt(pi) erf(X) -> sqrt(pi)
  \addplot[cgblue, line width=1.0pt] coordinates {
    (0.00,0.0000)(0.10,0.1993)(0.20,0.3947)(0.30,0.5825)(0.40,0.7593)
    (0.50,0.9226)(0.60,1.0703)(0.70,1.2014)(0.80,1.3153)(0.90,1.4125)
    (1.00,1.4936)(1.10,1.5601)(1.20,1.6135)(1.30,1.6555)(1.40,1.6879)
    (1.50,1.7124)(1.60,1.7305)(1.70,1.7437)(1.80,1.7531)(1.90,1.7597)
    (2.00,1.7642)(2.10,1.7672)(2.20,1.7692)(2.30,1.7704)(2.40,1.7712)
    (2.50,1.7717)(2.60,1.7720)(2.70,1.7722)(2.80,1.7723)(2.90,1.7724)
    (3.00,1.7724)(3.20,1.7724)(3.40,1.7725)(3.60,1.7725)(3.80,1.7725)
    (4.00,1.7725)(4.30,1.7725)};

  % f = 1/sqrt(1+|x|) : F(X) = 2 ln(1+X), logarithmisch divergent
  \addplot[birrenorange, line width=1.0pt, densely dashed,
           domain=0:4.3, samples=120] {2*ln(1+x)};

  % f = 1 : F(X) = 2X, das Nicht-Beispiel der Vorlesung
  \addplot[birrenred, line width=1.0pt,
           dash pattern=on 4pt off 2pt on 1pt off 2pt,
           domain=0:2.7, samples=2] {2*x};

  % Beschriftung direkt an den Kurven
  \node[lbl, text=cgblue, align=left, anchor=west]
    at (axis cs:1.60,0.92) {$f=e^{-x^{2}/2}$\\[1pt]$\in L^2(\R)$};
  \node[lbl, text=birrenorange, align=left, anchor=west]
    at (axis cs:2.42,4.05) {$f=\frac{1}{\sqrt{1+|x|}}$\\[1pt]$\notin L^2(\R)$\\[1pt]
                          {\scriptsize$\sim 2\ln X$}};
  \node[lbl, text=birrenred, align=left, anchor=west]
    at (axis cs:0.10,4.35) {$f=1$\\[1pt]$\notin L^2(\R)$};

  % Grenzwert des Gauss-Integrals, direkt über seiner Referenzlinie
  \node[lbl, text=birrengray, font=\scriptsize, anchor=south]
    at (axis cs:3.45,1.82) {$\sqrt{\pi}$};

\end{axis}

\end{tikzpicture}
\end{center}


Achten Sie auf die mittlere Kurve in beiden Bildern: bei $v_j = 1/\sqrt{j}$ und
$f(x) = 1/\sqrt{1+|x|}$ gehen die Einträge sehr wohl gegen null, die Norm aber
logarithmisch gegen unendlich. Nicht das Verschwinden entscheidet, sondern die
Rate.

\begin{greybox}[Drei Räume, drei Systeme]
	Damit haben wir praktisch alle Hilberträume beisammen, die wir brauchen werden: $\mathbb{C}^2$ gehört zum \emph{Qubit}, $\ell^2$ zum \emph{harmonischen Oszillator} (etwa einem Photon in einer Glasfaser) und $L^2(\mathbb{R})$ zu einem \emph{Teilchen in einer Dimension}.
\end{greybox}



%!TEX root = TheoretischePhysikC.tex
% Kern von content/c03-hilbertraum, c03.s4 bis c03.s7.1. Quellen: docs/plan/script-v0.3/sync-ch3.md

\section{Vorlesung: Dirac-Notation, Hilbertraum, Basis}

\subsection{Dirac-Notation}
Sei $V$ ein komplexer Vektorraum; denken Sie vorerst an $\mathbb{C}^n$ aus der letzten Vorlesung. In der Quantenmechanik wird die folgende Schreibweise universell verwendet: Ab sofort schreiben wir jeden Vektor $\vb{v}\in V$ als \emph{Ket} $\ket{v}$. Den Fettdruck lassen wir im Ket weg; Strich und spitze Klammer sagen schon, dass ein Vektor gemeint ist. Mehr bedeutet die Notation zunächst nicht, und warum wir sie einführen, wird bald klar. Die eigentliche Hürde am Anfang ist nicht die Physik, sondern die Notation.

\begin{beispiel}
	Sei $V = \mathbb{C}^2$, der Raum, mit dem wir die Zustände eines Qubits beschreiben werden. Aus den beiden Vektoren
	\begin{equation}
		\ket{e_0} = \begin{pmatrix} 1 \\ 0 \end{pmatrix}, \qquad
		\ket{e_1} = \begin{pmatrix} 0 \\ 1 \end{pmatrix}
	\end{equation}
	bilden wir durch Addieren und Multiplizieren mit Zahlen zwei neue,
	\begin{equation}
		\ket{e_\pm} = \tfrac{1}{\sqrt{2}}\ket{e_0} \pm \tfrac{1}{\sqrt{2}}\ket{e_1}
		= \frac{1}{\sqrt{2}} \begin{pmatrix} 1 \\ \pm 1 \end{pmatrix}.
	\end{equation}
	Die Vektoren $\ket{e_\pm}$ tauchen immer wieder auf, beim Drehimpuls, beim Wasserstoffatom und besonders bei Spinsystemen.
\end{beispiel}

Dabei steht $\ket{v + w}$ für $\ket{v} + \ket{w}$ und $\ket{zv}$ für $z\ket{v}$, wenn $z$ eine Zahl ist.

Ergänzt wird die Ket-Notation durch die \emph{Bra}-Vektoren. Mit dem Skalarprodukt $\vb{v}\adj\vb{w}$ ist der Bra $\bra{v}$ die konjugierte Zeile, die die Kreuz-Operation aus der Spalte $\ket{v}$ macht. Das Produkt von Bra und Ket heißt \emph{Bra-Ket}; es ist das Skalarprodukt aus der letzten Vorlesung:
\begin{equation}
	\bra{v} = (\ket{v})\adj = \begin{pmatrix} v_1^* & \cdots & v_n^* \end{pmatrix}, \qquad
	\braket{v}{w} = (\ket{v})\adj\ket{w} = \vb{v}\adj\vb{w}.
\end{equation}
Das ist keine neue Definition, sondern eine neue Notation für dieselben Objekte.

\begin{historical}
	Die Notation stammt von Paul Dirac. Englisch \emph{bracket} heißt Klammer: links steht der ,,bra``, rechts der ,,ket``, und das ,,c`` ist irgendwo verlorengegangen.
\end{historical}

Ihren Vorteil spielt die Notation mit einer Regel für Linearkombinationen aus. Für komplexe Zahlen $\alpha,\beta$ gilt
\begin{equation}
	(\alpha\ket{v} + \beta\ket{w})\adj = \alpha^*\bra{v} + \beta^*\bra{w},
\end{equation}
denn die Kreuz-Operation konjugiert jeden Eintrag. Damit rechnen wir Skalarprodukte aus, ohne die Vektoren als Spalten hinzuschreiben. In unendlich vielen Dimensionen passt ein Vektor nicht in eine endliche Spalte, und dort rechnen wir genau so.

\begin{beispiel}
	Seien $\ket{v} = \ket{e_0} + i\ket{e_1}$ und $\ket{w} = \ket{e_0} + \ket{e_1}$. Nach der Regel ist $\bra{v} = \bra{e_0} - i\bra{e_1}$; beachten Sie das Minuszeichen. Mit $\braket{e_0}{e_0} = \braket{e_1}{e_1} = 1$ und $\braket{e_0}{e_1} = \braket{e_1}{e_0} = 0$ ist
	\begin{equation}
		\begin{split}
			\braket{v}{w} &= (\bra{e_0} - i\bra{e_1})(\ket{e_0} + \ket{e_1}) \\
			&= \braket{e_0}{e_0} + \braket{e_0}{e_1} - i\braket{e_1}{e_0} - i\braket{e_1}{e_1} = 1 - i.
		\end{split}
	\end{equation}
	Mit Spalten gerechnet ebenso: $1\cdot 1 + (-i)\cdot 1 = 1 - i$.
\end{beispiel}

Seien wir ehrlich: In zwei Dimensionen ist der erste Weg aufwendiger. In höheren Dimensionen werden Sie aber ganz schnell eine Liebe für die Ket-Notation entdecken.

\subsection{Skalarprodukte}
Auf $\mathbb{C}^n$ gibt es neben $\vb{v}\adj\vb{w}$ andere Skalarprodukte, sogar unendlich viele; die folgende Definition erfasst sie alle. Bei Rechnungen mit Spalten verwenden wir weiterhin $\vb{v}\adj\vb{w}$. Ein Skalarprodukt brauchen wir, weil die Quantenmechanik probabilistisch ist: Es ist unser Werkzeug für Wahrscheinlichkeiten.

\begin{definition}
	Ein \emph{Skalarprodukt} auf einem komplexen Vektorraum $V$ ordnet je zwei Vektoren $\ket{\phi},\ket{\psi}\in V$ eine komplexe Zahl $\braket{\phi}{\psi}$ zu. Es ist additiv in jedem Argument, $\braket{\phi_1 + \phi_2}{\psi} = \braket{\phi_1}{\psi} + \braket{\phi_2}{\psi}$ und ebenso im zweiten, und für alle $z\in\mathbb{C}$ gilt
	\begin{equation}
		\braket{\phi}{z\psi} = z\braket{\phi}{\psi}, \qquad \braket{z\phi}{\psi} = z^*\braket{\phi}{\psi}.
	\end{equation}
	Dazu kommt die \emph{Positivität}: $\braket{\psi}{\psi}$ ist reell und $\geq 0$, und null nur für $\ket{\psi} = 0$.
\end{definition}

Das Skalarprodukt $\vb{v}\adj\vb{w}$ erfüllt alle drei Bedingungen. Der Stern im ersten Argument ist wichtig: Wäre das Skalarprodukt dort linear, so wäre $\braket{i\psi}{i\psi} = -\braket{\psi}{\psi}$, und beide Seiten könnten nur für $\braket{\psi}{\psi} = 0$ nichtnegativ sein, auch wenn $\ket{\psi} \neq 0$ ist.

Die Positivität erlaubt uns, die Länge eines Vektors zu definieren, die \emph{Norm}
\begin{equation}
	\norm{\psi} = \sqrt{\braket{\psi}{\psi}} \geq 0.
\end{equation}
Mit dem Skalarprodukt $\vb{v}\adj\vb{w}$ ist $\norm{v}^2 = \vb{v}\adj\vb{v} = \sum_j\abs{v_j}^2$.

Aus den drei Bedingungen folgen die konjugierte Symmetrie und die \emph{Schwarz'sche Ungleichung}, aus dieser die Dreiecksungleichung für endlich viele Vektoren:
\begin{equation}
	\begin{gathered}
		\braket{\phi}{\psi} = \braket{\psi}{\phi}^*, \qquad
		\abs{\braket{\phi}{\psi}} \leq \norm{\phi}\,\norm{\psi}, \\
		\norm{\sum_{j=1}^{N}\ket{\psi_j}} \leq \sum_{j=1}^{N}\norm{\psi_j}.
	\end{gathered}
\end{equation}
Beide Ungleichungen sind Arbeitsgeräte. Die Schwarz'sche steckt hinter der Heisenbergschen Unschärferelation (Vorlesung~9); man beweist beide mit demselben Schritt: Das Skalarprodukt eines geschickt gewählten Vektors mit sich selbst ist nie negativ. Mit der Dreiecksungleichung schätzen wir Größen ab.

Zwei Vektoren heißen \emph{orthogonal}, $\ket{\phi} \perp \ket{\psi}$, wenn $\braket{\phi}{\psi} = 0$ ist. Achtung: Das heißt \emph{nicht}, dass einer der Vektoren null ist; das erzwingt die Positivität nur für $\ket{\phi} = \ket{\psi}$. Ein System von Vektoren $\ket{\phi_j}$ mit $j$ aus einer Indexmenge $X$ heißt \emph{orthonormal}, wenn
\begin{equation}
	\braket{\phi_j}{\phi_k} = \delta_{jk} \quad \text{für alle } j,k\in X.
\end{equation}
Dabei ist $\delta_{jk}$ das Kronecker-Delta, $1$ für $j = k$ und $0$ sonst; es ist nicht die Diracsche Delta-Funktion, die später vorkommt.

\subsection{Hilbertraum}
Jetzt haben wir genug beisammen, um Hilberträume zu definieren.

\begin{definition}
	Ein \emph{Hilbertraum} $\Hil$ ist ein komplexer Vektorraum mit Skalarprodukt, der bezüglich der Norm $\norm{\psi}^2 = \braket{\psi}{\psi}$ vollständig ist. \emph{Vollständig} heißt: Kommen die Glieder einer Folge $\ket{\psi_1},\ket{\psi_2},\dots$ in $\Hil$ einander beliebig nahe, $\norm{\psi_m - \psi_k} \to 0$ für $m,k \to \infty$, so hat die Folge einen Grenzwert $\ket{\psi}\in\Hil$, also $\norm{\psi_m - \psi} \to 0$.
\end{definition}

Die Vollständigkeit ist wichtig, aber subtil; sie hat mit Konvergenz in unendlichdimensionalen Räumen zu tun. In $\mathbb{C}^n$ mit $\vb{v}\adj\vb{w}$ ist sie automatisch erfüllt, und für Rechnungen dort ist ein Hilbertraum schlicht ein komplexer Vektorraum mit Skalarprodukt; für das Qubit ist das $\mathbb{C}^2$. In unendlich vielen Dimensionen gilt diese Abkürzung nicht: Dort braucht man die Vollständigkeit, etwa für die Entwicklung nach einer Basis mit unendlich vielen Vektoren.

\subsection{Basen}
Basen brauchen wir aus einem praktischen Grund: Die meisten Rechnungen der Quantenmechanik werden erst explizit, wenn eine Basis gewählt ist. Wir betrachten hier nur endliche Indexmengen $X$. Ein orthonormales System $\{\ket{\phi_j}\}_{j\in X}$ in $\Hil$ heißt \emph{Basis} (genauer \emph{Orthonormalbasis}) von $\Hil$, wenn es \emph{vollständig} ist:
\begin{equation}
	\norm{\psi}^2 = \sum_{j\in X}\abs{\braket{\phi_j}{\psi}}^2 \quad \text{für alle } \ket{\psi}\in\Hil.
\end{equation}
Verwirrenderweise sagt man auch hier vollständig, meint aber etwas anderes als bei der Vollständigkeit des Hilbertraums.

Die Paare $\ket{e_0},\ket{e_1}$ und $\ket{e_+},\ket{e_-}$ sind Basen von $\mathbb{C}^2$. Orthonormal ist auch das zweite Paar: $\braket{e_+}{e_-} = \tfrac12(1 - 1) = 0$ und $\norm{e_\pm}^2 = \tfrac12(1 + 1) = 1$. Für $\ket{\psi} = \alpha\ket{e_0} + \beta\ket{e_1}$ ist $\norm{\psi}^2 = \abs{\alpha}^2 + \abs{\beta}^2$; die rechte Seite ergibt dasselbe, mit $\braket{e_0}{\psi} = \alpha$ und $\braket{e_1}{\psi} = \beta$ ebenso wie mit $\braket{e_\pm}{\psi} = (\alpha \pm \beta)/\sqrt{2}$, denn $\abs{\alpha + \beta}^2 + \abs{\alpha - \beta}^2 = 2(\abs{\alpha}^2 + \abs{\beta}^2)$. In $\mathbb{C}^n$ mit $\vb{v}\adj\vb{w}$ gibt es immer eine Basis, die Spalten mit einer $1$ im $j$-ten Eintrag und sonst lauter Nullen. Im unendlichdimensionalen Fall ist alles schwieriger: Dass es eine Basis gibt, muss man beweisen, und der Beweis ist schwer. Basen mit unendlich vielen Vektoren behandelt der Webkurs.

Die wichtigste Eigenschaft einer Basis ist, dass man jeden Vektor $\ket{\Psi}\in\Hil$ nach ihr entwickeln kann, mit den Koeffizienten $c_j = \braket{\phi_j}{\Psi}$:
\begin{equation}
	\eqbox{\ket{\Psi} = \sum_{j\in X}\braket{\phi_j}{\Psi}\,\ket{\phi_j}.}
\end{equation}
Denn für $\ket{\xi} = \ket{\Psi} - \sum_j c_j\ket{\phi_j}$ ergibt die Orthonormalität $\braket{\xi}{\xi} = \norm{\Psi}^2 - \sum_j\abs{c_j}^2$, nach der Vollständigkeit null, und nach der Positivität ist $\ket{\xi} = 0$. Diese Gleichung ist eine Form der sogenannten \emph{Vollständigkeitsrelation}: Jeder Vektor ist die Summe seiner Komponenten längs der Basisvektoren, und die Koeffizienten sind Skalarprodukte. Sie sieht harmlos aus, wird aber ständig angewandt.

% fig_vl03_vollst.tex — Vollständigkeitsrelation: ein Vektor, zerlegt in
% zwei verschiedenen Orthonormalbasen desselben Raumes.
%
% Follows the conventions of fig_vl11_kasten.tex (the reference plate).
%
%   * geometry comes from the expansion derived in the lecture,
%         |Psi> = sum_j <phi_j|Psi> |phi_j>,
%     evaluated for the two orthonormal bases the lecture has already
%     introduced in C^2,
%         {|e_0>, |e_1>}   and   |e_±> = (|e_0> ± |e_1>)/sqrt(2),
%     on the concrete unit vector
%         |Psi> = (2|e_0> + |e_1>)/sqrt(5)   (real, so it fits in a plane).
%     Every coordinate below is that expansion, not a sketch:
%         c_0 = 2/sqrt5,  c_1 = 1/sqrt5   ->  |c|^2 = 4/5, 1/5
%         c_± = (2 ± 1)/sqrt10            ->  |c|^2 = 9/10, 1/10
%     and the two component rectangles therefore share |Psi> as their
%     diagonal exactly, because sum_j c_j |phi_j> = |Psi> in both bases.
%
%   * colour encodes *which basis*, and nothing else: cgblue for
%     {|e_0>,|e_1>}, birrenorange for {|e_+>,|e_->}.  The state |Psi>
%     itself is neutral navy — it belongs to no basis, which is the whole
%     point.  The coordinate rays are thin neutral grey: reference lines,
%     not data.
%
%   * within one basis the two members are told apart by tint (full
%     colour / 35%), not by a second hue, so the bars survive greyscale
%     printing; the numbers are printed as well, so nothing depends on
%     colour at all.  In (a) the same redundancy is carried by the dash
%     pattern — the first basis closes its rectangle with dashes, the
%     second with dots — and every arrow carries its own label, so the
%     two decompositions stay separable with the colour thrown away.
%     The neutral rays are held at birrengray!60: at full strength they
%     print at the same grey value as birrenorange and swallowed the
%     orange component vectors lying on top of them.
%
% Panel (a) the two decompositions of the same |Psi>.  Two different
%           rectangles, one diagonal.
% Panel (b) the weights |c_j|^2 those rectangles imply.  4/5 : 1/5 against
%           9/10 : 1/10 — the split is basis-dependent, the total is not.
%           This is the content of the completeness axiom
%           ||psi||^2 = sum_j |<phi_j|psi>|^2, and the reason the numbers
%           can later be read as probabilities.

\begin{center}
\begin{tikzpicture}[
    projvec/.style   ={-{Stealth[length=5pt,width=3.5pt]}, line width=1.3pt},
    statevec/.style  ={-{Stealth[length=6pt,width=4pt]}, line width=1.6pt,
                       spacecadet},
    basisray/.style  ={-{Stealth[length=4pt,width=3pt]}, line width=0.5pt,
                       birrengray!60},
    projdash/.style  ={line width=0.6pt, densely dashed},
    projdot/.style   ={line width=0.7pt, densely dotted},
]

% ─────────────────────────────────────────────────────────────
%  (a)  |Psi> resolved in two orthonormal bases
% ─────────────────────────────────────────────────────────────
\begin{axis}[
    name=geom,
    tpcplot,
    scale only axis,
    width=5.4cm, height=5.4cm,
    axis lines=none, grid=none,
    xmin=-0.38, xmax=1.38,
    ymin=-0.64, ymax=1.12,     % same span as x: 45 degrees looks like 45
    clip=false,
    title={(a) \tpcifenglish{One vector, two bases}%
                            {Ein Vektor, zwei Basen}},
]

  % ---- coordinate rays of the two bases (reference lines, neutral) ----
  \draw[basisray] (axis cs:-0.18,0)      -- (axis cs:1.24,0);
  \draw[basisray] (axis cs:0,-0.18)      -- (axis cs:0,1.00);
  \draw[basisray] (axis cs:-0.127,-0.127) -- (axis cs:0.721,0.721);
  \draw[basisray] (axis cs:-0.127,0.127)  -- (axis cs:0.424,-0.424);

  \node[lbl, anchor=west]       at (axis cs:1.27,0)      {$|e_0\rangle$};
  \node[lbl, anchor=south]      at (axis cs:0,1.03)      {$|e_1\rangle$};
  \node[lbl, anchor=south west] at (axis cs:0.74,0.74)   {$|e_+\rangle$};
  \node[lbl, anchor=north west] at (axis cs:0.44,-0.44)  {$|e_-\rangle$};

  % ---- the two component rectangles: same diagonal, different sides ----
  % basis {e_0,e_1}:  feet at (2/sqrt5, 0) and (0, 1/sqrt5)
  \draw[projdash, cgblue] (axis cs:0.894427,0)
                       -- (axis cs:0.894427,0.447214);
  \draw[projdash, cgblue] (axis cs:0,0.447214)
                       -- (axis cs:0.894427,0.447214);
  % basis {e_+,e_-}:  feet at 3/sqrt10 along e_+ and 1/sqrt10 along e_-
  \draw[projdot, birrenorange] (axis cs:0.670820,0.670820)
                            -- (axis cs:0.894427,0.447214);
  \draw[projdot, birrenorange] (axis cs:0.223607,-0.223607)
                            -- (axis cs:0.894427,0.447214);

  % ---- the component vectors c_j |phi_j> ----
  \draw[projvec, cgblue]       (axis cs:0,0) -- (axis cs:0.894427,0);
  \draw[projvec, cgblue]       (axis cs:0,0) -- (axis cs:0,0.447214);
  \draw[projvec, birrenorange] (axis cs:0,0) -- (axis cs:0.670820,0.670820);
  \draw[projvec, birrenorange] (axis cs:0,0) -- (axis cs:0.223607,-0.223607);

  % ---- the state itself ----
  \draw[statevec] (axis cs:0,0) -- (axis cs:0.894427,0.447214);
  \node[lbl, text=spacecadet, font=\small, anchor=south west]
    at (axis cs:0.91,0.47) {$|\Psi\rangle$};

  % ---- component labels, each beside its own arrow ----
  \node[lbl, text=cgblue, font=\scriptsize, anchor=north]
    at (axis cs:0.72,-0.045) {$c_0|e_0\rangle$};
  \node[lbl, text=cgblue, font=\scriptsize, anchor=east]
    at (axis cs:-0.04,0.22)  {$c_1|e_1\rangle$};
  \node[lbl, text=birrenorange, font=\scriptsize, anchor=south east]
    at (axis cs:0.58,0.685)  {$c_+|e_+\rangle$};
  \node[lbl, text=birrenorange, font=\scriptsize, anchor=north east]
    at (axis cs:0.155,-0.255) {$c_-|e_-\rangle$};

  \node[lbl, text=birrengray, font=\scriptsize, anchor=north east]
    at (axis cs:1.38,1.12) {$c_j \equiv \langle \phi_j|\Psi\rangle$};
\end{axis}

% ─────────────────────────────────────────────────────────────
%  (b)  The weights those two decompositions carry
% ─────────────────────────────────────────────────────────────
\begin{axis}[
    tpcplot,
    scale only axis,
    at={(geom.north east)}, anchor=north west, xshift=22mm,
    width=5.0cm, height=4.2cm,
    axis lines=left, grid=none,
    xmin=0, xmax=1.58, ymin=0.42, ymax=2.55,
    xtick={0,0.5,1}, xticklabels={$0$,$1/2$,$1$},
    ytick={1,2},
    yticklabels={{$\{|e_+\rangle,|e_-\rangle\}$},{$\{|e_0\rangle,|e_1\rangle\}$}},
    ytick style={draw=none},
    xlabel={$|c_j|^2$},
    clip=false,
    title={(b) \tpcifenglish{Weights}{Gewichte}},
]

  % the common total, drawn first so it does not sit on the bar borders
  \draw[levellinegrey] (axis cs:1,0.50) -- (axis cs:1,2.46);

  % basis {e_0,e_1}:  4/5 and 1/5
  \fill[cgblue]     (axis cs:0,1.74)  rectangle (axis cs:0.8,2.26);
  \fill[cgblue!35]  (axis cs:0.8,1.74) rectangle (axis cs:1.0,2.26);
  \draw[spacecadet, line width=0.5pt]
                    (axis cs:0,1.74)  rectangle (axis cs:1.0,2.26);

  % basis {e_+,e_-}:  9/10 and 1/10
  \fill[birrenorange]    (axis cs:0,0.74)  rectangle (axis cs:0.9,1.26);
  \fill[birrenorange!35] (axis cs:0.9,0.74) rectangle (axis cs:1.0,1.26);
  \draw[spacecadet, line width=0.5pt]
                         (axis cs:0,0.74) rectangle (axis cs:1.0,1.26);

  \node[font=\footnotesize, text=white] at (axis cs:0.40,2.00)
    {$|c_0|^2=\tfrac45$};
  \node[lbl, text=cgblue, font=\footnotesize, anchor=west]
    at (axis cs:1.04,2.00) {$|c_1|^2=\tfrac15$};

  \node[font=\footnotesize, text=white] at (axis cs:0.45,1.00)
    {$|c_+|^2=\tfrac{9}{10}$};
  \node[lbl, text=birrenorange, font=\footnotesize, anchor=west]
    at (axis cs:1.04,1.00) {$|c_-|^2=\tfrac{1}{10}$};

  \node[lbl, text=birrengray, font=\scriptsize]
    at (axis cs:0.50,1.50)
    {$\textstyle\sum_j |c_j|^2 = \|\Psi\|^2 = 1$};
\end{axis}

\end{tikzpicture}
\end{center}


Die Abbildung rechnet das für $\ket{\Psi} = (2\ket{e_0} + \ket{e_1})/\sqrt{5}$ in beiden Basen vor. Die Koeffizienten sind hier reell; nur deshalb passt (a) in eine Ebene. Jede Zerlegung schließt sich zu einem eigenen Rechteck: verschiedene Seiten, dieselbe Diagonale. In (b) stehen die Gewichte $\abs{c_j}^2$. Ihre Aufteilung hängt von der Basis ab, ihre Summe nicht: Sie ist $\norm{\Psi}^2$, für einen Vektor der Norm eins also eins.

Wie bildet man in der Praxis Basen? Mit dem Orthonormierungsverfahren von Gram und Schmidt. Aus einer beliebigen Liste $\ket{\Phi_1},\dots,\ket{\Phi_N}$ bestimmt man nacheinander
\begin{equation}
	\ket{\Psi_n} = \ket{\Phi_n} - \sum_{k=1}^{n-1}\frac{\braket{\Psi_k}{\Phi_n}}{\norm{\Psi_k}^2}\,\ket{\Psi_k},
\end{equation}
wobei Summanden mit $\ket{\Psi_k} = 0$ ausgelassen werden, weil sonst durch null geteilt würde. Dann ist $\braket{\Psi_k}{\Psi_n} = 0$ für $k<n$, und die Vektoren $\ket{\Psi_n}/\norm{\Psi_n}$ mit $\ket{\Psi_n} \neq 0$ bilden immer ein orthonormales System. Gibt es außer dem Nullvektor keinen Vektor, der auf allen $\ket{\Phi_n}$ senkrecht steht, so ist dieses System eine Basis von $\Hil$.

\begin{beispiel}
	In $\mathbb{C}^2$ starten wir mit zwei ausgedachten Vektoren,
	\begin{equation}
		\ket{\Phi_1} = \begin{pmatrix} 2 \\ 0 \end{pmatrix}, \qquad
		\ket{\Phi_2} = \begin{pmatrix} 1 \\ 1 \end{pmatrix}.
	\end{equation}
	Im ersten Schritt läuft die Summe von $1$ bis $0$; es wird gar nicht summiert, und $\ket{\Psi_1} = \ket{\Phi_1}$. Im zweiten Schritt ist mit $\braket{\Psi_1}{\Phi_2} = 2$ und $\norm{\Psi_1}^2 = 4$
	\begin{equation}
		\ket{\Psi_2} = \ket{\Phi_2} - \frac{\braket{\Psi_1}{\Phi_2}}{\norm{\Psi_1}^2}\,\ket{\Psi_1}
		= \begin{pmatrix} 1 \\ 1 \end{pmatrix} - \frac{1}{2}\begin{pmatrix} 2 \\ 0 \end{pmatrix}
		= \begin{pmatrix} 0 \\ 1 \end{pmatrix}.
	\end{equation}
	Geteilt durch die Normen $\norm{\Psi_1} = 2$ und $\norm{\Psi_2} = 1$ wird aus der ausgedachten Liste die Basis $\ket{e_0},\ket{e_1}$. Die Bedingung ist erfüllt, denn aus $\braket{\Phi_1}{\psi} = 2\psi_1 = 0$ und $\braket{\Phi_2}{\psi} = \psi_1 + \psi_2 = 0$ folgt $\ket{\psi} = 0$.
\end{beispiel}

%!TEX root = TheoretischePhysikC.tex
% Kern von content/c03-hilbertraum, c03.s8 bis c03.s12. Quellen: docs/plan/script-v0.3/sync-ch4.md

\section{Vorlesung: Operatoren}

\subsection{Operatoren und der Hut}
Das zweite Standbein der Theorie sind die Operatoren: Mit ihnen beschreibt die Quantenmechanik Zustände, Messungen und Dynamik.
\begin{definition}
	Sei $\Hil$ ein Hilbertraum. Ein \emph{Operator} $\op{A}$ auf $\Hil$ ist eine lineare Abbildung $\op{A}\colon\Hil \to \Hil$, das heißt
\begin{equation}
	\op{A}(\ket{\phi} + \ket{\psi}) = \op{A}\ket{\phi} + \op{A}\ket{\psi}, \qquad
	\op{A}(\alpha\ket{\phi}) = \alpha\,\op{A}\ket{\phi}
\end{equation}
für alle $\alpha\in\mathbb{C}$ und alle $\ket{\phi},\ket{\psi}\in\Hil$.
\end{definition}

In dieser Vorlesung ist $\Hil$ endlichdimensional: Es hat eine Basis mit endlicher Indexmenge, wie $\mathbb{C}^n$ mit $\vb{v}\adj\vb{w}$. Mit einer Basis ist ein Operator nichts anderes als eine $n\times n$-Matrix; die Regeln für Matrix mal Spalte sind genau die beiden Bedingungen. Wenn Sie das Wort ,,Operator`` hören, denken Sie an eine Tabelle. Ein Produkt wird von rechts nach links ausgeführt: $\op{A}\op{B}$ heißt erst $\op{B}$, dann $\op{A}$.

Trivial, aber sehr nützlich sind der \emph{Einsoperator} $\mathbb{I}$ mit $\mathbb{I}\ket{\psi} = \ket{\psi}$ und der \emph{Nulloperator} $\widehat{0}$ mit $\widehat{0}\,\ket{\psi} = 0$ für alle $\ket{\psi}\in\Hil$; im endlichdimensionalen Fall sind das die Einheitsmatrix und die Matrix mit lauter Nullen. Einer hingeschriebenen Null sieht man nicht an, was gemeint ist. Wir vereinbaren deshalb:

\begin{merksatz}[Notation]
	Ein \emph{Hut} bedeutet immer einen Operator, $\op{A}$; der \emph{Strich} der Ket-Notation einen Vektor, $\ket{\psi}$; ein Symbol ohne Verzierung eine Zahl. Ohne Hut bleiben der Einsoperator $\mathbb{I}$, die Ket-Bras weiter unten und zwei Sorten von Operatoren, die erst später vorkommen: die Zustände und die Pauli-Matrizen.
\end{merksatz}
Die Vorlesungen ab Vorlesung~5 sind noch in der alten Fassung und schreiben Operatoren oft ohne Hut.

\subsection{Der adjungierte Operator}
Die Kreuz-Operation übertragen wir jetzt auf Operatoren. Der \emph{adjungierte Operator} $\op{A}\adj$ ist durch die folgende Gleichung festgelegt:
\begin{equation}
	\braket{\op{A}\adj\phi}{\psi} = \braket{\phi}{\op{A}\psi} \quad \text{für alle } \ket{\phi},\ket{\psi}\in\Hil.
\end{equation}
Er ist also nicht durch eine Rechenvorschrift bestimmt, sondern durch sein Verhalten unter dem Skalarprodukt. In endlicher Dimension gibt es genau einen solchen Operator. In $\mathbb{C}^n$ mit $\vb{v}\adj\vb{w}$ heißt das Kreuz: erst transponieren, dann alle Einträge komplex konjugieren. Wie der Bra hängt $\op{A}\adj$ aber vom Skalarprodukt ab: Mit einem anderen Skalarprodukt gilt diese Regel für die Einträge in $\mathbb{C}^n$ im Allgemeinen nicht mehr. Die Rechenregeln, die wir ständig verwenden, sind
\begin{equation}
	\begin{gathered}
		(\op{A} + \op{B})\adj = \op{A}\adj + \op{B}\adj, \qquad (z\op{A})\adj = z^*\op{A}\adj, \\
		(\op{A}\op{B})\adj = \op{B}\adj\op{A}\adj, \qquad (\op{A}\adj)\adj = \op{A}.
	\end{gathered}
\end{equation}

\subsection{Klassen von Operatoren}
Mit $\mathbb{I}$ und $\op{A}\adj$ allein definieren wir sechs wichtige Klassen. Ein Operator $\op{A}$ heißt
\begin{enumerate}
	\item \emph{hermitesch}, wenn $\op{A}\adj = \op{A}$,
	\item \emph{positiv}, wenn $\braket{\psi}{\op{A}\psi} \geq 0$ für alle $\ket{\psi}\in\Hil$,
	\item \emph{Projektor}, wenn $\op{A}^2 = \op{A}\adj = \op{A}$,
	\item \emph{isometrisch}, wenn $\op{A}\adj\op{A} = \mathbb{I}$,
	\item \emph{unitär}, wenn $\op{A}\adj\op{A} = \op{A}\op{A}\adj = \mathbb{I}$, und
	\item \emph{normal}, wenn $\op{A}\adj\op{A} = \op{A}\op{A}\adj$.
\end{enumerate}
Jede Klasse hat in der Quantenmechanik ihre eigene Rolle: Die hermiteschen Operatoren beschreiben die Observablen, die positiven, mit einer Normierung, die Zustände, die Projektoren Ja-Nein-Messungen (Ereignis oder kein Ereignis), und die isometrischen und unitären die Dynamik.

Hermitesche Operatoren heißen oft auch selbstadjungiert; in endlicher Dimension ist das dasselbe. Für Operatoren, die nur auf einem Teil des Raumes definiert sind, muss man symmetrisch und selbstadjungiert unterscheiden; das behandelt der Webkurs, ebenso die Isometrien in unendlich vielen Dimensionen, die nicht unitär sind.

$\op{A}$ ist genau dann hermitesch, wenn $\braket{\psi}{\op{A}\psi}$ für jedes $\ket{\psi}$ reell ist; deshalb ist jeder positive Operator hermitesch. Operatoren der Form $\op{B}\adj\op{B}$ sind positiv, weil das Skalarprodukt eines Vektors mit sich selbst nie negativ ist:
\begin{equation}
	\braket{\psi}{\op{B}\adj\op{B}\psi} = \braket{\op{B}\psi}{\op{B}\psi} \geq 0.
\end{equation}
Wegen $\op{P} = \op{P}\adj\op{P}$ ist damit auch jeder Projektor $\op{P}$ positiv. Mit demselben Schritt beginnt der Beweis der Unschärferelation in Vorlesung~9. Ein isometrischer Operator $\op{V}$ lässt das Skalarprodukt invariant, $\braket{\op{V}\phi}{\op{V}\psi} = \braket{\phi}{\psi}$. In endlicher Dimension ist jede Isometrie unitär, und hermitesche, positive und unitäre Operatoren, Projektoren und Isometrien sind normal.

\subsection{Matrixelemente und Ket-Bras}
Sei $\{\ket{\phi_x}\}_{x\in X}$ eine Basis von $\Hil$ mit endlicher Indexmenge. Die \emph{Matrixelemente} von $\op{A}$ in dieser Basis sind die Zahlen $[\op{A}]_{xy}$ in
\begin{equation}
	[\op{A}]_{xy} = \braket{\phi_x}{\op{A}\phi_y}, \qquad
	\op{A}\ket{\Psi} = \sum_{x,y\in X}\ket{\phi_x}\,[\op{A}]_{xy}\,\braket{\phi_y}{\Psi};
\end{equation}
dabei sagt $y$, welcher Basisvektor hineingeht (die Spalte), und $x$, welche Komponente herauskommt (die Zeile). Die zweite Gleichung folgt, wenn man die Entwicklung nach der Basis zweimal einsetzt; über $x$ und $y$ wird unabhängig summiert. Darin kommt keine Matrix vor, nur Basisvektoren und Zahlen: Ein Operator ist bekannt, sobald man seine Matrixelemente kennt. Auch der adjungierte Operator lässt sich ablesen, $[\op{A}\adj]_{xy} = [\op{A}]_{yx}^*$: Für Matrixelemente in einer Orthonormalbasis gilt ,,transponieren und konjugieren`` bei jedem Skalarprodukt.

Kets und Bras kann man selbst als lineare Abbildungen lesen. Ab sofort ist das Ket $\ket{\Psi}$ die Abbildung, die eine Zahl frisst und ein Vielfaches des Vektors $\vb{\Psi}$ produziert, für den es bisher stand. Die Definition des Adjungierten gilt ebenso für lineare Abbildungen zwischen zwei Hilberträumen, hier $\mathbb{C}$ mit dem Skalarprodukt $a^*b$ und $\Hil$. Der adjungierte Operator zum Ket ist der Bra; er frisst einen Vektor und produziert eine Zahl:
\begin{equation}
	\ket{\Psi}\colon \mathbb{C} \to \Hil,\ z \mapsto z\,\vb{\Psi}, \qquad
	\bra{\Psi} = (\ket{\Psi})\adj\colon \Hil \to \mathbb{C},\ \vb{\Phi} \mapsto \braket{\Psi}{\Phi}.
\end{equation}
Bra mal Ket ist dann die Abbildung $z \mapsto \braket{\Phi}{\Psi}\,z$. Wir unterscheiden sie nicht von der Zahl $\braket{\Phi}{\Psi}$; für das Rechnen ändert sich also nichts. In der anderen Reihenfolge, Ket mal Bra, entsteht ein Operator auf $\Hil$, das \emph{Ket-Bra}
\begin{equation}
	\ketbra{\Psi}{\Phi}\,\ket{\chi} = \braket{\Phi}{\chi}\,\ket{\Psi} \quad \text{für alle } \ket{\chi}\in\Hil.
\end{equation}
Ist die Basis aus dem Kontext klar, so trägt der Buchstabe $\phi$ keine Information mehr, und wir schreiben $\ket{x}$ und $\bra{x}$. Damit werden die Entwicklungsformeln zu Operatorgleichungen:
\begin{equation}
	(\ketbra{\Psi}{\Phi})\adj = \ketbra{\Phi}{\Psi}, \qquad
	\mathbb{I} = \sum_{x\in X}\ketbra{x}{x}, \qquad
	\op{A} = \sum_{x,y\in X}\melem{x}{\op{A}}{y}\,\ketbra{x}{y}.
\end{equation}
Prüfen Sie sie, indem Sie beide Seiten auf einen beliebigen Vektor anwenden. Die Gleichung für $\mathbb{I}$ ist eine der hilfreichsten der Quantenmechanik. Physikalisch sagt sie nur, dass Eins gleich Eins ist, aber man darf diese komplizierte Eins an jeder Stelle einer Rechnung einschieben.

\subsection{Drei Operatoren auf \texorpdfstring{$\mathbb{C}^3$}{C3}}
\begin{beispiel}
	Sei $\Hil = \mathbb{C}^3$ mit $\vb{v}\adj\vb{w}$, und seien
	\begin{equation}
		\op{A} = \begin{pmatrix} 0 & 1 & 0 \\ 1 & 0 & 1 \\ 0 & 1 & 0 \end{pmatrix}, \qquad
		\op{B} = \frac13\begin{pmatrix} 1 & 1 & 1 \\ 1 & 1 & 1 \\ 1 & 1 & 1 \end{pmatrix}, \qquad
		\op{C} = \begin{pmatrix} 0 & 1 & 0 \\ 0 & 0 & i \\ -1 & 0 & 0 \end{pmatrix}.
	\end{equation}
	Transponieren und Konjugieren gibt bei $\op{A}$ und $\op{B}$ dieselbe Matrix; beide sind hermitesch. Jeder Eintrag von $\op{B}^2$ ist $\tfrac19(1 + 1 + 1) = \tfrac13$, also ist $\op{B}$ ein Projektor und damit positiv. $\op{A}$ ist dagegen \emph{nicht} positiv: Für $\ket{\psi}$ mit den Einträgen $\tfrac12$, $-\tfrac{1}{\sqrt{2}}$, $\tfrac12$ hat $\op{A}\ket{\psi}$ die Einträge $-\tfrac{1}{\sqrt{2}}$, $1$, $-\tfrac{1}{\sqrt{2}}$, also ist
	\begin{equation}
		\melem{\psi}{\op{A}}{\psi}
		= \begin{pmatrix} \frac12 & -\frac{1}{\sqrt{2}} & \frac12 \end{pmatrix}
		  \begin{pmatrix} - \frac{1}{\sqrt{2}} \\ 1 \\ -\frac{1}{\sqrt{2}} \end{pmatrix}
		= -\sqrt{2}.
	\end{equation}
	$\op{C}$ ist unitär, denn
	\begin{equation}
		\op{C}\adj\op{C}
		= \begin{pmatrix} 0 & 0 & -1 \\ 1 & 0 & 0 \\ 0 & -i & 0 \end{pmatrix}
		  \begin{pmatrix} 0 & 1 & 0 \\ 0 & 0 & i \\ -1 & 0 & 0 \end{pmatrix}
		= \begin{pmatrix} 1 & 0 & 0 \\ 0 & 1 & 0 \\ 0 & 0 & 1 \end{pmatrix},
	\end{equation}
	und ebenso ist $\op{C}\op{C}\adj = \mathbb{I}$. Mit der kanonischen Basis $\ket{0},\ket{1},\ket{2}$ (die Spalten mit der $1$ im ersten, zweiten, dritten Eintrag) ist
	\begin{equation}
		\op{C} = \ketbra{0}{1} + i\,\ketbra{1}{2} - \ketbra{2}{0}.
	\end{equation}
\end{beispiel}

Was die drei Rechnungen gemeinsam haben, sieht man an den Eigenwerten. Ein Eigenwert von $\op{A}$ ist, informell gesagt, eine Zahl $\lambda$, für die $\op{A}$ einen Vektor $\ket{v} \neq 0$ einfach mit $\lambda$ multipliziert, $\op{A}\ket{v} = \lambda\ket{v}$; die genaue Definition kommt in der nächsten Vorlesung. Für jeden Eigenwert, den ein Operator hat, überträgt sich die definierende Gleichung seiner Klasse auf die Zahl $\lambda$:
\begin{equation}
	\begin{aligned}
		&\op{A}\adj = \op{A} &&\Rightarrow\quad \lambda^* = \lambda, \\
		&\braket{\psi}{\op{A}\psi} \geq 0 \text{ für alle } \ket{\psi} &&\Rightarrow\quad \lambda \geq 0, \\
		&\op{A}^2 = \op{A}\adj = \op{A} &&\Rightarrow\quad \lambda^2 = \lambda, \\
		&\op{A}\adj\op{A} = \mathbb{I} &&\Rightarrow\quad \abs{\lambda} = 1.
	\end{aligned}
\end{equation}
Das sind Einschränkungen für die Eigenwerte, die es gibt; die Umkehrung behaupten wir nicht. Die Normalität allein schränkt die Eigenwerte nicht ein. Die Eigenwerte der drei Beispieloperatoren sind
\begin{equation}
	\op{A}\colon\ 0,\ \pm \sqrt{2}, \qquad
	\op{B}\colon\ 1,\ 0,\ 0, \qquad
	\op{C}\colon\ e^{-i\pi/6},\ i,\ e^{7i\pi/6} \quad (\lambda^3 = -i),
\end{equation}
und weitere gibt es nicht.

% fig_vl04_spektren.tex — Operatorklassen an ihren Spektren erkannt:
% wo in der komplexen Ebene die Eigenwerte einer Klasse liegen dürfen,
% und wo die Eigenwerte der drei Beispielmatrizen der Vorlesung landen.
%
% Follows the conventions of fig_vl11_kasten.tex.
%
%   * geometry comes from the defining relations of the lecture, not from
%     a sketch.  For a normal operator the defining equation passes
%     verbatim to each eigenvalue, and that is the whole content of
%     panel (a):
%         A^dag = A          =>  lambda^* = lambda   (the real axis)
%         A^dag A = I        =>  |lambda|  = 1       (the unit circle)
%         A^2 = A^dag = A    =>  lambda^2  = lambda  (the points 0 and 1)
%         <psi|A psi> >= 0   =>  lambda   >= 0       (the half-axis)
%     so the nesting projector < positive < Hermitian is not asserted,
%     it is visible: two points inside a half-line inside a line;
%   * panel (b) plots the spectra of exactly the three matrices of the
%     lecture's Beispiel, from their characteristic polynomials:
%         det(A - lambda) = -lambda(lambda^2 - 2)  =>  0, +-sqrt(2)
%         B = |u><u|, u = (1,1,1)/sqrt(3)          =>  1, 0, 0
%         det(C - lambda) = -(lambda^3 + i)        =>  lambda^3 = -i,
%             i.e. exp(-i pi/6), exp(i pi/2), exp(i 7pi/6) — three points
%             120 degrees apart on the unit circle.
%     The vector psi = (1/2, -1/sqrt2, 1/2) for which the lecture
%     computes <psi|A|psi> = -sqrt(2) is the eigenvector of that
%     eigenvalue, so the number computed in the text is the leftmost blue
%     point: that is why A fails to be positive, and the point is
%     labelled accordingly;
%   * colour encodes the operator class and nothing else — blue
%     Hermitian, green positive, red projector, orange unitary — and it
%     is the same code in both panels, which is what lets the eye carry
%     the map of (a) over to the data of (b);
%   * in (a) the loci ARE the statement, so they are drawn in colour; in
%     (b) they drop back to being reference lines and are drawn neutral
%     grey (`levellinegrey`), with the eigenvalues as the only coloured
%     objects.  Marker shape (disc / ring / diamond) repeats the class
%     distinction, so the plate survives greyscale printing — it also
%     lets the shared eigenvalue 0 of A and B be drawn where it belongs,
%     as a blue disc inside a red ring, instead of being nudged apart;
%   * the classification is a statement about NORMAL operators (the
%     spectral theorem); the grey note in (a) records that normality
%     itself constrains nothing, and the accompanying prose says it.

\begin{center}
\begin{tikzpicture}[x=1.45cm, y=1.45cm,
    % The axes of the complex plane are reference furniture, not data:
    % neutral grey, so the coloured loci of (a) carry the eye.
    cplane axis/.style={-{Stealth[length=4.5pt,width=3pt]},
                        birrengray, line width=0.5pt}]

% ─────────────────────────────────────────────────────────────
%  (a)  The loci: which eigenvalues each class allows
% ─────────────────────────────────────────────────────────────

  % Axes of the complex plane
  \draw[cplane axis] (-2.02,0) -- (2.08,0);
  \draw[cplane axis] (0,-1.50) -- (0,1.88);
  \node[lbl, text=birrengray, font=\scriptsize, anchor=south east]
    at (2.06,0.05) {$\real\lambda$};
  \node[lbl, text=birrengray, font=\scriptsize, anchor=west]
    at (0.09,1.78) {$\imag\lambda$};

  % positive:  lambda >= 0 — a highlighter stroke under the Hermitian line
  \draw[line width=6pt, birrengreen!55, line cap=butt] (0,0) -- (1.80,0);
  % Hermitian: lambda^* = lambda
  \draw[line width=1.4pt, cgblue] (-1.80,0) -- (1.80,0);
  % unitary:   |lambda| = 1
  \draw[line width=1.3pt, birrenorange] (0,0) circle (1);
  % projector: lambda^2 = lambda
  \fill[birrenred] (0,0) circle (3pt);
  \fill[birrenred] (1,0) circle (3pt);
  \node[lbl, text=birrenred, font=\scriptsize, anchor=north east] at (-0.04,-0.04) {$0$};
  \node[lbl, text=birrenred, font=\scriptsize, anchor=south west] at ( 1.05, 0.04) {$1$};

  % Class labels, each in its own colour, each outside the unit circle
  \node[lbl, text=birrenorange, font=\scriptsize, align=center, anchor=south]
    at (0,1.07) {\tpcifenglish{unitary}{unitär}\\$|\lambda|=1$};
  \node[lbl, text=birrenred, font=\scriptsize, align=right, anchor=south east]
    at (-1.04,0.12) {\tpcifenglish{projector}{Projektor}\\$\lambda\in\{0,1\}$};
  \node[lbl, text=cgblue, font=\scriptsize, align=right, anchor=north east]
    at (-1.04,-0.10) {\tpcifenglish{Hermitian}{hermitesch}\\$\lambda^{*}=\lambda$};
  \node[lbl, text=birrengreen, font=\scriptsize, align=left, anchor=north west]
    at (1.04,-0.10) {\tpcifenglish{positive}{positiv}\\$\lambda\ge 0$};
  \node[lbl, text=birrengray, font=\scriptsize, anchor=north west]
    at (0.12,-1.14) {normal: \tpcifenglish{all of}{ganz} $\C$};

  \node[lbl head, anchor=south] at (0,2.02)
    {(a) \tpcifenglish{where the eigenvalues lie}{wo die Eigenwerte liegen}};

% ─────────────────────────────────────────────────────────────
%  (b)  The same plane, with the spectra of the lecture's example
% ─────────────────────────────────────────────────────────────
\begin{scope}[xshift=6.95cm]

  % Reference loci, now neutral: the unit circle and the real axis
  \draw[levellinegrey] (0,0) circle (1);
  \draw[cplane axis] (-2.02,0) -- (2.08,0);
  \draw[cplane axis] (0,-1.50) -- (0,1.88);
  \node[lbl, text=birrengray, font=\scriptsize, anchor=south east]
    at (2.06,0.05) {$\real\lambda$};
  \node[lbl, text=birrengray, font=\scriptsize, anchor=west]
    at (0.09,1.78) {$\imag\lambda$};
  \node[lbl, text=birrengray, font=\scriptsize, anchor=north west]
    at (0.10,-1.06) {$|\lambda|=1$};

  % C:  lambda^3 = -i  — three points 120 degrees apart on the circle
  \foreach \a in {-30,90,210}{
    \fill[birrenorange] (\a:1) ++(0,3.4pt) -- ++(3.4pt,-3.4pt)
                        -- ++(-3.4pt,-3.4pt) -- ++(-3.4pt,3.4pt) -- cycle;
  }
  % A:  lambda = 0, +-sqrt(2)
  \fill[cgblue] ({-sqrt(2)},0) circle (2.5pt);
  \fill[cgblue] (0,0)          circle (2.5pt);
  \fill[cgblue] ({ sqrt(2)},0) circle (2.5pt);
  % B:  lambda = 1, 0, 0  — the ring at 0 encircles A's disc
  \draw[birrenred, line width=1.1pt] (0,0) circle (4.3pt);
  \draw[birrenred, line width=1.1pt] (1,0) circle (4.3pt);

  % Eigenvalue readings
  \node[lbl, font=\scriptsize, anchor=north]      at ({-sqrt(2)},-0.09) {$-\sqrt{2}$};
  \node[lbl, font=\scriptsize, anchor=north]      at ({ sqrt(2)},-0.09) {$\sqrt{2}$};
  \node[lbl, font=\scriptsize, anchor=north]      at (1,-0.19)          {$1$};
  \node[lbl, font=\scriptsize, anchor=north east] at (-0.05,-0.12)      {$0$};

  % Operator labels, in the colour of the class each one belongs to
  \node[lbl, text=cgblue, font=\scriptsize, align=right, anchor=south east]
    at (-1.06,0.44) {$\widehat A$\\\tpcifenglish{Hermitian}{hermitesch}};
  \draw[cgblue, line width=0.4pt] (-1.45,0.42) -- (-1.42,0.08);
  \node[lbl, text=birrenorange, font=\scriptsize, align=center, anchor=south]
    at (0,1.10) {$\widehat C$\\\tpcifenglish{unitary}{unitär}};
  \node[lbl, text=birrenred, font=\scriptsize, align=left, anchor=south west]
    at (1.06,0.34) {$\widehat B$\\\tpcifenglish{projector}{Projektor}};
  \draw[birrenred, line width=0.4pt] (1.14,0.32) -- (1.05,0.09);

  % Why the lecture bothers to compute <psi|A|psi> = -sqrt(2)
  \node[lbl, text=birrengray, font=\scriptsize, anchor=north]
    at ({-sqrt(2)},-0.70) {\tpcifenglish{not positive}{nicht positiv}};

  \node[lbl head, anchor=south] at (0,2.02)
    {(b) \tpcifenglish{the example's $\widehat A,\widehat B,\widehat C$}%
                      {$\widehat A$, $\widehat B$, $\widehat C$ aus dem Beispiel}};

\end{scope}
\end{tikzpicture}
\end{center}


Teil (a) der Abbildung zeigt, wo die Eigenwerte einer Klasse liegen dürfen, Teil (b) die Eigenwerte von $\op{A}$, $\op{B}$ und $\op{C}$. Dort liegt $-\sqrt{2}$ links der Null; genau daran scheitert die Positivität von $\op{A}$, denn $\ket{\psi}$ von oben ist ein Eigenvektor der Norm eins zu diesem Eigenwert.

%!TEX root = TheoretischePhysikC.tex

\section{Vorlesung: Position-Basis, Spur, Diagonalisierung}
\subsection{Position-Basis}
Dirac und mit ihm viele Physiker gehen mit der Ket-Notation noch einen Schritt weiter. Wir betrachten zuerst den Fall, dass $\mathcal{H} \equiv \ell^2(\mathbb{N})$. Das heisst, die Vektoren $|\Psi\rangle \in \mathcal{H}$ sind lineare Kombinationen 
\begin{equation}
	|\Psi\rangle = \sum_{n=1}^{\infty} \Psi(n)|\phi_n\rangle \equiv \sum_{n=1}^{\infty} \Psi(n)|n\rangle,
\end{equation}
wobei $\Psi(n)\in \mathbb{C}$ und
\begin{equation}
	\sum_{n=1}^{\infty} |\Psi(n)|^2 < \infty.
\end{equation}
Die Ket-Notation wird dann erweitert, indem man den Wert $\Psi(n)$ einer ,,Wellenfunktion`` $\Psi\in \ell^2(\mathbb{N})$ an der Stelle $n\in\mathbb{N}$ als
\begin{equation}
	\langle n|\Psi\rangle \equiv \Psi(n)
\end{equation}
definiert. Achten Sie hier auf eine Feinheit: $\ell^2(\mathbb{N})$ verlangt nur Quadratsummierbarkeit; die \emph{normierten} Wellenfunktionen mit $\sum_{n} |\Psi(n)|^2 = 1$ bilden darin nur eine Teilmenge.

Nun folgt die wichtigste und zugleich unangenehmste Definition dieses Kapitels. Alles Bisherige war mathematisch völlig in Ordnung; ab hier verlassen wir diesen sicheren Boden bewusst.

Sei jetzt $\mathcal{H}\equiv L^2(\mathbb{R})$. Dieser Hilbertraum besteht aus quadratintegrierbaren Funktionen $f:\mathbb{R}\rightarrow\mathbb{C}$, welche die Bedingung
\begin{equation}
	\int_{-\infty}^{\infty} |f(x)|^2\,dx < \infty
\end{equation}
erfüllen. Die Ket-Notation bedeutet in diesem Fall folgendes: $|\Psi\rangle$ steht für die quadratintegrierbare Funktion $\Psi$. Man führt jetzt das ,,uneigentliche`` Bra $\langle x|$ ein, indem man 
\begin{equation}
	\langle x|\Psi\rangle \equiv \Psi(x)
\end{equation}
definiert. Warnung: Das Bra $\langle x|$ ist nun beim besten Willen kein Vektor im Hilbertraum mehr. Trotzdem arbeiten wir mit solchen ,,Bra-Vektoren`` als ob sie tatsächlich Vektoren darstellen. Bei den Bra-Vektoren handelt es sich um eine Gratwanderung zwischen Stringenz und Unsinn. Um Kohärente Rechenregeln zu erstellen, führen wir die  sogenannten ,,Position-Basis`` ein: die Ket-Vektoren $|x\rangle$, $x\in \mathbb{R}$ werden mit Hilfe den Bra-Vektoren $\langle x|$ durch
\begin{equation}
	\langle x|y\rangle \equiv \delta(x-y)
\end{equation}
gegeben, wobei $\delta(z)$ die Dirac-Delta-Funktion darstellt.\footnote{An die Stelle des Kronecker-Deltas $\langle n|m\rangle = \delta_{nm}$ tritt hier eine Distribution, die im Sinne einer Funktion gar nicht existiert. Daher der typische Fehler: für $y=x$ steht rechts nicht Eins, sondern $\delta(0)=\infty$. Solche Ausdrücke muss man vermeiden.} Wir definieren auch $(|x\rangle)^\dag\equiv \langle x|$ und (kontinuierliche) lineare Kombinationen von $|x\rangle$-Vektoren sind explizit erlaubt. Daher kann man eine quadratintegrierbare Funktion $\Psi \in L^2(\mathbb{R})$ als eine (kontinuierliche) lineare Kombination von den $|x\rangle$ darstellen\footnote{Diese Schreibweise erinnert uns an die Darstellung, die in dem endlichen dimensionalen Fall gilt:
\begin{equation}
	|\Phi\rangle = \sum_{n\in X} \Phi(n)|n\rangle.
\end{equation}}:
\begin{equation}
	|\Psi\rangle \equiv \int_{-\infty}^{\infty} \Psi(x)|x\rangle \, dx.
\end{equation}
\begin{intuition}
	Was ist $|x\rangle$ physikalisch? Der Zustand eines Teilchens, das sich mit Sicherheit an der Stelle $x$ befindet. Bemerkenswert ist daran, dass eine Position in der Natur nie ganz scharf festliegt: es gibt immer eine Unschärfe.
\end{intuition}

Diese Notation erlaubt uns mehrere Rechnungen analog mit dem endlichen dimensionalen Fall durchzuführen. Zum Beispiel kann man die Norm von $|\Psi\rangle$ wie folgt erhalten:
\begin{equation}
	\begin{split}
		\|\Psi\|^2 &= \langle \Psi|\Psi\rangle \\
		&= \left(\int_{-\infty}^{\infty} \Psi(x)|x\rangle \, dx\right)^\dag \left(\int_{-\infty}^{\infty} \Psi(y)|y\rangle \, dy\right) \\
		&= \left(\int_{-\infty}^{\infty} \Psi^*(x) \langle x| \, dx\right)\left(\int_{-\infty}^{\infty} \Psi(y)|y\rangle \, dy\right) \\
		&= \int_{-\infty}^{\infty} \Psi^*(x)\Psi(y) \langle x|y\rangle\,dxdy \\
		&= \int_{-\infty}^{\infty} \Psi^*(x)\Psi(y) \delta(x-y)\,dxdy\\
		&= \int_{-\infty}^{\infty} \Psi^*(x)\Psi(x) \,dx \\
		&= \int_{-\infty}^{\infty} |\Psi(x)|^2 \,dx.
	\end{split}
\end{equation} 

\subsection{Spur}
Die Spur eines Operators spielt in der Quantenmechanik eine wichtige Rolle. Die Spur eines Operators $A$ auf einem endlich dimensionalen Hilbertraum $\mathcal{H}$ definieren wir als 
\begin{equation}
	\tr(A) \equiv \sum_{x\in X} \langle x|A|x\rangle,
\end{equation}
wobei $x$ über die Indizes einer Orthonormalbasis $\{|\phi_x\rangle\}_{x\in X}$ läuft. Gemeint ist eine gewöhnliche Orthonormalbasis, nicht die Positions-Basis des vorigen Abschnitts. Wenn $A : \mathcal{H}_1 \rightarrow \mathcal{H}_2$ und $B : \mathcal{H}_2 \rightarrow \mathcal{H}_1$ rechnet man nach, dass
\begin{equation}
	\tr_2(AB) = \tr_1(BA)
\end{equation}
wobei der Index an der Spur nur daran erinnert, auf welchem der beiden Hilberträume die Spur zu nehmen ist. Diese Zyklizität benutzen wir ständig. In der Positions-Basis schreibt man die Spur auch als
\begin{equation}
	\tr(A) = \int_{-\infty}^{\infty} \langle x|A|x\rangle \, dx;
\end{equation}
in der Literatur begegnet sie einem oft.

\subsection{Diagonalisierung}
Für viele Rechnungen in der Quantenmechanik müssen Operatoren \emph{diagonalisiert} werden. Wie allgemein in der linearen Algebra bedeutet dies die Transformation auf eine besonders einfache Gestalt. Der Grund für den Aufwand: allgemeine Matrizen kommutieren nicht und sind mühsam zu handhaben, diagonale Operatoren dagegen sind so gut wie Zahlen. Alle Operatoren gleichzeitig diagonal zu bekommen gelingt allerdings in der Regel nicht. Der einfachste Fall ist  wieder endliche Matrizen, also Operatoren in einem Raum $\mathcal{H}\equiv \mathbb{C}^N$ für endliches $N$. Ein \emph{diagonaler Operator} $A$ hat per Definition Matrixelemente der Form
\begin{equation}
	[A]_{nm} = \langle\phi_n|A|\phi_m\rangle = \delta_{nm} a_m,
\end{equation}
wobei $a_1, \ldots, a_N$ die diagonalen Matrixelemente sind. Der Operator $A$ wirkt auf Vektoren $|\Psi\rangle$ wie folgt:
\begin{equation}
	A|\Psi\rangle = \sum_{n,m} |n\rangle\langle n|A|m\rangle\langle m|\Psi\rangle = \sum_{m} a_{m}\Psi(m) |m\rangle,
\end{equation}
wobei $\Psi(m) = \langle m|\Psi\rangle$.  

\emph{Diagonalisierung} bedeutet nun das Auffinden einer invertierbaren Matrix $S$ so, dass $SAS^{-1}$ ein diagonaler Operator ist. Wir werden uns (fast) ausschließlich für den Fall interessieren, in dem diese Transformation sogar durch einen unitären Operator $S$ geschieht (dann ist $S^{−1} = S^\dag$). Allgemein bedeutet für uns also
\emph{diagonalisieren} von $A$ das Auffinden eines unitären Operators $U : \mathcal{H} \rightarrow L^2(\mathbb{R})$, so dass 
\begin{equation}
	\langle x|(UAU^\dag )|\Psi\rangle = a(x)\Psi(x).
\end{equation}
Jeder hermitesche und jeder unitäre Operator ist diagonalisierbar. Für beliebige Operatoren gilt das nicht; es gibt durchaus Matrizen, die sich nicht diagonalisieren lassen. In der Quantenmechanik haben wir insofern Glück.

\subsection{Eigenwert, Eigenvektor}
Ist $V$ ein komplexer Vektorraum und $A:V\rightarrow A$ eine lineare Abbildung von $V$ in sich selbst, so bezeichnet man als \emph{Eigenvektor} einen Vektor $\mathbf{v}\not=\mathbf{0}$, der durch $A$ auf ein Vielfaches 
$\alpha\mathbf{v}$ von sich selbst mit $\alpha\in\mathbb{C}$ abgebildet wird:
\begin{equation}
	A\mathbf{v} = \alpha \mathbf{v}.
\end{equation}
Den Faktor $\alpha$ nennt man dann den zugehörigen Eigenwert. 

Eine Diagonalisierung von einem Operator (Matrix) $A$ liefert uns die Eigenwerte und Eigenvektoren. Sei zum Beispiel $\mathcal{H} \equiv \mathbb{C}^N$. Angenommen wir haben eine unitäre Matrix $U$ gefunden, so dass $UAU^\dag$ diagonal ist. Man bildet die Basis 
\begin{equation}
	\phi_1 \equiv U^\dag \begin{pmatrix}
	1\\0\\ \vdots \\ 0
	\end{pmatrix}, \quad  \phi_2 \equiv U^\dag \begin{pmatrix}
	0\\1\\ \vdots \\ 0
	\end{pmatrix}, \quad \ldots, \quad \phi_N \equiv U^\dag \begin{pmatrix}
	0\\0\\ \vdots \\ 1
	\end{pmatrix}.
\end{equation}
(Übung: Bildet diese Reihenfolge eine Basis?) Ein Hinweis: $\langle\phi_j|\phi_k\rangle = e_j^\dag UU^\dag e_k = \delta_{jk}$ mit dem kanonischen Basisvektor $e_j$, und genau hier geht die Unitarität ein. Der Operator $A$ wirkt auf diesen Vektoren wie folgt:
\begin{equation}
	A\phi_1 \equiv AU^\dag \begin{pmatrix}
	1\\0\\ \vdots \\ 0
	\end{pmatrix} = U^\dag U AU^\dag \begin{pmatrix}
	1\\0\\ \vdots \\ 0
	\end{pmatrix} = U^\dag \begin{pmatrix}
		a_1 & 0 & 0 & \cdots & 0 \\
		0 & a_2 & 0 & \cdots & 0 \\
		\vdots & \vdots &  & \ddots & \vdots \\
		0 & 0 & 0 & \cdots & a_N
	\end{pmatrix} \begin{pmatrix}
	1\\0\\ \vdots \\ 0
	\end{pmatrix} = a_1 \phi_1.
\end{equation}
Genau analog wirkt $A$ auf $\phi_j$ als
\begin{equation}
	A\phi_j = a_j\phi_j.
\end{equation}
Nach dieser Art und Weise indentifiziert man $\phi_j$ als die Eigenvektoren von $A$ mit dazugehörigen Eigenwerten $a_j$. 

Ebenso wichtig ist die Rückrichtung: aus einer Liste von Eigenvektoren gewinnt man die unitäre Matrix. Sei $A$ hermitesch auf $\mathcal{H}\equiv\mathbb{C}^N$ mit
\begin{equation}
	A|\phi_j\rangle = a_j |\phi_j\rangle, \qquad j = 1,\ldots,N.
\end{equation}
Zwei Vorbereitungen. Erstens die Normierung: ein aus einem Gleichungssystem gewonnener Eigenvektor hat fast nie die Norm Eins; man teilt dann durch sie und erhält wieder einen Eigenvektor zum selben Eigenwert. Zweitens die Orthogonalität, die keineswegs automatisch gilt. Sind alle Eigenwerte verschieden, so berechnen wir $\langle\phi_j|A|\phi_k\rangle$ auf zwei Arten: nach rechts wirkend ist
\begin{equation}
	\langle\phi_j|A|\phi_k\rangle = a_k \langle\phi_j|\phi_k\rangle.
\end{equation}
Der Trick ist, $A$ nach links wirken zu lassen. Dabei gehen $A^\dag = A$ und die Reellwertigkeit der Eigenwerte ein --- letztere, weil $a_j = \langle\phi_j|A|\phi_j\rangle$ für hermitesches $A$ nach Vorlesung 4 reell ist. Damit ist
\begin{equation}
	\langle\phi_j|A = \left(A^\dag|\phi_j\rangle\right)^\dag = \left(a_j|\phi_j\rangle\right)^\dag = a_j^*\langle\phi_j| = a_j \langle\phi_j|,
\end{equation}
also $\langle\phi_j|A|\phi_k\rangle = a_j\langle\phi_j|\phi_k\rangle$. Der Vergleich ergibt
\begin{equation}
	(a_j - a_k)\langle\phi_j|\phi_k\rangle = 0,
\end{equation}
und da die Eigenwerte verschieden sind, bleibt nur $\langle\phi_j|\phi_k\rangle = 0$ für $j\neq k$: die Eigenvektoren bilden eine Orthonormalbasis. Damit setzen wir
\begin{equation}
	U \equiv \sum_{j} |j\rangle\langle\phi_j|,
\end{equation}
mit dem kanonischen Basisvektor $|j\rangle$, der eine Eins in der $j$-ten Zeile trägt. Orthonormalität und Vollständigkeitsrelation liefern
\begin{equation}
	UU^\dag = \sum_{j,k} |j\rangle\langle\phi_j|\phi_k\rangle\langle k| = \sum_{j,k} \delta_{jk}|j\rangle\langle k| = \sum_{j} |j\rangle\langle j| = \mathbb{I},
\end{equation}
und dieselbe Rechnung mit $A$ in der Mitte ergibt
\begin{equation}
	UAU^\dag = \sum_{j,k} |j\rangle\langle\phi_j|A|\phi_k\rangle\langle k| = \sum_{j} a_j |j\rangle\langle j|,
\end{equation}
also einen diagonalen Operator mit den Eigenwerten auf der Diagonalen. (Übung: Zeigen Sie ebenso $U^\dag U = \mathbb{I}$.)

\begin{beispiel}
	Sei $\mathcal{H}\equiv\mathbb{C}^2$, in der Quanteninformationstheorie ein \emph{Qubit}, und
	\begin{equation}
		A \equiv \begin{pmatrix} 0 & 1 \\ 1 & 0 \end{pmatrix}.
	\end{equation}
	Das charakteristische Polynom $\det(A - \lambda\mathbb{I}) = \lambda^2 - 1$ liefert $\lambda = \pm 1$ mit den normierten Eigenvektoren
	\begin{equation}
		|\phi_1\rangle = \frac{1}{\sqrt{2}}\begin{pmatrix} 1 \\ 1 \end{pmatrix}, \qquad |\phi_2\rangle = \frac{1}{\sqrt{2}}\begin{pmatrix} 1 \\ -1 \end{pmatrix}.
	\end{equation}
	Ein Ket ist ein Spalten-, ein Bra ein Zeilenvektor, ihr Produkt also eine Matrix:
	\begin{equation}
		U = |1\rangle\langle\phi_1| + |2\rangle\langle\phi_2| = \frac{1}{\sqrt{2}}\begin{pmatrix} 1 & 1 \\ 0 & 0 \end{pmatrix} + \frac{1}{\sqrt{2}}\begin{pmatrix} 0 & 0 \\ 1 & -1 \end{pmatrix} = \frac{1}{\sqrt{2}}\begin{pmatrix} 1 & 1 \\ 1 & -1 \end{pmatrix},
	\end{equation}
	die \emph{Hadamard-Matrix}; man rechnet $UAU^\dag = \diag(1,-1)$ nach. Vertauscht man die beiden Eigenvektoren, so liefert dasselbe Rezept eine andere unitäre Matrix $U'$ mit $U'AU'^\dag = \diag(-1,1)$: die Diagonalisierung ist nicht eindeutig.
\end{beispiel}

Dieses Verfahren funktioniert auch im unendlichen dimensionalen Fall, ist jedoch aufwendiger, allerdings werden wir immer nur von Fall zu Fall unendlich dimensionale Hilberträume behandeln.

\subsection{Funktionalkalkül}
Wenn $H$ ein hermitescher Operator ist, dann ist die diagonalisierung $SHS^{-1} = SHS^{\dag}$ von $H$ auch hermitesch, also die Multiplikation mit einer reellen Funktion $a_x$ (bezüglich der Eigenvektor-Basis $\{|\phi_x\rangle\}_{x\in X}$):
\begin{equation}
	SHS^\dag |\phi_x\rangle = a_x |\phi_x\rangle,
\end{equation}
wobei $|\phi_x\rangle$ die Eigenvektoren von $A$ mit dazugehörigen Eigenwerten $a_x$ bilden. 
Dass die $a_x$ reell sind, sieht man sofort: für unitäres $S$ ist $(SHS^\dag)^\dag = SH^\dag S^\dag = SHS^\dag$, die Diagonalform ist also wieder hermitesch, und das heißt für eine Diagonalmatrix $a_x^* = a_x$.

Wir definieren von diesem Operator die Funktion 
\begin{equation}
	U_t \equiv e^{\frac{−it}{\hbar}H}
\end{equation}
die durch
\begin{equation}
	U_t |\phi_x\rangle \equiv e^{\frac{−it}{\hbar}a_x} |\phi_x\rangle, \quad \forall x\in X,
\end{equation}
gegeben wird. (Wenn $A$ auf einen Eigenvektor angewandt wird, wird der Operator $H$ mit der reellen Zahl $a_x$ ersetzt.) Der Operator $U_t$ ist auch (bezüglich die Eigenvektor-Basis) ein Multiplikationsoperator, allerdings wird der Skalar $a_x$ durch eine (x-abhängige) \emph{Phase} $e^{\frac{−it}{\hbar}a_x}$ ersetzt. Weil die $a_x$ reell sind, hat jede Phase den Betrag Eins; genau deshalb ist $U_t$ unitär und lässt die Norm eines Zustands unverändert. 

Im Fall, dass $\mathcal{H}$ endlich-dimensional ist, ist die Operator-Funktion $e^{\frac{−it}{\hbar}H}$ das gleiche wie eine Potenzreihe 
\begin{equation}
	e^{\frac{−it}{\hbar}H} \equiv \sum_{n=0}^\infty \frac{1}{n!} \left(\frac{−it}{\hbar}\right)^n H^n.
\end{equation}
Übrigens ist die Definition der Exponentialfunktion durch Diagonalisierung auch noch bei unbeschränkten selbst-adjungierten Operatoren $H$ vernünftig, während die Potenzen eines unendlich-dimensionalen Operators und erst recht die Potenzreihe nur unter sehr einschränkenden Bedingungen einen Sinn ergeben.

\begin{merksatz}
	Der Operator $U_t$ beschreibt die Zeitentwicklung eines Quantensystems. Damit steht die Lösung der Schrödingergleichung schon hier, bevor die Gleichung selbst aufgetaucht ist.
\end{merksatz}

%!TEX root = TheoretischePhysikC.tex
\section{Vorlesung: Postulate der Quantenmechanik}
Die Zahlen, die die Quantenmechanik als Ergebnis liefert, sind Wahrscheinlichkeiten oder, etwas allgemeiner Erwartungswerte. Da in den bisherigen Theorievorlesungen von solchen Größen wohl wenig die Rede war, müssen wir kurz darauf eingehen, wie solche Ergebnisse zu interpretieren, d.h., mit dem Experiment zu vergleichen sind. Die allgemeine Vorschrift für statistische Experimente ist, dass ein solches Experiment genügend häufig zu wiederholen
ist.

Dass die Theorie überhaupt nur Wahrscheinlichkeiten liefert, ist der eigentliche Bruch mit der klassischen Physik. Klassisch ist eine Wahrscheinlichkeit stets Ausdruck fehlender Information; im Prinzip kann dort alles beliebig präzise sein. Quantenmechanisch ist sie fundamental, und es gibt Messungen, die sich nicht gemeinsam mit voller Präzision durchführen lassen.

Für ein Wahrscheinlichkeitsexperiment muss dabei im Voraus festgelegt sein, welche Ergebnisse als ,,Treffer`` zu werten sind. Die von der Theorie gelieferte \emph{Wahrscheinlichkeit} ist dann zu vergleichen mit der Anzahl der Treffer, dividiert durch die Gesamtzahl der Wiederholungen. Entsprechend werden Erwartungswerte von Experimenten gemessen, bei denen jedes einzelne Experiment einen Zahlenwert liefert: Der Erwartungswert ist zu vergleichen mit der Summe der Messwerte, dividiert durch die Gesamtzahl der Wiederholungen. Was dabei als ,,genügend häufige`` Wiederholung für eine angestrebte Genauigkeit gilt, wird in dem als Statistik bekannten Zweig der Wahrscheinlichkeitstheorie behandelt. Die Berechnung der statistischen Konfidenzintervalle entbindet aber den Experimentalphysiker nicht von der Pflicht, sich zu vergewissern, dass die Theorie überhaupt anwendbar ist, dass also die so gewonnenen Häufigkeiten tatsächlich reproduzierbar und frei von systematischen Fluktuationen
sind, also zum Beispiel nicht abhängig von irgendwelchen Aktivitäten im Nachbarlabor schwanken. Das ist bei Wahrscheinlichkeitsmessungen nicht anders als bei allen anderen Messungen auch.

Zu erläutern wäre noch, was denn nun eigentlich als Wiederholung ,,des gleichen`` Experiments gilt. Hier gelten die Kriterien der Experimentalphysik: zwei innerhalb gegebener Toleranzen nach den gleichen Plänen aufgebaute Experimente gelten als gleich. Es muss also nicht etwa dieselbe Apparatur verwendet werden. Entscheidend ist also das experimentelle \emph{Verfahren}. Dies muss natürlich genügend präzise spezifiziert sein, wenn Reproduzierbarkeit von Wahrscheinichkeiten gewährleistet werden soll.

In den Experimenten der Quantenmechanik sind grundsätzlich zwei Teile zu unterscheiden: die \emph{Präparation} und die \emph{Messung}. Zum Beispiel für die Polarisationsexperimente bestand die Präparation in der Herstellung des Strahls, und die Messung in der Anwendung der Filter, zusammen mit dem Photonenzähler. Das Ziel der Theorie ist nun die Berechnung der Wahrscheinlichkeiten für praktisch beliebige Kombinationen von Präparation und Messung. Die Theorie muss also mindestens enthalten: (1) Größen zur Beschreibung der Präparierverfahren, (2) Größen zur Beschreibung der Messverfahren und (3) eine Formel zur Berechnung der Wahrscheinlichkeiten für jede Kombination.

\subsection{Die Postulate}

\begin{greybox}[Zum Status der Postulate]
Man hört oft, wir verstünden die Quantenmechanik nicht. Das trifft es nicht: Wir beherrschen den Formalismus, und wir haben seit über fünfzig Jahren Postulate, die uns die besten Modelle des Universums liefern. Was wir nicht können, ist, sie weiter zu \emph{begründen} --- sie sind uns von der Natur gegeben.
\end{greybox}

\begin{center}
\parbox[c][5cm][c]{14cm}{\begin{tcolorbox}[colback=banana,arc=0cm,frame hidden,enhanced]{\textbf{Postulat 1}. Jeder Sorte quantenmechanischer Systeme wird ein Hilbertraum $\mathcal{H}$ zugeordnet.}\end{tcolorbox}}
\end{center}
Dies ist der Raum der sogenannten \emph{Wellenfunktionen} $\Psi$. Für ein \emph{Qubit}, bzw.\ \emph{Qudit}, ist $\mathcal{H}\equiv \mathbb{C}^2$, bzw.\ $\mathcal{H}\equiv \mathbb{C}^d$. Für ein freies Teilchen in drei Dimensionen ist $\mathcal{H} \equiv L^2(\mathbb{R}^{3})$, der Raum der quadratintegriebare Funktionen $\Psi:\mathbb{R}^3\rightarrow \mathbb{C}$, mit
\begin{equation}
	\iiint |\Psi(x,y,z)|^2\, dxdydz  <\infty.
\end{equation} 
Das Skalarprodukt zweier Vektoren $|\Phi\rangle, |\Psi\rangle \in \mathcal{H}$ lautet
\begin{equation}
	\langle \Phi|\Psi\rangle \equiv \iiint \Phi^*(x,y,z)\Psi(x,y,z)\, dxdydz.
\end{equation}

Zum Hilbertraum wird $L^2(\mathbb{R}^3)$ erst mit diesem Skalarprodukt, und mit der Vollständigkeit, die wir hier weglassen. Beachten Sie außerdem: Postulat 1 beschreibt \emph{ein} Teilchen --- den Hilbertraum mehrerer Teilchen zu bauen verlangt ein weiteres Rezept, das wir später nachreichen.

\begin{center}
\parbox[c][5cm][c]{14cm}{\begin{tcolorbox}[colback=banana,arc=0cm,frame hidden,enhanced]{\textbf{Postulat 2}. Ein Präparierverfahren wird charakterisiert durch einen \emph{Zustand}, bzw. eine \emph{Dichtematrix}. Das ist ein positiver Operator $\rho$ auf $\mathcal{H}$, mit $\tr(\rho) = 1$. Speziell für reine Zustände ist $\rho$ ein eindimensionaler Projektor $\rho = |\Psi\rangle\langle \Psi|$, also gegeben durch einen Vektor $|\Psi\rangle \in \mathcal{H}$.}\end{tcolorbox}}
\end{center}

Hier ist eine Warnung angebracht: Der Hilbertraum beschreibt \emph{nicht} unmittelbar die Zustände. Viele Bücher sagen, die Zustände seien die Wellenfunktionen; das stimmt so nicht. Er ist der Raum, \emph{in dem} wir Zustände definieren --- die Zustände selbst sind Operatoren. Reine Zustände beschreiben geschlossene Systeme, und die gibt es in der Physik nicht: Das elektromagnetische Feld koppelt immer an. Gemischte Zustände sind der Normalfall.

Das Wort ,,eindimensional`` meint dabei keine Raumdimension, sondern den Rang: $\tr(P)$ zählt die Ket-Bra-Terme eines Projektors, und wegen $\tr(\rho)=1$ steht dort genau einer. Daraus folgt auch die Normierung der Wellenfunktionen, denn mit $\tr(AB) = \tr(BA)$ ist
\begin{equation}
	1 = \tr(\rho) = \tr(|\Psi\rangle\langle \Psi|) = \tr(\langle \Psi|\Psi\rangle) = \langle \Psi|\Psi\rangle.
\end{equation}
Nur normierte Vektoren stellen also reine Zustände dar.

\begin{beispiel}[Zustände eines Qubits]
Für $\mathcal{H} = \mathbb{C}^2$ ist $\rho$ eine positive $2\times 2$-Matrix mit Spur eins. Positivität liefert $\rho_{00} \equiv p \ge 0$, die Spurbedingung erzwingt $\rho_{11} = 1-p$, und weil auch $\rho_{11}\ge 0$ sein muss, liegt $p$ zwischen null und eins. Hermitezität macht schließlich die Nebendiagonaleinträge zueinander komplex konjugiert:
\begin{equation}\label{eq:qubitzustand}
	\rho = \begin{pmatrix} p & x \\ x^* & 1-p\end{pmatrix}, \qquad 0 \le p \le 1, \quad x \in \mathbb{C}.
\end{equation}
Ausgeschöpft ist die Positivität damit nicht; sie verlangt, dass \emph{alle} Eigenwerte nichtnegativ sind, hier also
\begin{equation}
	\det \rho = p(1-p) - |x|^2 \ge 0.
\end{equation}
Ein Qubit-Zustand hängt somit von genau drei reellen Zahlen ab.
\end{beispiel}

\begin{übung}
Berechnen Sie die beiden Eigenwerte der Matrix in Gl.~(\ref{eq:qubitzustand}) und leiten Sie daraus die Bedingung $p(1-p) - |x|^2 \ge 0$ her.
\end{übung}

\begin{center}
\parbox[c][5cm][c]{14cm}{\begin{tcolorbox}[colback=banana,arc=0cm,frame hidden,enhanced]{\textbf{Postulat 3}. Ein Detektor, das heisst ein Messapparat mit den beiden Ansprechmöglichkeiten ,,ja`` (Treffer) und ,,nein``, ist gegeben durch einen Operator $F$ mit $0 \le F \le \mathbb{I}$.}\end{tcolorbox}}
\end{center}

Die Schreibweise $0 \le F \le \mathbb{I}$ ist erklärungsbedürftig: Für Operatoren bedeutet $A \le B$, dass $B-A$ positiv ist. Postulat 3 verlangt also \emph{zwei} Bedingungen, $F \ge 0$ \emph{und} $\mathbb{I} - F \ge 0$, und rechts steht der Einsoperator, nicht die Zahl eins.

\begin{beispiel}[Die Z-Messung am Qubit]
Die Standardmessung auf heutigen Quantenrechnern ist die Z-Messung,
\begin{equation}
	F \equiv |0\rangle\langle 0| = \begin{pmatrix} 1 & 0 \\ 0 & 0 \end{pmatrix}.
\end{equation}
Beide Bedingungen sind zu prüfen. Erstens ist $F$ positiv, denn ein Operator ist genau dann positiv, wenn er sich als $A^\dag A$ schreiben lässt, und hier ist $F = F^\dag F$. Zweitens --- getrennt zu zeigen --- ist
\begin{equation}
	\mathbb{I} - F = \begin{pmatrix} 0 & 0 \\ 0 & 1\end{pmatrix} = |1\rangle\langle 1|
\end{equation}
von derselben Form, also ebenfalls positiv. Erst jetzt dürfen wir $F$ einen Detektor nennen; physikalisch fragt er, ob der Drehimpuls in die $z$-Richtung zeigt.
\end{beispiel}

\begin{center}
\parbox[c][5cm][c]{14cm}{\begin{tcolorbox}[colback=banana,arc=0cm,frame hidden,enhanced]{\textbf{Postulat 4}. Die Ansprechhäufigkeit eines Detektors $F$ an gemäß $\rho$ präparierten System ist
$\tr(\rho F)$. In einem reinen Zustand $|\Psi\rangle$ spezialisiert sich dieser Ausdruck zu $\langle \Psi|F|\Psi\rangle$.}\end{tcolorbox}}
\end{center}

\begin{beispiel}[Ansprechhäufigkeit der Z-Messung]
Mit dem Zustand aus Gl.~(\ref{eq:qubitzustand}) und $F = |0\rangle\langle 0|$ ist
\begin{equation}
	\rho F = \begin{pmatrix} p & x \\ x^* & 1-p\end{pmatrix}\begin{pmatrix} 1 & 0 \\ 0 & 0\end{pmatrix} = \begin{pmatrix} p & 0 \\ x^* & 0\end{pmatrix}, \qquad \mathbb{P}[\text{,,ja``}] = \tr(\rho F) = p.
\end{equation}
Wir erhalten also mit Wahrscheinlichkeit $p$ einen Treffer --- womit gerechtfertigt ist, die Diagonaleinträge einer Dichtematrix als Wahrscheinlichkeiten zu lesen.
\end{beispiel}

Der Zusatz über reine Zustände ist keine weitere Annahme, sondern eine Rechnung: mit $\rho = |\Psi\rangle\langle \Psi|$ und der Zyklizität der Spur ist
\begin{equation}
	\tr(\rho F) = \tr(|\Psi\rangle\langle \Psi|F) = \tr(\langle \Psi|F|\Psi\rangle) = \langle \Psi|F|\Psi\rangle.
\end{equation}
Machen Sie es sich am Anfang zur Gewohnheit, bei jeder Rechnung ausdrücklich zu sagen, welches Postulat Sie gerade benutzen.

Im Allgemeinen kann ein Messapparat in jedem Einzelexperiment ein komplexeres Ergebnis liefern als eine einzelne ja/nein-Entscheidung. Diesen Fall kann man aber auf ja/nein Entscheidungen zurückführen. Nennen wir die Menge der möglichen Ausgänge einer Messung $X$. Für jede Teilmenge $S \subset X$ können wir nun die Frage stellen: Liegt der Messwert $x$ in $S$? Dazu muss es also nach Postulat 3 einen Detektoroperator $F(S)$ geben. Die
Menge der Operatoren $F(S)$ nennen wir dann eine \emph{Observable}. Es gilt für disjunkte Mengen $S_1, S_2 \subset X$ 
\begin{equation}
	F(S_1 \cup S_2) = F(S_1) + F(S_2)
\end{equation}
das heißt, $F$ ist ein ,,Maß`` mit Werten in den positiven Operatoren. Wir setzen normalerweise $F(X) = \mathbb{I}$, das heißt, auch Misserfolg wird als ein mögliches Ergebnis betrachtet. Häufig reicht es, endliche Ausgangsmengen $X$ zu betrachten. Dann ist wegen der Maßeigenschaft alles durch die Operatoren $F_x = F(\{x\})$ bestimmt, die einfach die Relationen 
\begin{equation}
	F_x \ge 0, \quad \text{und}\quad \sum_{x\in X} F_x = \mathbb{I}.
\end{equation}
erfüllen müssen.

Wenn jedes $F(S)$ ein Projektor ist, also $F(S)^2 = F(S)$ erfüllt, nennen wir die Observable \emph{projektionswertig}. Ein weiterer wichtiger Spezialfall sind Observable, die als Ergebnisse reelle Zahlen liefern. Dann können wir den Erwartungswert $\langle A\rangle$ der Verteilung als
\begin{equation}
	\langle A \rangle = \tr(\rho A)
\end{equation}
schreiben, wobei 
\begin{equation}\label{eq:observableA}
	A = \sum_{x\in X} x F_x
\end{equation}
ein hermitescher Operator ist. Im projektionswertigen Fall ist nach dem sogenannten Spektralsatz die ganze Observable schon aus $A$ bestimmt. Man liest deshalb häufig als Definition, dass die Observablen genau die hermiteschen Operatoren sind. 

Ein Zustand charakterisiert ja nicht ein einzelnes System mit irgendwie veränderlichen Eigenschaften, sondern ein Präparierverfahren, das (schon allein zur Wiederholung des statistischen Häufigkeitsexperiments) zu jeder beliebigen Zeit in Gang gesetzt werden können muss. Nach der Annahme der Reproduzierbarkeit muss es auch egal sein, ob sowohl Präparation als auch Messung etwas später durchgeführt werden. Das Einzige was zählt ist ein \emph{relativer} Zeitversatz.

Mit den Postulaten 1 bis 4 wären wir im Prinzip fertig; weil aber ständig Zeitverschiebungen auftreten, wollen wir ein bequemes Hilfspostulat.

\begin{center}
\parbox[c][5cm][c]{14cm}{\begin{tcolorbox}[colback=banana,arc=0cm,frame hidden,enhanced]{\textbf{Postulat 5}. Jede Systemsorte ist weiter charakterisiert durch einen \emph{Hamiltonoperator} $H$, der den \emph{Zeitentwicklungsoperator} $U_t = e^{-\frac{it}{\hbar} H}$ bestimmt. Wendet man in einem Experiment mit gegebenem Präparier- und Messverfahren die Messung mit einer Verzögerung von $t$ Sekunden an, so erhält man den Erwartungswert
\begin{equation}
	\langle A\rangle (t) \equiv \tr(U_t\rho U_t^\dag A) = \tr(\rho U_t^\dag A U_t).
\end{equation}
}\end{tcolorbox}}
\end{center}
Diese Gleichung folgt einfach aus der Eigenschaft $\tr(AB) = \tr(BA)$ der Spur. Sie ist hier aufgenommen, weil die beiden Seiten unterschiedlichen, aber offensichtlich äquivalenten Auffassungen der Zeitentwicklung entsprechen: auf
der linken Seite liegt es nahe, $\rho_t \equiv U_t \rho U_t^\dag$ als zeitabhängige Dichtematrix einzuführen, und die Observablen als zeitunabhängig zu betrachten. Dies nennt sich das \emph{Schrödingerbild}. Für reine Zustände ergibt sich dann auch für Wellenfunktionen: $|\Psi_t\rangle = U_t |\Psi_0\rangle$. Die rechte Seite legt die Einführung zeitabhängiger Observablen $A_t\equiv U_t^\dag A U_t$, bei festgehaltenen Zuständen nahe. Dies nennt sich
das \emph{Heisenbergbild}. Die mathematisch ohnehin evidente Äquivalenz dieser Bilder bedeutet physikalisch, dass nur die relative Zeit zwischen den Teilen des Experiments wichtig ist.

Im Postulat fehlt eine Voraussetzung, die wir hier nachtragen: Der Hamiltonoperator $H$ ist stets \emph{hermitesch}. Nur dann ist $U_t$ unitär, und nur dann ist $U_t^\dag F U_t$ wieder ein zulässiges Messverfahren mit $0 \le U_t^\dag F U_t \le \mathbb{I}$. Ein Experiment mit verzögerter Messung ist damit einfach ein Experiment mit einer anderen Messung; insofern ist Postulat 5 mit den übrigen verträglich. Beweisen lässt es sich aus ihnen trotzdem nicht, und wesentlich bleibt: $H$ ist eine \emph{zusätzliche} Information über die Systemsorte.

\begin{beispiel}[Zeitentwicklung eines Qubits]
Rechnen wir $U_t$ für ein Qubit mit
\begin{equation}
	H \equiv \begin{pmatrix} 0 & 1 \\ 1 & 0 \end{pmatrix}
\end{equation}
aus; dass $H$ hermitesch ist, sieht man sofort. Im Allgemeinen berechnet man Operatorfunktionen über Eigenwerte und Eigenvektoren, hier ist die Taylorreihe bequemer:
\begin{equation}
	U_t = e^{-\frac{it}{\hbar}H} = \sum_{n=0}^\infty \frac{1}{n!}\left(\frac{-it}{\hbar}\right)^n H^n, \qquad H^0 \equiv \mathbb{I}.
\end{equation}
Der Trick ist $H^2 = \mathbb{I}$, also $H^{2n} = \mathbb{I}$ und $H^{2n+1} = H$. Sortiert man die Reihe nach Termen mit $\mathbb{I}$ und mit $H$, erkennt man wegen $(-i)^{2n} = (-1)^n$ die Reihen aus dem ersten Semester wieder:
\begin{equation}
	U_t = \left[\sum_{n=0}^\infty \frac{(-1)^n}{(2n)!}\left(\frac{t}{\hbar}\right)^{2n}\right]\mathbb{I} - i\left[\sum_{n=0}^\infty \frac{(-1)^n}{(2n+1)!}\left(\frac{t}{\hbar}\right)^{2n+1}\right] H,
\end{equation}
also
\begin{equation}
	\eqbox{U_t = \cos\left(\frac{t}{\hbar}\right)\mathbb{I} - i\sin\left(\frac{t}{\hbar}\right) H}
\end{equation}
beziehungsweise
\begin{equation}
	U_t = \begin{pmatrix} \cos\left(\frac{t}{\hbar}\right) & -i\sin\left(\frac{t}{\hbar}\right) \\ -i\sin\left(\frac{t}{\hbar}\right) & \cos\left(\frac{t}{\hbar}\right)\end{pmatrix}.
\end{equation}
Eine Warnung: $H^2 = \mathbb{I}$ ist eine Besonderheit dieses Beispiels und gilt keineswegs allgemein.
\end{beispiel}

% fig_vl06_messung.tex — Was die Postulate 2, 4 und 5 zusammen festlegen:
% der Zustandsraum eines Qubits ist eine Scheibe, die Zeitentwicklung dreht
% darin, und der Detektor liest die waagerechte Koordinate ab.
%
% Chosen over the obvious "measurement cycle" flow chart because these
% postulates contain NO state-update rule -- there is no projection
% postulate in this lecture, so a "state after the measurement" panel would
% illustrate something the text does not claim.  What the text does claim
% has a shape, and the shape is nowhere drawn:
%
%   * Postulat 2 + Hermitezitaet + Spur eins give the lecture's
%         rho = ( p , x ; x^* , 1-p ),
%     and positivity is exactly det rho = p(1-p) - |x|^2 >= 0, i.e.
%         (p - 1/2)^2 + |x|^2 <= (1/2)^2 :
%     a disc of radius 1/2 about (1/2, 0).  Panel (a) is that inequality,
%     drawn in the slice Re x = 0 -- the plane the lecture's own
%     Hamiltonian preserves.  The boundary det rho = 0 is the pure states,
%     the centre is I/2, everything outside is a Hermitian trace-one
%     matrix that is NOT a state: hatched, not filled, so it survives
%     greyscale printing (and hatching is the right mark for "excluded").
%   * Postulat 5 with the lecture's H = ( 0,1 ; 1,0 ) and
%         U_t = cos(t/hbar) I - i sin(t/hbar) H
%     gives, for rho_t = U_t rho U_t^dag,
%         p_t  - 1/2 = r cos(2t/hbar + phi),   Im x_t = r sin(2t/hbar + phi),
%     with (p-1/2, Im x) = r(cos phi, sin phi).  So U_t is a rotation about
%     the centre of the disc at angular velocity 2/hbar: the radius r --
%     equivalently det rho -- is conserved.  The circles in (a) are those
%     orbits, not decoration.
%   * Postulat 4 with the lecture's detector F = |0><0| gives
%         tr(rho_t F) = p_t = 1/2 + r cos(2t/hbar + phi),
%     the horizontal coordinate of the rotating point.  Panel (b) is that
%     function.  Hence the point of the plate: the DISTANCE FROM THE CENTRE
%     in (a) is the AMPLITUDE of the measured frequency in (b).  The pure
%     state swings the full 0..1, the maximally mixed state answers 1/2
%     forever -- it is deaf to this experiment.
%
% The three marked states, all in the slice Re x = 0 (numbers used below):
%   |0><0|   r = 1/2, phi =  0     -> (p, Im x) = (1     , 0     )
%   rho      r = 0.3, phi = 70 deg -> (p, Im x) = (0.6026, 0.2819)
%   I/2      r = 0                 -> (p, Im x) = (0.5   , 0     )
% and their curves in (b) are 1/2 + r cos(2t/hbar + phi) with the same
% r and phi, so tr(rho F) at t = 0 is the p of the dot in (a) -- the
% lecture's own "Ansprechhaeufigkeit der Z-Messung" example.
%
% Colour encodes one thing only: the state, ordered by purity r
% (spacecadet r=0, munsell r=0.3, birrenorange r=1/2, the ramp of
% `tpclevels`).  A state has the same colour in both panels, and its orbit
% in (a) has the colour of the state that runs on it.  Construction lines
% (radius, projection onto p) are neutral grey.
%
% Both axes of (a) carry the same unit length (scale only axis with matched
% width/height), otherwise the circle -- the whole content of the panel --
% would render as an ellipse.

\begin{center}
\begin{tikzpicture}

% ─────────────────────────────────────────────────────────────
%  (a)  The state disc  det rho = p(1-p) - |x|^2 >= 0
% ─────────────────────────────────────────────────────────────
\begin{axis}[
    name=disc,
    tpcplot,
    scale only axis,
    width=4.61cm, height=5.11cm,     % 1.28 x 1.42 units at 3.6 cm/unit
    xmin=-0.14, xmax=1.14,
    ymin=-0.68, ymax=0.74,
    grid=none, clip=false, axis on top,
    xtick={0,0.5,1}, xticklabels={$0$,$\tfrac12$,$1$},
    ytick={-0.5,0,0.5}, yticklabels={$-\tfrac12$,$0$,$\tfrac12$},
    xlabel={$p$}, ylabel={$\imag x$},
    title={(a) \tpcifenglish{state disc}{Zustandsscheibe}},
]

  % Hermitesch, Spur eins, aber nicht positiv: kein Zustand.
  \fill[forbidden] (axis cs:-0.14,-0.68) rectangle (axis cs:1.14,0.74);

  % Das Innere der Scheibe freistellen: det rho >= 0.
  \addplot[fill=white, draw=none, domain=0:360, samples=181]
    ({0.5+0.5*cos(x)}, {0.5*sin(x)});

  % Bahn des gemischten Zustands unter U_t: det rho bleibt erhalten.
  \addplot[munsell, densely dotted, line width=0.8pt,
           domain=0:360, samples=181]
    ({0.5+0.3*cos(x)}, {0.3*sin(x)});
  \addplot[munsell, densely dotted, line width=0.8pt,
           -{Stealth[length=5pt]}, domain=118:152, samples=20]
    ({0.5+0.3*cos(x)}, {0.3*sin(x)});

  % Rand det rho = 0: die reinen Zustaende, ebenfalls eine Bahn.
  \addplot[birrenorange, line width=1.2pt, domain=0:360, samples=181]
    ({0.5+0.5*cos(x)}, {0.5*sin(x)});
  \addplot[birrenorange, line width=1.2pt, -{Stealth[length=5pt]},
           domain=175:205, samples=20]
    ({0.5+0.5*cos(x)}, {0.5*sin(x)});

  % Postulat 4: der Detektor liest die waagerechte Koordinate ab.
  \draw[-{Stealth[length=4pt]}, birrengray!75, line width=0.5pt, densely dotted]
    (axis cs:0.6026,0.2819) -- (axis cs:0.6026,-0.66);
  \node[lbl, font=\scriptsize, text=birrengray, anchor=north west,
        fill=white, fill opacity=0.9, text opacity=1, inner sep=1pt]
    at (axis cs:0.635,-0.52) {$\tr(\rho F)$};

  % Radius r = Abstand vom maximal gemischten Zustand.
  \draw[-{Stealth[length=5pt]}, birrengray, line width=0.7pt]
    (axis cs:0.5,0) -- (axis cs:0.6026,0.2819);
  \node[lbl, font=\scriptsize, text=birrengray, anchor=east]
    at (axis cs:0.525,0.175) {$r$};

  % Die drei Zustaende
  \node[circle, fill=birrenorange, inner sep=1.7pt] at (axis cs:1,0) {};
  \node[circle, fill=munsell,      inner sep=1.7pt] at (axis cs:0.6026,0.2819) {};
  \node[circle, fill=spacecadet,   inner sep=1.7pt] at (axis cs:0.5,0) {};

  \node[lbl, anchor=east]  at (axis cs:0.975,0) {$\ketbra{0}{0}$};
  \node[lbl, anchor=west]  at (axis cs:0.645,0.295) {$\rho$};
  \node[lbl, anchor=north east] at (axis cs:0.49,-0.05) {$\tfrac{\mathbb{I}}{2}$};

  % Rand und Umgebung benennen
  \draw[birrengray, line width=0.5pt] (axis cs:0.235,0.545) -- (axis cs:0.185,0.405);
  \node[lbl, font=\scriptsize, anchor=south west, inner sep=1pt]
    at (axis cs:-0.02,0.55) {\tpcifenglish{pure}{rein}: $\det\rho=0$};
  \node[lbl, font=\scriptsize, text=birrengray, anchor=south west,
        align=left, inner sep=1pt]
    at (axis cs:-0.10,-0.655) {$\det\rho<0$\\ \tpcifenglish{no state}{kein Zustand}};
  \node[lbl, font=\scriptsize, text=birrengray, anchor=north east]
    at (axis cs:1.14,0.72) {$\real x=0$};
  \node[lbl, font=\scriptsize, text=birrengray, anchor=north west]
    at (axis cs:0.125,0.315) {$U_t$};
\end{axis}

% ─────────────────────────────────────────────────────────────
%  (b)  Postulat 4 + 5:  tr(rho_t F) = 1/2 + r cos(2t/hbar + phi)
% ─────────────────────────────────────────────────────────────
\begin{axis}[
    tpcplot,
    at={(disc.south east)}, anchor=south west, xshift=15mm,
    scale only axis,
    width=5.7cm, height=5.11cm,
    xmin=0, xmax=6.2832, ymin=-0.06, ymax=1.10,
    clip=false,
    xtick={0,1.5708,3.1416,4.7124,6.2832},
    xticklabels={$0$,$\tfrac{\pi}{2}$,$\pi$,$\tfrac{3\pi}{2}$,$2\pi$},
    ytick={0,0.5,1}, yticklabels={$0$,$\tfrac12$,$1$},
    xlabel={$t/\hbar$},
    ylabel={$\tr(\rho_t F)$},
    every axis y label/.style={at={(axis description cs:-0.17,0.5)},
                               rotate=90, anchor=south, font=\small},
    title={(b) $Z$-\tpcifenglish{measurement}{Messung}, $F=\ketbra{0}{0}$},
]
  % tr(rho_t F) = 1/2 + r cos(2t/hbar + phi), Amplitude r = Radius aus (a)
  \addplot[birrenorange, line width=1.1pt, domain=0:6.2832, samples=241]
    {0.5+0.5*cos(deg(2*x))};
  \addplot[munsell, line width=1.1pt, domain=0:6.2832, samples=241]
    {0.5+0.3*cos(deg(2*x)+70)};
  \addplot[spacecadet, line width=1.1pt, domain=0:6.2832, samples=2]
    {0.5};

  \node[lbl, anchor=west] at (axis cs:6.40,1.00)  {$r=\tfrac12$};
  \node[lbl, anchor=west] at (axis cs:6.40,0.603) {$r=0.3$};
  \node[lbl, anchor=west] at (axis cs:6.40,0.50)  {$r=0$};
\end{axis}

\end{tikzpicture}
\end{center}


Die Abbildung fasst die Postulate 2, 4 und 5 zusammen: Positivität schneidet aus den Matrizen~(\ref{eq:qubitzustand}) eine Scheibe heraus, $U_t$ dreht darin, und der Detektor liest die waagerechte Koordinate ab. Der Abstand vom Zentrum ist die Amplitude des Signals.

Eine bedeutende Lücke bleibt in diesen Postulaten: nirgends ist gesagt, welche Operatoren $\rho$ oder $F$ man einem gegebenem Experiment zuordnen soll. Dies ist eine Frage, die man tatsächlich nicht durch ein paar allgemeine Postulate klären kann. Der häufigste Zugang ist, ein paar ,,Standardobservablen`` auszuzeichnen, in der Regel durch ihr Transformationsverhalten unter raumzeitlichen Symmetrien. Dies sind dann Größen mit engem Bezug zu klassischen Größen, und heißen auch so: Ort, Impuls, Energie, Drehimpuls und so weiter. Wer eine genauere Zuordnung von Messgeräten und Operatoren wünscht, muss auch den konkreten Messprozess genauer analysieren, besonders die quantenmechanische Wechselwirkung zwischen dem gegebenen System und Teilen des Apparates, wie zum Beispiel ionisierbaren Testteilchen im Apparat. Für fast alle Fragen in dieser Vorlesung werden wir aber mit der Untersuchung der Standardobservablen auskommen.


%!TEX root = TheoretischePhysikC.tex
\section{Vorlesung: ,,Quantisierung``}

Wie indentifiziert man den richtigen Hilbertraum, Hamiltonoperator und Observablen für ein gegebenes System? Leider gibt es keine systematische Methode, um diese Information für ein physikalisches System zu bestimmen. Das bedeutet, dass es bei einer klassischen Beschreibung eines Systems keine mathematisch gut definierte Abbildung gibt, die ihm automatisch das richtige Quantensystem zuordnet. Der Grund dafür ist der folgende. Die klassische Physik stellt eine gute Annäherung an die Physik quantenmechanischer Systeme dar, die Wechselwirkungen mit der Umgebung ausgesetzt sind. Diese schädlichen Auswirkungen der Umwelt auf ein Quantensystem sind als Dekohärenz bekannt. Die Dekohärenz ist der Grund dafür, dass die klassische Mechanik eine so gute Annäherung an die Alltagsphysik darstellt. Geschlossene Systeme gibt es nämlich praktisch nicht: Alle geladenen Objekte koppeln an das elektromagnetische Feld, das selbst ein Quantensystem ist. Bewegen wir ein solches Objekt hin und her, so strahlt es, und die Strahlung trägt die Information mit Lichtgeschwindigkeit ins Unendliche.

Im Idealfall möchte man eine Möglichkeit haben, das ursprüngliche Quantensystem allein aufgrund der Beschreibung des Systems durch die klassische Mechanik aufzudecken. Das heißt, wir wollen eine Möglichkeit, die Dekohärenz umzukehren. Leider ist dies im Allgemeinen unmöglich. Ein solches fiktives Umkehrverfahren ist unter dem Begriff ,,Quantisierung`` bekannt geworden.

\begin{greybox}[Ein inverses Problem]
	Die Dekohärenz bildet verschiedene Quantensysteme auf dasselbe klassische System ab -- zwei Punkte landen auf einem. Ebenso wenig lässt sich aus dem verrauschten Bild einer Überwachungskamera ein Gesicht rekonstruieren.
\end{greybox}

Wir umgehen die Schwierigkeiten der Dekohärenz durch Versuch und Irrtum. Jahrzehntelange Erfahrung hat uns eine Liste von Hilberträumen, Hamiltonoperatoren und Observablen beschert, von denen wir inzwischen wissen, dass sie die korrekte physikalische Beschreibungen für breite Klassen von Systemen liefern. Tatsächlich sind alle fundamentalen Wechselwirkungen (mit Ausnahme der Gravitation) auf diese Weise verstanden worden. Ziel dieser Vorlesung ist es, eine Liste der grundlegenden Quantensysteme zu erstellen, von denen bekannt ist, dass sie bestimmte ideale Systeme korrekt modellieren. Denken Sie dabei nie, es gäbe ein einfaches Rezept dafür -- wäre es so leicht, bräuchten wir keine Experimentalphysik.


\subsection{Ein Teilchen in einer Dimension}
Das wichtigste System, das erfolgreich mit der Quantenmechanik modelliert wurde, ist das Teilchen in einer Dimension. Wie immer gehen wir dabei pflichtgemäß alle fünf Postulate durch. Der Hilbertraum (Postulat 1) lautet in diesem Fall $\mathcal{H}\equiv L^2(\mathbb{R})$. Das heißt, wir beschreiben Präparationen (Postulat 2) von reinen Zuständen eines Teilchens mit Hilfe von komplexen Funktionen -- \emph{Wellenfunktionen} -- $\Psi:\mathbb{R}\rightarrow \mathbb{C}$. Wellenfunktionen werden mit der Ket-Notation nach
\begin{equation}
	|\Psi\rangle \equiv \int_{-\infty}^\infty \Psi(x) |x\rangle \, dx
\end{equation}
beschrieben. Die Beschränkung auf reine Zustände ist keine Willkür: Jede Dichtematrix ist ein Gemisch reiner Zustände. Erlaubt sind nur normierte Vektoren, $\langle \Psi|\Psi\rangle = 1$.

Die Wellenfunktion ist bislang nur Notation; Bedeutung bekommt sie erst durch die Messungen. Eine Theorie braucht beides, Zustände \emph{und} Detektoren.

Gemäß Postulat 3 werden Detektoren mit Operatoren modelliert. Wir beschreiben hier nur Orts-Detektoren. Solche Detektoren liefern ,,ja``/,,nein``-Ergebnisse, d.H.\ das Teilchen befindet sich in einem gegebenem Ort oder nicht. Um solche Detektoren zu beschreiben benötigen wir ein Intervall $[a,b]\subset \mathbb{R}$, wobei $a<b$. Der entsprechende Detektor wird durch den Operator
\begin{equation}
	F_{[a,b]} \equiv \int_{a}^b |x\rangle\langle x|\, dx
\end{equation}
definiert. Solche Operatoren decken zwar nicht alle physikalisch möglichen Detektoren ab, dienen aber als physikalisch relevante Grundlage. Warum ein Bereich und nicht ein Punkt? Weil die Frage nach der genauen Stelle immer die Antwort ,,Wahrscheinlichkeit null`` hat. Auf Intervalle sind wir nicht festgelegt: Für kompliziertere Bereiche sorgt die Additivität -- für disjunkte Mengen $S_1$, $S_2$ gilt $F(S_1\cup S_2) = F(S_1)+F(S_2)$.

Angenommen, das System befindet sich in einem reinen Zustand $|\Psi\rangle$ und wir werden die Messung $F_{[a,b]}$ durchführen. Mit welcher Wahrscheinlichkeit werden wir das Teilchen im Bereich $a \le x \le b$ finden? Postulat 4 liefert uns die Ansprechhäufigkeit des Detektors $F_{[a,b]}$ gemäß einer Präparation $\rho \equiv |\Psi\rangle\langle \Psi|$ nach 
\begin{equation}
	\begin{split}
		\mathbb{P}[\text{,,$x\in [a,b]$``}] &= \tr(\rho F_{[a,b]}) \\
		&= \tr (|\Psi\rangle\langle \Psi|F_{[a,b]})\\
		&= \tr (\langle \Psi|F_{[a,b]}|\Psi\rangle) \\
		&= \langle \Psi|F_{[a,b]}|\Psi\rangle \\
		&= \int_{-\infty}^\infty dx \int_{-\infty}^\infty dy\, \Psi^*(x)\Psi(y) \langle x|F_{[a,b]}|y\rangle   \\
		&= \int_{-\infty}^\infty dx \int_{-\infty}^\infty dy \int_{a}^b dz\, \Psi^*(x)\Psi(y) \langle x|z\rangle\langle z|y\rangle \\
		&= \int_{-\infty}^\infty dx \int_{-\infty}^\infty dy \int_{a}^b dz\, \Psi^*(x)\Psi(y) \delta(x-z)\delta(z-y) \\
		&= \int_{a}^b dz\, \Psi^*(z)\Psi(z) \\
		&= \int_{a}^b |\Psi(z)|^2 \, dz.
	\end{split}
\end{equation}
Im Fall, dass $a\rightarrow -\infty$ und $b\rightarrow \infty$ erhalten wir die Bedingung
\begin{equation}
	1 = \mathbb{P}[\text{,,$x\in \mathbb{R}$``}].
\end{equation} 
Weil die Größe $|\Psi(x)|^2$ ist nicht negativ ist, können wir daher die Größe 
\begin{equation}
	\eqbox{p(x) \equiv |\Psi(x)|^2}
\end{equation}
als eine Wahrscheinlichtkeitsdichte interpretieren.

\begin{intuition}
	Die Wellenfunktion ordnet jeder möglichen Position eine \emph{komplexe} Zahl zu und liefert die Position nur indirekt: Die Wahrscheinlichkeit, das Teilchen zwischen $a$ und $b$ zu finden, ist die Fläche unter der Kurve $p(x)$. Der Versuchung, sich schon $\Psi$ selbst als Dichte vorzustellen, soll man widerstehen.
\end{intuition}

Warnung: Das Objekt
\begin{equation}
	F_{\{x\}} \equiv |x\rangle \langle x|,
\end{equation}
stellt einen Operator \emph{nicht} dar. Dies hat mit der Interpretation von $p(x) \, (\equiv \tr(F_{\{x\}}\rho))$ zu tun: Die Wahrscheinlichkeit, dass wir das Teilchen \emph{genau} an der Stelle $x$ finden, ist immer gleich Null. Man muss immer ein Intervall angeben um eine nicht verschwindende Wahrscheinlichkeit zu erhalten.

Mathematisch sieht man dasselbe daran, dass $F_{\{x\}}$ die Projektoreigenschaft $F^2 = F$ verfehlt:
\begin{equation}
	F_{\{x\}}^2 = |x\rangle \langle x|x\rangle \langle x| = \langle x|x\rangle\, |x\rangle\langle x| = \delta(0)\, F_{\{x\}},
\end{equation}
und $\delta(0)$ ist per Definition unendlich: Statt $F_{\{x\}}$ kommt $\infty\cdot F_{\{x\}}$ heraus. Für ein Intervall dagegen steht in der Mitte $\langle x|y\rangle = \delta(x-y)$:
\begin{equation}
	\begin{split}
		F_{[a,b]}^2 &= \int_a^b \int_a^b |x\rangle \langle x|y\rangle \langle y|\, dx\, dy \\
		&= \int_a^b \int_a^b |x\rangle\, \delta(x-y)\, \langle y|\, dx\, dy \\
		&= \int_a^b |x\rangle\langle x|\, dx = F_{[a,b]}.
	\end{split}
\end{equation}
Die Delta-Funktion erledigt eines der Integrale, und übrig bleibt der Ausgangsoperator: $F_{[a,b]}$ ist wirklich ein Projektor.
\begin{merksatz}
	Alleinstehende Ket-Bra-Kombinationen der Positionsbasis sind zu vermeiden -- sie sind keine Operatoren. Unter einem Integral sind dieselben Ausdrücke harmlos.
\end{merksatz}
Trotzdem arbeitet man manchmal rein informell mit solchen ,,Projektoren``. Insbesondere ist es hilfreich, sich den ,,Projektor`` $|x\rangle\langle x|$ so vorzustellen, dass er das Ereignis darstellt, dass sich das Teilchen an der Position $x$ befindet. Nach dieser Art und Weise führt man die Positionsoperator (gemäß Gl.~(\ref{eq:observableA}))
\begin{equation}
	\eqbox{\widehat{x} \equiv \int_{-\infty}^{\infty} x |x\rangle \langle x|\, dx}
\end{equation}
ein. Man liest ihn als \emph{erstes Moment} der Ortsmessung; das Hütchen erinnert daran, dass hier ein Operator steht und keine Zahl. Er wirkt als Multiplikationsoperator,
\begin{equation}
	\widehat{x}|\Psi\rangle = \int_{-\infty}^{\infty} x\, |x\rangle \langle x|\Psi\rangle\, dx = \int_{-\infty}^{\infty} x\, \Psi(x)\, |x\rangle\, dx,
\end{equation}
aus $\Psi(x)$ wird also $x\,\Psi(x)$.

Eine wichtige Bermerkung: Es ist physikalisch unmöglich, das Teilchen in dem ,,klassischem Zustand`` $|\Psi\rangle = |x\rangle$ zu präparieren. Man bräuchte ein Gerät mit unendlicher Präzision. Dasselbe sagt die Suche nach Eigenvektoren: Ein $|u\rangle$ mit $\widehat{x}|u\rangle = \lambda|u\rangle$ hätte als Wellenfunktion eine Delta-Funktion, wäre also ein Positionsbasiselement. Von ,,Positionseigenvektoren`` spricht man trotzdem ständig -- nur sind es idealisierte Objekte, keine Elemente von $\mathcal{H}$. Wir werden auf dieses Thema in einer späteren Vorlesung zurückkommen. Natürlich vorkommende reine Zustände befinden sich im Allgemeinen immer in Superposition:
\begin{equation}
	\Psi(x) \not= \delta(x-y).
\end{equation}


Der Erwartungswert von $\widehat{x}$ stellt die durchschnittliche Position des Teilchens dar:
\begin{equation}
	\begin{split}
		\langle \widehat{x}\rangle &= \tr(\widehat{x}\rho) \\
		&= \int x|\Psi(x)|^2\, dx,\\
		&= \int x p(x)\, dx.
	\end{split}
\end{equation}
(Diese Gleichung gilt nur im Fall, dass sich das Teilchen in dem reinen Zustand $\rho \equiv |\Psi\rangle\langle \Psi|$ sich befindet.) Hinter dem Symbol $\langle \widehat{x}\rangle$ steckt ein ganzes Verfahren: präparieren, messen, notieren -- viele Male, bis der Mittelwert steht.

Der Impulsoperator eines eindimensionalen quantenmechanischen Teilchens wird durch
\begin{equation}
	\eqbox{\widehat{p} \equiv -i\hbar \frac{d}{dx}}
\end{equation}
definiert. Warum? Man kann zwar viele heuristische Argumente für diese Identifizierung anführen, aber die ehrliche Antwort ist, dass die Natur es einfach so gewollt hat.  Wie wirkt $\widehat{p}$ auf einen Zustand $|\Psi\rangle$? Diese Wirkung wird durch die folgende ,,Matrixelemente`` gegeben (diese gelten für alle $|\Psi\rangle \in \mathcal{H}$):
\begin{equation}
	\langle x|\widehat{p}|\Psi\rangle \equiv -i\hbar \frac{d}{dx}\Psi(x).
\end{equation}

Postulat 5: Zeitentwicklung eines Teilchens wird mit Hilfe des Hamiltonoperators (entspricht der \emph{kinetischen Energie})
\begin{equation}
	\widehat{H} \equiv \frac{\widehat{p}^2}{2m} = -\frac{\hbar^2}{2m}\frac{d^2}{dx^2}
\end{equation}
mit Matrixelemente
\begin{equation}
	\langle x|\widehat{H}|\Psi\rangle = -\frac{\hbar^2}{2m}\frac{d^2}{dx^2}\Psi(x)
\end{equation}
gegeben. Hierbei bezeichnet $m$ die Masse des Teilchens. Woher kommt dieser Ausdruck? Aus der einzigen Quantisierungsregel, die wir haben: Der Hamiltonoperator ist die klassische Energie, mit Hütchen. Klassisch ist $E = p^2/2m$; wir setzen das Hütchen auf. Das ist eine Faustregel, kein Verfahren -- jede Quantisierung bleibt eine Vermutung, über die erst das Experiment entscheidet.

\begin{beispiel}
	Angenommen, ein Teilchen wurde in \emph{gleicher Superposition} über ein Intervall $[-1,1]$ präpariert: Der entsprechende Zustand/Wellenfunktion lautet
	\begin{equation}
		\Psi(x) \equiv \begin{cases}
			\frac{1}{\sqrt2}, \quad x\in [-1,1] \\
			0, \quad \text{sonst}.
		\end{cases}
	\end{equation}
	Bevor wir rechnen, raten wir: Aus Symmetriegründen sollte der Mittelwert genau null sein. Die Normierung $\tr(\rho) = \tr(|\Psi\rangle\langle \Psi|) = \langle \Psi|\Psi\rangle$ des Zustandes wird durch
	\begin{equation}
		\begin{split}
			\langle \Psi|\Psi\rangle &= \int_{-\infty}^\infty |\Psi(x)|^2 \, dx \\
			&= \int_{-1}^1  \frac{1}{2} \, dx \\
			&= 1
		\end{split}
	\end{equation} 
	gegeben. Den Erwartungswert des Positionsoperators $\widehat{x}$ berechnet man wie folgt:
	\begin{equation}
		\begin{split}
			\langle \widehat{x}\rangle = \tr(\widehat{x}\rho) &= \langle \Psi|\widehat{x}|\Psi\rangle \\
		&= \int_{-1}^1 x \times \frac{1}{2}\, dx = \frac12\left[\frac{x^2}{2}\right]_{x= -1}^1 = 0.
		\end{split}
	\end{equation}
\end{beispiel}

%!TEX root = TheoretischePhysikC.tex

\section{Vorlesung: Die Schrödingergleichung}



\subsection{Die Schrödingergleichung}
Von den fünf Postulaten sind die ersten vier meist schnell identifiziert; die eigentliche Arbeit steckt in Postulat 5, und analytisch lösbar ist dort nur eine Handvoll Systeme -- etwa fünf.

Gemäß Postulat 5 wird die Zeitentwicklung eines Systems durch den Zeitentwicklungsoperator
\begin{equation}
	U_t \equiv e^{-\frac{it}{\hbar}\widehat{H}}
\end{equation}
berechnet. Angenommen, dass sich das System in einem reinen Zustand $|\Psi\rangle$ mit dazugehörigem Wellenfunktion $\Psi$  befindet: Wie wird die Zeitentwicklung 
\begin{equation}
	|\Psi(t)\rangle \equiv U_t|\Psi\rangle
\end{equation}
berechnet? Im Allgemeinen wird der Zeitentwicklungsoperator $\widehat{H}$ durch einen Operator, der auf einen unendlichen dimensionalen Hilbertraum wirkt, gegeben. Solche Operatoren sind nur schwer zu diagonalisieren. Deshalb wird eine indirekte Methode verwendet, um $|\Psi(t)\rangle$ zu berechnen: Man bildet die erste Ableitung von $|\Psi(t)\rangle$ bezüglich der Zeit:
\begin{equation}
	\frac{\partial}{\partial t} |\Psi(t)\rangle = \frac{\partial}{\partial t} e^{-\frac{it}{\hbar}\widehat{H}} |\Psi\rangle = -\frac{i}{\hbar}\widehat{H} e^{-\frac{it}{\hbar}\widehat{H}} |\Psi\rangle = -\frac{i}{\hbar}\widehat{H}|\Psi(t)\rangle.
\end{equation}
Dass man so rechnen darf, liegt an der Potenzreihe: Nur die Exponentialfunktion hängt von der Zeit ab, und in ihrer Reihe verhält sich $\widehat{H}$ wie eine Zahl. Umgestellt steht dort
\begin{equation}
	i\hbar\frac{\partial}{\partial t}|\Psi(t)\rangle = \widehat{H}|\Psi(t)\rangle,
\end{equation}
die Schrödingergleichung in allgemeiner Form -- eine für jedes System. Fundamental ist bei uns aber Postulat 5; die Gleichung ist bloß das Mittel, $U_t|\Psi\rangle$ auszurechnen.

Um die Berechnungen zu vereinfachen, nehmen wir das Beispeil eines Teilchens in einer Dimension. Weil die Schrödingergleichung eine Vektorgleichung ist, dürfen wir beide Seiten mit demselben Bra multiplizieren; links vertauscht $\langle x|$ dabei mit der Zeitableitung. Wir berechnen nun das Skalarprodukt von dieser Gleichung und $|x\rangle$:
\begin{equation}
	\begin{split}
		i\hbar \langle x|\frac{\partial}{\partial t} |\Psi(t)\rangle  &= \langle x|\frac{\widehat{p}^2}{2m}|\Psi(t)\rangle \\
		&= \frac{1}{2m}\langle x|\widehat{p}^2|\Psi(t)\rangle 
	\end{split}
\end{equation}
Wir wenden jetzt die von
\begin{equation}
	\langle x|\widehat{p}^2|\Psi\rangle = -\hbar^2\frac{d^2}{dx^2}\Psi(x)
\end{equation}
gegebenen Matrixelemente von $\widehat{H}$ an und setzen die Wellenfunktion-Darstellung ein, um 
\begin{equation}\label{eq:1dfpschr}
	\eqbox{i\hbar\frac{\partial}{\partial t}\Psi(x,t) = -\frac{\hbar^2}{2m}\frac{\partial^2}{\partial x^2}\Psi(x,t)}
\end{equation}
zu erhalten. Diese Gleichung ist als die \emph{Schrödingergleichung} (für ein freies Teilchen in einer Dimension) bekannt.

Es lohnt sich, kurz bei ihrer Form zu verweilen: Streicht man das $i$, so steht dort -- bis auf ein Vorzeichen -- die Wärmeleitungsgleichung. Solche Gleichungen sind gründlich untersucht -- darin liegt der Gewinn der Herleitung: Wir dürfen den ganzen alten Werkzeugkasten anwenden.

Die Lösung der Schrödingergleichung beschreibt die Zeitentwicklung des Zustandes eines Teilchens. 
\begin{beispiel}
	Angenommen, ein Teilchen befindet sich zum Zeitpunkt $t=0$ in einem reinen Zustand -- einem \emph{Gauß'schen Wellenpaket}:
	\begin{equation}
		\Psi(x,t=0) = \sqrt{\frac{1}{\sigma\sqrt{2\pi}}}e^{-\frac14 \frac{(x-a)^2}{\sigma^2}}.
	\end{equation}
	Wie entwickelt sich die Wellenfunktion bezüglich der Zeit? Die Lösung der Schrödingergleichung (\ref{eq:1dfpschr}) in diesem Fall lautet
	\begin{equation}
		\Psi(x,t) = \sqrt{\frac{1}{\sigma(1+i(t/\tau))\sqrt{2\pi}}}e^{-\frac14 \frac{(x-a)^2}{\sigma^2(1+i(t/\tau))}},
	\end{equation}
	wobei
	\begin{equation}
		\tau \equiv \frac{2m\sigma^2}{\hbar}.
	\end{equation}	
\end{beispiel}

Wir fangen mit einer reellen Wellenfunktion an; der Faktor $1+i(t/\tau)$ macht sie komplex. Das Betragsquadrat bleibt eine Glockenkurve, nur eine immer breitere -- das Wellenpaket \emph{zerfließt}.

Bei Problemen mit Teilchen, die sich in Potentialen bewegen, ist es oft einfacher, ihre Dynamik zu berechnen, indem man, wie oben beschrieben, die Schrödingergleichung löst.


\subsection{Fourier Transformation}
Um die Dynamik freien Teilchen im Allgemeinen zu verstehen, ist es hilfreich eine neue mathematische Methode einzuführen. Rein mathematisch liefert die Fourier-Transformation eine \emph{alternative Darstellung} einer Funktion: dieselbe Information, andere Form -- zum Lösen von Gleichungen oft bequemer. In der Musik heißt die Transformierte eines Signals sein \emph{Spektrum}. Angenommen, wir haben eine Funktion $f:\mathbb{R}\rightarrow\mathbb{C}$. Die \emph{Fourier-Transformierte} der Funktion $f(x)$ wird definiert als 
\begin{equation}
	\eqbox{F(k)\equiv \frac{1}{\sqrt{2\pi}}\int_{-\infty}^\infty f(x)e^{-ikx}\, dx.}
\end{equation}
Die \emph{inverse Transformation} lautet
\begin{equation}\label{eq:invfourier}
	\eqbox{f(x)\equiv \frac{1}{\sqrt{2\pi}}\int_{-\infty}^\infty F(k)e^{ikx}\, dk.}
\end{equation}
Die Funktionen $f(x)$ und $F(k)$ sind \emph{Fourier-Paare}. Beide Funktionen beinhalten dieselbe Information. Um diese Äquivalenz zu zeigen, setzen wir die beiden Gleichungen ineinander ein:
\begin{equation}
	\begin{split}
		f(x) &= \frac{1}{\sqrt{2\pi}}\int_{-\infty}^\infty \left(\frac{1}{\sqrt{2\pi}}\int_{-\infty}^\infty f(y)e^{-iky}\, dy\right)e^{ikx}\, dk \\
		&= \int_{-\infty}^\infty f(y)\left(\frac{1}{2\pi}\int_{-\infty}^\infty e^{ik(x-y)}\, dk \right)dy \\
		&= \int_{-\infty}^\infty f(y)\delta(x-y)\,dy.
	\end{split}
\end{equation}
Hierbei haben wir in der zweiten Reihe die Darstellung der Dirac'sche Deltafunktion als 
\begin{equation}
	\delta(x) = \frac{1}{2\pi}\int_{-\infty}^\infty e^{ikx}\, dk
\end{equation}
eigesetzt. Diese Integraldarstellung der Deltafunktion ist nur \emph{eine} von mehreren möglichen; wir übernehmen sie hier ohne Beweis.

\subsection{Lösung der Schrödingergleichung für ein freies Teilchen in $1D$}
Die Schrödingergleichung eines Teilchens in einer Dimension lautet
\begin{equation}
	i\hbar\frac{\partial}{\partial t}\Psi(x,t) = -\frac{\hbar^2}{2m}\frac{\partial^2}{\partial x^2}\Psi(x,t).
\end{equation}
Wie löst man eine Gleichung dieser Form? Der erste Schritt ist bei Differentialgleichungen immer derselbe: Man probiert die Exponentialfunktion. Weil die Schrödingergleichung linear ist, suchen wir nach Lösungen der Form ebener Wellen: $\Psi(x,t) \propto e^{i(kx-\omega t)}$. Einsetzen in die Schrödingergleichung liefert
\begin{equation}
	i\hbar (-i\omega)e^{i(kx-\omega t)} = -\frac{\hbar^2}{2m}(-k^2)e^{i(kx-\omega t)}.
\end{equation}
Wir erhalten dabei eine Ebene-Welle Lösung nur dann, wenn
\begin{equation}
	\eqbox{\hbar\omega = \frac{\hbar^2k^2}{2m}.}
\end{equation} 
Diese Bedingung wird \emph{Dispersionsrelation} bezeichnet. Wenn die Dispersionsrelation erfüllt ist, dann bildet 
\begin{equation}
	\Psi_{k}(x,t) = e^{i\left(kx-\frac{\hbar k^2}{2m} t\right)}
\end{equation}
eine Familie von Lösungen, die mit $k$ bezeichnet sind. Für jedes $k$ gibt es eine solche Lösung, also unendlich viele; addiert werden sie, indem man die Summe durch ein Integral ersetzt -- dabei entdecken wir die Fourier-Transformation quasi wieder.

Wegen der Linearität der Schrödingergleichung können wir neue Lösungen nach dem Superpositionsprinzip bilden:
\begin{equation}
	\eqbox{\Psi(x,t) = \int_{-\infty}^{\infty} \Phi(k)e^{i\left(kx-\frac{\hbar k^2}{2m} t\right)}\, dk.}
\end{equation}
Um die Schrödingergleichung vollständig zu lösen, ist es nötig, die Anfangsbedingungen zu beachten:
\begin{equation}
	\Psi(x,t=0) = \int_{-\infty}^{\infty} \Phi(k)e^{ikx}\, dk.
\end{equation}
Hierbei wird $\Psi(x,0)$ vorgegeben und wir müssen für $\Phi(k)$ lösen. Um dies zu tun vergleichen wir diese Gleichung mit der inversen Transformation Gl.~(\ref{eq:invfourier}): Wir betrachten $\Phi(k)$ als die Fourier-Transformierte von $\Psi(x,0)$ und wenden die Fourier Transformation an, um
\begin{equation}
	\eqbox{\Phi(k) = \frac{1}{2\pi}\int_{-\infty}^{\infty} \Psi(x,0)e^{-ikx}\, dx}
\end{equation}
zu erhalten. Beim Vergleich stört zunächst der Vorfaktor $\frac{1}{\sqrt{2\pi}}$, der in Gl.~(\ref{eq:invfourier}) vor dem Integral steht, in der Anfangsbedingung dagegen fehlt. Die Zuordnung lautet deshalb $f(x)\equiv \Psi(x,0)$ und $F(k)\equiv \sqrt{2\pi}\,\Phi(k)$, und die Hintransformation liefert
\begin{equation}
	\sqrt{2\pi}\,\Phi(k) = \frac{1}{\sqrt{2\pi}}\int_{-\infty}^{\infty} \Psi(x,0)e^{-ikx}\, dx.
\end{equation}
Der Vorfaktor $\frac{1}{2\pi}$ oben ist also nicht geraten, sondern Buchhaltung. Für die Fourier-Transformation sind viele Konventionen im Umlauf; prüfen Sie in fremden Texten zuerst nach, welche gilt.

Warnung: Eine ebene Welle ist \emph{nicht} ein Element von $L^2(\mathbb{R})$, weil die Wellenfunktion $\Psi_{k}(x,t)$ nicht quadratintegriebar ist. Die Funktion $\Psi_{k}(x,t)$ bildet zwar eine Lösung, ist aber selbst nicht ein Element von dem Hilbertraum $\mathcal{H}$. Diese Beobachtung schließt nicht die Möglichkeit aus, dass wir durch lineare Kombinationen von ebenen Wellen quadratintegriebare Funktionen bilden können. Man sieht das sofort an der Wahrscheinlichkeitsdichte: $|\Psi_{k}(x,t)|^2 = 1$, überall gleich groß, und das Integral darüber divergiert -- auch keine Normierungskonstante macht es endlich.

Obwohl die ,,stationäre`` ebene Welle $\Psi_{k}(x) \equiv e^{ikx}$ keinem Vektor im Hilbertraum entspricht,  können wir dennoch eine Ket-Notation einführen: Der ,,Ket-Vektor`` $|\Psi_{k}\rangle$ steht für die ebene Welle. Die dazugehörige Wellenfunktion wird durch
\begin{equation}
	\langle x| \Psi_{k}\rangle = e^{ikx}
\end{equation} 
gegeben.


\subsection{Impuls-Raum}
Wir haben bisher nur Messungen von der Teilchenposition berücksichtigt. Wie kann aus der Wellenfunktion $\Psi(x,t)$ die Statistik für \emph{Impulsmessungen} oder für Größen, die nur vom Impuls abhängen, gewonnen werden? Für eine stationäre ebene Welle $\Psi(x,t) = e^{ikx}$ ist der Impuls durch
\begin{equation}
	\langle x|\widehat{p}|\Psi_k\rangle \overset{!}{=}  -i\hbar \frac{d}{dx} e^{ikx}  \overset{!}{=} \hbar k   e^{ikx} \overset{!}{=} \hbar k \langle x|\Psi_k\rangle
\end{equation}
zu interpretieren: Obwohl der Ket $|\Psi_k\rangle$ nicht als Element in $\mathcal{H}$ existiert, behandeln wir $|\Psi_k\rangle$ trotzdem als einen fiktiven Vektoren. Nach dieser Berechnung ist es sinnvoll $|\Psi_k\rangle$ als eine Art ,,Eigenvektor`` für $\widehat{p}$ zu interpretieren. Daraus ergibt sich, dass eine ebene Welle $\Psi(x,t) = e^{ikx}$ einen bestimmen Impuls von $\hbar k$ hat. Der Vektorgröße $|\Psi_k\rangle$ spielt für $\widehat{p}$ genau diesselbe Rolle, dass $|x\rangle$ für $\widehat{x}$ spielt.

\begin{intuition}
Denken Sie an einen gewöhnlichen Vektor, zerlegt nach zwei verschiedenen Basen: derselbe Vektor, zwei völlig verschiedene Zahlenlisten. Genauso hier -- die Koeffizienten von $|\Psi\rangle$ heißen in der Positionsbasis \emph{Wellenfunktion}, in der Basis der ebenen Wellen \emph{Impulswellenfunktion}.
\end{intuition}

Mit Hilfe dieser Überlegungen führen wir jetzt die folgende Wellenfunktion im Impulsraum ein:
\begin{equation}
	\Phi(p,t) \equiv \frac{1}{\sqrt{2\pi\hbar}} \int_{-\infty}^\infty \Psi(x,t)e^{-i\frac{px}{\hbar}}\,dx.	
\end{equation} 
Man kann sehen, dass diese neue \emph{Impulswellenfunktion} einfach durch die Fourier-Transformation der ursprünglichen Ortswellenfunktion gegeben wird. Daraus schließen wir, dass die Impulswellenfunktion die gleiche Information enthält wie die Ortswellenfunktion; sie beschreiben in unterschiedlichen Darstellungen den gleichen Zustand des Teilchens. Angenommen, wir haben nur die Impulswellenfunktion, können wir die Ortswellenfunktion  mit der inversen Fourier transformation rekonstruieren:
\begin{equation}
	\Psi(x,t) \equiv \frac{1}{\sqrt{2\pi\hbar}} \int_{-\infty}^\infty \Phi(p,t)e^{i\frac{px}{\hbar}}\,dp.	
\end{equation} 
Der Vorfaktor $\frac{1}{\sqrt{2\pi\hbar}}$ und das $\hbar$ im Exponenten wirken unmotiviert; sie sind der Preis für eine schöne Analogie: Mit dieser Konvention steht der Impulsmittelwert nachher genau dort, wo im Ortsraum $\langle\widehat{x}\rangle = \int x|\Psi(x,t)|^2\,dx$ steht.

% fig_vl08_fourier.tex — Das Fourier-Paar: dasselbe Wellenpaket im Orts-
% und im Impulsraum, und die reziproke Breite der beiden Darstellungen.
%
% Follows the conventions of fig_vl11_kasten.tex.
%
%   * geometry from the closed form, not freehand.  The lecture's packet
%         Psi(x,0) = (sigma sqrt(2 pi))^{-1/2} exp(-(x-a)^2 / (4 sigma^2))
%     has the momentum wave function (Fourier transform, p = hbar k)
%         Phi(p)  proportional to exp(-sigma^2 p^2 / hbar^2),
%     so the two densities are Gaussians of standard deviation
%         sigma_x = sigma        and      sigma_p = hbar / (2 sigma).
%     Every curve, every bar length and every peak height in the plate is
%     evaluated from those two numbers:
%         |Psi|^2 peak = 1/(sigma_x sqrt(2 pi)),
%         |Phi|^2 peak = 1/(sigma_p sqrt(2 pi)),
%     which is why the tallest curve of (a) is the flattest of (b).  The
%     three states are sigma_x = sigma/2, sigma, 2 sigma, spaced by a
%     factor of two so that the reciprocal factor of two in (b) can be
%     read off the ruler rather than believed.
%
%   * plotted are the *densities* |Psi|^2 and |Phi|^2, because those are
%     what the lecture integrates in <p> = int p |Phi(p,t)|^2 dp;
%   * colour encodes the state (which packet), nothing else — the same
%     three colours in the same order appear in both panels, so the eye
%     can carry a packet across the Fourier transform;
%   * the +/- sigma rulers sit in a strip above the curves, clear of the
%     data: their lengths are the whole point of the plate, and their
%     order reverses from (a) to (b);
%   * dotted grey drop lines take the middle packet's ruler down onto its
%     own flank in both panels, so the bars are visibly *measured* off a
%     curve rather than asserted; that packet is the worked example, the
%     other two are the trend;
%   * the grey double arrow in the gap carries the only thing the two
%     panels do not say separately: they are one state, not two;
%   * axis units are sigma and hbar/sigma, so the plate carries no
%     arbitrary numbers, and the y axes are unnumbered — the message is
%     the width, not the height.
%
% Panel (a) three packets in position space, narrow to wide.
% Panel (b) the same three in momentum space, wide to narrow, with the
%           product sigma_x sigma_p = hbar/2 that the two panels measure.

\begin{center}
\begin{tikzpicture}

% Dimension ruler: bar–arrow–arrow–bar, the width of one packet.
\tikzset{
  dimbar/.style={line width=0.7pt, >={Stealth[length=4pt,width=3pt]}, |<->|},
}

% ─────────────────────────────────────────────────────────────
%  (a)  Position space
% ─────────────────────────────────────────────────────────────
\begin{axis}[
    name=posspace,
    tpcplot,
    width=6.4cm, height=5.9cm,
    grid=none,
    xmin=-5, xmax=5, ymin=0, ymax=2.62,
    xtick={-4,-2,0,2,4}, ytick=\empty,
    xlabel={$x/\sigma$},
    ylabel={$|\Psi(x)|^2$},
    every axis x label/.style={at={(axis description cs:0.5,-0.13)},
                               anchor=north, font=\small},
    every axis y label/.style={at={(axis description cs:-0.045,0.5)},
                               rotate=90, anchor=south, font=\small},
    title={(a) \tpcifenglish{position space}{Ortsraum}},
    domain=-5:5, samples=201,
    clip=false,
]
  % |Psi|^2 = exp(-x^2/(2 sigma_x^2)) / (sigma_x sqrt(2 pi))
  \addplot[wavefn, spacecadet]   {1.59577*exp(-2*x^2)};       % sigma_x = sigma/2
  \addplot[wavefn, cgblue]       {0.79788*exp(-0.5*x^2)};     % sigma_x = sigma
  \addplot[wavefn, birrenorange] {0.39894*exp(-0.125*x^2)};   % sigma_x = 2 sigma

  % Width rulers, +/- sigma_x, in the state's own colour
  \draw[dimbar, spacecadet]   (axis cs:-0.5,2.40) -- (axis cs:0.5,2.40);
  \draw[dimbar, cgblue]       (axis cs:-1,  2.15) -- (axis cs:1,  2.15);
  \draw[dimbar, birrenorange] (axis cs:-2,  1.90) -- (axis cs:2,  1.90);

  \node[lbl, font=\scriptsize, text=spacecadet,   anchor=west]
    at (axis cs:0.65,2.40) {$\sigma/2$};
  \node[lbl, font=\scriptsize, text=cgblue,       anchor=west]
    at (axis cs:1.15,2.15) {$\sigma$};
  \node[lbl, font=\scriptsize, text=birrenorange, anchor=west]
    at (axis cs:2.15,1.90) {$2\sigma$};

  % The ruler really is read off the curve: drop the middle packet's
  % +/- sigma_x onto its own flank, where |Psi|^2 has fallen to e^{-1/2}
  % of its peak.  The same packet is worked in panel (b).
  \draw[classicallim] (axis cs:-1,2.15) -- (axis cs:-1,0.484);
  \draw[classicallim] (axis cs: 1,2.15) -- (axis cs: 1,0.484);

  \node[lbl, font=\scriptsize, text=birrengray, anchor=west]
    at (axis cs:-4.9,2.54) {$\pm\sigma_x$};
\end{axis}

% ─────────────────────────────────────────────────────────────
%  (b)  Momentum space — the same three packets
% ─────────────────────────────────────────────────────────────
\begin{axis}[
    name=momspace,
    tpcplot,
    at={(posspace.south east)}, anchor=south west, xshift=19mm,
    width=6.4cm, height=5.9cm,
    grid=none,
    xmin=-2.5, xmax=2.5, ymin=0, ymax=2.62,
    xtick={-2,-1,0,1,2}, ytick=\empty,
    xlabel={$p\,\sigma/\hbar$},
    ylabel={$|\Phi(p)|^2$},
    every axis x label/.style={at={(axis description cs:0.5,-0.13)},
                               anchor=north, font=\small},
    every axis y label/.style={at={(axis description cs:-0.045,0.5)},
                               rotate=90, anchor=south, font=\small},
    title={(b) \tpcifenglish{momentum space}{Impulsraum}},
    domain=-2.5:2.5, samples=201,
    clip=false,
]
  % |Phi|^2 with sigma_p = hbar/(2 sigma_x): the order of the widths flips
  \addplot[wavefn, spacecadet]   {0.39894*exp(-0.5*x^2)};     % sigma_p = hbar/sigma
  \addplot[wavefn, cgblue]       {0.79788*exp(-2*x^2)};       % sigma_p = hbar/(2 sigma)
  \addplot[wavefn, birrenorange] {1.59577*exp(-8*x^2)};       % sigma_p = hbar/(4 sigma)

  \draw[dimbar, spacecadet]   (axis cs:-1,   2.40) -- (axis cs:1,   2.40);
  \draw[dimbar, cgblue]       (axis cs:-0.5, 2.15) -- (axis cs:0.5, 2.15);
  \draw[dimbar, birrenorange] (axis cs:-0.25,1.90) -- (axis cs:0.25,1.90);

  \node[lbl, font=\scriptsize, text=spacecadet,   anchor=west]
    at (axis cs:1.15,2.40) {$\hbar/\sigma$};
  \node[lbl, font=\scriptsize, text=cgblue,       anchor=west]
    at (axis cs:0.65,2.15) {$\hbar/(2\sigma)$};
  \node[lbl, font=\scriptsize, text=birrenorange, anchor=west]
    at (axis cs:0.40,1.90) {$\hbar/(4\sigma)$};

  \draw[classicallim] (axis cs:-0.5,2.15) -- (axis cs:-0.5,0.484);
  \draw[classicallim] (axis cs: 0.5,2.15) -- (axis cs: 0.5,0.484);

  \node[lbl, font=\scriptsize, text=birrengray, anchor=west]
    at (axis cs:-2.45,2.54) {$\pm\sigma_p$};

  % What the two rulers measure together
  \node[lbl, font=\small, text=cgblue, anchor=west]
    at (axis cs:-2.45,1.30) {$\sigma_x\,\sigma_p = \hbar/2$};
\end{axis}

% The two panels are one state, not two: the transform lives in the gap.
\coordinate (fmid) at ($(posspace.north east)!0.5!(momspace.north west)+(0,-11mm)$);
\draw[{Stealth[length=4pt,width=3pt]}-{Stealth[length=4pt,width=3pt]},
      birrengray, line width=0.7pt]
  ($(fmid)+(-4mm,0)$) -- ($(fmid)+(4mm,0)$);
\node[lbl, font=\scriptsize, text=birrengray, anchor=south]
  at ($(fmid)+(0,0.7mm)$) {$\mathcal{F}$};

\end{tikzpicture}
\end{center}


Vergleichen Sie in der Abbildung die drei Balken über den Kurven: Halbiert man die Breite im Ortsraum, so verdoppelt sie sich im Impulsraum. Gleiche Farbe heißt dasselbe Paket -- ihre Reihenfolge kippt, das Produkt $\sigma_x\sigma_p$ nicht.


\subsubsection{Statistik von Impulsmessungen}
Die Statistik der Impulsmessungen kann von der Positionsraum-Wellenfunktion abgelesen werden:
\begin{align}
	\langle \widehat{p}\rangle &= \int_{-\infty}^{\infty} \Psi^*(x,t) \left(-i\hbar \frac{\partial}{\partial x} \Psi(x,t)\right)\,dx \\
	\langle \widehat{p}^2\rangle &= \int_{-\infty}^{\infty} \Psi^*(x,t) \left(-\hbar^2 \frac{\partial^2}{\partial x^2} \Psi(x,t)\right)\,dx.
\end{align}
Alternativ können wir die Impulsstatistik mit Hilfe der Impulswellenfunktion berechnen, z.B.:
\begin{equation}
	\begin{split}
		\langle \widehat{p}\rangle &= \frac{1}{2\pi\hbar}\int_{-\infty}^{\infty}\int_{-\infty}^{\infty}\int_{-\infty}^{\infty} \Phi(p',t)e^{-i\frac{p'x}{\hbar}} \left(-i\hbar \frac{\partial}{\partial x} \Phi^*(p,t)e^{i\frac{px}{\hbar}}\right)\,dxdpdp' \\
		&= \frac{1}{2\pi\hbar}\int_{-\infty}^{\infty}\int_{-\infty}^{\infty}\int_{-\infty}^{\infty} \Phi^*(p',t)e^{-i\frac{p'x}{\hbar}} \left(p \Phi(p,t)e^{i\frac{px}{\hbar}}\right)\,dxdpdp' \\
		&= \int_{-\infty}^{\infty}\int_{-\infty}^{\infty} p \Phi^*(p',t) \Phi(p,t) \left(\frac{1}{2\pi\hbar}\int_{-\infty}^{\infty} e^{i\frac{(p-p')x}{\hbar}}\,dx\right)\,dpdp' \\
		&= \int_{-\infty}^{\infty}\int_{-\infty}^{\infty} p \Phi^*(p',t) \Phi(p,t) \left(\frac{1}{2\pi}\int_{-\infty}^{\infty} e^{i{(p-p')x}}\,dx'\right)\,dpdp' \\
		&= \int_{-\infty}^{\infty}\int_{-\infty}^{\infty} p \Phi^*(p',t) \Phi(p,t) \delta(p-p') \,dpdp'\\
		&= \int_{-\infty}^{\infty} p |\Phi(p,t)|^2 \,dp.
	\end{split}
\end{equation}
Nach einem ähnlichen Verfahren erhalten wir:
\begin{equation}
		\langle \widehat{p}^2\rangle = \int_{-\infty}^{\infty} p^2 |\Phi(p,t)|^2 \,dp.
\end{equation}
Das heißt, die Statistik der Impulsmessungen ist wesentlich einfacher bezüglich Impulswellenfunktionen zu berechnen. Vergleichen Sie den Aufwand: erst ableiten, mit $\Psi^*$ multiplizieren, über $x$ integrieren -- oder nur $|\Phi(p,t)|^2$ mit dem Gewicht $p$ integrieren.


\subsection{Ein Teilchen in drei Dimensionen}
Eigentlich ist das Problem schon gelöst: Ein freies Teilchen in drei Dimensionen ist nichts anderes als drei Kopien eines Teilchens in einer Dimension und über den Separationsansatz mitgelöst. Ein Teilchen in drei Dimensionen wird in der Quantenmechanik wie folgt beschrieben
\begin{enumerate}
	\item Hilbertraum: $L^2(\mathbb{R}^3)$.
	\item Präparationen von reinen Zuständen $\rho \equiv |\Psi\rangle\langle\Psi|$: Wellenfunktionen $\Psi:\mathbb{R}^3\rightarrow \mathbb{C}$, wobei
	\begin{equation}
		|\Psi\rangle \equiv \iiint_{-\infty}^\infty \Psi(x,y,z) |x,y,z\rangle \, dxdydz
	\end{equation}
	und
	\begin{equation}
		\langle x',y',z'|x,y,z\rangle = \delta(x'-x)\delta(y'-y)\delta(z'-z).
	\end{equation}
	Damit eine solche Funktion als Wellenfunktion durchgeht, muss sie -- genau wie in einer Dimension -- auf eins normiert sein:
	\begin{equation}
		\langle\Psi|\Psi\rangle = \iiint_{-\infty}^\infty |\Psi(x,y,z)|^2 \, dxdydz = 1.
	\end{equation}
	\item Detektoren: Sei $S\subset \mathbb{R}^3$ ein Bereich. Der Messoperator, der die Messung, ob sich das Teilchen in einer Region S befindet, modelliert, ist durch 
	\begin{equation}
		F_{S} \equiv \iiint_S  	|x,y,z\rangle \langle x,y,z|\, dxdydz
	\end{equation} 
	gegeben.
	\item Ansprechhäufigkeit des Detektors $F_S$: 
	\begin{equation}
		\mathbb{P}[\text{,,$(x,y,z)\in S$``}] \equiv \tr(\rho F_S) = \iiint_S |\Psi(x,y,z)|^2 \, dxdydz.
	\end{equation}
	\item Hamiltonoperator:
	\begin{equation}
		\begin{split}
			\widehat{H} \equiv \frac{\widehat{P}^2_{\text{ges.}}}{2m} &= \frac{\widehat{p}_x^2}{2m}+\frac{\widehat{p}_y^2}{2m}+\frac{\widehat{p}_z^2}{2m} \\ &= -\frac{\hbar^2}{2m}\frac{\partial^2}{\partial x^2}-\frac{\hbar^2}{2m}\frac{\partial^2}{\partial y^2}-\frac{\hbar^2}{2m}\frac{\partial^2}{\partial z^2}.
		\end{split}
	\end{equation}
\end{enumerate}


\subsection{Die Quantisierungsregeln für das Teilchen}
Die Ableitung des Hamiltonoperators für das freie Teilchen ist ein Beispiel für ein Verfahren, wie
aus klassischen Größen ,,entsprechende`` quantenmechanische gewonnen werden können: Dabei werden aus Funktionen der klassischen Variablen Ort und Impuls Operatoren im Raum der Wellenfunktionen. Die folgende Korrespondenz-Tabelle stellt ein Paar Quantisierungsregeln für das Teilchen (in drei Dimensionen) dar
\begin{center}
\begin{tabular}{lcc}    \toprule
{\bf klassisch} & & {\bf quantenmechanischer Operator wirkt nach}  \\\midrule
Ort ($j$-Koord.): &$x_j$ & $\widehat{x}_j \Psi(x,y,z) = x_j \Psi(x,y,z)$  \\ 
\rowcolor{isabelline} Impuls ($j$-Koord.):& $p_j$ & $\widehat{p}_j \Psi(x,y,z) = -i \hbar \frac{\partial}{\partial x_j}\Psi(x,y,z)$  \\ 
kinetische Energie: & $K$ & $\widehat{K}\Psi(x,y,z) = -\frac{\hbar^2}{2m} \nabla^2 \Psi(x,y,z)$ \\
\rowcolor{isabelline} potentielle Energie & $V(x,y,z)$ & $\widehat{V} \Psi(x,y,z) = V(x,y,z)\Psi(x,y,z)$ \\
gesamte Energie & $E$ & $\widehat{H}\Psi(x,y,z) = \left(-\frac{\hbar^2}{2m} \nabla^2+ V(x,y,z)\right)\Psi(x,y,z)$ \\
\rowcolor{isabelline} Drehimpuls & $\mathbf{L}$ & $\widehat{L}_1\Psi(x,y,z) = -i \hbar (y \frac{\partial}{\partial z} - z \frac{\partial}{\partial y})\Psi(x,y,z)$  \\
allgemein & $f$ & ??? \\
\bottomrule
\end{tabular}
\end{center}

Kurz gesagt also, lautet die Regel: Ersetze in jedem klassischen Ausdruck die Ortskoordinaten durch Multiplikationsoperatoren, und die Impulskoordinaten durch ($-i\hbar \times$) Differentiationsoperatoren. Diese Regel ist recht nützlich, und man sollte sie sich unbedingt merken. Davor steht immer ein physikalischer Schritt: Zuerst muss man die \emph{klassische} Energie verstehen, erst daraus gewinnt man den Operator. Und die Regel ist kein wohldefiniertes Verfahren, sondern eine Faustregel -- jedes Mal eine Vermutung, über die erst der Vergleich mit dem Experiment entscheidet. Es gibt aber eine entscheidende Einschränkung: Klassisch ist der Ausdruck 
\begin{equation*}
	p_j x_k - x_k p_j = 0.
\end{equation*}
Mit den oben eingeführten Operatoren dagegen erhält man
\begin{equation}
	\begin{split}
		-i \hbar \left(\frac{\partial}{\partial x_j}x_k - x_k \frac{\partial}{\partial x_j}\right)\Psi(x,y,z) &= -i \hbar \left(x_k\frac{\partial}{\partial x_j}  + \delta_{jk} - x_k \frac{\partial}{\partial x_j} \right)\Psi(x,y,z) \\
		&= -i\hbar \delta_{jk} \Psi(x,y,z).
	\end{split}
\end{equation}
Das heißt, der entsprechende quantenmechanische Operator wirkt wie $-i\hbar\delta_{jk}\mathbb{I}$. Zusammengefasst schreiben wir
\begin{equation}
	[\widehat{p}_j,\widehat{x}_k]\equiv\widehat{p}_j\widehat{x}_k - \widehat{x}_k  \widehat{p}_j = -i\hbar\delta_{jk}\mathbb{I},
\end{equation} 
wobei $[\widehat{A},\widehat{B}] = \widehat{A}\widehat{B} - \widehat{B}\widehat{A}$ den \emph{Kommutator} zweier Operatoren bezeichnet. Diese Operatorgleichung wird als kanonische Vertauschungsrelationen bezeichnet. Der ganze Unterschied ist dabei die Leibniz-Regel: Einmal multipliziert man zuerst mit $x_k$ und leitet dann das Produkt ab, einmal leitet man zuerst ab und multipliziert danach -- der so entstehende Term $\delta_{jk}$ ist, bis auf den Faktor $-i\hbar$, der ganze Kommutator.

Für das Ausfüllen der obigen Korrespondenztabelle bedeuten sie einfach, dass es keine allgemein gültige ,,Quantisierungsvorschrift`` für klassische Größen gibt. Die konkreten Einträge in der Tabelle sind übrigens gerade solche Funktionen von $x$ und $p$, bei denen solche \emph{Operatorordnungsprobleme} keine Rolle spielen. Quantenmechanisch stehen für die klassische Größe $x\,p_x$ dagegen drei Kandidaten nebeneinander -- $\widehat{x}\widehat{p}_x$, $\widehat{p}_x\widehat{x}$ und das symmetrisierte Mittel --, und nichts zeichnet einen davon aus. Die Differenz $\widehat{x}\widehat{p}_x-\widehat{p}_x\widehat{x}$ ist klassisch exakt null, quantenmechanisch mit $\hbar\approx 10^{-34}\,\mathrm{J\,s}$ nur winzig -- aber eben nicht null, und dieser Unterschied ist beobachtbar. Er läuft unter dem Stichwort \emph{Unschärfe}.


%!TEX root = TheoretischePhysikC.tex
\section{Vorlesung: Unschärfe}
Wir betrachten nun freie Teilchen, die im Hilbertraum $\mathcal{H}\equiv L^2(\mathbb{R}^3)$ beschrieben werden. Der Raum $\mathbb{R}^3$ ist hier als Ortsraum interpretiert, das heißt, in dieser Darstellung des Hilbertraums sind die Funktionen der Ortsobservable als Multiplikationsoperatoren dargestellt. In einem reinen Zustand $\rho = |\Psi\rangle\langle \Psi|$ ist daher die Wahrscheinlichkeitsdichte für die Ergebnisse von Ortsmessungen $p(\mathbf{x})\equiv |\Psi(\mathbf{x})|^2$.

Eine Ortsmessung ist dabei ein Gitter von Ja/Nein-Detektoren, im Grunde eine Kamera; der Operator $\widehat{x}_j$ selbst wird nie direkt gemessen, er steht für dieses Verfahren.

Der Erwartungswert der Ortskomponente $\widehat{x}_j$ ergibt sich durch Integration:
\begin{equation}
	\mathbb{E}[x_j] = \iiint x_j |\Psi(\mathbf{x})|^2\, dxdydz = \tr(\rho \widehat{x}_j),
\end{equation}
und entsprechend für höhere Momente. Speziell ist also die Varianz der Verteilung von $x_j$:
\begin{equation}
	\text{var}(x_j) = \tr(\rho \widehat{x}_j^2) - \tr(\rho \widehat{x}_j)^2.
\end{equation}

Für den Impulsoperator erhalten wir analog
\begin{align}
	\mathbb{E}[p_j] &= \tr(\rho \widehat{p}_j) \\
	\text{var}(p_j) &= \tr(\rho \widehat{p}_j^2) - \tr(\rho \widehat{p}_j)^2
\end{align}

Das sind zwei getrennte Experimente: Entweder misst man den Ort, oder man misst den Impuls. Bis hierher unterscheiden sich die Vorhersagen von klassischer Physik und Quantenmechanik nicht; der Bruch kommt jetzt.

Es stellte sich heraus, dass nicht beide Verteilungen zugleich scharf konzentriert sein können. Das heißt: beide Varianzen können nicht beliebig klein werden. Seien $A$ und $B$ beliebige hermitesche Operatoren. Dann ist für alle 
$\lambda \in \mathbb{R}$ der Operator $M = (A+i\lambda B)^\dag (A+i\lambda B)$ (per definition) positiv,
hat also positive Erwartungswerte. Dass $\lambda$ reell sein muss, ist wesentlich: Die Adjungierung
konjugiert komplexe Zahlen, und für komplexes $\lambda$ ginge die Rechnung schief. Der Operator $M$
ist ein Kunstgriff; entscheidend ist allein, dass er die Form $C^\dag C$ mit $C=A+i\lambda B$ hat, denn diese Form garantiert Positivität. Damit ist für $\lambda\in \mathbb{R}$ und eine beliebige
Dichtematrix $\rho$:
\begin{equation}
	\tr(\rho M) = \tr(\rho A^2)+\lambda^2\tr(\rho B^2) + \lambda\tr(\rho i[A,B])\ge 0.
\end{equation}
Setzt man hier für $\lambda$ den Wert ein, bei dem die linke Seite am kleinsten wird (Übung),
bekommt man:
\begin{equation}
	\tr(\rho A^2)\tr(\rho B^2) \ge \frac14 \tr(\rho i[A,B])^2.
\end{equation}
Das Rezept dafür: nach $\lambda$ ableiten, die Ableitung null setzen. Die Ungleichung gilt für beliebige hermitesche Operatoren.
Setzt man hier ein: $A=\widehat{p}_j - \tr(\rho \widehat{p}_j )\mathbb{I}$ und $B=\widehat{x}_j - \tr(\rho \widehat{x}_j )\mathbb{I}$, und benutzt die kanonischen Vertauschungsrelationen für die rechte Seite, so bekommt man
\begin{equation}
	\eqbox{\text{var}(p_j)\text{var}(x_j) \ge \frac{\hbar^2}{4}, \quad\text{oder}\quad (\Delta p_j)(\Delta x_j) \ge \frac{\hbar}{2},}
\end{equation}
wenn man die $(\Delta A)$ als Wurzel der Varianz (Standardabweichung) einführt. Dies ist die \emph{Heisenberg'sche Unschärferelation}, die in dieser Form von Kennard stammt.

Die Interpretation der Heisenberg'schen Unschärferelation lautet: Es ist unmöglich, quantenmechanische Teilchen so zu präparieren, dass die Varianzen bei Impuls-und Ortsmessungen beide beliebig klein ausfallen. Die Größen $(\Delta p_j)$ und $(\Delta x_j)$ beziehen sich also auf völlig verschiedene statistische Experimente. Eine Standardabweichung entsteht erst durch viele Wiederholungen: präparieren, messen, wieder genauso präparieren, wieder messen; $(\Delta p_j)$ bestimmt man in einem davon getrennten Durchgang.

\begin{merksatz}[Was die Unschärferelation nicht sagt]
Sie sagt nichts über sukzessive Messungen. Die verbreitete Darstellung, man messe erst den Ort und dann am selben Teilchen den Impuls, ist falsch, auch wenn fast alle Autoren sie so bringen. Für Messfolgen gilt zwar ebenfalls eine Ungleichung, aber eine andere und kompliziertere. Diese Unterscheidung ist prüfungsrelevant.
\end{merksatz}
%!TEX root = TheoretischePhysikC.tex


\section{Vorlesung: Zerfließen des Wellenpakets}
Der Ausdruck ,,Wellenpaket`` wird für räumlich begrenzte Funktionen verwendet, wie etwa $\Psi\in L^2(\mathbb{R}^3)$, nicht für ebene Wellen. Vom Standpunkt der statistischen Interpretation gibt es also nur solche ,,Pakete``, und die Diskussion ebener Wellen ist bestenfalls eine etwas verkürzte Art, über die Fourier-Transformation zu sprechen. In der Frühzeit der Quantenmechanik galten gut lokalisierte Wellenpakete außerdem als eine Art Ersatz für Teilchen. Gemessen
daran war eine Eigenschaft der Wellenpakete etwas störend: sie ,,flossen auseinander``,  was für Teilchen ja wirklich ein ungebührliches Betragen ist. Das Zerfließen des Wellenpakets galt danach als typisch quantenmechanischer Effekt. Wenn man allerdings $\Psi$ konsequent statistisch deutet, ist nichts Paradoxes an diesem Vorgang: warum soll schließlich eine Wahrscheinlichkeitsverteilung
nicht mit der Zeit breiter werden? Wir werden hier sehen, dass das Zerfließen der Wellenpakete gerade ein klassischer Effekt ist: er wird asymptotisch genau durch die Impulsverteilung bestimmt, wie bei einem Schwarm klassischer Teilchen.

\begin{greybox}[Warum reale Teilchen trotzdem nicht zerfließen]
	Allein ist ein Teilchen selten. Materie besteht aus geladenen Bestandteilen, und das elektromagnetische Feld fragt sie unablässig ,,bist du da?`` -- das Teilchen antwortet, indem es ein wenig Feld erzeugt, und diese pausenlose Ortsmessung lokalisiert es immer wieder neu. Stichwort: \emph{Dekohärenz}.
\end{greybox}

Wir betrachten also jetzt Teilchen in $\mathbb{R}^3$ mit Hamiltonoperator
\begin{equation}
	H = \frac{\widehat{p}_\text{ges.}^2}{2m} = \frac{\widehat{p}_{x}^2}{2m}+\frac{\widehat{p}_{y}^2}{2m}+\frac{\widehat{p}_{z}^2}{2m}.
\end{equation}
In diesem Fall entspricht wie üblich in der QM $H$ der Energie eines freien Teilchens. 

Wir interessieren uns für die Wahrscheinlichkeitsverteilung des Ortes zu verschiedenen Zeiten, besonders zu großer Zeit, das heißt, für die Ergebnisse von
Ortsmessungen, wenn zwischen Präparation und Messung eine große Zeit $t$ verstrichen ist. Wir haben zwei äquivalente Alternativen, um dieses Problem anzugehen: das Heisenbergbild und das Schrödingerbild. Im Schrödingerbild würden wir die Schrödingergleichung 
\begin{equation}
	i\hbar \frac{\partial}{\partial t}\Psi = \widehat{H}\Psi
\end{equation}
lösen, und dann das Ergebnis als 
\begin{equation}
	p_t(x) \equiv |\Psi(x,t)|^2
\end{equation}
ablesen -- eine Wolke im Raum, Detektorklick für Detektorklick aufgesammelt. Der Unterschied der Bilder ist reine Buchhaltung: dort bewegt sich der Zustand, hier die Messung, und wegen $\tr(\rho_t\widehat{F}) = \tr(\rho\,\widehat{F}(t))$ liefern beide dieselben Zahlen. Wir werden die Rechnung im Heisenbergbild durchführen, d.h. wir halten $\Psi$ (und dabei
die Dichtematix $\rho$) fest und berechnen die Operatoren
\begin{equation}
	\widehat{x}(t) \equiv e^{\frac{it}{\hbar} \widehat{H}}\widehat{x} e^{-\frac{it}{\hbar} \widehat{H}}, \quad \widehat{y}(t) \equiv e^{\frac{it}{\hbar} \widehat{H}}\widehat{y} e^{-\frac{it}{\hbar} \widehat{H}} \quad \text{und} \quad \widehat{y}(t) \equiv e^{\frac{it}{\hbar} \widehat{H}}\widehat{y} e^{-\frac{it}{\hbar} \widehat{H}}.
\end{equation}

Als nächstes suchen wir passende ,,Koordinaten``, in denen wir unsere Rechnung duchführen können. Wegen des Auftretens der Exponentialfunktion ist es naheliegend, eine Darstellug des Hilbertraums zu wählen, in der $H$ ein Multiplikationsoperator ist, also die Impulsdarstellung:
\begin{equation}
	\Phi(\mathbf{p})\equiv \Phi(p_x,p_y,p_z) \equiv \frac{1}{\sqrt{2\pi\hbar}} \iiint_{-\infty}^\infty \Psi(x)e^{-\frac{i}{\hbar}(p_x x+ p_y y+ p_z z)}\,dxdydz.	
\end{equation} 
Zu jedem $\mathbf{p}$ gehört dabei die ebene Welle $|\Psi_{\mathbf{p}}\rangle$ aus Vorlesung 8, hier so normiert, dass $\langle\Psi_{\mathbf{p}}|\Psi_{\mathbf{p}'}\rangle = \delta^{(3)}(\mathbf{p}-\mathbf{p}')$ gilt:
\begin{equation}
	\langle \mathbf{x}|\Psi_{\mathbf{p}}\rangle = \frac{1}{(2\pi\hbar)^{3/2}}e^{\frac{i}{\hbar}(p_x x + p_y y + p_z z)},
\end{equation}
und diese uneigentlichen Kets liefern eine Zerlegung der Eins:
\begin{equation}\label{eq:eins-impuls}
	\mathbb{I} = \iiint_{-\infty}^\infty |\Psi_{\mathbf{p}}\rangle\langle\Psi_{\mathbf{p}}|\,dp_xdp_ydp_z.
\end{equation}
Die Koeffizienten $\langle\Psi_{\mathbf{p}}|\Psi\rangle$ in dieser Basis sind gerade $\Phi(\mathbf{p})$.\footnote{Der Vorfaktor oben in $\Phi(\mathbf{p})$ ist der eindimensionale; dreidimensional gehört dorthin $(2\pi\hbar)^{-3/2}$.} Solches Einschieben einer Eins ist unser wichtigstes Rechenwerkzeug; wir brauchen es gleich wieder.

Gewählt ist diese Basis, weil die ebenen Wellen Eigenvektoren von $\widehat{H}$ sind,
\begin{equation}\label{eq:heigen}
	\widehat{H}|\Psi_{\mathbf{p}}\rangle = \frac{|\mathbf{p}|^2}{2m}|\Psi_{\mathbf{p}}\rangle,
\end{equation}
so dass in der Potenzreihe jedes $\widehat{H}$ zu einer Zahl wird:
\begin{equation}
	e^{-\frac{it}{\hbar}\widehat{H}}|\Psi_{\mathbf{p}}\rangle = \sum_{n=0}^{\infty}\frac{1}{n!}\left(-\frac{it}{\hbar}\right)^n \widehat{H}^n|\Psi_{\mathbf{p}}\rangle = e^{-\frac{it}{\hbar}\frac{|\mathbf{p}|^2}{2m}}|\Psi_{\mathbf{p}}\rangle.
\end{equation}
Im Exponenten stand ein Operator, jetzt eine Zahl -- daher unsere Besessenheit von Eigenvektoren. Erlaubt ist das \emph{nur} auf einem Eigenvektor: Eine Matrix durch eine Zahl zu ersetzen, das muss man sich verdienen.

Wir brauchen gleich außerdem, dass $\widehat{x}$ in der Impulsdarstellung als $i\hbar\frac{\partial}{\partial p_x}$ wirkt. Eindimensional schiebt man dazu die Ortszerlegung der Eins ein,
\begin{equation}
	\langle \Psi_p|\widehat{x}|\Psi\rangle = \int_{-\infty}^\infty \langle \Psi_p|x\rangle\, x\, \Psi(x)\,dx = \frac{1}{\sqrt{2\pi\hbar}}\int_{-\infty}^\infty x\, e^{-\frac{i}{\hbar}px}\,\Psi(x)\,dx,
\end{equation}
erzeugt den Faktor $x$ als Ableitung nach $p$,
\begin{equation}
	i\hbar\frac{\partial}{\partial p}e^{-\frac{i}{\hbar}px} = i\hbar\left(-\frac{ix}{\hbar}\right)e^{-\frac{i}{\hbar}px} = x\,e^{-\frac{i}{\hbar}px},
\end{equation}
und zieht diese vor das Integral:
\begin{equation}\label{eq:xinp}
	\langle \Psi_p|\widehat{x}|\Psi\rangle = i\hbar\frac{\partial}{\partial p}\Phi(p).
\end{equation}
Für $\widehat{y}$ und $\widehat{z}$ läuft es genauso -- eine schöne Symmetrie zur Ortsdarstellung, wo umgekehrt $\widehat{p}$ als $-i\hbar\frac{\partial}{\partial x}$ wirkt: dieselbe Struktur mit vertauschten Rollen, bis auf das Vorzeichen.

Dann ist
\begin{equation}\label{eq:timedepx}
	\begin{split}
		(\widehat{x}(t)\Phi)(\mathbf{p}) &= e^{\frac{it}{\hbar} \frac{|\mathbf{p}|^2}{2m}}\left(\widehat{x}(t)e^{-\frac{it}{\hbar}\widehat{H}}\Phi\right)(\mathbf{p}) \\
		&= e^{\frac{it}{\hbar} \frac{|\mathbf{p}|^2}{2m}}\left(i\hbar \frac{\partial}{\partial p_x}\right)\left(e^{-\frac{it}{\hbar}\frac{|\mathbf{p}|^2}{2m}}\Phi(\mathbf{p})\right) \\
		&= t \frac{p_x}{m}\Phi(\mathbf{p}) + i\hbar \frac{\partial}{\partial p_x} \Phi(\mathbf{p}), \\ 
	\end{split}
\end{equation}
also
\begin{equation}
	\widehat{x}(t) = t \frac{\widehat{p}_x}{m} + \widehat{x}(0),
\end{equation}
wobei $\frac{\widehat{p}_x}{m}$ den \emph{Geschwindigkeitsoperator} (in der $x$-Richtung) darstellt. Es gelten ähnliche Ausdrücke für $\widehat{y}(t)$ und $\widehat{z}(t)$. Den Ausdruck $\frac{p_\mu}{m} = \frac{\partial}{\partial p_\mu} \frac{|\mathbf{p}|^2}{2m}$, $\mu = x,y,z$, bezeichnet man auch als \emph{Gruppengeschwindigkeit}, weil er die Geschwindigkeit von \emph{Gruppen}, also Wellenpaketen beschreibt. Die Rechnung ist aufwendig; merken sollte man sich ihr Ergebnis, nicht die Schrittfolge.

Um die Zeitabhängigkeit des Erwartungswertes von $\widehat{x}(t)$ zu bekommen, braucht man nun nur noch auf beiden Seiten von (\ref{eq:timedepx}) die Spur mit $\rho =  |\Psi\rangle\langle\Psi|$ zu nehmen. Wir erhalten dabei
\begin{equation}
	\begin{split}
		\langle\widehat{x}\rangle(t) &= \tr(\rho \widehat{x}(t)) \\
		&= \tr\left(\rho \left(t \frac{\widehat{p}_x}{m} + \widehat{x}(0)\right)\right)\\
		&= t \frac{\langle\widehat{p}_x \rangle(0)}{m} + \langle\widehat{x}\rangle(0).
	\end{split}
\end{equation}
Das ist genau die newtonsche Regel ,,Ort gleich Anfangsort plus Geschwindigkeit mal Zeit``.

\begin{intuition}
	Dieselbe Formel, andere Physik: Klassisch \emph{hat} ein Teilchen eine Position, Erwartungswerte braucht man gar nicht. Hier läuft nur der Mittelwert klassisch, während die Breite wächst -- eine klassische Wolke mit scharfem Ort und Impuls zerfließt nicht, eine quantenmechanische muss es. Klassische Mechanik ist der Grenzfall $\hbar\to 0$.
\end{intuition}

Die Operatorgleichung enthält aber viel mehr Information als nur das: Die Diagonalisierung der Operatoren $\widehat{x}(t)$ enthält ja alle Information über die Wahrscheinlichkeitsverteilungen des Ortes zur Zeit $t$ in beliebigen Zuständen. Das Aussehen der Ortswahrscheinlichkeitsverteilung zu großen Zeiten ist also völlig von der gemeinsamen Verteilung der Operatoren $\frac{\widehat{p}_\mu}{m}$ bestimmt. Am Verhalten von $\widehat{x}$, $\widehat{y}$ und $\widehat{z}$  ist also überhaupt kein Unterschied zwischen klassischen und quantenmechanischen Teilchen zu sehen.

% fig_vl10_zerfliessen.tex — Zerfliessen des Wellenpakets: the packet
% drifts uniformly while it spreads, and its width settles onto exactly
% the straight line a swarm of classical particles with the same
% momentum spread would follow.
%
% GEOMETRY FROM THE FORMULA.  Everything here follows from the operator
% identity this lecture derives,
%
%       x(t) = x(0) + (t/m) p ,
%
% by taking mean and variance in a state with no initial x-p
% correlation:
%
%       <x>(t)  = <x>(0) + <p> t/m ,
%       Dx(t)^2 = Dx(0)^2 + (Dp t/m)^2 .
%
% With tau = m Dx(0)/Dp — and for a minimum-uncertainty packet,
% Dx(0) Dp = hbar/2, hence tau = 2 m Dx(0)^2/hbar — the second line is
%
%       Dx(t) = Dx(0) sqrt(1 + (t/tau)^2),
%
% the hyperbola of panel (b).  Its asymptote Dp t/m is precisely how a
% classical swarm with momentum spread Dp comes apart, so the claim of
% the lecture — that the spreading is asymptotically fixed by the
% momentum distribution alone — is the statement that the two curves in
% (b) merge.  Panel (a) plots the Gaussian carrying those two moments,
%
%       |Psi(x,t)|^2 = exp(-(x-<x>(t))^2/(2 Dx(t)^2)) / (sqrt(2 pi) Dx(t)),
%
% so every centre, width and height is forced by the formula rather than
% drawn by eye.  Lengths in units of Dx(0), times in units of tau, drift
% fixed by <p> = 3 Dp: the centres then sit at 0, 3, 6, 9 and the widths
% at 1, sqrt2, sqrt5, sqrt10.
%
% DESIGN DECISIONS
%   * colour encodes TIME and nothing else — entries 1, 2, 4, 5 of the
%     ordered `tpclevels` ramp (navy, blue, orange, red), the same four
%     in both panels, so a curve in (a) can be found again as a dot in
%     (b).  Entry 3 (munsell) is skipped on purpose: next to cgblue it
%     is the same teal to the eye, and here the two would be neighbours
%     in time rather than well separated as in fig_vl11;
%   * the width curve in (b) is green because it is a different
%     QUANTITY, not a fifth time; its formula is set in the same green,
%     which spares the panel a legend;
%   * grey is reference furniture only, and every grey annotation is
%     set in grey too: the verticals at <x>(t), whose EQUAL SPACING is
%     uniform motion; the dashed envelope through the peaks, which
%     falls as 1/Dx(t) exactly as fast as the widths grow, so the
%     conserved area is shown rather than asserted, and whose label is
%     set in the empty strip immediately above it rather than in the
%     right margin, where it crowded the t=3tau label; the initial width
%     and the classical asymptote in (b);
%   * time labels sit to the RIGHT of each peak rather than above it,
%     which is the only side on which the descending envelope leaves
%     room;
%   * panel (b) uses `scale only axis` with a 4.2 x 5.4 cm box on ranges
%     4.2 x 5.4, i.e. 1 cm per unit on both axes.  The asymptote is then
%     a true 45-degree line and its label can be set along it.  The y
%     range runs to 5.4 although the curve reaches only sqrt(1+4.2^2) =
%     4.32 inside the frame: that headroom is where the three annotation
%     lines sit, so none of them is written over the curve.  The widest
%     of them, the green formula, is 3.8 cm in the fonts of the notes —
%     wider than in the cmfonts preview, which is why the panel is sized
%     against the former;
%   * both panels are 5.4 cm tall so that titles and x axes line up;
%   * the orders of magnitude for tau — 10^-16 s for an electron
%     localised to an angstrom, 10^13 s for a microgram grain — are
%     NOT in the panel.  They were, at its foot, and in the fonts of the
%     notes (wider than the cmfonts preview) the grey asymptote ran
%     straight through them; a 4.2 cm panel crossed by two reference
%     lines has no free region 3.4 cm wide, so the numbers now stand in
%     the sentence below the figure, where the definition of tau in (b)
%     hands them straight over.

\begin{center}
\begin{tikzpicture}

% ─────────────────────────────────────────────────────────────
%  (a)  |Psi(x,t)|^2 at four times
% ─────────────────────────────────────────────────────────────
\begin{axis}[
    name=packets,
    tpcplot,
    scale only axis, width=6.4cm, height=5.4cm,
    grid=none,
    xmin=-3.2, xmax=16.8,
    ymin=0,    ymax=0.47,
    xtick={0,3,6,9,12,15},
    ytick={0,0.1,0.2,0.3,0.4},
    xlabel={$x/\Delta x(0)$},
    ylabel={$\Delta x(0)\,|\Psi(x,t)|^2$},
    every axis x label/.style={at={(axis description cs:0.5,-0.15)},
                               anchor=north, font=\small},
    every axis y label/.style={at={(axis description cs:-0.16,0.5)},
                               rotate=90, anchor=south, font=\small},
    title={(a) \tpcifenglish{The packet spreads}{Das Paket zerfließt}},
    samples=401,
]

  % Centres <x>(t) = 3 t/tau, each rule drawn only up to the peak it
  % belongs to.  Their equal spacing is the newtonian rule.
  \draw[levellinegrey] (axis cs:0,0) -- (axis cs:0,0.3989);
  \draw[levellinegrey] (axis cs:3,0) -- (axis cs:3,0.2821);
  \draw[levellinegrey] (axis cs:6,0) -- (axis cs:6,0.1784);
  \draw[levellinegrey] (axis cs:9,0) -- (axis cs:9,0.1262);

  % Envelope of the peaks: 1/(sqrt(2 pi) Dx(t)) with t/tau = s/3.  It
  % passes exactly through the four maxima.
  \addplot[envelope, domain=0:12.6, samples=140]
      {1/(sqrt(2*pi)*sqrt(1+x^2/9))};

  % The packets.  Widths sqrt(1+(t/tau)^2) = 1, sqrt2, sqrt5, sqrt10.
  % Each runs across the whole panel: a curve stopped short leaves a
  % blunt end lying on the axis, which reads as a cut-off rather than
  % as a tail that has died away.
  \addplot[wavefn, spacecadet,   domain=-3.2:16.8]
      {exp(-(x)^2/2)/sqrt(2*pi)};
  \addplot[wavefn, cgblue,       domain=-3.2:16.8]
      {exp(-(x-3)^2/4)/sqrt(2*pi*2)};
  \addplot[wavefn, birrenorange, domain=-3.2:16.8]
      {exp(-(x-6)^2/10)/sqrt(2*pi*5)};
  \addplot[wavefn, birrenred,    domain=-3.2:16.8]
      {exp(-(x-9)^2/20)/sqrt(2*pi*10)};

  % Direct labels beside the peaks, each in the colour of its curve and
  % set just above the peak height, where the envelope has already
  % dropped away.
  \node[lbl, text=spacecadet,   anchor=south west]
    at (axis cs:0.55,0.410) {$t=0$};
  \node[lbl, text=cgblue,       anchor=south west]
    at (axis cs:3.45,0.293) {$t=\tau$};
  \node[lbl, text=birrenorange, anchor=south west]
    at (axis cs:6.45,0.189) {$t=2\tau$};
  \node[lbl, text=birrenred,    anchor=south west]
    at (axis cs:9.45,0.136) {$t=3\tau$};

  % What the grey verticals say, and what the grey envelope says.
  \node[lbl, text=birrengray, font=\scriptsize, anchor=north east,
        align=right] at (axis cs:16.6,0.468)
    {\tpcifenglish{centres}{Zentren}:
     $\langle\widehat{x}\rangle(t)=\langle\widehat{p}\rangle\,t/m$\\[1pt]
     \tpcifenglish{area}{Fl\"ache}: $\int|\Psi|^2dx=1$};
  \node[lbl, text=birrengray, font=\scriptsize, anchor=south west]
    at (axis cs:4.7,0.232)
    {\tpcifenglish{peak height}{H\"ohe} $\propto 1/\Delta x(t)$};
\end{axis}

% ─────────────────────────────────────────────────────────────
%  (b)  The width: hyperbola against its classical asymptote
% ─────────────────────────────────────────────────────────────
\begin{axis}[
    name=widths,
    tpcplot,
    at={(packets.south east)}, anchor=south west, xshift=17mm,
    scale only axis, width=4.2cm, height=5.4cm,
    xmin=0, xmax=4.2, ymin=0, ymax=5.4,
    xtick={0,1,2,3,4}, ytick={0,1,2,3,4},
    grid=major,
    xlabel={$t/\tau$},
    ylabel={$\Delta x(t)/\Delta x(0)$},
    every axis x label/.style={at={(axis description cs:0.5,-0.15)},
                               anchor=north, font=\small},
    every axis y label/.style={at={(axis description cs:-0.21,0.5)},
                               rotate=90, anchor=south, font=\small},
    title={(b) \tpcifenglish{Width}{Breite} $\Delta x(t)$},
]

  % Reference furniture: the classical asymptote Dp t/m and the width
  % the packet started with.  They cross at t = tau.
  \addplot[levellinegrey, domain=0:4.2, samples=2] {x};
  \draw[levellinegrey] (axis cs:0,1) -- (axis cs:1,1);

  % Dx(t)/Dx(0) = sqrt(1+(t/tau)^2)
  \addplot[birrengreen, line width=1.1pt, domain=0:4.2, samples=140]
      {sqrt(1+x^2)};

  % The four times of panel (a), in their colours.
  \addplot[only marks, mark=*, mark size=2.0pt, spacecadet]
      coordinates {(0,1)};
  \addplot[only marks, mark=*, mark size=2.0pt, cgblue]
      coordinates {(1,1.41421)};
  \addplot[only marks, mark=*, mark size=2.0pt, birrenorange]
      coordinates {(2,2.23607)};
  \addplot[only marks, mark=*, mark size=2.0pt, birrenred]
      coordinates {(3,3.16228)};

  % The curve names itself, in its own colour; the grey lines below it
  % state the asymptotics and fix the time scale.
  \node[text=birrengreen, font=\scriptsize, anchor=north west]
    at (axis cs:0.08,5.32)
    {$\Delta x(t)=\Delta x(0)\sqrt{1+(t/\tau)^{2}}$};
  \node[lbl, text=birrengray, font=\scriptsize, anchor=north west]
    at (axis cs:0.08,4.68)
    {$\longrightarrow\ \Delta p\,t/m$ \tpcifenglish{for}{f\"ur} $t\gg\tau$};
  \node[lbl, text=birrengray, font=\scriptsize, anchor=north west]
    at (axis cs:0.08,4.14) {$\tau=2m\,\Delta x(0)^2/\hbar$};

  % The asymptote, labelled along itself — the panel is drawn at
  % 1 cm per unit on both axes, so the line really does run at 45 deg.
  \node[lbl, text=birrengray, font=\scriptsize, rotate=45,
        anchor=north west, inner sep=1.5pt] at (axis cs:1.65,1.50)
    {\tpcifenglish{classical swarm}{klassischer Schwarm}};
\end{axis}

\end{tikzpicture}
\end{center}


In~(a) rücken die grauen Marken der Zentren in gleichen Abständen vor, während
die Höhe genau so weit fällt, wie die Breite wächst. In~(b) achte man darauf,
wie früh die Kurve auf die klassische Gerade einschwenkt: schon nach wenigen
$\tau$ ist der Unterschied nicht mehr zu sehen. Ob man das Zerfließen je zu
Gesicht bekommt, entscheidet die Zeitskala $\tau$: für ein Elektron mit
$\Delta x(0) = 10^{-10}$\,m ist $\tau \sim 10^{-16}$\,s, für ein Staubkorn von
$1\,\mu$g mit $\Delta x(0) = 10^{-6}$\,m dagegen $\tau \sim 10^{13}$\,s, also
rund $300\,000$ Jahre.

Die interessanteren Unterschiede zwischen klassischer Mechanik und Quantenmechanik (und die wichtigeren Anwendungen der Quantenmechanik) finden sich allerdings bei Teilchen in äußeren Potentialen, oder in Wechselwirkung miteinander.

Der nächste Schritt dorthin ist bescheiden: ein Teilchen in einem \emph{endlich} großen Kasten -- unser erstes Beispiel dafür, was ,,Quanten`` in Quantenmechanik heißt.
%!TEX root = TheoretischePhysikC.tex

\section{Vorlesung: Teilchen im Kasten}
Wenn die Bewegungsfreiheit eines Teilchens durch äußere Wände eingeschränkt wird, äußert sich das zunächst einmal in der Einschränkung in den möglichen Werten der Ortsvariablen $x$. Wenn sich das Teilchen nur in einem Gebiet $G \subset \mathbb{R}^3$ bewegen kann, modifizieren wir den Hilbertraum erst einmal zu $\mathcal{H}\equiv L^2(G)$, d.h.\ Wellenfunktionen sind nur noch auf $G$ definiert und im
Skalarprodukt wird die Integration 
\begin{equation}
	\langle \Phi|\Psi\rangle \equiv \iiint_G \Phi^*(x,y,z)\Psi(x,y,z)\, dxdydz
\end{equation}
ebenfalls auf $G$ eingeschränkt. 

Den Hamiltonoperator versuchen wir erst einmal wieder als 
\begin{equation}
	H = -\frac{\hbar^2}{2m}\nabla^2
\end{equation}
anzusetzen. Das geht allerdings schief ($H$ wird nicht hermitesch sein). Man kann die Hermitizität von $H$ wieder herstellen, wenn man den Definitionsbereich von $H$ auf solche Funktionen einschränkt, die gewisse \emph{Randbedingungen} erfüllen. Diese Randbedingungen sind letztlich eine idealisierte Beschreibung der Wechselwirkung zwischen Teilchen und Wand. Eine häufig benutzte Randbedingung für den Hamiltonoperator ist die so genannte \emph{Dirichlet-Randbedingung}: im Definitionsbereich von $H$ werden nur Funktionen zugelassen, die auf dem Rand von $G$ verschwinden.

Um die Dynamik von $H$ zu lösen, ist es extrem hilfreich, die Eigenwerte und dazugehörigen Eigenvektoren von $H$ zu berechnen. (Dann ist $U_t=e^{-\frac{it}{\hbar}\widehat{H}}$ durch nur einen Multiplikationsoperator gegeben.) Wir beginnen mit der Berechnung der Eigenfunktionen für den einfachsten Fall: ein ,,eindimensionales`` Teilchen im Intervall $[0,L] \subset \mathbb{R}$. Das heißt, wir berechnen zuerst die Lösungen der Gleichung
\begin{equation}
	\widehat{H}|\Psi\rangle = E|\Psi\rangle
\end{equation}
für alle mögliche Eigenwerte $E$. Um weiter zu kommen werden wir diese Gleichung als eine Wellenfunktion-Gleichung umschreiben 
\begin{equation}
	\langle x|\widehat{H}|\Psi\rangle = -\frac{\hbar^2}{2m}\frac{d}{dx}\Psi(x) = E\Psi(x) = E\langle x|\Psi\rangle.
\end{equation}
Ohne weitere Begrenzungen ist der Hamiltonoperator nicht hermitesch. Wir müssen auch die folgende Randbedingung 
\begin{equation}
	\Psi(0) = \Psi(L) = 0
\end{equation} 
einführen. Dies ist eine lineare Differentialgleichung mit konstanten Koeffizienten, wird
also gelöst durch den Ansatz 
\begin{equation}
	\Psi(x) = e^{\mu x}, \quad\text{mit}\quad \mu\in \mathbb{C}.
\end{equation}
Einsetzen ergibt die Gleichung
\begin{equation}
	-\frac{\hbar^2}{2m} \mu^2 e^{\mu x} = E e^{\mu x},
\end{equation}
also
\begin{equation}
	\mu = \pm i \sqrt{\frac{2mE}{\hbar^2}}, 
\end{equation}
und damit zwei linear unabhängige Lösungen.

Die allgemeine Lösung der Differentialgleichung ist damit: 
\begin{equation}
	\Psi(x) = A e^{\mu x} + B e^{-\mu x},
\end{equation}
wobei $\mu$ eine der beiden Lösungen ist, und $A,B \in \mathbb{C}$. Die Randbedingung $\Psi(0)= 0$ erzwingt $B = −A$. Es bleibt also die Gleichung 
\begin{equation}
	\Psi(L) =A e^{\mu L} + B e^{-\mu L} = 0,
\end{equation}
oder äquivalent: 
\begin{equation}
	e^{2\mu L} = 1,
\end{equation}
aus der sich die möglichen Werte für $\mu$ (und damit für $E$) ergeben. Zunächst scheiden reelle $\mu$ aus,
also ist $E \ge 0$. Für imaginäre $\mu$ gibt es Lösungen, nämlich 
\begin{equation}
	2\mu L = 2\pi i n,
\end{equation}
wobei $n \in \mathbb{Z}$. Also ist $\mu = n\pi i/L$, und die Eigenfunktion $\Psi(x) = A \sin(n\pi x/L)$.
Offensichtlich geben $n$ und $−n$ die gleiche Eigenfunktion, und $n = 0$ scheidet aus (Dort ist $\Psi'' =0$, also $
\Psi(x) = ax+b$ linear, was sich nicht mit den Randbedingungen verträgt). Nach Normierung bekommen wir also die Eigenfunktionen und Eigenwerte
\begin{equation}
	\phi_n(x) = \sqrt{\frac{2}{L}}\sin\left(\frac{n\pi}{L} x\right)\quad \text{und}\quad E_n = \frac{\hbar^2 \pi^2}{2mL^2}n^2\quad \text{für $n=1,2,\ldots$.}
\end{equation}

Die $\phi_n$ bilden eine Basis von unseren Hilbertraum $\mathcal{H} = L^2([0,L])$ und damit einen unitären
Operator $U:\mathcal{H} \rightarrow \ell^2(\mathbb{N})$, nämlich $(U\Psi)(n) = \langle \phi_n|\Psi\rangle$, der $H$ diagonalisiert. Den Zeitentwicklungsoperator $U_t = e^{−itH/\hbar}$ bekommt man also durch Entwickeln nach der Basis $\phi_n$ (=Anwendung von $U$), triviale Anwendung von $U_t$ auf die Basisvektoren (=Multiplikation mit einer Funktion von $n$) und Aufsummation der Reihe mit den neuen Komponenten (=Anwendung von $U^\dag$):
\begin{equation}
	U_t |\Psi\rangle = \sum_{n=1}^{\infty} e^{-\frac{it}{\hbar} E_n} |\phi_n\rangle \langle \phi_n|\Psi\rangle. 
\end{equation}
Als Merkregel in der Dirac'schen Bra-Ket Schreibweise gilt also wiederum
\begin{equation}
	\eqbox{U_t = \sum_{n=1}^{\infty} e^{-\frac{it}{\hbar} E_n} |\phi_n\rangle \langle \phi_n|.}
\end{equation}

Für ein dreidimensionales Teilchen in einer quaderförmigen Kiste ist damit die
wesentliche Arbeit schon geleistet. Wir betrachten also das Gebiet
\begin{equation}
	G = \left\{(x,y,z)\,\middle| \,  0\le x \le L_x, 0\le y \le L_y, 0\le z \le L_z  \right\},
\end{equation}
wobei $L_x$, $L_y$ und $L_z$ die Kantenlängen des Quaders sind. Der Hamiltonoperator ist
nun 
\begin{equation}
	\widehat{H} = -\frac{\hbar^2}{2m}\nabla^2,
\end{equation}
wieder mit der Randbedingung, dass Funktionen im Definitionsbereich von $H$ auf dem Rand des Gebietes verschwinden. Zur Lösung der Eigenwertgleichung macht man einen \emph{Separationsansatz}
\begin{equation}
	\Psi(x,y,z) = \Phi_x(x)\Phi_y(y)\Phi_z(z)
\end{equation}
und erfährt, dass die Funktion $\Phi_x''/\Phi_x$ nicht mehr von $x$ abhängt (analog für $\Phi_y$ und $\Phi_z$), also eine Eigenfunktion für das eben diskutierte Problem eines eindimensionalen Intervalls sein muss. Damit bekommen wir also Eigenfunktionen und Eigenwerte
\begin{align}
	\phi_{n_x,n_y,n_z}(x,y,z) &= \sqrt{\frac{8}{L_xL_yL_z}} \sin\left(\frac{n_x\pi}{L_x} x\right)\sin\left(\frac{n_y\pi}{L_y} y\right)\sin\left(\frac{n_z\pi}{L_z} z\right) \\
	E_{n_x,n_y,n_z}&=\frac{\hbar^2\pi^2}{2m} \left(\frac{n_x^2}{L_x^2} + \frac{n_y^2}{L_y^2} + \frac{n_z^2}{L_z^2} \right).
\end{align}
Was kann man über anders geformte Kisten sagen? Ein Separationsansatz funktioniert nur in wenigen Fällen, in denen man Koordinaten finden kann, in denen sich die Form der Kiste leicht beschreiben lässt, und in der der
Laplace-Operator außerdem in eine Summe von zweiter Ableitungen zerfällt. Außer dem Quader kann man so noch die Kugel behandeln und auch noch Ellipsoide. Für allgemeine Gebiete $G$ ist die Bestimmung der Eigenfunktionen
und Eigenwerte des Laplace-Operators nicht geschlossen möglich, aber ein Problem, das bei verschiedenen Wellenphänomenen in der Physik eine Rolle spielt (Hohlraum-Resonanz).
%!TEX root = TheoretischePhysikC.tex

\section{Vorlesung: Harmonischer Oszillator}
Der Energie eines \emph{klassischen} eindimensionalen Oszillators ist gegeben durch
\begin{equation}
	E = \frac{1}{2m}p^2 + \frac{m\omega^2}{2}x^2,
\end{equation}
wobei $\omega$ die Frequenz des klassischen Oszillators mit dem gleichen Potential
ist.

Nach unserer Quantisierungstabelle lautet der Hamiltonoperator des quantumechanischen Oszillators:
\begin{equation}
	\widehat{H} = \frac{1}{2m}\widehat{p}^2 + \frac{m\omega^2}{2}\widehat{x}^2.
\end{equation}	


Der Oszillator tritt in vielen Kontexten der Physik auf: nämlich immer dann, wenn man kleine Bewegungen um eine Gleichgewichtslage näherungsweise untersucht. Ein freies Quantenfeld lässt sich als System von unendlich
vielen Oszillatoren analysieren. Es lohnt sich jedenfalls, dieses System genau kennenzulernen. Deshalb werden wir auch mehrere Betrachtungsweisen untersuchen.

Die direkteste Methode, das Spektrum zu erhalten, basiert auf der Beobachtung, dass sich $H$ (im wesentlichen) als Quadrat einer komplexen Linearkombination von $\widehat{p}$ und $\widehat{x}$ schreiben lässt. Wir setzen
\begin{equation}
	\widehat{a} = \sqrt{\frac{m\omega}{2\hbar}}\widehat{x} + i \sqrt{\frac{1}{2m\omega\hbar}}\widehat{p},
\end{equation}
sodass
\begin{align}
	\widehat{a} \widehat{a}^\dag &= \frac{m\omega}{2\hbar}\widehat{x}^2 + \frac{1}{2m\omega\hbar}\widehat{p}^2 + \frac{i}{2\hbar} (\widehat{p}\widehat{x} - \widehat{x}\widehat{p}) \\
	\widehat{a}^\dag \widehat{a} &= \frac{m\omega}{2\hbar}\widehat{x}^2 + \frac{1}{2m\omega\hbar}\widehat{p}^2 - \frac{i}{2\hbar} (\widehat{p}\widehat{x} - \widehat{x}\widehat{p}).
\end{align}
Unter Verwendung der kanonischen Vertauschungsrelationen $[\widehat{p},\widehat{x}] = -i\hbar\mathbb{I}$ ist damit
\begin{equation}\label{eq:ccra}
	\eqbox{ [\widehat{a}, \widehat{a}^\dag] = \mathbb{I}}
\end{equation}
und
\begin{equation}\label{eq:qho}
	\eqbox{\widehat{H} = \hbar \omega \left(\widehat{a}^\dag \widehat{a} + \tfrac{1}{2}\mathbb{I}\right).}
\end{equation}
Zunächst ist aus dieser Form von $H$ klar, dass $H \ge \hbar\omega/2 \mathbb{I}$. Der Oszillator ist
also niemals ,,in Ruhe``: $x=0$ und $p = 0$ können wegen der Unschärferelation nicht beide beliebig genau präpariert werden. 

Ist die ,,Nullpunktsenergie`` $\hbar \omega/2$ selbst ein Eigenwert? Der zugehörige Eigenvektor $|\phi_0\rangle$ müsste die Gleichung 
\begin{equation}
	\widehat{a}^\dag \widehat{a}|\phi_0\rangle = 0,
\end{equation}
also $\|\widehat{a}|\phi_0\rangle\|=0$, also $\widehat{a}|\phi_0\rangle = 0$ erfüllen. In der Ortsdarstellung ist
dies die gewöhnliche Differentialgleichung $\langle x|\widehat{a}|\phi_0\rangle = 0$:
\begin{equation}
	\sqrt{\frac{m\omega}{2\hbar}}x\phi_0(x) + i\sqrt{\frac{1}{2m\omega\hbar}} \frac{\hbar}{i}\frac{d}{dx} \phi_0(x) = 0,
\end{equation}
also 
\begin{equation}
	\frac{d\phi_0}{dx}  = -\frac{m\omega}{\hbar} \phi_0,
\end{equation}
mit der Lösung
\begin{equation}
	\phi_0(x) = Ne^{-\frac{m\omega}{\hbar} \frac{x^2}{2}},
\end{equation}
wobei $N$ die Normierung 
\begin{equation}
	N = \left(\frac{m\omega}{\pi\hbar}\right)^{\frac{1}{4}}
\end{equation}
ist.


In der Ableitung der Unschärferelation hatten wir genau wie hier die \emph{Positivität} das Operators $\widehat{a}^\dag \widehat{a}$ benutzt. Der Fall der Gleichheit führt auch dort auf die Bedingung $\widehat{a}|\Psi\rangle = 0$. Der Oszillator-Grundzustand ist also genau ein Zustand minimaler Unschärfe.

Für die weitere Rechnung brauchen wir die Relationen
\begin{align}
	[\widehat{H},\widehat{a}^\dag] &= \hbar \omega \widehat{a}^\dag \\
	[\widehat{H},\widehat{a}] &= -\hbar \omega \widehat{a}, 
\end{align}
die aus (\ref{eq:ccra}) und (\ref{eq:qho}) folgen. Wenn nun $\widehat{H}|\Psi\rangle = E|\Psi\rangle$ ein Eigenvektor ist, dann ist 
\begin{equation}
	\widehat{H}\widehat{a}^\dag|\Psi\rangle = (E  + \hbar \omega)\widehat{a}^\dag|\Psi\rangle
\end{equation}
Also erfüllt $\widehat{a}^\dag|\Psi\rangle$ die Eigenwertgleichung zum Eigenwert $E + \hbar \omega$. Wegen 
\begin{equation}
	\|\widehat{a}^\dag|\Psi\rangle\|^2 = \langle \Psi|\widehat{a}\widehat{a}^\dag|\Psi\rangle = \langle \Psi|\widehat{a}^\dag\widehat{a}+\mathbb{I}|\Psi\rangle \ge \langle \Psi|\Psi\rangle \ge 0
\end{equation}
kann $\widehat{a}^\dag|\Psi\rangle$ auch nicht Null sein: wir bekommen
also wirklich einen neuen Eigenvektor. Startend von $|\phi_0\rangle$ erhalten wir so die
Eigenvektoren 
\begin{equation}\label{eq:qhoeigenstates}
	\eqbox{|\phi_n\rangle = \frac{(\widehat{a}^\dag)^n|\phi_0\rangle}{\|(\widehat{a}^\dag)^n|\phi_0\rangle \|}}
\end{equation}
und Eigenwerte
\begin{equation}
	\eqbox{E_n = \hbar\omega(n+\tfrac{1}{2}).}
\end{equation}

Wegen dieses Verfahrens heißen übrigens die Operatoren $\widehat{a}$ und $\widehat{a}^\dag$ auch \emph{Leiteroperatoren}, bzw. ,,Absteigeoperator`` und ,,Aufsteigeoperator``. Den Absteigeoperator $\widehat{a}$ verwenden wir, um zu zeigen, dass es auch keine weiteren Eigenwerte gibt. Denn wenn $|\Psi\rangle$ ein Eigenvektor zum Eigenwert $E$ ist, erfüllt $(\widehat{a})^n|\Psi\rangle$ die Eigenwertgleichung zum Eigenwert $E − n\hbar \omega$. Für genügend große $n$ wird dieser Ausdruck negativ, kann also kein Eigenwert des positiven Operators
$H$ sein. Für große $n$ muss also $(\widehat{a})^n|\Psi\rangle=0$ sein. Sei $n$ die kleinste Zahl, für die dies gilt. Dann ist $|\Phi\rangle = (\widehat{a})^{n-1}|\Psi\rangle$ eine von Null verschiedene Lösung
der Gleichung $\widehat{a}|\Phi\rangle = 0$, also nach unserer obigen Rechnung proportional zu $|\phi_0\rangle$. Nach der Vertauschungsrelation für $\widehat{a}$ und $\widehat{a}^\dag$ ist $|\Psi\rangle$ dann ein Vielfaches von
$(\widehat{a}^\dag)^n|\Phi_0\rangle$. Die angegebene vollständige Liste von Eigenfunktionen ist also eine Basis des Hilbertraums.

Die Norm $\|(\widehat{a}^\dag)^n|\Phi_0\rangle\|$ lässt sich ebenfalls allein aus der Vertauschungsrelation
berechnen. Es ist
\begin{equation}
	\|\widehat{a}^\dag|\phi_n\rangle\|^2 = \langle\phi_n|\widehat{a}\widehat{a}^\dag|\phi_n\rangle =  \langle\phi_n|\left(\frac{1}{\hbar\omega}\widehat{H}+\frac12\mathbb{I}\right)|\phi_n\rangle=(n+1)\|\phi_n\|^2 = (n+1).
\end{equation}
Damit ist 
\begin{equation}
	\widehat{a}^\dag|\phi_n\rangle = \sqrt{n+1}|\phi_{n+1}\rangle,
\end{equation}
und (durch Induktion)
\begin{equation}
	|\phi_n\rangle = \frac{1}{\sqrt{n!}}(\widehat{a}^\dag)^n|\Phi_0\rangle.
\end{equation}

Da die $|\phi_n\rangle$ eine Basis des Hilbertraumes bilden, lassen sich alle Operatoren durch ihre Matrixelemente bezüglich dieser Basis ausdrücken. Das ist immer dann sinnvoll, wenn zeitabhängige Größen untersucht werden sollen, da ja dann die Zeitentwicklungsoperatoren Multiplikationsoperatoren sind. Wir bekommen
(mit der Matrixnotation $[\widehat{C}]_{nm} = \langle\phi_n|\widehat{C}|\phi_m\rangle$)
\begin{equation}
	[\widehat{a}^\dag]_{nm} = \sqrt{n}\delta_{n,m+1} \quad \text{und} \quad [\widehat{a}]_{nm} = \sqrt{m}\delta_{n+1,m}.
\end{equation}
Mit 
\begin{equation}
	\widehat{p} = \frac{\sqrt{2m\omega\hbar}}{2i}(\widehat{a}-\widehat{a}^\dag) 
\end{equation}
bekommt man daraus unmittelbar die Matrix für $\widehat{p}$ und entsprechend die Matrix für $\widehat{x}$. Hingemalt sind es:
\begin{equation}
	\widehat{p} = \frac{\sqrt{2m\omega\hbar}}{2i}
	\begin{pmatrix}
		0 & \sqrt{1} & 0 & 0 & 0 & 0 & \hdots \\
		-\sqrt{1} & 0 & \sqrt{2} & 0 & 0 & 0 & \\
		0 & -\sqrt{2} & 0 & \sqrt{3} & 0 & 0 & \\
		0 & 0 & -\sqrt{3} & 0 & \sqrt{4} & 0 & \\
		0 & 0 & 0 & -\sqrt{4} & 0 & \sqrt{5} & \\
		0 & 0 & 0 & 0 & -\sqrt{5} & 0 & \\
		\vdots & & & & & & \ddots \\
	\end{pmatrix}
\end{equation}
und
\begin{equation}
	\widehat{x} = \sqrt{\frac{\hbar}{2m\omega}}
	\begin{pmatrix}
		0 & \sqrt{1} & 0 & 0 & 0 & 0 & \hdots \\
		\sqrt{1} & 0 & \sqrt{2} & 0 & 0 & 0 & \\
		0 & \sqrt{2} & 0 & \sqrt{3} & 0 & 0 & \\
		0 & 0 & \sqrt{3} & 0 & \sqrt{4} & 0 & \\
		0 & 0 & 0 & \sqrt{4} & 0 & \sqrt{5} & \\
		0 & 0 & 0 & 0 & \sqrt{5} & 0 & \\
		\vdots & & & & & & \ddots \\
	\end{pmatrix}.
\end{equation}

Ein wichtiger Vorteil der oben erwähnten \emph{algebraischen Methode} der Leiteroperatoren ist,
dass die Ergebnisse sofort in jeder Darstellung des Hilbertraums zur Verfügung stehen. Zum Beispiel liefert die Formel (\ref{eq:qhoeigenstates}) die Wellenfunktionen/Eigenfunktionen $\phi_n(x) \in L^2(\mathbb{R})$ in der Ortsdarstellung, sobald $\widehat{a}$ konkret in dieser Darstellung eingesetzt wird. Durch Induktion sieht man, dass 
\begin{equation}
	\phi_n(x) = c_n \mathrm{H}_n(\lambda x)  e^{-\frac{\lambda^2 x^2}{2}}
\end{equation}
mit
\begin{equation}
	\lambda = \sqrt{\frac{m\omega}{\hbar}}\quad \text{und} \quad c_n = \sqrt{\frac{\lambda}{2^n n!\sqrt{\pi}}},
\end{equation}
wobei
$\mathrm{H}_n(\lambda x)$ die sogenannten \emph{Hermite-Polynome} sind. Die ersten $3$ Hermite-Polynome
lauten
\begin{align}
	\mathrm{H}_0(z) &= 1,\\
	\mathrm{H}_1(z) &= 2z,\\
	\mathrm{H}_2(z) &= 4z^2-2,
\end{align}
usw.

Mit der Diagonalisierung des Hamiltonoperators haben wir natürlich auch detaillierte Informationen über die Zeitentwicklung des Oszillators gewonnen. Wie der klassische Oszillator ist auch der quantenmechanische periodisch mit der Periode $T = 2\pi/\omega$. Das folgt allein schon aus der Liste der Eigenwerte: die Eigenwerte des Zeitentwicklungsoperators $U_T = e^{-\frac{iT}{\hbar} \widehat{H}}$ sind ja nach Funktionalkalkül
\begin{equation}
	e^{-\frac{iT}{\hbar} \hbar\omega(n+\frac12)} = e^{-2\pi i (n+\frac12)} = -1,
\end{equation}
also unabhängig von $n$, und $U_T = −\mathbb{I}$ ist ein Vielfaches des $\mathbb{I}$-Operators.

Damit gilt für die Zeitabhängigkeit aller Erwartungswerte, nach Postulat 5 also für die Funktionen der Form $f(t) = \tr(U_t \rho U_t^\dag A)$:
\begin{equation}
	f(t+T) = \tr(U_tU_T \rho U_T^\dag U_t^\dag A) =f(t).
\end{equation}	
Die Analogie mit der klassischen Entwicklung geht aber noch viel weiter. Am einfachsten sieht man das hier, indem man die Zeitabhängigkeit des Leiteroperators $\widehat{a}$ berechnet. Es gilt 
\begin{equation}
	\frac{d}{dt} U_t^\dag \widehat{a} U_t = U_t^\dag \left(\frac{i}{\hbar} [\widehat{H},\widehat{a}] \right) U_t = -i\omega U_t^\dag \widehat{a} U_t
\end{equation}
also
\begin{equation}
	U_t^\dag \widehat{a} U_t = e^{-i\omega t}\widehat{a}.
\end{equation}
Dabei findet man, dass die quantenmechanischen Orts-und Impulsoperatoren genau den klassischen Bewegungsgleichungen folgen:
\begin{equation}
	\begin{split}
		\widehat{p}(t) &= U_t^\dag \widehat{p} U_t  = \cos(\omega t) \widehat{p}  - m\omega\sin(\omega t) \widehat{x} \\
		\widehat{x}(t) &= U_t^\dag \widehat{x} U_t  = \frac{1}{m\omega}\sin(\omega t) \widehat{p}  + \cos(\omega t) \widehat{x}
	\end{split}
\end{equation}
%!TEX root = TheoretischePhysikC.tex

\section{Vorlesung: Endlich hoher Potentialtopf}
Wir betrachten nun ein Teilchen in einem endlich hohen Potentialtopf mit Tiefe $V_0$ und Breite $L$:
\begin{equation}
	V(x) = \begin{cases}
		-V_0, \quad &|x|\le \frac{L}{2}\\
		0, \quad &|x| > \frac{L}{2}.
	\end{cases}
\end{equation}
Wir werden uns mit den folgenden Fragen beschäftigen: (a) Gibt es gebundene Zustände (d.h.\ Energieeigenvektoren mit Eigenenergie $E\le 0$)? (b) Welche Unterschiede zum unendlich hohen Potentialtopf sind zu beobachten?

\begin{intuition}
	Schießen wir ein Wellenpaket von außen auf den Topf, so fließt später ein Teil hindurch, ein Teil wird reflektiert, und ein Rest bleibt im Topf sitzen. Ein Streuprozess also, im Prinzip derselbe wie am LHC, wo die gebundenen Zustände neue Teilchen sind.
\end{intuition}

Die volle Dynamik $U_t$ mischt gebundene und freie Zustände und ist zu kompliziert; ein gebundener Zustand aber bewegt sich zeitlich nicht und ist damit ein Eigenvektor von $\widehat{H}$.

Es ist hilfreich, um die Lösung der Eigenvektorgleichung  
\begin{equation}
	\widehat{H}|\Psi\rangle = E|\Psi\rangle
\end{equation}
zu erhalten, die Konstanten
\begin{equation}
	\kappa = \frac{\sqrt{2m|E|}}{\hbar}\quad \text{und}\quad q = \frac{\sqrt{2m(E+V_0)}}{\hbar}
\end{equation}
zu definieren. 

Die dazugehörigen Wellenfunktionen erhalten wir mittels
\begin{equation}
	\langle x|\widehat{H}|\Psi\rangle = \left(-\frac{\hbar^2}{2m}\frac{d^2}{dx^2} + V(x)\right)\Psi(x) = E\Psi(x) = E\langle x|\Psi\rangle.
\end{equation}
Nach Umformung und Einsetzen erhalten wir
\begin{equation}
	\begin{cases}
		\frac{d^2}{dx^2}\Psi(x) = -q^2 \Psi(x), \quad &\text{für $|x|<\frac{L}{2}$}\\
		\frac{d^2}{dx^2}\Psi(x) = \kappa^2 \Psi(x), \quad &\text{für $|x|>\frac{L}{2}$}.
	\end{cases}
\end{equation}
Wegen Symmetrie können die Eigenfunktionen gerade oder ungerade gewählt werden. Gerade heißt: ersetzt man $x$ durch $-x$, kommt dieselbe Funktion zurück. Wir beschränken uns auf $E<0$; die unendlich vielen Lösungen mit $E>0$ sind \emph{Streuzustände}, außerhalb ebene Wellen, und bleiben den Übungen überlassen. Die geraden Lösungen lauten
\begin{equation}
	\Psi(x) = \begin{cases}
		A_1 e^{\kappa x}, \quad & x< -\frac{L}{2} \\
		A_2 \cos(qx), \quad & |x| \le \frac{L}{2} \\
		A_1 e^{-\kappa x}, \quad & x> \frac{L}{2}
	\end{cases}
\end{equation}
und die ungeraden Lösungen
\begin{equation}
	\Psi(x) = \begin{cases}
		-B_1 e^{\kappa x}, \quad & x< -\frac{L}{2} \\
		B_2 \sin(qx), \quad & |x| \le \frac{L}{2} \\
		B_1 e^{-\kappa x}, \quad & x> \frac{L}{2}.
	\end{cases}
\end{equation}
Außerhalb fallen sie exponentiell ab, verschwinden aber nicht: das Teilchen ist mit kleiner Wahrscheinlichkeit im klassisch verbotenen Bereich zu finden, und das in einem zeitlich unveränderlichen Zustand.

Alle Wellenfunktionen werden als \emph{stetig} und \emph{differenzierbar} angenommen. Dies liefert Anschlussbedingungen bei $x = \pm L/2$. Die Begründung ist physikalisch, nicht mathematisch: ein Sprung in $\Psi$ bedeutet eine unendliche erste Ableitung, und die hängt am Impulsoperator -- das Teilchen hätte dort unendlichen Impuls und unendliche kinetische Energie. Unmöglich ist das nicht -- es gibt solche Systeme --, aber wir verlangen endliche kinetische Energie. Genau darin steckt die Quantisierung. Für die geraden Lösungen erhalten wir wegen
\begin{enumerate}
	\item Stetigkeit: $A_1 e^{-\frac{\kappa L}{2}} = A_2 \cos\left(\frac{qL}{2}\right)$,
	\item Differenzierbarkeit: $-A_1\kappa e^{-\frac{\kappa L}{2}} = -A_2 q\sin\left(\frac{qL}{2}\right)$.
\end{enumerate}
Dividieren wir die zweite Bedingung durch die erste, erhalten wir
\begin{equation}\label{eq:symmcond}
	\kappa = q\tan\left(\frac{qL}{2}\right).
\end{equation}
Für die ungeraden Lösungen erfolgt die Bedingung  
\begin{enumerate}
	\item Stetigkeit: $B_1 e^{-\frac{\kappa L}{2}} = B_2 \sin\left(\frac{qL}{2}\right)$,
	\item Differenzierbarkeit: $-B_1\kappa e^{-\frac{\kappa L}{2}} = B_2 q\cos\left(\frac{qL}{2}\right)$.
\end{enumerate}
Dabei wird analog die Folgerung erhalten werden:
\begin{equation}\label{eq:unsymm}
	\kappa = -q\cot\left(\frac{qL}{2}\right).
\end{equation}

Wir betrachten nun die Bedingung (\ref{eq:symmcond}) für gerade Lösungen: Die Quadratur beider Seiten ergibt
\begin{equation}
	\kappa^2 = q^2\tan^2\left(\frac{qL}{2}\right).
\end{equation}
Nach der Definition $\kappa$ erhalten wir
\begin{equation}
	\kappa^2 = \frac{2m|E|}{\hbar^2} = \frac{2mV_0}{\hbar^2} - \frac{2m(V_0+E)}{\hbar^2} = \frac{2mV_0}{\hbar^2} - q^2.
\end{equation}
Führen wir die Konstanten
\begin{equation}
	u\equiv \frac{qL}{2}\quad \text{und}\quad \alpha^2 \equiv \frac{m V_0 L^2}{2\hbar^2}
\end{equation}
ein, ergibt sich
\begin{equation}
	\alpha^2-u^2 = u^2\tan^2(u).
\end{equation}
Nach Umformung:
\begin{equation}
	\frac{\alpha^2}{u^2} = 1+\tan^2(u) = \frac{1}{\cos^2(u)}.
\end{equation}
sodass
\begin{equation}
	|\cos(u)| = \frac{u}{\alpha}.
\end{equation}
Die Lösung dieser \emph{transzendenten} Gleichung erfolgt graphisch. Wegen der Bedingungung $\kappa>0$, $q>0$ und $\kappa = q\tan(u)$ muss $\tan(u)>0$ gelten, sodass nur Lösungen in $0<u<\frac{\pi}{2}$ und $\pi < u < \frac{3\pi}{2}$ korrekt sind. Zusammengefasst: Es gibt mindestens einen gebundenen Zustand und insgesamt endlich viele Zustände mit gerader Wellenfunktion.

\begin{merksatz}
	Auf $\alpha$ allein kommt es an: klein bei leichtem Teilchen, flachem oder schmalem Topf, groß bei schwerem Teilchen, tiefem oder breitem Topf. Je größer $\alpha$, desto mehr gebundene Zustände -- einen geraden gibt es immer.
\end{merksatz}


Die Bedingung (\ref{eq:unsymm}) für ungerade Lösungen wird analog behandelt: Die Quadratur beider Seiten ergibt in diesem Fall
\begin{equation}
	\kappa^2 = q^2\cot^2\left(u\right).
\end{equation}
Daher ergibt sich
\begin{equation}
	\alpha^2-u^2 = u^2\cot^2(u).
\end{equation}
Nach Umformung:
\begin{equation}
	\frac{\alpha^2}{u^2} = 1+\cot^2(u) = \frac{1}{\sin^2(u)}.
\end{equation}
sodass
\begin{equation}
	|\sin(u)| = \frac{u}{\alpha}.
\end{equation}
Lösung dieser Gleichung erfolgt auch graphisch: Wegen $\kappa>0$, $q>0$ und $\kappa = -q\cot(u)$ muss $\cot(u)<0$ gelten, sodass nur Lösungen in $\frac{\pi}{2}<u< \pi$ und $\frac{3\pi}{2} < u < 2\pi$ in Frage kommen. Anders als $|\cos(u)|$ startet $|\sin(u)|$ im Ursprung bei Null; für kleines $\alpha$ kann es deshalb gar keine ungerade Lösung geben. Der Schnittpunkt bei $u=0$ zählt nicht: dort wäre $q=0$ und die Sinuslösung identisch Null -- die typische Falle beim graphischen Lösen.

% fig_vl13_topf.tex — Endlich hoher Potentialtopf: die gebundenen
% Zustaende mit ihren exponentiellen Auslaeufern, und die graphische
% Loesung, die abzaehlt, wie viele es davon gibt.
%
% GEOMETRY FROM THE FORMULA.  Nothing here is drawn by eye.  With the
% abbreviations of the lecture,
%
%       u = qL/2 ,   alpha^2 = m V_0 L^2 / (2 hbar^2) ,
%
% the matching conditions kappa = q tan(u) and kappa = -q cot(u) reduce
% to the two transcendental equations
%
%       |cos u| = u/alpha   (even)      |sin u| = u/alpha   (odd),
%
% each restricted to the branch on which the sign condition holds:
% tan u > 0 on (0,pi/2) and (pi,3pi/2), cot u < 0 on (pi/2,pi) and
% (3pi/2,2pi).  Those four intervals are exactly the four DESCENDING
% half-arches of |cos u| and |sin u|, alternating even/odd — which is
% why panel (b) can be read off as "count the arches the straight line
% u/alpha still reaches".  The ascending halves are the branches the
% sign conditions throw away, the crossing at u = 0 among them.
%
% For alpha = 4 there are ceil(2 alpha/pi) = 3 bound states, with roots
%
%       u_1 = 1.25235   u_2 = 2.47458   u_3 = 3.59530 ,
%
% and with v = kappa L/2 = sqrt(alpha^2 - u^2),
%
%       v_1 = 3.79890   v_2 = 3.14269   v_3 = 1.75322 .
%
% Panel (a) then follows without any further choice:
%
%       E_n / V_0 = u_n^2/alpha^2 - 1  = -0.90198, -0.61728, -0.19211 ,
%       psi_n(x)  = cos(2 u_n x/L)  or  sin(2 u_n x/L)      inside,
%                   (+-) (u_n/alpha) exp(-2 v_n (|x|/L - 1/2))  outside,
%
% the amplitude at the wall being u_n/alpha = |cos u_n| resp. |sin u_n|
% — the transcendental equation itself.  So the tail of the topmost
% state starts almost three times as high AND decays 2.2x more slowly
% than the ground state's (1/kappa_3 = 0.285 L against 1/kappa_1 =
% 0.132 L): the weaker the binding, the further the particle leaks out.
% That comparison is what the picture is for; prose cannot do it.
%
% DESIGN DECISIONS
%   * colour encodes the quantum number n and nothing else, and it is
%     the SAME code in both panels, so the arch a state comes from in
%     (b) can be found again as a wavefunction in (a).  Entries 1, 2, 4,
%     5 of the ordered `tpclevels` ramp (navy, teal, orange, red);
%     entry 3, munsell, is skipped on purpose — beside cgblue it is the
%     same teal to the eye, and arches 2 and 3 in (b) are neighbours;
%   * the fourth colour appears only in (b): arch 4 is the state that
%     alpha = 4 does not yet support.  It is drawn, not omitted, because
%     the empty arch is the point — increase alpha, the line flattens,
%     and the fourth state appears;
%   * everything that is reference rather than data is neutral grey: the
%     energy levels, the -V_0 guide, the straight line u/alpha, and the
%     dotted ascending half-arches;
%   * classically forbidden region hatched, not filled, so the leakage
%     of psi into it survives greyscale printing;
%   * amplitude 0.12 V_0 in (a) is the largest that keeps psi_2 clear of
%     E_1, the tightest gap (E_2 - E_1 = 0.285 V_0).

\begin{center}
\begin{tikzpicture}

% Roots of |cos u| = u/alpha and |sin u| = u/alpha at alpha = 4, and the
% quantities that follow from them.  Local to this tikzpicture.
\def\uA{1.25235}\def\vA{3.79890}\def\EA{-0.90198}\def\bA{0.31309}
\def\uB{2.47458}\def\vB{3.14269}\def\EB{-0.61728}\def\bB{0.61864}
\def\uC{3.59530}\def\vC{1.75322}\def\EC{-0.19211}\def\bC{-0.89883}
\def\amp{0.12}

% ─────────────────────────────────────────────────────────────
%  (a)  The three bound states, each on its own level
% ─────────────────────────────────────────────────────────────
\begin{axis}[
    name=well,
    tpcplot,
    width=7.6cm, height=6.6cm,
    axis lines=none,
    xmin=-1.30, xmax=1.30,
    ymin=-1.24, ymax=0.24,
    clip=false,
    title={(a) \tpcifenglish{Bound states, $\alpha=4$}%
                            {Gebundene Zust\"ande, $\alpha=4$}},
    domain=-1.24:1.24, samples=401,
]

  % Classically forbidden for every bound state: outside the well V = 0
  % lies above E < 0.  Hatched, so psi can be drawn straight over it.
  \addplot[forbidden, draw=none, domain=-1.24:-0.5, samples=2] {-1.24} \closedcycle;
  \addplot[forbidden, draw=none, domain=0.5:1.24,   samples=2] {-1.24} \closedcycle;

  % The well itself
  \draw[potential] (axis cs:-1.24,0)  -- (axis cs:-0.5,0);
  \draw[potential] (axis cs:-0.5,0)   -- (axis cs:-0.5,-1);
  \draw[potential] (axis cs:-0.5,-1)  -- (axis cs:0.5,-1);
  \draw[potential] (axis cs:0.5,-1)   -- (axis cs:0.5,0);
  \draw[potential, -{Stealth[length=5pt]}] (axis cs:0.5,0) -- (axis cs:1.30,0);

  % Energy levels E_n/V_0 = u_n^2/alpha^2 - 1, and the floor -V_0.
  % Reference rules, hence grey: colour is already carrying n.
  \addplot[levellinegrey, domain=-1.24:1.24, samples=2] {\EA};
  \addplot[levellinegrey, domain=-1.24:1.24, samples=2] {\EB};
  \addplot[levellinegrey, domain=-1.24:1.24, samples=2] {\EC};
  \addplot[levellinegrey, domain=-1.24:-0.5, samples=2] {-1};

  % Eigenfunctions about their levels.  Inside: cos/sin(2 u_n x/L).
  % Outside: amplitude at the wall +-u_n/alpha, times exp(-2 v_n (|x|/L - 1/2)).
  \addplot[wavefn, spacecadet] {\EA + \amp*(
       (abs(x)<0.5)*cos(deg(2*\uA*x))
     + (abs(x)>=0.5)*\bA*exp(-2*\vA*(abs(x)-0.5)) )};
  \addplot[wavefn, cgblue]     {\EB + \amp*(
       (abs(x)<0.5)*sin(deg(2*\uB*x))
     + (abs(x)>=0.5)*((x>0)*2-1)*\bB*exp(-2*\vB*(abs(x)-0.5)) )};
  \addplot[wavefn, birrenorange] {\EC + \amp*(
       (abs(x)<0.5)*cos(deg(2*\uC*x))
     + (abs(x)>=0.5)*\bC*exp(-2*\vC*(abs(x)-0.5)) )};

  % Level labels
  \node[lbl, anchor=east] at (axis cs:-1.30,\EA) {$E_{1}$};
  \node[lbl, anchor=east] at (axis cs:-1.30,\EB) {$E_{2}$};
  \node[lbl, anchor=east] at (axis cs:-1.30,\EC) {$E_{3}$};
  \node[lbl, anchor=east] at (axis cs:-1.30,-1)  {$-V_{0}$};

  % Axis furniture
  \node[lbl, anchor=west]       at (axis cs:1.32,0)     {$x$};
  \node[lbl, anchor=south east] at (axis cs:1.24,0.03)  {$V=0$};
  \node[lbl, anchor=south]      at (axis cs:-0.5,0.04)  {$-\tfrac{L}{2}$};
  \node[lbl, anchor=south]      at (axis cs:0.5,0.04)   {$\tfrac{L}{2}$};

  % Well width
  \draw[{Stealth[length=4pt]}-{Stealth[length=4pt]}, birrengray, thin]
    (axis cs:-0.5,-1.14) -- (axis cs:0.5,-1.14)
    node[midway, lbl, fill=white, inner sep=1pt] {$L$};

  % Penetration depth 1/kappa_n = L/(2 v_n): 0.132 L for the ground state
  % against 0.285 L for the topmost one.  White underlay so the measure
  % reads over the hatching.
  \draw[white, line width=2.6pt] (axis cs:0.5,-0.96) -- (axis cs:0.632,-0.96);
  \draw[white, line width=2.6pt] (axis cs:0.5,-0.40) -- (axis cs:0.785,-0.40);
  \draw[{Stealth[length=3pt]}-{Stealth[length=3pt]}, birrengray!60!black, thin]
    (axis cs:0.5,-0.96) -- (axis cs:0.632,-0.96);
  \draw[{Stealth[length=3pt]}-{Stealth[length=3pt]}, birrengray!60!black, thin]
    (axis cs:0.5,-0.40) -- (axis cs:0.785,-0.40);
  \node[lbl, font=\scriptsize, anchor=west, fill=white, inner sep=1pt]
    at (axis cs:0.66,-0.96) {$1/\kappa_{1}$};
  \node[lbl, font=\scriptsize, anchor=west, fill=white, inner sep=1pt]
    at (axis cs:0.82,-0.40) {$1/\kappa_{3}$};

  % The hatched zone, named once
  \node[lbl, text=birrengray, font=\scriptsize, align=center,
        fill=white, inner sep=1pt]
    at (axis cs:-0.87,-1.13)
    {\tpcifenglish{classically\\forbidden}{klassisch\\verboten}};
\end{axis}

% ─────────────────────────────────────────────────────────────
%  (b)  |cos u| = u/alpha and |sin u| = u/alpha: counting the states
% ─────────────────────────────────────────────────────────────
\begin{axis}[
    tpcplot,
    at={(well.south east)}, anchor=south west, xshift=12mm,
    width=5.9cm, height=6.6cm,
    xlabel={$u=\tfrac{qL}{2}$},
    xmin=0, xmax=5.45, ymin=0, ymax=1.55,
    xtick={0,1.5707963,3.1415927,4.712389},
    xticklabels={$0$,$\tfrac{\pi}{2}$,$\pi$,$\tfrac{3\pi}{2}$},
    ytick={0,0.5,1},
    grid=major,
    samples=120,
    title={(b) \tpcifenglish{Graphical solution}{Grafische L\"osung}},
]

  % Branches killed by the sign conditions tan u > 0 / cot u < 0: the
  % ascending half-arches.  Dotted grey — they are not solutions.
  \addplot[birrengray!75, line width=0.7pt, densely dotted,
           domain=0:1.5707963]          {abs(sin(deg(x)))};
  \addplot[birrengray!75, line width=0.7pt, densely dotted,
           domain=1.5707963:3.1415927]  {abs(cos(deg(x)))};
  \addplot[birrengray!75, line width=0.7pt, densely dotted,
           domain=3.1415927:4.712389]   {abs(sin(deg(x)))};
  \addplot[birrengray!75, line width=0.7pt, densely dotted,
           domain=4.712389:5.45]        {abs(cos(deg(x)))};

  % The straight line u/alpha.  Reference, not data: grey.
  \addplot[birrengray!55!black, line width=0.9pt,
           domain=0:5.45, samples=2] {x/4};

  % The four admissible branches, coloured by the state they carry.
  \addplot[spacecadet,   line width=1.0pt, domain=0:1.5707963]         {abs(cos(deg(x)))};
  \addplot[cgblue,       line width=1.0pt, domain=1.5707963:3.1415927] {abs(sin(deg(x)))};
  \addplot[birrenorange, line width=1.0pt, domain=3.1415927:4.712389]  {abs(cos(deg(x)))};
  \addplot[birrenred,    line width=1.0pt, domain=4.712389:5.45]       {abs(sin(deg(x)))};

  % The three roots
  \addplot[only marks, mark=*, mark size=2pt, spacecadet]
    coordinates {(1.25235,0.31309)};
  \addplot[only marks, mark=*, mark size=2pt, cgblue]
    coordinates {(2.47458,0.61864)};
  \addplot[only marks, mark=*, mark size=2pt, birrenorange]
    coordinates {(3.59530,0.89883)};

  \node[lbl, text=spacecadet,   font=\scriptsize, fill=white, inner sep=1pt]
    at (axis cs:1.25235,0.17) {$n=1$};
  \node[lbl, text=cgblue,       font=\scriptsize, fill=white, inner sep=1pt]
    at (axis cs:2.47458,0.47) {$n=2$};
  \node[lbl, text=birrenorange, font=\scriptsize, fill=white, inner sep=1pt]
    at (axis cs:3.30000,0.56) {$n=3$};

  % Curve identities, at the top of the first two arches
  \node[lbl, anchor=south west] at (axis cs:0.06,1.02) {$|\cos u|$};
  \node[lbl, anchor=south west] at (axis cs:1.63,1.02) {$|\sin u|$};

  % The straight line, named where it runs alone under the top border
  \node[lbl, anchor=east] at (axis cs:4.95,1.30) {$u/\alpha$, $\alpha=4$};

  % The excluded branches, named once, on the row above
  \node[lbl, text=birrengray, font=\scriptsize, anchor=west]
    at (axis cs:0.10,1.46)
    {\tpcifenglish{dotted: excluded branches}{gepunktet: verworfene Zweige}};
\end{axis}

\end{tikzpicture}
\end{center}


Rechts schneidet die Gerade $u/\alpha$ nur die absteigenden Bögen: bei $\alpha=4$ dreimal, den vierten erreicht sie nicht mehr. Links dieselben drei Zustände im Topf -- je schwächer gebunden, desto weiter kriecht der Ausläufer in den verbotenen Bereich.

Aus den gefundenen Lösungen für $u$ folgt $q = \frac{2u}{L}$ und $E = \frac{\hbar^2q^2}{2m}-V_0$. Wir können den Grenzübergang zum unendlich tiefen Potentialtopf erhalten, indem wir den Grenzwert
\begin{equation}
	V_0\rightarrow \infty\quad \text{d.h.} \quad \alpha\rightarrow\infty
\end{equation}
nehmen. In diesem Limes werden wir unendlich viele Schnittpunkte für (\ref{eq:symmcond}) und (\ref{eq:unsymm}) beobachten. Diese liefern uns unendlich viele gebundene Lösungen.


%!TEX root = TheoretischePhysikC.tex
\section{Vorlesung: Drehimpuls}

\subsection{Bahndrehimpuls}
Nach den drei lösbaren Modellen folgt das Wasserstoffatom: ein Teilchen im Zentralpotential.

Wir untersuchen nun Quantensysteme, die einen Bahndrehimpuls in dem klassischen Limes aufweisen. Das kanonische Beispiel ist das eines Teilchens, das sich in drei räumlichen Dimensionen bewegen kann. Daher werden wir den Hilbertraum $\mathcal{H} \equiv L^2(\mathbb{R}^3)$ zuordnen um die Zustände zu modellieren. Unser Ziel ist es, die richtigen Quantenobservablen und ihre Eigenschaften zu untersuchen.

Nach unserer Quantisierungstabelle identifizieren wir die folgenden Operatoren als Observablen, die zu der klassischen Drehimpulskomponente gehören.
\begin{definition}
	Die \emph{Bahndrehimpulsoperatoren} lauten $\widehat{\mathbf{L}} \equiv \widehat{\mathbf{r}}\times \widehat{\mathbf{p}}$, wobei 
	\begin{equation}
		\mathbf{r} \equiv \begin{pmatrix}
			\widehat{x} \\
			\widehat{y} \\
			\widehat{z} 
		\end{pmatrix}\quad
		\mathbf{p} \equiv \begin{pmatrix}
			\widehat{p}_x \\
			\widehat{p}_y \\
			\widehat{p}_z 
		\end{pmatrix}\quad\text{und}\quad \widehat{\mathbf{L}} \equiv \begin{pmatrix}
			\widehat{L}_x \\
			\widehat{L}_y \\
			\widehat{L}_z 
		\end{pmatrix} \equiv \begin{pmatrix}
			\widehat{L}_1 \\
			\widehat{L}_2 \\
			\widehat{L}_3 
		\end{pmatrix}.
	\end{equation}
	Die Positionsbasis-Matrixelemente lauten  
	\begin{align}
		\langle x,y,z|\widehat{L}_x|\Psi\rangle &= -i{\hbar} \left(y\frac{\partial}{\partial z} - z\frac{\partial}{\partial y} \right)\Psi(x,y,z) \\
		\langle x,y,z|\widehat{L}_y|\Psi\rangle &=  i{\hbar} \left(x\frac{\partial}{\partial z} - z\frac{\partial}{\partial x} \right)\Psi(x,y,z) \\
		\langle x,y,z|\widehat{L}_z|\Psi\rangle &=  -i{\hbar} \left(x\frac{\partial}{\partial y} - y\frac{\partial}{\partial x} \right)\Psi(x,y,z),
	\end{align}
	wobei $\Psi(x,y,z) = \langle x,y,z|\Psi\rangle$ eine beliebige Wellenfunktion ist. 
\end{definition}

Man kann folgende Kommutatoren zeigen:
\begin{equation}
	\eqbox{[\widehat{L}_1,\widehat{L}_2] = i\hbar \widehat{L}_3, \quad [\widehat{L}_2,\widehat{L}_3] = i\hbar \widehat{L}_1\quad\text{und}\quad [\widehat{L}_3,\widehat{L}_1] = i\hbar \widehat{L}_2.}
\end{equation}
Es folgt auch, dass
\begin{equation}
	\eqbox{[\widehat{\mathbf{L}}^2, \widehat{L}_j]=0, \quad \text{für alle $j=1,2,3$,}}
\end{equation}
wobei
\begin{equation}
	\widehat{\mathbf{L}}^2 \equiv \widehat{L}_1^2+\widehat{L}_2^2+\widehat{L}_3^2.
\end{equation}

Der ganze Kommutator hängt an einer Stelle: $\frac{\partial}{\partial z}$ vertauscht mit $x$, nicht mit $z$; wer die Produktregel vergisst, erhält null.

Die Kommutatoren werden wie folgt bewiesen: Sei $\Psi(x,y,z)$ eine beliebige Wellenfunktion und betrachten
\begin{equation}
	\begin{split}
		\langle x,y,z|\widehat{L}_1\widehat{L}_2|\Psi\rangle &= {\hbar}^2 \left(y\frac{\partial}{\partial z} - z\frac{\partial}{\partial y} \right)\left(x\frac{\partial}{\partial z} - z\frac{\partial}{\partial x} \right)\Psi(x,y,z) \\
		&= {\hbar}^2 \left(xy\frac{\partial^2}{\partial z^2} - y\frac{\partial}{\partial x} - yz\frac{\partial^2}{\partial x \partial z} - xz\frac{\partial^2}{\partial y\partial z}  + z^2\frac{\partial^2}{\partial x\partial y} \right)\Psi(x,y,z)
	\end{split}
\end{equation}
und 
\begin{equation}
	\begin{split}
		\langle x,y,z|\widehat{L}_2\widehat{L}_1|\Psi\rangle &= {\hbar}^2 \left(x\frac{\partial}{\partial z} - z\frac{\partial}{\partial x} \right)\left(y\frac{\partial}{\partial z} - z\frac{\partial}{\partial y} \right)\Psi(x,y,z) \\
		&= {\hbar}^2 \left(xy\frac{\partial^2}{\partial z^2} - x\frac{\partial}{\partial y} - xz\frac{\partial^2}{\partial y \partial z} - yz\frac{\partial^2}{\partial x\partial z}  + z^2\frac{\partial^2}{\partial x\partial y} \right)\Psi(x,y,z).
	\end{split}
\end{equation}
Daher berechnen wir den Kommutator
\begin{equation}
	\langle x,y,z|\widehat{L}_1\widehat{L}_2 - \widehat{L}_2\widehat{L}_1|\Psi\rangle = {\hbar}^2 \left(- y\frac{\partial}{\partial x} + x\frac{\partial}{\partial y} \right)\Psi(x,y,z) = i\hbar \langle x,y,z|\widehat{L}_3|\Psi\rangle.
\end{equation}
Da dies für jedes $|\Psi\rangle$ gilt, ist es eine Operatoridentität.

Dass $\widehat{\mathbf{L}}^2$ und $\widehat{L}_3$ gemeinsame Eigenzustände besitzen, ist ein Spezialfall einer allgemeinen Tatsache.
\begin{lemma}
	Sind $\widehat{A},\widehat{B}:\mathcal{H}\to\mathcal{H}$ hermitesch mit $[\widehat{A},\widehat{B}]=0$, so besitzen sie gemeinsame Eigenvektoren.
\end{lemma}
\begin{proof}
	Endliche Dimension, keine Entartung: aus $\widehat{A}|\phi_j\rangle = a_j|\phi_j\rangle$ und $|\psi_j\rangle \equiv \widehat{B}|\phi_j\rangle$ folgt
	\begin{equation}
		\widehat{A}|\psi_j\rangle = \widehat{A}\widehat{B}|\phi_j\rangle = \widehat{B}\widehat{A}|\phi_j\rangle = a_j\widehat{B}|\phi_j\rangle = a_j|\psi_j\rangle.
	\end{equation}
	Der Eigenraum zu $a_j$ ist eindimensional, also $|\psi_j\rangle \propto |\phi_j\rangle$, d.h.\ $\widehat{B}|\phi_j\rangle = b_j|\phi_j\rangle$. Bei Entartung diagonalisiert man $\widehat{B}$ im Eigenraum; allgemein gilt die Aussage ebenso, ist aber schwerer zu begründen.
\end{proof}
Die gemeinsamen Eigenzustände von $\widehat{\mathbf{L}}^2$ und $\widehat{L}_3$ werden $|\ell,m\rangle$ bezeichnet und erfüllen die folgenden Eigenwertgleichungen 
\begin{equation}
	\eqbox{
	\begin{aligned}
		\widehat{\mathbf{L}}^2|\ell,m\rangle &= \hbar^2\ell(\ell+1)|\ell,m\rangle, \\
		\widehat{L}_3|\ell,m\rangle &= \hbar m|\ell,m\rangle,
	\end{aligned}
	}
\end{equation}
wobei $\ell = 0, \frac12, 1, \frac32, \ldots$ und $m = -\ell, -\ell + 1, \ldots, \ell$. Hierbei bezeichnet $\ell$ die \emph{Drehimpulsquantenzahl} und $m$ die \emph{magnetische Quantenzahl}.

Die Eigenwerte erhält man wie beim harmonischen Oszillator. Mit
\begin{equation}
	\eqbox{\widehat{L}_\pm \equiv \widehat{L}_1 \pm i\widehat{L}_2}
\end{equation}
folgt aus den obigen Vertauschungsrelationen
\begin{equation}
	[\widehat{L}_3,\widehat{L}_\pm] = \pm\hbar\widehat{L}_\pm, \quad [\widehat{L}_+,\widehat{L}_-] = 2\hbar\widehat{L}_3 \quad\text{und}\quad [\widehat{\mathbf{L}}^2,\widehat{L}_\pm] = 0.
\end{equation}
Die $\widehat{L}_\pm$ wirken als Leiteroperatoren: sie schieben den Eigenwert von $\widehat{L}_3$ um $\pm\hbar$ und lassen den von $\widehat{\mathbf{L}}^2$ unberührt. Da $\widehat{L}_3^2$ nicht größer als $\widehat{\mathbf{L}}^2$ werden kann, bricht die Leiter oben und unten ab, und das erzwingt die obige Liste.

% fig_vl14_multiplett.tex — Das Drehimpuls-Multiplett: die Leiter in m und
% der Kegel fester Länge, beide für ell = 2.
%
%   * geometry comes from the eigenvalue equations derived in the lecture,
%     never from a freehand sketch:
%         L^2 |ell,m> = hbar^2 ell(ell+1) |ell,m>,
%         L_3 |ell,m> = hbar m |ell,m>,   m = -ell, ..., +ell,
%         [L_3, L_pm] = +- hbar L_pm,     L_+|ell,ell> = L_-|ell,-ell> = 0.
%     Both panels are drawn in units of hbar, with ell = 2, so
%     |L|/hbar = sqrt(ell(ell+1)) = sqrt6 = 2.4495 is a fixed length and
%     the rungs sit at the integers -2,...,2.  The rung spacing is 1
%     because [L_3,L_pm] = +- hbar L_pm; nothing is spaced by eye.
%
%   * the point of the plate is the number sqrt6 = 2.4495 against the
%     number 2.  The top rung of the ladder falls SHORT of the length of
%     the vector, and that gap is what forces the ladder to break off and
%     what keeps the vector off the axis.  Panel (a) shows the gap as a
%     distance, panel (b) shows the same gap as an angle:
%         cos theta_min = ell / sqrt(ell(ell+1)) = sqrt(ell/(ell+1)),
%     i.e. 35 deg at ell = 2 and only 6 deg at ell = 100 — the classical
%     limit, as an order of magnitude rather than as an assertion.
%
%   * in panel (b) the vector of state m is drawn to the right rim of its
%     cone, at (sqrt(ell(ell+1)-m^2), m).  That point lies exactly on the
%     circle of radius sqrt6: the arrows END on the dashed circle by
%     construction, they are not fitted to it.  The cones are ellipses
%     with vertical semi-axis 0.15 times the horizontal one — the flattest
%     perspective that still reads as a cone, chosen because at 0.24 the
%     m = 0 and m = 1 cones would already overlap near the axis.
%
%   * colour encodes the magnetic quantum number m and nothing else: one
%     m, one colour, both panels, so the top rung in (a) and the top cone
%     in (b) are the same state.  The ramp runs m = -2,...,+2 through
%     `tpclevels` (spacecadet, cgblue, munsell, birrenorange, birrenred).
%     Ladder arrows and construction arcs are neutral: the operators are
%     not quantum numbers.
%
%   * the region |L_3| > |L|, which no state can reach, is hatched rather
%     than filled, so it survives greyscale printing; the bound itself is
%     a neutral grey rule, a reference line rather than data.  Following
%     fig_vl11_kasten, the hatch is a band beside that rule rather than a
%     full-bleed field, so the two L_pm|ell,+-ell> = 0 labels stand on
%     clean paper instead of on the hatch.
%
% Panel (a) the 2ell+1 rungs with L_+ and L_- between them, and the two
%           blocked arrows at the ends where the ladder terminates.
% Panel (b) the same five states as cones of fixed slant length: the
%           vector never lies along L_3, and the minimum tilt shrinks
%           only slowly with ell.

\begin{center}
\begin{tikzpicture}

% ─────────────────────────────────────────────────────────────
%  (a)  The ladder in m:  L_3 |ell,m> = hbar m |ell,m>
% ─────────────────────────────────────────────────────────────
\begin{axis}[
    name=ladder,
    tpclevelplot,
    scale only axis,
    width=5.4cm, height=5.8cm,
    xmin=-1.75, xmax=1.45,
    ymin=-3.6,  ymax=3.6,
    title={(a) $\ell=2$: \tpcifenglish{multiplet and ladder}%
                                      {Multiplett und Leiter}},
]

  % |L_3| > |L| is unreachable for every state: hatched, not filled.
  \fill[forbidden] (axis cs:-1.75,2.4495)  rectangle (axis cs:1.45,3.16);
  \fill[forbidden] (axis cs:-1.75,-2.4495) rectangle (axis cs:1.45,-3.16);

  % The bound itself, a neutral reference rule at |L|/hbar = sqrt6.
  \draw[levellinegrey] (axis cs:-1.75,2.4495)  -- (axis cs:1.45,2.4495);
  \draw[levellinegrey] (axis cs:-1.75,-2.4495) -- (axis cs:1.45,-2.4495);

  % The 2ell+1 = 5 rungs, at L_3/hbar = m.  Colour = m.
  \draw[line width=1.1pt, birrenred]    (axis cs:-0.75,2)  -- (axis cs:0.65,2);
  \draw[line width=1.1pt, birrenorange] (axis cs:-0.75,1)  -- (axis cs:0.65,1);
  \draw[line width=1.1pt, munsell]      (axis cs:-0.75,0)  -- (axis cs:0.65,0);
  \draw[line width=1.1pt, cgblue]       (axis cs:-0.75,-1) -- (axis cs:0.65,-1);
  \draw[line width=1.1pt, spacecadet]   (axis cs:-0.75,-2) -- (axis cs:0.65,-2);

  \node[lbl, anchor=east] at (axis cs:-0.82,2)  {$m=+2$};
  \node[lbl, anchor=east] at (axis cs:-0.82,1)  {$m=+1$};
  \node[lbl, anchor=east] at (axis cs:-0.82,0)  {$m=0$};
  \node[lbl, anchor=east] at (axis cs:-0.82,-1) {$m=-1$};
  \node[lbl, anchor=east] at (axis cs:-0.82,-2) {$m=-2$};

  % L_+ raises m by one:  [L_3,L_+] = +hbar L_+
  \draw[-{Stealth[length=5pt,width=3.5pt]}, spacecadet, line width=0.9pt]
    (axis cs:0.42,-1.86) -- (axis cs:0.42,-1.14);
  \draw[-{Stealth[length=5pt,width=3.5pt]}, spacecadet, line width=0.9pt]
    (axis cs:0.42,-0.86) -- (axis cs:0.42,-0.14);
  \draw[-{Stealth[length=5pt,width=3.5pt]}, spacecadet, line width=0.9pt]
    (axis cs:0.42,0.14)  -- (axis cs:0.42,0.86);
  \draw[-{Stealth[length=5pt,width=3.5pt]}, spacecadet, line width=0.9pt]
    (axis cs:0.42,1.14)  -- (axis cs:0.42,1.86);
  \node[lbl, anchor=west] at (axis cs:0.50,1.5) {$\widehat{L}_+$};

  % L_- lowers m by one:  [L_3,L_-] = -hbar L_-
  \draw[-{Stealth[length=5pt,width=3.5pt]}, spacecadet, line width=0.9pt]
    (axis cs:-0.52,1.86)  -- (axis cs:-0.52,1.14);
  \draw[-{Stealth[length=5pt,width=3.5pt]}, spacecadet, line width=0.9pt]
    (axis cs:-0.52,0.86)  -- (axis cs:-0.52,0.14);
  \draw[-{Stealth[length=5pt,width=3.5pt]}, spacecadet, line width=0.9pt]
    (axis cs:-0.52,-0.14) -- (axis cs:-0.52,-0.86);
  \draw[-{Stealth[length=5pt,width=3.5pt]}, spacecadet, line width=0.9pt]
    (axis cs:-0.52,-1.14) -- (axis cs:-0.52,-1.86);
  \node[lbl, anchor=east] at (axis cs:-0.60,-1.5) {$\widehat{L}_-$};

  % Where the ladder breaks off: the next rung would need |m| > sqrt6.
  \draw[-{Stealth[length=5pt,width=3.5pt]}, birrengray, line width=0.9pt,
        densely dashed] (axis cs:0.42,2.14) -- (axis cs:0.42,2.70);
  \draw[birrengray, line width=0.9pt]
    ([xshift=-1.4mm,yshift=-1.4mm]axis cs:0.42,2.70) --
    ([xshift=1.4mm,yshift=1.4mm]axis cs:0.42,2.70);
  \draw[birrengray, line width=0.9pt]
    ([xshift=-1.4mm,yshift=1.4mm]axis cs:0.42,2.70) --
    ([xshift=1.4mm,yshift=-1.4mm]axis cs:0.42,2.70);

  \draw[-{Stealth[length=5pt,width=3.5pt]}, birrengray, line width=0.9pt,
        densely dashed] (axis cs:-0.52,-2.14) -- (axis cs:-0.52,-2.70);
  \draw[birrengray, line width=0.9pt]
    ([xshift=-1.4mm,yshift=-1.4mm]axis cs:-0.52,-2.70) --
    ([xshift=1.4mm,yshift=1.4mm]axis cs:-0.52,-2.70);
  \draw[birrengray, line width=0.9pt]
    ([xshift=-1.4mm,yshift=1.4mm]axis cs:-0.52,-2.70) --
    ([xshift=1.4mm,yshift=-1.4mm]axis cs:-0.52,-2.70);

  \node[lbl, anchor=south] at (axis cs:0.42,3.24)
    {$\widehat{L}_+\ket{\ell,\ell}=0$};
  \node[lbl, anchor=north] at (axis cs:-0.52,-3.24)
    {$\widehat{L}_-\ket{\ell,-\ell}=0$};

  % What the grey rules are
  \node[lbl, font=\scriptsize, anchor=south west, fill=white,
        fill opacity=1, inner sep=2pt]
    at (axis cs:-1.73,2.53) {$|\mathbf{L}|/\hbar=\sqrt{\ell(\ell+1)}=\sqrt{6}$};

  % The multiplet is 2ell+1 states long
  \draw[{Stealth[length=4pt]}-{Stealth[length=4pt]}, birrengray,
        line width=0.6pt] (axis cs:0.95,-2) -- (axis cs:0.95,2);
  \node[lbl, font=\scriptsize, text=birrengray, rotate=90, anchor=south]
    at (axis cs:1.05,0) {$2\ell+1$ \tpcifenglish{states}{Zustände}};
\end{axis}

% ─────────────────────────────────────────────────────────────
%  (b)  The vector model: |L| fixed, L_3 quantised, L_1,L_2 undecided
% ─────────────────────────────────────────────────────────────
\begin{axis}[
    name=cones,
    tpclevelplot,
    scale only axis,
    width=4.78cm, height=5.8cm,      % 6.1 x 7.4 units, equal unit lengths
    at={(ladder.south east)}, anchor=south west, xshift=14mm,
    xmin=-3.05, xmax=3.05,
    ymin=-3.9,  ymax=3.5,
    title={(b) \tpcifenglish{vector model}{Vektormodell}},
]

  % The cones: |L| = sqrt6 hbar with L_3 = m hbar fixed, so the tip runs
  % on a circle of radius sqrt(6 - m^2) at height m.
  \addplot[line width=0.6pt, densely dotted, birrenred,
           domain=0:360, samples=121] ({1.41421*cos(x)}, {2+0.21213*sin(x)});
  \addplot[line width=0.6pt, densely dotted, birrenorange,
           domain=0:360, samples=121] ({2.23607*cos(x)}, {1+0.33541*sin(x)});
  \addplot[line width=0.6pt, densely dotted, munsell,
           domain=0:360, samples=121] ({2.44949*cos(x)}, {0.36742*sin(x)});
  \addplot[line width=0.6pt, densely dotted, cgblue,
           domain=0:360, samples=121] ({2.23607*cos(x)}, {-1+0.33541*sin(x)});
  \addplot[line width=0.6pt, densely dotted, spacecadet,
           domain=0:360, samples=121] ({1.41421*cos(x)}, {-2+0.21213*sin(x)});

  % The sphere of fixed length: every arrow below ends on it.
  \addplot[levellinegrey, domain=0:360, samples=181]
    ({2.44949*cos(x)}, {2.44949*sin(x)});

  % The axis the multiplet is quantised along
  \draw[axisstyle] (axis cs:0,-2.66) -- (axis cs:0,3.30);
  \node[lbl, anchor=west] at (axis cs:0.10,3.22) {$L_3/\hbar$};

  % The second generator of the top cone, so it reads as a cone
  \draw[line width=0.6pt, densely dotted, birrenred]
    (axis cs:0,0) -- (axis cs:-1.41421,2);

  % One representative vector per cone, ending on the circle at
  % (sqrt(6-m^2), m).  Colour = m, as in panel (a).
  \draw[-{Stealth[length=5pt,width=3.5pt]}, birrenred, line width=1.0pt]
    (axis cs:0,0) -- (axis cs:1.41421,2);
  \draw[-{Stealth[length=5pt,width=3.5pt]}, birrenorange, line width=1.0pt]
    (axis cs:0,0) -- (axis cs:2.23607,1);
  \draw[-{Stealth[length=5pt,width=3.5pt]}, munsell, line width=1.0pt]
    (axis cs:0,0) -- (axis cs:2.44949,0);
  \draw[-{Stealth[length=5pt,width=3.5pt]}, cgblue, line width=1.0pt]
    (axis cs:0,0) -- (axis cs:2.23607,-1);
  \draw[-{Stealth[length=5pt,width=3.5pt]}, spacecadet, line width=1.0pt]
    (axis cs:0,0) -- (axis cs:1.41421,-2);

  \node[lbl, anchor=west] at (axis cs:1.52,2)  {$m=\ell$};
  \node[lbl, anchor=west] at (axis cs:1.52,-2) {$m=-\ell$};

  % The tilt that never vanishes: cos theta_min = ell/sqrt(ell(ell+1)),
  % drawn in the gap between the m = 1 and m = 2 cones.
  \addplot[birrengray, line width=0.7pt, domain=54.7356:90, samples=40]
    ({1.68*cos(x)}, {1.68*sin(x)});
  \node[lbl, font=\scriptsize, anchor=west] at (axis cs:0.07,1.13)
    {$\theta_{\min}$};

  % Its pole: the length the top cone falls short of, sqrt6 against m = 2
  \draw[birrengray, line width=0.5pt]
    (axis cs:-0.09,2.4495) -- (axis cs:0.09,2.4495);
  \node[lbl, font=\scriptsize, text=birrengray, anchor=south east,
        inner sep=1pt] at (axis cs:-0.11,2.53) {$\sqrt{6}$};

  % What the dashed circle is
  \node[lbl, anchor=north] at (axis cs:0,-2.72)
    {$|\mathbf{L}|=\hbar\sqrt{\ell(\ell+1)}=\sqrt{6}\,\hbar$};

  % The classical limit as a number, not as an assertion
  \node[lbl, font=\scriptsize, text=birrengray, anchor=north, align=center]
    at (axis cs:0,-3.28)
    {\tpcifenglish{never along the axis}{nie entlang der Achse}:\\[-1pt]
     $\ell=2$: $35^\circ$\quad $\ell=10$: $18^\circ$\quad
     $\ell=100$: $6^\circ$};
\end{axis}

\end{tikzpicture}
\end{center}


Man beachte in (a) den Abstand zwischen der obersten Sprosse $m=\ell$ und der
gestrichelten Schranke $\sqrt{\ell(\ell+1)}$: genau dort bricht die Leiter ab.
In (b) ist dieselbe Lücke ein Winkel -- der Vektor liegt nie auf der Achse.

\begin{merksatz}
	Position und Impuls sind nicht quantisiert, der Drehimpuls schon: es gibt eine kleinste Stufe, von der Größe $\hbar$. Im Alltag unsichtbar, spektroskopisch nachgewiesen.
\end{merksatz}

Nun werden wir die Orstdarstellung der Eigenfunktionen herleiten. Mit Hilfe von den Kugelkoordinaten $r$, $\theta$, $\phi$ erhalten wir (Übung)
\begin{equation}
	\begin{split}
		\langle r,\theta,\phi|\widehat{L}_3|\Psi\rangle &= -i\hbar \frac{\partial}{\partial \phi} \Psi(r,\theta,\phi), \\
		\langle r,\theta,\phi|\widehat{\mathbf{L}}^2|\Psi\rangle &= -\hbar^2\frac{1}{\sin^2(\theta)}\left(\sin(\theta)\frac{\partial}{\partial \theta}\left(\sin(\theta)\frac{\partial}{\partial \theta}\right) + \frac{\partial^2}{\partial \phi^2}\right) \Psi(r,\theta,\phi),
	\end{split}
\end{equation}
wobei $|\Psi\rangle \in L^2(\mathbb{R}^3)$ einen beliebigen reinen Zustand ist. 

Die Eigenfunktionen von $\widehat{L}_3$ sind durch $e^{im\phi}$ gegeben:
\begin{equation}
	-i\hbar \frac{\partial}{\partial \phi}e^{im\phi} = \hbar m e^{im\phi}.
\end{equation}
Die Eindeutigkeit der Wellenfunktion $\Psi(r,\theta,\phi)$ unter die Verschiebung $\phi\mapsto \phi+2\pi$ fordert $m\in \mathbb{Z}$, sodass nur $\ell = 0,1,2, \ldots$ auftritt. (Die Eigenvektoren für Halbzahlige $\ell$ wird später mit eine andere Methode aufgebaut.) Die gemeinsamen Eigenfunktionen von $\widehat{\mathbf{L}}^2$ und $\widehat{L}_3$ lauten
\begin{equation}
	\eqbox{\text{Y}_\ell^m (\theta,\phi) \equiv (-1)^m\left(\frac{2\ell + 1}{4\pi}\frac{(\ell-m)!}{(\ell+m)!}\right)^{\frac12} \text{P}_\ell^m(\cos(\theta))e^{im\phi},}
\end{equation}
wobei für $m\ge 0$
\begin{equation}
	\eqbox{\text{P}^m_\ell(x) = (1-x^2)^{\frac{m}{2}}\frac{d^m}{dx^m}\text{P}_\ell(x)}
\end{equation}
die \emph{assoziierten Legendrefunktionen} die mit Hilfe von den \emph{Legendre-Polynomen}
\begin{equation}
	\eqbox{\text{P}_\ell(x)=\frac{1}{2^\ell \ell!} \frac{d^\ell}{dx^\ell}(x^2-1)^\ell}
\end{equation}
definiert sind. Die Funktionen $\text{Y}_\ell^m$ sind als die \emph{Kugelflächenfunktionen} bekannt. Sie hängen nicht von $r$ ab -- der Drehimpuls sieht die Radialkomponente nicht, was später den Separationsansatz trägt.

Die ersten drei Kugelflächenfunktionen lauten
\begin{align}
	\text{Y}^0_0(\theta,\phi) &= \frac{1}{\sqrt{4\pi}}, \\
	\text{Y}^0_1(\theta,\phi) &= \sqrt{\frac{3}{4\pi}}\cos(\theta), \\
	\text{Y}^1_1(\theta,\phi) &= -\sqrt{\frac{3}{8\pi}}\sin(\theta)e^{i\phi}.
\end{align}
$\text{Y}^0_0$ ist konstant: das Teilchen besitzt dann gar keinen Drehimpuls.

%!TEX root = TheoretischePhysikC.tex

\section{Vorlesung: Zentralpotential}
Man betrachte Kugelsymmetrisches Potential in drei Dimensionen:
\begin{equation}
	V(\mathbf{r}) = V(r) \quad\text{mit $r = |\mathbf{r}|$.}
\end{equation}
Das Ziel ist es, die Kinematik und Dynamik eines Teilchens, das sich in dem Potential bewegt, zu lösen. Mittels Schrödingergleichung ist
\begin{equation}
	\widehat{H}|\Psi\rangle = E|\Psi\rangle
\end{equation}
zu lösen. Diese Eigenvektorgleichung liefert uns die Eigenwerte und Eigenvektoren $\widehat{H}$, die man verwenden kann, um den Unitäroperator $U_t$ zu berechnen. Eigenwerte sind hier mehr als Technik: gegenüber den atomaren Anregungsenergien ist Raumtemperatur winzig, das Elektron sitzt fast immer im Grundzustand.

\subsection{Separation der Variablen}
Die Positionsbasis $x,y,z$ wäre nicht falsch, nur schlecht gewählt: man richtet die Darstellung nach der Symmetrie.

Die Matrixelemente des Impulsoperators in Kugelkoordinaten lauten $\forall |\Psi\rangle \in \mathcal{H}$:
\begin{equation}
	\begin{split}
		\langle r,\theta,\phi|\frac{\widehat{\mathbf{p}}^2}{2m}|\Psi\rangle &= -\frac{\hbar^2}{2m}\nabla^2 \Psi(r,\theta,\phi) \\
		&= -\frac{\hbar^2}{2m} \left[\left(\frac{\partial}{\partial r} + \frac1r\right)^2 + \frac{1}{r^2\sin(\theta)}\left(\sin(\theta)\frac{\partial}{\partial \theta}\right)^2 +  \frac{1}{r^2\sin^2(\theta)}\frac{\partial^2}{\partial \phi^2}\right]\Psi(r,\theta,\phi) \\
		&= \langle r,\theta,\phi|\frac{1}{2m}\left(\widehat{p_r}^2 + \frac{\widehat{\mathbf{L}}^2}{r^2}\right)|\Psi\rangle,
	\end{split}
\end{equation}
wobei $\Psi(r,\theta,\phi) \equiv \langle r,\theta,\phi|\Psi\rangle$ und $\langle r,\theta,\phi|\widehat{p_r}|\Psi\rangle = -i\hbar \left(\frac{\partial}{\partial r} + \frac1r\right)\Psi(r,\theta,\phi)$ den radialen Impulsoperator darstellt. Der Zusatzterm $\frac1r$ ist wesentlich -- der radiale Impuls ist \emph{nicht} $-i\hbar\frac{\partial}{\partial r}$.

Die dritte Zeile ist der Angelpunkt des Kapitels: der Winkelanteil der kinetischen Energie \emph{ist} $\widehat{\mathbf{L}}^2/(2mr^2)$.

Daher werden die Eigenwerte und Eigenvektoren $\widehat{H}$ durch
\begin{equation}\label{eq:centralpotsse}
	\langle r,\theta,\phi|\widehat{H}|\Psi\rangle = \left(\frac{\widehat{p_r}^2}{2m} + \frac{\widehat{\mathbf{L}}^2}{2mr^2} + V(\widehat{\mathbf{r}})\right)\Psi(r,\theta,\phi) = E\Psi(r,\theta,\phi)
\end{equation}
gegeben. Achten Sie auf den Hut: im Matrixelement wird aus dem Operator $V(\widehat{\mathbf{r}})$ die Zahl $V(r)$.

\begin{greybox}[Symmetrie und Erhaltungsgröße]
	Eine Drehung führt das System in sich selbst über, und zu jeder Symmetrie gehört eine Erhaltungsgröße -- der Satz von Noether, den Sie aus Theoretische Physik B als Ladungserhaltung kennen. Die Drehimpulsoperatoren \emph{erzeugen} die Rotationen; weil $\widehat{H}$ rotationsinvariant ist, muss der Drehimpuls erhalten sein.
\end{greybox}

Es gilt:
\begin{equation}
	[\widehat{H}, \widehat{\mathbf{L}}^2] = [\widehat{H}, \widehat{L}_z] = [\widehat{\mathbf{L}}^2,\widehat{L}_z] = 0.
\end{equation}
Diese Beobachtung hat eine wichtige Auswirkung: Es gibt gemeinsame Eigenvektoren für die Operatoren $\widehat{H}$, $\widehat{L}_z$ und $\widehat{\mathbf{L}}^2$ (dies ist eine allgemeine Implikation, die für alle Mengen von kommutierenden Operatoren gilt).

Ziel ist es, die gemeinsame Eigenfunktionen von $\widehat{H}$ von $\widehat{H}$, $\widehat{L}_z$ und $\widehat{\mathbf{L}}^2$ zu finden. Wenn auf diese Weise eine eindeutige Wahl von Basiszuständen des Hilbertraums erreicht wird, spricht man von einem vollständigen Satz kommutierender Observablen.

\begin{intuition}
	Normalerweise wird ein Problem schwerer, wenn man mehr fordert; hier wird es leichter. Kommutierende Observablen sind gleichzeitig scharf messbar und verhalten sich, grob gesagt, wie unabhängige Teilsysteme -- deshalb zerfällt die Lösung. Ein Separationsansatz ist kein Glücksfall, sondern die Rechenform einer Symmetrie.
\end{intuition}

Um die Eigenvektoren und Eigenwerte von $\widehat{H}$ zu berechnen, wenden wir einen Separationsansatz an:
\begin{equation}
	\Psi(r,\theta,\phi) = \text{Y}_{\ell}^m(\theta,\phi) R(r) 
\end{equation}
mit
\begin{equation}
	\widehat{\mathbf{L}}^2\text{Y}_{\ell}^m(\theta,\phi) = \hbar^2 \ell(\ell+1)\text{Y}_{\ell}^m(\theta,\phi).
\end{equation}
Wegen Normierung der $\text{Y}_{\ell}^m(\theta,\phi)$ gilt
\begin{equation}
	\iiint |\Psi(\mathbf{r})|^2\,dxdydz = \int_0^{\infty} |R(r)|^2 r^2\, dr.
\end{equation}
Setzen wir den Separationsansatz in (\ref{eq:centralpotsse}) ein, erhalten wir
\begin{equation}
	\left(\frac{\widehat{p_r}^2}{2m} + \frac{\hbar^2\ell(\ell+1)}{2mr^2} + V(r)\right)R(r) = E R(r).
\end{equation}
Der Operator $\widehat{p}_r^2$ hat die Matrixelemente 
\begin{equation}
	\begin{split}
		\langle r,\theta,\phi|\widehat{p}_r^2|\Psi\rangle &= -\hbar^2\left(\frac{\partial}{\partial r} + \frac{1}{r}\right)\left(\frac{\partial}{\partial r} + \frac{1}{r}\right)\Psi(r,\theta,\phi) \\
		&= -\hbar^2\left(\frac{1}{r}\frac{\partial}{\partial r} r\right)\left(\frac{1}{r}\frac{\partial}{\partial r} r\right)\Psi(r,\theta,\phi) \\ 
		&= -\hbar^2 \frac{1}{r}\frac{\partial^2}{\partial r^2}\left(r \Psi(r,\theta,\phi)\right).
	\end{split}
\end{equation}
Mit Hilfe von der Definition $u(r)\equiv rR(r)$ können wir die radiale Wirkung wie folgt erhalten
\begin{equation}
	\left(-\frac{\hbar^2}{2m}\frac{d^2}{dr^2} + \frac{\hbar^2\ell(\ell+1)}{2mr^2} + V(r) - E\right)u(r) = 0.
\end{equation}
Das ist ein eindimensionales Problem: $r$ spielt die Rolle der Position, der Zusatzterm $\hbar^2\ell(\ell+1)/(2mr^2)$ ist die Zentrifugalbarriere -- wie bei den Planetenbahnen.

Man kann zeigen, dass nur Lösungen mit $u(0)=0$ physikalisch sinnvoll sind, da sonst entweder das Normquadrat 
\begin{equation}
	\int_0^\infty |u(r)|^2\, dr = \infty
\end{equation}
beträgt, oder den Erwartungswert der kinetischen Energie wegen $\langle \widehat{p_r}^2\rangle = \infty$ unendlich werden lässt. Daher ist die Radialgleichung äquivalent zu einem eindimensionalen Problem auf $x>0$ mit Randbedingung bei $x=0$. Festgelegt ist $V(r)$ damit noch nicht; die wichtigste Wahl, das Coulombpotential, folgt in der nächsten Vorlesung.


%!TEX root = TheoretischePhysikC.tex

\section{Vorlesung: Das Coulombpotential (Wasserstoffatom)}

\begin{center}
	\begin{tikzpicture}
		\begin{axis}[width=\columnwidth, height=8cm,hide axis, tick style={color=black}, domain=0:12, xtick=\empty, ytick=\empty, samples =200, xmin = -0.5, xmax = 8, ymin = -3, ymax = 2]
		\addplot[cgblue, dashed] {1/x^2-2/x};
		\addplot[cgblue] {-2/x};
		\draw[->] (axis cs: -0.5, 0) -- (axis cs: 8, 0);
		\draw (axis cs: 7.8, -0.2) node[below] {$r$};
		\draw (axis cs: 6, 1.5) node[below] {$U_\ell(r) = \frac{\hbar^2 \ell(\ell+1)}{2m_{\text{e}}r^2}+V(r)$};
		\draw[->] (axis cs: 0, -3) -- (axis cs: 0, 2);
    	\draw (axis cs: 0.5,1.5) node[right] {$\ell = 1$};
   		\draw (axis cs:3.5,-1.25) node[above] {$\ell=0$};
		\end{axis}
	\end{tikzpicture}
\end{center} 

Das Wasserstoffatom besteht aus einem Proton und einem Elektron. Ein erstes grobes Modell für das Atom erfolgt indem man das Proton als eine unendlichen massiven Punktmasse annähert. Nach dieser Annäherung kann nun das Wasserstoffatom als ein Teilchen (das Elektron), das sich in einem Zentralpotential 
\begin{equation}
	V(r) = -\frac{Ze^2}{4\pi \epsilon_0 r}
\end{equation}
bewegen, modelliert werden.

Das Elektron ist dabei kein Punkt auf einer Bahn, sondern eine Wolke; das Proton
ebenso, nur lässt der Massenunterschied es als Punktmasse durchgehen. Die Reduktion
auf ein Teilchen ist keine Näherung: Relativkoordinaten und die reduzierte Masse
$\mu = m_{\text{e}}m_{\text{p}}/(m_{\text{e}}+m_{\text{p}})$ machen sie exakt;
genähert wird erst mit $\mu \approx m_{\text{e}}$. Aus solchen Rechnungen wurden
die Wasserstofflinien bestimmt, mit denen man Entfernungen zu fernen Galaxien
vermaß.

Die Radialgleichung ist ein eindimensionales Problem im effektiven Potential
$U_\ell(r)$ des Bildes oben. Lösen lässt sie sich mit dem Potenzreihenansatz, der
Frobenius-Methode; das ist so aufwendig, dass wir die Lösungen einfach
hinschreiben und durch Einsetzen prüfen.

Die Lösungen der Radialgleichung 
\begin{equation}
	\left(-\frac{\hbar^2}{2m_{\text{e}}}\frac{d^2}{dr^2} + \frac{\hbar^2\ell(\ell+1)}{2m_{\text{e}}r^2} + V(r) - E\right)u(r) = 0.
\end{equation}
sind der Form
\begin{equation}
	R_{n\ell}(r) = \sqrt{\frac{(n-\ell-1)!(2\kappa)^3}{2n(n+\ell)!}}(2\kappa r)^\ell e^{-\kappa r}\text{L}_{n-\ell-1}^{2\ell +1}(2\kappa r),	
\end{equation}
wobei 
\begin{equation}
	\kappa = \frac{\sqrt{2m_{\text{e}}|E|}}{\hbar} = \frac{Z}{n a_0}
\end{equation}
und $L_r^s$ die Laguerre-Polynomen sind. Die dazugehörige Eigenwerte lauten
\begin{equation}
	\eqbox{E_n = -\frac{(Ze)^2}{4\pi\epsilon_0}\frac{1}{a_0}\frac{1}{2n^2},}
\end{equation}
mit \emph{Hauptquantenzahl} $n=1,2,\ldots$ und \emph{Bohrradius} 
\begin{equation}
	a_0 = \frac{4\pi\epsilon_0\hbar^2}{m_{\text{e}}e^2} = 5.29177210903(80) \times 10^{-11} \text{m}.
\end{equation}
Die Zählung beginnt bei $n=1$, ein $n=0$ gibt es nicht; im Termschema häufen sich
die unendlich vielen Niveaus an der Ionisationsschwelle $E=0$.
\begin{merksatz}
Beobachtbar sind nicht die Eigenwerte, sondern ihre Differenzen: beim Übergang
$n \to n'$ mit $n'<n$ emittiert das Elektron ein Photon der Energie $E_n-E_{n'}$,
und das sind die Spektrallinien.
\end{merksatz}

% fig_vl16_termschema.tex — Das Termschema des Wasserstoffs und die
% Spektrallinien, die aus seinen Niveaudifferenzen folgen.
%
% Follows the conventions of fig_vl11_kasten (the reference plate).
%
%   * geometry comes from the closed form derived in the lecture, never
%     from a freehand sketch.  With the Rydberg energy
%         E_R = e^2/(4 pi eps_0 a_0) / 2 = 13.6057 eV,
%     the lecture's eigenvalue (Z = 1)
%         E_n = -E_R / n^2
%     fixes every level height in panel (a):
%         n  = 1      2       3       4       5       6      ...
%         E  = -13.61 -3.401 -1.512  -0.850  -0.544  -0.378  ... -> 0
%     so the 1/n^2 accumulation at the ionisation threshold is forced by
%     the formula and not drawn by eye.  Levels up to n = 20 are drawn;
%     from n ~ 8 they can no longer be resolved on the plate, which is
%     precisely the statement "infinitely many levels pile up at E = 0".
%   * the transition arrows carry the Merksatz of the lecture: observable
%     are not the eigenvalues but their differences.  Hence the levels are
%     neutral grey reference rules and the *arrows* are the coloured data.
%     Arrow length is literally the photon energy E_n - E_n', to scale:
%     the Lyman arrows are ~4x the Balmer arrows, which is the whole
%     reason the one series is deep UV and the other visible light.
%   * colour encodes the final state n' of the series, and nothing else:
%     n' = 1 (Lyman) cgblue, n' = 2 (Balmer) birrenorange — the same two
%     colours in both panels, so a stick in (b) can be traced back to an
%     arrow in (a).
%   * panel (b) plots the wavelengths that follow from the same formula,
%         1/lambda = R_inf (1/n'^2 - 1/n^2),  1/R_inf = 91.127 nm,
%     on a logarithmic axis, because only a log axis holds Lyman
%     (91-122 nm) and Balmer (365-656 nm) in one frame.  Consequences
%     that are therefore forced: the two series never overlap, the band
%     122-365 nm carries no hydrogen line at all, and the lines of each
%     series crowd at their series limit lambda = n'^2 / R_inf exactly as
%     the levels crowd at E = 0.  The visible window 380-750 nm is drawn
%     as a plain grey bracket, not a coloured band: it is reference
%     furniture and must survive greyscale printing.
%
% Panel (a) the term scheme with the Lyman and Balmer transitions.
% Panel (b) the same transitions as spectral lines on a wavelength axis —
%           the order of magnitude prose cannot deliver: Balmer lands in
%           the visible, Lyman does not.

\begin{center}
\begin{tikzpicture}

\tikzset{
  trans/.style={-{Stealth[length=5pt,width=3.5pt]}, line width=0.9pt},
}

% ─────────────────────────────────────────────────────────────
%  (a)  Term scheme E_n = -13.6 eV / n^2
% ─────────────────────────────────────────────────────────────
\begin{axis}[
    name=terme,
    tpcplot,
    width=6.2cm, height=6.9cm,
    axis y line=left, axis x line=none,
    grid=none, clip=false,
    xmin=0, xmax=1.20,
    ymin=-16.4, ymax=1.6,
    ytick={0,-3.4014,-13.6057},
    yticklabels={$0$,$-3.4$,$-13.6$},
    ylabel={$E$ [eV]},
    every axis y label/.style={at={(axis description cs:-0.20,0.5)},
                               rotate=90, anchor=south, font=\small},
    title={(a) \tpcifenglish{Term scheme}{Termschema}},
]

  % Bound levels E_n = -13.6057/n^2.  Written out rather than looped:
  % \foreach defers expansion inside a pgfplots axis.
  \draw[levellinegrey] (axis cs:0.02,-13.6057) -- (axis cs:0.88,-13.6057);
  \draw[levellinegrey] (axis cs:0.02, -3.4014) -- (axis cs:0.88, -3.4014);
  \draw[levellinegrey] (axis cs:0.02, -1.5117) -- (axis cs:0.88, -1.5117);
  \draw[levellinegrey] (axis cs:0.02, -0.8504) -- (axis cs:0.88, -0.8504);
  \draw[levellinegrey] (axis cs:0.02, -0.5442) -- (axis cs:0.88, -0.5442);
  \draw[levellinegrey] (axis cs:0.02, -0.3779) -- (axis cs:0.88, -0.3779);
  \draw[levellinegrey] (axis cs:0.02, -0.2777) -- (axis cs:0.88, -0.2777);
  \draw[levellinegrey] (axis cs:0.02, -0.2126) -- (axis cs:0.88, -0.2126);
  \draw[levellinegrey] (axis cs:0.02, -0.1680) -- (axis cs:0.88, -0.1680);
  \draw[levellinegrey] (axis cs:0.02, -0.1361) -- (axis cs:0.88, -0.1361);
  \draw[levellinegrey] (axis cs:0.02, -0.0945) -- (axis cs:0.88, -0.0945);

  % The ionisation threshold: a real boundary, so heavy ink.  It needs no
  % "E = 0" tag at its right end — the y tick already reads 0 — and that
  % end is where the accumulating levels have to be labelled.
  \draw[potential] (axis cs:0.02,0) -- (axis cs:0.88,0);
  \node[lbl, text=birrengray, font=\scriptsize, anchor=south]
    at (axis cs:0.58,0.10)
    {\tpcifenglish{ionisation threshold}{Ionisationsschwelle}};

  % Lyman series, n -> 1.  Arrow length = E_n - E_1, to scale.
  \draw[trans, cgblue] (axis cs:0.16, -3.4014) -- (axis cs:0.16,-13.6057);
  \draw[trans, cgblue] (axis cs:0.23, -1.5117) -- (axis cs:0.23,-13.6057);
  \draw[trans, cgblue] (axis cs:0.30, -0.8504) -- (axis cs:0.30,-13.6057);

  % Balmer series, n -> 2
  \draw[trans, birrenorange] (axis cs:0.60,-1.5117) -- (axis cs:0.60,-3.4014);
  \draw[trans, birrenorange] (axis cs:0.67,-0.8504) -- (axis cs:0.67,-3.4014);
  \draw[trans, birrenorange] (axis cs:0.74,-0.5442) -- (axis cs:0.74,-3.4014);

  % Series labels, directly under the level each series ends on
  \node[lbl, text=cgblue, anchor=north, align=center]
    at (axis cs:0.30,-13.90) {Lyman ($n\to1$)\\[-1pt] UV};
  \node[lbl, text=birrenorange, anchor=north, align=center]
    at (axis cs:0.70,-3.72)
    {Balmer ($n\to2$)\\[-1pt] \tpcifenglish{visible}{sichtbar}};

  % Quantum numbers, only where the levels are still resolved
  \node[lbl, text=birrengray, font=\scriptsize, anchor=west]
    at (axis cs:0.90,-13.6057) {$n=1$};
  \node[lbl, text=birrengray, font=\scriptsize, anchor=west]
    at (axis cs:0.90, -3.4014) {$n=2$};
  \node[lbl, text=birrengray, font=\scriptsize, anchor=west]
    at (axis cs:0.90, -1.5117) {$n=3$};

  % The accumulation at the threshold: labelled exactly like the resolved
  % levels, in the same column, because that is what it is — the rest of
  % the same sequence, no longer separable on the plate.
  \node[lbl, text=birrengray, font=\scriptsize, anchor=west]
    at (axis cs:0.90,-0.30) {$n\to\infty$};
\end{axis}

% ─────────────────────────────────────────────────────────────
%  (b)  The same transitions as spectral lines
% ─────────────────────────────────────────────────────────────
\begin{axis}[
    tpcplot,
    at={(terme.north east)}, anchor=north west, xshift=13mm,
    width=7.4cm, height=5.7cm,
    axis y line=none, axis x line=bottom,
    grid=none, clip=false,
    xmode=log, log ticks with fixed point,
    xmin=82, xmax=900,
    ymin=0,   ymax=1.42,
    xtick={100,200,400,800},
    minor xtick={90,300,500,600,700},
    xlabel={$\lambda$ [nm]},
    every axis x label/.style={at={(axis description cs:0.5,-0.26)},
                               anchor=north, font=\small},
    title={(b) \tpcifenglish{Spectral lines}{Spektrallinien}},
]

  % Lyman, n = 2..12 -> 1.  lambda = 91.127 nm / (1 - 1/n^2)
  \draw[cgblue, line width=0.7pt] (axis cs:121.50,0) -- (axis cs:121.50,1);
  \draw[cgblue, line width=0.7pt] (axis cs:102.52,0) -- (axis cs:102.52,1);
  \draw[cgblue, line width=0.7pt] (axis cs: 97.20,0) -- (axis cs: 97.20,1);
  \draw[cgblue, line width=0.7pt] (axis cs: 94.92,0) -- (axis cs: 94.92,1);
  \draw[cgblue, line width=0.7pt] (axis cs: 93.73,0) -- (axis cs: 93.73,1);
  \draw[cgblue, line width=0.7pt] (axis cs: 93.03,0) -- (axis cs: 93.03,1);
  \draw[cgblue, line width=0.7pt] (axis cs: 92.57,0) -- (axis cs: 92.57,1);
  \draw[cgblue, line width=0.7pt] (axis cs: 92.27,0) -- (axis cs: 92.27,1);
  \draw[cgblue, line width=0.7pt] (axis cs: 92.05,0) -- (axis cs: 92.05,1);
  \draw[cgblue, line width=0.7pt] (axis cs: 91.76,0) -- (axis cs: 91.76,1);

  % Balmer, n = 3..12 -> 2.  lambda = 91.127 nm / (1/4 - 1/n^2)
  \draw[birrenorange, line width=0.7pt] (axis cs:656.11,0) -- (axis cs:656.11,1);
  \draw[birrenorange, line width=0.7pt] (axis cs:486.01,0) -- (axis cs:486.01,1);
  \draw[birrenorange, line width=0.7pt] (axis cs:433.94,0) -- (axis cs:433.94,1);
  \draw[birrenorange, line width=0.7pt] (axis cs:410.07,0) -- (axis cs:410.07,1);
  \draw[birrenorange, line width=0.7pt] (axis cs:396.91,0) -- (axis cs:396.91,1);
  \draw[birrenorange, line width=0.7pt] (axis cs:388.81,0) -- (axis cs:388.81,1);
  \draw[birrenorange, line width=0.7pt] (axis cs:383.44,0) -- (axis cs:383.44,1);
  \draw[birrenorange, line width=0.7pt] (axis cs:379.69,0) -- (axis cs:379.69,1);
  \draw[birrenorange, line width=0.7pt] (axis cs:374.92,0) -- (axis cs:374.92,1);

  % Series limits lambda = n'^2 / R_inf, where n -> infinity accumulates —
  % the image of the pile-up at E = 0 in panel (a).  Annotated once, from
  % the empty band between the series, so no label sits on a line.
  \draw[levellinegrey] (axis cs: 91.13,0) -- (axis cs: 91.13,1.20);
  \draw[levellinegrey] (axis cs:364.51,0) -- (axis cs:364.51,1.20);
  \node[lbl, text=birrengray, font=\scriptsize, anchor=south west]
    at (axis cs:93.5,1.20) {\tpcifenglish{series limit}{Seriengrenze}};
  \node[lbl, text=birrengray, font=\scriptsize, anchor=south east]
    at (axis cs:358,1.02) {\tpcifenglish{series limit}{Seriengrenze}};

  % The band between the two series is empty, and that emptiness is a
  % result, not a gap in the drawing — hence it is labelled.
  \node[lbl, text=birrengray, font=\scriptsize, align=center]
    at (axis cs:210,0.46)
    {\tpcifenglish{no hydrogen}{keine}\\[-1pt]
     \tpcifenglish{lines here}{Wasserstofflinien}};

  % The one hydrogen line everybody has seen: red H-alpha, n = 3 -> 2.
  \node[lbl, text=birrenorange, font=\scriptsize, anchor=south]
    at (axis cs:656.11,1.00) {$\mathrm{H}\alpha$};

  % Spectral windows, as neutral grey brackets under the axis: reference
  % furniture, so no colour, and greyscale-safe.  The visible window
  % starts just longward of the Balmer limit — which is why that whole
  % series, and only that series, is a laboratory-bench phenomenon.
  \draw[birrengray, line width=0.5pt]
    (axis cs:84,-0.145) -- (axis cs:84,-0.18) -- (axis cs:378,-0.18)
      -- (axis cs:378,-0.145);
  \node[lbl, text=birrengray, font=\scriptsize, anchor=north]
    at (axis cs:178,-0.19) {UV};
  \draw[birrengray, line width=0.5pt]
    (axis cs:384,-0.145) -- (axis cs:384,-0.18) -- (axis cs:750,-0.18)
      -- (axis cs:750,-0.145);
  \node[lbl, text=birrengray, font=\scriptsize, anchor=north]
    at (axis cs:537,-0.19) {\tpcifenglish{visible}{sichtbar}};
\end{axis}

\end{tikzpicture}
\end{center}


Farbig sind hier nicht die Niveaus, sondern die Übergänge: die Pfeillänge ist
maßstäblich die Photonenenergie, die Lyman-Pfeile sind rund viermal so lang wie
die Balmer-Pfeile. Rechts kehrt die Häufung bei $E=0$ als Gedränge an der
Seriengrenze wieder.

Die vollständige Eigenfunktionen lauten
\begin{equation}
	\psi_{n\ell m}(r,\theta,\phi) = R_{n\ell}(r)\text{Y}_{\ell}^m(\theta,\phi).
\end{equation}
Zu jedem $n$ gibt es mögliche Drehimpulsquantenzahlen $\ell = 0,1,2, \ldots, n-1$ ($\equiv \text{s}, \text{p}, \text{d}, \text{f}, \ldots $). Zu jedem $\ell$ gehören $2\ell +1$ mögliche Werte der magnetischen Quantenzahahl $m = -\ell, -\ell+1, \ldots, \ell$. Eine wichtige Bemerkung: die Energie ist von $m$ und $\ell$ unabhängig. 

Die ersten Laguerre-Polynomen lauten
\begin{equation}
	\begin{split}
		L_0^s(x) &= 1 \\
		L_1^s(x) &= -x+(s+1) \\
		L_2^s(x) &= \frac{x^2}{2}-(s+2)x + \frac{1}{2}(s+1)(s+2) \\
		&\vdots
	\end{split}
\end{equation}
Daher erhalten wir die folgende Eigenwellenfunktionen $\psi_{n\ell m}(r,\theta,\phi)$ bis $n=2$:
\begin{align*}
	\psi_{100}(r,\theta,\phi) &= \frac{1}{\sqrt{\pi}}\left(\frac{Z}{a_0}\right)^{\frac{3}{2}}e^{-\frac{Zr}{a_0}} &\quad & 1\text{s} \\
	\psi_{200}(r,\theta,\phi) &= \frac{1}{4\sqrt{2\pi}}\left(\frac{Z}{a_0}\right)^{\frac{3}{2}}\left(2-\frac{Zr}{a_0}\right)e^{-\frac{Zr}{2a_0}} &\quad & 2\text{s} \\
	\psi_{210}(r,\theta,\phi) &= \frac{1}{4\sqrt{2\pi}}\left(\frac{Z}{a_0}\right)^{\frac{3}{2}}\frac{Zr}{a_0}e^{-\frac{Zr}{2a_0}}\cos(\theta) &\quad & 2\text{p} \quad (m=0) \\
	\psi_{21\pm1}(r,\theta,\phi) &= \mp\frac{1}{8\sqrt{\pi}}\left(\frac{Z}{a_0}\right)^{\frac{3}{2}}\frac{Zr}{a_0}e^{-\frac{Zr}{2a_0}}\sin(\theta)e^{\pm i\phi} &\quad & 2\text{p} \quad (m=\pm1). \\
\end{align*}

Die üblichen bunten Orbitalbilder zeigen nur $|\psi_{n\ell m}|^2$, nicht die
Phase.

%!TEX root = TheoretischePhysikC.tex

\section{Vorlesung: Spin}
Die Drehimpulsoperatoren erfüllen die folgende Vertauschungsrelationen:
\begin{equation}\label{eq:amcr}
	[\widehat{L}_1,\widehat{L}_2] = i\hbar \widehat{L}_3, \quad [\widehat{L}_2,\widehat{L}_3] = i\hbar \widehat{L}_1\quad\text{und}\quad [\widehat{L}_3,\widehat{L}_1] = i\hbar \widehat{L}_2.
\end{equation}
Es folgt, dass
\begin{equation}
	[\widehat{\mathbf{L}}^2, \widehat{L}_j]=0, \quad \text{für alle $j=1,2,3$,}
\end{equation}
wobei
\begin{equation}
	\widehat{\mathbf{L}}^2 \equiv \widehat{L}_1^2+\widehat{L}_2^2+\widehat{L}_3^2.
\end{equation}
Wir haben gesehen, dass die Bahndrehimpulsoperatoren $\widehat{\mathbf{L}} \equiv \widehat{\mathbf{r}}\times \widehat{\mathbf{p}}$ diese Kommutatoren erfüllen. Man sagt, dass $\widehat{\mathbf{r}}\times \widehat{\mathbf{p}}$ uns eine \emph{Darstellung} dieser Kommutatoren liefert. Eine wichtige Eigenschaft dieser Darstellung ist, dass den Gesamtdrehimpuls \emph{quantisiert} ist, im Sinne, dass die Eigenwerte des Gesamtdrehimpulsoperators $\widehat{\mathbf{L}}^2$ nur die Werte $\hbar^2 \ell(\ell+1)$ annehmen können, wobei $\ell = 0, \frac12, 1, \ldots$. Allerdings für die Bahndrehimpulsoperatoren $\widehat{\mathbf{r}}\times \widehat{\mathbf{p}}$ tauchen nur ganzzahlige Werte auf: $\ell = 0,1,2,\ldots$.

Zu jeder Darstellung gehört ein Hilbertraum, auf dem die Operatoren wirken; ausführlich diskutiert haben wir bisher allein $\mathcal{H}=L^2(\mathbb{R}^3)$.

Eine faszinierende Frage bleibt offen: Gibt es Quantensysteme, die halb-ganzzahlige Eigenwerte aufweisen? Um diese Frage zu bejahen benötigen wir einen Hilbertraum $\mathcal{H}$ und drei Operatoren $\widehat{S}_j:\mathcal{H}\rightarrow \mathcal{H}$, $j=1,2,3$, die die Relationen (\ref{eq:amcr}) erfüllen:
\begin{equation}\label{eq:amcr-spin}
	\eqbox{[\widehat{S}_1,\widehat{S}_2] = i\hbar \widehat{S}_3, \quad [\widehat{S}_2,\widehat{S}_3] = i\hbar \widehat{S}_1\quad\text{und}\quad [\widehat{S}_3,\widehat{S}_1] = i\hbar \widehat{S}_2.}
\end{equation}
Falls ein solcher Hilbertraum und solche Operatoren existieren, können wir zumindest feststellen, dass es keine \emph{mathematische} Hindernisse gibt die die Existenz von Quantensystemen mit halb-ganzzahligem gesamten Drehimpuls verhindern. Selbst wenn ein solches Quantensystem mathematisch existieren kann, ist es noch ein weiterer Schritt, festzustellen, ob ein solches System in der Natur vorkommen kann.

Es zeigt sich, dass Quantensysteme mit halbzahligen Eigenwerten für den gesamten Drehimpulsoperator tatsächlich existieren. Wir untersuchen nun den Fall $\ell = \frac12$. Der dazugehörige Hilbertraum ist durch $\mathcal{H} \equiv \mathbb{C}^2$ gegeben. Dass wir die Operatoren jetzt $\widehat{S}_j$ statt $\widehat{L}_j$ nennen, ist eine reine Umbenennung. Wir definieren dazu die Matrizen
\begin{equation}
	\widehat{S}_1 = \frac{\hbar}{2} \begin{pmatrix}
		0 & 1 \\ 1 & 0
	\end{pmatrix}, \quad \widehat{S}_2 = \frac{\hbar}{2} \begin{pmatrix}
		0 & -i \\ i & 0
	\end{pmatrix} \quad \text{und} \quad \widehat{S}_3 = \frac{\hbar}{2} \begin{pmatrix}
		1 & 0 \\ 0 & -1
	\end{pmatrix}.
\end{equation}
Bis auf den Faktor $\frac{\hbar}{2}$ sind das die \emph{Pauli-Matrizen} $\sigma_1,\sigma_2,\sigma_3$. Zu prüfen bleibt, ob sie (\ref{eq:amcr-spin}) erfüllen:
\begin{equation}
	[\widehat{S}_1,\widehat{S}_2] = \frac{\hbar^2}{4}\left(\sigma_1\sigma_2-\sigma_2\sigma_1\right) = \frac{\hbar^2}{4}\cdot 2i \begin{pmatrix}
		1 & 0 \\ 0 & -1
	\end{pmatrix} = i\hbar\, \widehat{S}_3.
\end{equation}
Die beiden übrigen Relationen folgen nach demselben Muster.

\begin{historical}
	Damit ist erst die \emph{mathematische} Frage beantwortet: Man kann wilde Hilberträume definieren, die alle Bedingungen erfüllen und trotzdem kein System in der Natur beschreiben. Beim Spin geht es gut -- im Stern-Gerlach-Experiment, das einen Silberatomstrahl in genau zwei Teilstrahlen aufspaltete, wurde eines gefunden: Das Elektron besitzt einen \emph{internen} Freiheitsgrad mit Gesamtdrehimpuls $\frac12$.
\end{historical}

Der Gesamtdrehimpulsoperator $\widehat{\mathbf{S}}^2$ lautet
\begin{equation}
	\widehat{\mathbf{S}}^2 = \frac{3\hbar^2}{4} \mathbb{I}.
\end{equation}
Die Rechnung passt in eine Zeile, denn $\sigma_1^2 = \sigma_2^2 = \sigma_3^2 = \mathbb{I}$. Dieser Operator besitzt zwei (entartete) Eigenwerte, nämlich, $\frac{3\hbar^2}{4}$ und $\frac{3\hbar^2}{4}$. Vergleichen wir die Wirkung von $\widehat{\mathbf{S}}^2$ mit
\begin{equation}
	\widehat{\mathbf{S}}^2|\ell,m\rangle = \hbar^2\ell(\ell+1)|\ell,m\rangle
\end{equation}
stellen wir fest, dass $\ell=\frac12$ sein muss, denn $\frac34 = \frac12\left(\frac12+1\right)$.

Dass $[\widehat{\mathbf{S}}^2,\widehat{S}_j]=0$ gilt, ist hier trivial: Ein Vielfaches des Einsoperators vertauscht mit allem. Bezüglich dieser Darstellung ist $\widehat{S}_3$ bereits diagonal. Das heisst, die Eigenbasis von $\widehat{S}_3$ darf mit $|\ell,m\rangle$ identifiziert werden:
\begin{equation}
	|\ell=\tfrac12,m=\tfrac12\rangle \equiv \begin{pmatrix}
		1 \\ 0
	\end{pmatrix}\quad\text{und}\quad 
	|\ell=\tfrac12,m=-\tfrac12\rangle \equiv\begin{pmatrix}
		0 \\ 1
	\end{pmatrix}.
\end{equation}
Man schreibt dafür auch kurz $|\uparrow\rangle$ und $|\downarrow\rangle$.

Darstellungen zu $\ell=\frac32,\frac52,\ldots$ lassen sich ebenso angeben, \emph{fundamentale} Teilchen mit $\ell=\frac32$ hat aber noch niemand entdeckt. \emph{Aufbauen} lassen sich solche Drehimpulse allerdings: Ein System mit $\ell=\frac12$ und eines mit $\ell=1$ zusammengebracht sieht von fern aus wie $\ell=\frac32$. Damit sind wir bei zusammengesetzten Systemen.

\section*{Zusammengesetzte Systeme}
Oft müssen wir bei der Modellierung eines quantenmechanischen Systems das System aus mehreren Teilsystemen aufbauen: aus einem Elektron und einem Proton wird dann ein Wasserstoff-Atom und so weiter; oder wir verfeinern unsere Beschreibung durch Hinzunahme weiterer Freiheitsgrade: Aus einem ``Elektron ohne Spin'', wird zusammen mit einer Beschreibung eines isolierten Spinfreiheitsgrads eine Theorie für Elektronen, die den Spin einschließt.

In all diesen Fällen verfahren wir nach der Regel, dass jeder Sorte quantenmechanischer
Systeme ein Hilbertraum zugeordnet wird. Sei also $\mathcal{H}_1$ der Hilbertraum der Systemsorte ``1'' und ebenso $\mathcal{H}_2$. Die Frage lautet also: Wie und in welchem Hilbertraum ``$\mathcal{H}_{\text{$1$\&$2$}}$'' sollen wir die Systemsorte ``1 \& 2'' von Systemen beschreiben, die aus einem Paar von Teilchen/Freiheitsgraden der Sorten 1 und 2 besteht? Die Antwort ist mit Hilfe durch das \emph{Tensorprodukt} $\mathcal{H}_1\otimes \mathcal{H}_2$ von $\mathcal{H}_1$ und $\mathcal{H}_2$ gegeben.

Aus zwei Hilberträumen lässt sich auf mehrere Weisen ein dritter bauen; nur eine davon ist die physikalisch richtige. Die Forderung lautet: Die Teilsysteme müssen sich \emph{unabhängig} voneinander präparieren lassen -- ob neben dem einen Wasserstoffatom noch ein zweites steht, darf keine Rolle spielen. Also muss jede Kombination von Zuständen einem Vektor in $\mathcal{H}_{\text{$1$\&$2$}}$ entsprechen, und daraus folgt
\begin{equation}
	\dim\left(\mathcal{H}_{\text{$1$\&$2$}}\right) = \dim(\mathcal{H}_1)\cdot \dim(\mathcal{H}_2):
\end{equation}
Die Dimensionen \emph{multiplizieren} sich, sie addieren sich nicht -- wie bei klassischen Bits, wo drei Bits acht Zustände besitzen.

\begin{center}
\parbox[c][5cm][c]{14cm}{\begin{tcolorbox}[colback=banana,arc=0cm,frame hidden,enhanced]{Das Tensorprodukt $\mathcal{H}_1\otimes \mathcal{H}_2$ zweier Hilberträume $\mathcal{H}_1$ und $\mathcal{H}_2$ ist \emph{der} Hilbertraum, der Vektoren $\phi_1\otimes \phi_2$ mit $\phi_1\in \mathcal{H}_1$ und $\phi_2\in \mathcal{H}_2$ enthält, und von ihnen erzeugt wird, und in dem die Rechenregel 
\begin{equation}
	\langle \phi_1\otimes \phi_2|\psi_1\otimes\psi_2\rangle \equiv \langle\phi_1| \psi_1\rangle\langle\phi_2| \psi_2\rangle
\end{equation}
für Skalarprodukte gilt.
}\end{tcolorbox}}
\end{center}

Betrachten wir also als nächstes das Tensorprodukt von Operatoren. Für $A_1:\mathcal{H}_1 \rightarrow \mathcal{H}_1$ und $A_2:\mathcal{H}_2 \rightarrow \mathcal{H}_2$ definieren wir den Operator $A_1\otimes A_2:\mathcal{H}_1 \otimes\mathcal{H}_2 \rightarrow \mathcal{H}_1 \otimes\mathcal{H}_2$ durch die Gleichung
\begin{equation}
	\eqbox{(A_1\otimes A_2)(\phi_1\otimes \phi_2) \equiv (A_1\phi_1)\otimes (A_2\phi_2).}
\end{equation}
Da $A_1\otimes A_2$ ja ein linearer Operator sein soll, ist er damit auch auf allen
Linearkombinationen von Produktvektoren definiert.

Ein Tensorprodukt von Operatoren lässt sich natürlich auch durch seine Matrixelemente angeben. Direkt aus der Definition folgt
\begin{equation}
	\begin{split}
			[A_1\otimes A_2]_{jj',kk'} &\equiv \langle \phi_j\otimes \phi_{j'}|\left(A_1\otimes A_2\right) \phi_k\otimes\phi_{k'}\rangle \\
			&= \langle \phi_j\otimes \phi_{j'}| (A_1\phi_k)\otimes (A_2\phi_{k'})\rangle \\
			&= \langle \phi_j|A_1\phi_k\rangle \langle \phi_{j'}|A_2\phi_{k'}\rangle \\
			&= [A_1]_{j,k}[A_2]_{j',k'}.
	\end{split}
\end{equation}
Die Indizes an der Produktmatrix, sind also Paare von Indizes. Das Tensorprodukt einer $3\times 3$-Matrix mit einer $7\times 7$-Matrix ergibt also eine $21\times 21$-Matrix. Zur numerischen Verarbeitung solcher Produkte ist es nötig, aus den Paaren von Indizes wieder einfache Indizes zu machen, also alle möglichen Paare irgendwie durchzunummerieren.

Rechenregeln für das Tensorprodukt von Operatoren folgen meist leichter aus der Definition. Wir bekommen neben der Linearität in jedem Faktor, die für alle Tensorprodukte (auch von Vektoren) gilt:
\begin{align}
	(A_1\otimes A_2)(B_1\otimes B_2) &= A_1B_1\otimes A_2B_2 \\
	(A_1\otimes A_2)^\dag &= A_1^\dag \otimes A_2^\dag \\
	\text{$A_1\ge 0$ und $A_2\ge 0$} &\Rightarrow A_1\otimes A_2 \ge 0 \\
	\|A_1\otimes A_2\| &= \|A_1\|\| A_2\| \\
	\tr(A_1\otimes A_2) &= \tr(A_1)\tr(A_2).
\end{align}

Schließlich sollte noch erwähnt werden, dass alles, was hier über das Tensorprodukt aus zwei Faktoren gesagt wurde, sofort auf endlich viele Faktoren übertragbar ist.

\section*{Dynamik zusammengesetzter Systeme}
Laufen beide Teilsysteme unabhängig voneinander, so ist die Zeitentwicklung des Gesamtsystems durch $U_t = U^{(1)}_t\otimes U^{(2)}_t$ gegeben. Beim Hamiltonoperator lauert dagegen eine Falle: Man würde vermuten, er sei das Produkt $\widehat{H}_1\otimes\widehat{H}_2$ -- das ist falsch. Richtig ist
\begin{equation}\label{eq:H-tensor}
	\eqbox{\widehat{H} = \widehat{H}_1\otimes \mathbb{I} + \mathbb{I}\otimes \widehat{H}_2.}
\end{equation}
Der Grund ist die Produktregel: Ableiten von $U_t$ nach der Zeit gibt
\begin{equation}
	\frac{d}{dt}\left(U^{(1)}_t\otimes U^{(2)}_t\right) = \left(\frac{dU^{(1)}_t}{dt}\right)\otimes U^{(2)}_t + U^{(1)}_t\otimes\left(\frac{dU^{(2)}_t}{dt}\right),
\end{equation}
und daraus liest man mit $i\hbar\frac{dU^{(j)}_t}{dt} = \widehat{H}_j U^{(j)}_t$ und $i\hbar\frac{dU_t}{dt} = \widehat{H}U_t$ genau (\ref{eq:H-tensor}) ab. Energien addieren sich, sie multiplizieren sich nicht. Mit gelösten Teilproblemen $\widehat{H}_1|E_j\rangle = E_j|E_j\rangle$ und $\widehat{H}_2|E_k\rangle = E_k|E_k\rangle$ ist auch $\widehat{H}$ diagonalisiert:
\begin{equation}
	\begin{split}
		\widehat{H}\left(|E_j\rangle\otimes |E_k\rangle\right) &= \left(\widehat{H}_1|E_j\rangle\right)\otimes |E_k\rangle + |E_j\rangle \otimes \left(\widehat{H}_2|E_k\rangle\right) \\
		&= (E_j+E_k)\, |E_j\rangle\otimes |E_k\rangle.
	\end{split}
\end{equation}
\begin{übung}
	Zeigen Sie, dass diese Produktvektoren eine \emph{vollständige} Eigenbasis von $\widehat{H}$ bilden.
\end{übung}

\section*{Verschränkung}
Eine letzte Bemerkung, und sie ist die folgenreichste. Für $\mathcal{H} = \mathbb{C}^2\otimes\mathbb{C}^2$ mit Orthonormalbasen $\{\psi_1,\psi_2\}$, $\{\phi_1,\phi_2\}$ ist jeder reine Zustand eine Linearkombination der vier Produktvektoren $\psi_j\otimes\phi_k$ -- aber nicht jede ist selbst ein Produkt. Ein Beispiel ist der Einstein-Podolsky-Rosen-Zustand ($\psi_1=\phi_1=|\uparrow\rangle$, $\psi_2=\phi_2=|\downarrow\rangle$):
\begin{equation}
	|\text{EPR}\rangle = \frac{1}{\sqrt{2}}\left(\psi_1\otimes\phi_1 + \psi_2\otimes\phi_2\right) = \frac{1}{\sqrt{2}}\left(|\uparrow\uparrow\rangle + |\downarrow\downarrow\rangle\right).
\end{equation}
Er lässt sich nicht als $\psi\otimes\phi$ schreiben: Ein solches Produkt $(a_1\psi_1+a_2\psi_2)\otimes(b_1\phi_1+b_2\phi_2)$ hat die Koeffizienten $a_jb_k$, und aus $a_1b_2 = a_2b_1 = 0$ folgt $a_1b_1 = 0$ oder $a_2b_2 = 0$ -- beide können nicht $\tfrac{1}{\sqrt{2}}$ sein. Also lässt er sich nicht als unabhängige Präparation zweier Systeme auffassen -- obwohl es gerade die unabhängige Präparierbarkeit war, die uns zum Tensorprodukt geführt hat. Das Tensorprodukt liefert also mehr, als wir von ihm verlangt haben.

% fig_vl17_bloch.tex — Spin 1/2: die Blochkugel eines einzelnen Spins, und
% was aus dem Blochpfeil wird, sobald der Spin zu einem verschränkten Paar
% gehört.
%
% Follows the conventions of fig_vl11_kasten (the reference plate).
%
%   * geometry comes from the formulae of the lecture, never freehand.
%
%     Panel (a).  Every pure state of H = C^2 is, up to a global phase,
%         |psi(theta,phi)> = cos(theta/2)|up> + e^{i phi} sin(theta/2)|down>,
%     and with the S_j of the lecture its vector of expectation values is
%         (2/hbar)(<S_1>,<S_2>,<S_3>)
%             = (sin th cos ph, sin th sin ph, cos th),
%     a unit vector.  The panel therefore draws the unit sphere of that
%     vector, and every marked point sits where the formula puts it:
%     |up>,|down> at th = 0,pi — antipodal, not perpendicular, because the
%     half angle th/2 turns the 90 degrees of orthogonality in C^2 into
%     180 degrees in R^3; |+>,|-> at th = pi/2, ph = 0,pi;  |+i>,|-i> at
%     th = pi/2, ph = +-pi/2.  The projection is a genuine orthographic
%     one, camera azimuth psi = 25 deg and elevation eps with sin eps = 0.3:
%         x^ -> (-0.42262, -0.27189)
%         y^ -> ( 0.90631, -0.12679)
%         z^ -> ( 0.00000,  0.95394)
%     Consequences that are therefore forced rather than drawn by eye: the
%     equator projects to an ellipse of axis ratio 0.3; its *near* half is
%     the lower one, hence solid, the far half dashed; and the poles land
%     at 0.954 R, just inside the outline, which is what an orthographic
%     projection of a sphere actually does.  \tpcbx/\tpcby implement the
%     projection, so no coordinate in the panel was typed by hand.
%
%     Panel (b).  The one-parameter family
%         |psi_a> = cos(a)|up up> + sin(a)|down down>
%     runs from a product state through the EPR state of the lecture
%     (a = pi/4) to another product state.  Both curves are exact:
%         (2/hbar)|<S (x) 1>| = |cos^2 a - sin^2 a| = |cos 2a|,
%             the length of spin 1's arrow in panel (a) — for this family
%             <S_1 (x) 1> = <S_2 (x) 1> = 0 and <S_3 (x) 1> = (hbar/2)cos 2a,
%             all three computable with the tensor-product rules of the
%             lecture alone;
%         2|c_11 c_22 - c_12 c_21| = 2|cos a sin a| = |sin 2a|,
%             exactly the obstruction the lecture computes: a product needs
%             a_1 b_2 = a_2 b_1 = 0, i.e. a vanishing determinant of the
%             coefficient matrix.
%     Since cos^2 + sin^2 = 1 the two readings trade off exactly, and the
%     crossing points at a = pi/8, 3pi/8 are at 1/sqrt(2).  At the EPR
%     point the arrow of panel (a) has length zero: it has retreated from
%     the surface of the ball to its centre.  That is the whole content of
%     the plate — the two panels are one picture, read twice.
%
%   * colour carries the observable and nothing else: eigenstates of S_3
%     navy, of S_1 blue, of S_2 teal (the ordered ramp of `tpclevels`);
%     the one state that is not a basis state, the general |psi>, is
%     orange, and so is the entanglement curve in (b) — the blue curve in
%     (b) is the length of a blue-and-navy object, the arrow.  Axes,
%     equator and the EPR rule are neutral grey: reference structure.
%   * near/far is coded by fill, not by colour: a point of the near
%     hemisphere is a filled disc, one of the far hemisphere an open
%     circle, and the far halves of the axes are dashed.  This survives
%     greyscale printing and makes the antipodal pairs read as pairs.
%   * (b) separates its curves by dash pattern as well as by colour, for
%     the same reason.
%
% Panel (a) the state space of one spin 1/2.
% Panel (b) what becomes of that arrow when the spin is one of a pair.

% Orthographic projection of the sphere of radius #1, see the header.
% #2 = polar angle theta, #3 = azimuth phi, both in degrees.
\def\tpcbx#1#2#3{#1*(-0.42262*sin(#2)*cos(#3)+0.90631*sin(#2)*sin(#3))}
\def\tpcby#1#2#3{#1*(-0.27189*sin(#2)*cos(#3)-0.12679*sin(#2)*sin(#3)%
                     +0.95394*cos(#2))}
% \tpcbc{name}{radius}{theta}{phi} — name a point on that sphere.
\newcommand{\tpcbc}[4]{%
  \coordinate (#1) at ({\tpcbx{#2}{#3}{#4}},{\tpcby{#2}{#3}{#4}});}
\def\tpcRb{1.7}

\begin{center}
\begin{tikzpicture}[
    % near hemisphere: filled disc;  far hemisphere: open circle
    sdot/.style ={circle, inner sep=0pt, minimum size=3.4pt,
                  fill=#1, draw=#1},
    sdot/.default=cgblue,
    sdotb/.style={circle, inner sep=0pt, minimum size=3.4pt,
                  fill=white, draw=#1, line width=0.65pt},
    sdotb/.default=cgblue,
    blochaxis/.style={-{Stealth[length=5pt,width=3.5pt]}, birrengray,
                      line width=0.7pt},
    blochback/.style={birrengray, line width=0.7pt, densely dashed},
  ]

% ─────────────────────────────────────────────────────────────
%  (a)  The Bloch sphere of one spin 1/2
% ─────────────────────────────────────────────────────────────
\node[font=\sf\small\bfseries, text=cgblue] at (0,3.00)
  {(a) \tpcifenglish{one spin: the Bloch sphere}{ein Spin: die Blochkugel}};

% Every named point is placed by the parametrisation, nothing by eye.
\coordinate (O) at (0,0);
\tpcbc{Sup}  {\tpcRb}{0}  {0}      % |up>      = eigenstate of S_3, +
\tpcbc{Sdn}  {\tpcRb}{180}{0}      % |down>    = eigenstate of S_3, -
\tpcbc{Splu} {\tpcRb}{90} {0}      % |+>       = eigenstate of S_1, +
\tpcbc{Smin} {\tpcRb}{90} {180}    % |->       = eigenstate of S_1, -
\tpcbc{Spi}  {\tpcRb}{90} {90}     % |+i>      = eigenstate of S_2, +
\tpcbc{Smi}  {\tpcRb}{90} {270}    % |-i>      = eigenstate of S_2, -
\tpcbc{Spsi} {\tpcRb}{55} {-55}    % a general |psi(theta,phi)>
\tpcbc{Sfoo} {\tpcRb*sin(55)}{90}{-55}   % its foot in the equatorial plane
% Axis tips, all three at the same distance on the page (1.95 cm).
\tpcbc{Axa}  {3.880}{90}{0}
\tpcbc{Axb}  {2.131}{90}{90}
\tpcbc{Axc}  {2.044}{0} {0}
% Far ends of the three axes, just past the far surface.
\tpcbc{Bxa}  {1.05*\tpcRb}{90} {180}
\tpcbc{Bxb}  {1.05*\tpcRb}{90} {270}
\tpcbc{Bxc}  {1.05*\tpcRb}{180}{0}

% Far halves of the axes: dashed, they run behind the ball
\draw[blochback] (O) -- (Bxa);
\draw[blochback] (O) -- (Bxb);
\draw[blochback] (O) -- (Bxc);

% Equator: far (upper) half dashed, near (lower) half solid.  The axis
% ratio 0.3 = sin(eps) is forced by the projection, not chosen.
\draw[birrengray, line width=0.55pt, densely dashed]
  (-\tpcRb,0) arc (180:0:{\tpcRb} and {0.3*\tpcRb});
\draw[birrengray, line width=0.55pt]
  (-\tpcRb,0) arc (180:360:{\tpcRb} and {0.3*\tpcRb});

% Outline
\draw[birrengray, line width=0.9pt] (0,0) circle (\tpcRb);

% Near halves of the axes
\draw[blochaxis] (O) -- (Axa);
\draw[blochaxis] (O) -- (Axb);
\draw[blochaxis] (O) -- (Axc);
\node[lbl, anchor=north east, xshift=-1pt] at (Axa) {$\widehat{S}_1$};
\node[lbl, anchor=west, xshift=2pt]        at (Axb) {$\widehat{S}_2$};
\node[lbl, anchor=south, yshift=1pt]       at (Axc) {$\widehat{S}_3$};

% The general state: its arrow, its foot in the equatorial plane, and the
% two angles of the parametrisation.
\draw[tpc thin, birrengray, densely dotted] (O)    -- (Sfoo);
\draw[tpc thin, birrengray, densely dotted] (Sfoo) -- (Spsi);
\draw[-{Stealth[length=6pt,width=4pt]}, birrenorange, line width=1.1pt]
  (O) -- (Spsi);

\draw[birrengray, line width=0.55pt]
  plot[domain=0:55, samples=30, variable=\s]
    ({\tpcbx{0.62}{\s}{-55}},{\tpcby{0.62}{\s}{-55}});
\node[lbl, font=\scriptsize] at (-0.37,0.66) {$\theta$};

\draw[birrengray, line width=0.55pt]
  plot[domain=0:-55, samples=30, variable=\p]
    ({\tpcbx{0.55}{90}{\p}},{\tpcby{0.55}{90}{\p}});
\node[lbl, font=\scriptsize] at (-0.70,-0.25) {$\varphi$};

% The six eigenstates of S_1, S_2, S_3.  Colour = observable;
% filled = near hemisphere, open = far hemisphere.
\node[sdot=spacecadet]  at (Sup)  {};
\node[sdotb=spacecadet] at (Sdn)  {};
\node[sdot=cgblue]      at (Splu) {};
\node[sdotb=cgblue]     at (Smin) {};
\node[sdot=munsell]     at (Spi)  {};
\node[sdotb=munsell]    at (Smi)  {};

\node[lbl, anchor=west, xshift=3pt]       at (Sup)  {$\ket{\uparrow}$};
\node[lbl, anchor=west, xshift=3pt]       at (Sdn)  {$\ket{\downarrow}$};
\node[lbl, anchor=east, xshift=-4pt]      at (Splu) {$\ket{+}$};
\node[lbl, anchor=south west, xshift=2pt] at (Smin) {$\ket{-}$};
\node[lbl, font=\scriptsize, text=birrengray, anchor=north, yshift=-4pt]
                                          at (Spi)  {$\ket{{+}i}$};
\node[lbl, font=\scriptsize, text=birrengray, anchor=south, yshift=4pt]
                                          at (Smi)  {$\ket{{-}i}$};
\node[lbl, text=birrenorange, anchor=south east, yshift=1pt]
                                          at (Spsi) {$\ket{\psi}$};

% The formula the whole panel is drawn from, and how to read it.
\node[lbl, font=\scriptsize] at (0,-2.38)
  {$\ket{\psi}=\cos\tfrac{\theta}{2}\,\ket{\uparrow}
    +e^{i\varphi}\sin\tfrac{\theta}{2}\,\ket{\downarrow}$};
\node[lbl, font=\scriptsize, text=birrengray] at (0,-2.82)
  {\tpcifenglish{half angle: orthogonal $\Rightarrow$ antipodal}%
                {halber Winkel: orthogonal $\Rightarrow$ antipodal}};

% ─────────────────────────────────────────────────────────────
%  (b)  The same arrow, for one spin of an entangled pair
% ─────────────────────────────────────────────────────────────
\begin{axis}[
    tpcplot,
    at={(4.30cm,-1.70cm)}, anchor=south west,
    width=7.0cm, height=4.9cm,
    clip=false,
    xmin=0, xmax=90, ymin=0, ymax=1.42,
    xtick={0,22.5,45,67.5,90},
    xticklabels={$0$,$\pi/8$,$\pi/4$,$3\pi/8$,$\pi/2$},
    ytick={0,0.5,1},
    xlabel={$\alpha$},
    title={(b) \tpcifenglish{a pair: arrow vs.\ entanglement}%
                            {ein Paar: Pfeil gegen Verschränkung}},
    legend style={font=\scriptsize, draw=light-gray,
                  at={(0.5,-0.26)}, anchor=north, legend columns=2},
]

  % The EPR point alpha = pi/4: a neutral reference rule
  \addplot[levellinegrey] coordinates {(45,0) (45,1)};

  % Length of spin 1's Bloch arrow
  \addplot[cgblue, line width=1.0pt, domain=0:90, samples=181]
    {abs(cos(2*x))};
  \addlegendentry{$\tfrac{2}{\hbar}\bigl|\langle\widehat{\mathbf{S}}\otimes\mathbb{I}\rangle\bigr|=|\cos 2\alpha|$}

  % Failure of the product condition c_11 c_22 = c_12 c_21
  \addplot[birrenorange, line width=1.0pt, densely dashed,
           domain=0:90, samples=181] {abs(sin(2*x))};
  \addlegendentry{$2|c_{11}c_{22}-c_{12}c_{21}|=|\sin 2\alpha|$}

  \addplot[only marks, mark=*, mark size=2pt, cgblue, draw=none]
    coordinates {(45,0)};
  \addplot[only marks, mark=*, mark size=2pt, birrenorange, draw=none]
    coordinates {(45,1)};

  % The family the panel is drawn from
  \node[lbl, font=\scriptsize] at (axis cs:45,1.26)
    {$\ket{\psi_\alpha}=\cos\alpha\,\ket{\uparrow\uparrow}
      +\sin\alpha\,\ket{\downarrow\downarrow}$};

  % What the three marked values of alpha are
  \node[lbl, font=\scriptsize, text=birrengray, anchor=north]
    at (axis cs:0,-0.15)  {$\ket{\uparrow\uparrow}$};
  \node[lbl, font=\scriptsize, text=birrengray, anchor=north]
    at (axis cs:45,-0.15) {$\ket{\text{EPR}}$};
  \node[lbl, font=\scriptsize, text=birrengray, anchor=north]
    at (axis cs:90,-0.15) {$\ket{\downarrow\downarrow}$};
\end{axis}

\end{tikzpicture}
\end{center}


Der Pfeil links ist der Erwartungswert $\tfrac{2}{\hbar}\langle\widehat{\mathbf{S}}\rangle$ eines einzelnen Spins. Rechts sieht man ihn schrumpfen, sobald der Spin verschränkt wird: beim EPR-Zustand ist er in den Mittelpunkt der Kugel zurückgewichen -- dort, wo die gestrichelte Kurve ihr Maximum hat.

Solche Zustände heissen \emph{verschränkt}, und sie werden heute im Labor routinemässig erzeugt.

%!TEX root = TheoretischePhysikC.tex

\section{Vorlesung: Identische Teilchen}
Alle bekannten Elementarteilchen einer Sorte sind identisch: zwei Elektronen unterscheiden sich durch nichts. Das ist kein Ergebnis der Theorie, sondern ein empirischer Befund -- es müsste nicht so sein, aber es ist so. Zusammengesetzte Teilchen wie Atome sind dagegen in der Regel unterscheidbar, denn sie tragen innere Freiheitsgrade. Entscheidend ist nun, dass ,,identisch`` keine neue Wechselwirkung ist, sondern eine \emph{Einschränkung an die erlaubten Zustände}: am Ende bleiben zwei Unterräume von $\mathcal{H}^{\otimes N}$ übrig, die Bosonen und Fermionen beschreiben.

Sei $\mathcal{H}$ der Hilbertraum eines Teilchens. Nach unseren bisherigen Regeln ist
\begin{equation}
  \mathcal{H}^{\otimes N} \underbrace{\equiv \mathcal{H}\otimes \cdots \otimes \mathcal{H}}_{\text{$N$ Faktoren}}
\end{equation}
der Hilbertraum für $N$ solcher Teilchen. Um eine \emph{Vertauschung} darzustellen führen wir eine Symmetrieoperation für diese Vertauschung ein: Sei also $\pi$ eine \emph{Permutation} der $N$ Symbole $\{1,2, \ldots,N\}$, das heißt eine bijektive Abbildung dieser Menge in sich. Bis hierher ist das reine Mathematik. Physik wird daraus erst mit Postulat 5, nach dem physikalische Prozesse durch unitäre Operatoren beschrieben werden: Eine Vertauschung ist erst dann physikalisch realisierbar, wenn wir einen unitären Operator angeben können, der sie darstellt. Zu jeder solchen \emph{Umordnungsabbildung} definieren
wir einen unitären Operator $U_\pi$ auf $\mathcal{H}^{\otimes N}$ durch
\begin{equation}
  U_\pi (\phi_1\otimes\phi_2\otimes \cdots\otimes \phi_N) \equiv \phi_{\pi^{-1}(1)}\otimes\phi_{\pi^{-1}(2)}\otimes \cdots\otimes \phi_{\pi^{-1}(N)}.
\end{equation}
Wie üblich ist dieser Operator durch Linearität von den Produktvektoren auf ganz $\mathcal{H}^{\otimes N}$ fortzusetzen. Beachten Sie die \emph{inverse} Permutation in dieser Definition -- $U_\pi$ schiebt die Faktoren an neue Plätze, er bildet nicht die Indizes ab. Erfahrungsgemäß ist das hier die häufigste Verwirrungsquelle.

Jede Permutation lässt sich als Produkt von Paarvertauschungen $\theta_{jk}$ schreiben, wobei $\theta_{jk}$
die Vertauschung von $j$ und $k\not=j$ bezeichnet. Das heißt, wir können immer eine beliebige Permutation $\pi$ als Produkt $\pi = \theta_{j_mk_m}\cdots\theta_{j_2k_2}\theta_{j_1k_1}$ schreiben. Mit Hilfe von dieser Zerlegung führen wir die Funktion $\text{sign}(\pi)$ ein als das ,,Vorzeichen`` einer Permutation, also genau
dann $-1$, wenn $\pi$ das Produkt einer ungeraden Anzahl von Paarvertauschungen ist und $+1$, wenn diese Anzahl gerade ist. Dass das wohldefiniert ist, ist nicht selbstverständlich: dieselbe Permutation lässt sich verschieden zerlegen, eindeutig ist nur die \emph{Parität} der Anzahl.

\begin{beispiel}
  Für drei Qubits, also $\mathcal{H}\equiv \mathbb{C}^2$, hat $\mathcal{H}^{\otimes 3}$ die $2^3 = 8$ Basisvektoren
  \begin{equation}
    |000\rangle,\ |001\rangle,\ |010\rangle,\ |011\rangle,\ |100\rangle,\ |101\rangle,\ |110\rangle,\ |111\rangle.
  \end{equation}
  Die Inverse einer Paarvertauschung ist sie selbst, $\theta_{12}^{-1} = \theta_{12}$, also vertauscht $U_{\theta_{12}}$ die ersten beiden Faktoren:
  \begin{equation}
    \begin{array}{llll}
      |000\rangle \mapsto |000\rangle, & |001\rangle \mapsto |001\rangle, & |010\rangle \mapsto |100\rangle, & |011\rangle \mapsto |101\rangle,\\
      |100\rangle \mapsto |010\rangle, & |101\rangle \mapsto |011\rangle, & |110\rangle \mapsto |110\rangle, & |111\rangle \mapsto |111\rangle.
    \end{array}
  \end{equation}
  Vier Basisvektoren bleiben stehen, zwei Paare werden getauscht; als Matrix ist das
  \begin{equation}
    U_{\theta_{12}} = \begin{pmatrix}
      1&0&0&0&0&0&0&0\\
      0&1&0&0&0&0&0&0\\
      0&0&0&0&1&0&0&0\\
      0&0&0&0&0&1&0&0\\
      0&0&1&0&0&0&0&0\\
      0&0&0&1&0&0&0&0\\
      0&0&0&0&0&0&1&0\\
      0&0&0&0&0&0&0&1
    \end{pmatrix}.
  \end{equation}
  In jeder Zeile und jeder Spalte steht genau eine $1$ -- so sieht jede Permutationsmatrix aus. Die übrigen vier Permutationen sind eine gute Übung.
\end{beispiel}

Damit lässt sich sagen, welche Zustände identische Teilchen annehmen dürfen. Nach Postulat 2 werden Zustände mit Dichteoperatoren beschrieben; $\rho$ heißt \emph{invariant}, wenn
\begin{equation}
  U_\pi \rho\, U_\pi^\dagger = \rho \qquad \text{für alle Permutationen $\pi$}
\end{equation}
gilt, und genau diese sind die erlaubten. Im Drei-Qubit-Beispiel sieht man sofort zwei davon, nämlich $|000\rangle\langle 000|$ und $|111\rangle\langle 111|$ -- es wäre aber ein Fehler zu glauben, das seien schon alle.

\begin{beispiel}
  Wir nehmen zwei Qubits und darin den \emph{Singulett-Zustand}
  \begin{equation}
    |\psi^-\rangle \equiv \frac{1}{\sqrt{2}}\left(|01\rangle - |10\rangle\right),
  \end{equation}
  eine geschickt gewählte Superposition. Die Vertauschung $\theta_{12}$ tauscht $|01\rangle$ und $|10\rangle$, also
  \begin{equation}
    U_{\theta_{12}}|\psi^-\rangle = \frac{1}{\sqrt{2}}\left(|10\rangle - |01\rangle\right) = -|\psi^-\rangle.
  \end{equation}
  Der \emph{Vektor} ist demnach nicht invariant. Der \emph{Zustand} ist es trotzdem: im Dichteoperator tritt das Vorzeichen zweimal auf und kürzt sich weg:
  \begin{equation}
    U_{\theta_{12}}|\psi^-\rangle\langle \psi^-|U_{\theta_{12}}^\dagger = (-1)(-1)|\psi^-\rangle\langle \psi^-| = |\psi^-\rangle\langle \psi^-|.
  \end{equation}
\end{beispiel}

Diese kleine Rechnung ist der Kern der Vorlesung. Weil die globale Phase eines Vektors physikalisch bedeutungslos ist und allein der Dichteoperator zählt, gibt es \emph{zwei} Sorten erlaubter Zustände statt einer: exakt invariante und solche, die es nur bis auf ein Vorzeichen sind. Ohne dieses Argument gäbe es Fermionen nicht.

Man kann nun überprüfen, dass die Operatoren
\begin{align}
  P_{+} &\equiv \frac{1}{N!}\sum_{\pi} U_\pi \\
  P_{-} &\equiv \frac{1}{N!}\sum_{\pi} \text{sign}(\pi)U_\pi 
\end{align}
die Projektoren auf die Teilräume von Vektoren sind, die auf den
Raum 
\begin{equation}
  P_{+}\mathcal{H}^{\otimes N} \equiv \mathcal{H}^{\otimes N}_{+}
\end{equation}
von \emph{symmetrischen Vektoren} und auf den Raum
\begin{equation}
  P_{-}\mathcal{H}^{\otimes N} \equiv \mathcal{H}^{\otimes N}_{-}
\end{equation}
von antisymmetrischen Vektoren projektieren. Der Hilbertraum für $N$ gleicher Elementarteilchen ist entweder $\mathcal{H}^{\otimes N}_{+}$ für sogenannte \emph{Bosonen} oder $\mathcal{H}^{\otimes N}_{-}$ für \emph{Fermionen}.

Äquivalent, und meist anschaulicher, sind dieselben Räume direkt beschrieben:
\begin{align}
  \mathcal{H}^{\otimes N}_{+} &= \left\{|\psi\rangle \in \mathcal{H}^{\otimes N}\ \middle|\ U_\pi|\psi\rangle = |\psi\rangle\ \text{für alle $\pi$}\right\},\\
  \mathcal{H}^{\otimes N}_{-} &= \left\{|\psi\rangle \in \mathcal{H}^{\otimes N}\ \middle|\ U_\pi|\psi\rangle = \text{sign}(\pi)|\psi\rangle\ \text{für alle $\pi$}\right\}.
\end{align}
Für den Singulett-Zustand passt das exakt, denn $\text{sign}(\theta_{12}) = -1$. Rechnen Sie als Übung nach, dass für zwei Qubits $\dim \mathcal{H}^{\otimes 2}_{+} = 3$ und $\dim \mathcal{H}^{\otimes 2}_{-} = 1$ gilt.

\begin{center}
\parbox[c][5cm][c]{14cm}{\begin{tcolorbox}[colback=banana,arc=0cm,frame hidden,enhanced]{\textbf{Postulat 6}: Als Wellenfunktionen eines Vielteilchensystems kommen nur solche Funktionen in Frage, die gemeinsame Eigenfunktionen aller Vertauschungen von je zwei gleichen Elementarteilchen sind. Der Eigenwert zu einer solchen Vertauschung ist
der gleiche Faktor, mit dem die Wellenfunktion eines Teilchens dieser Sorte bei der Drehung um $2\pi$ multipliziert wird, also $+1$ für Teilchen mit ganzzahligem Spin und $-1$ für Teilchen mit halbzahligem Spin.}\end{tcolorbox}}
\end{center}

Lesen Sie das Postulat gegen das eben Erarbeitete: die Vektoren aus $\mathcal{H}^{\otimes N}_{\pm}$ sind gemeinsame Eigenvektoren aller $U_\pi$, mit Eigenwert $+1$ im bosonischen und $\text{sign}(\pi)$ im fermionischen Fall. Dass dieser Eigenwert ausgerechnet mit dem Spin zusammenhängt, ist keineswegs offensichtlich -- es ist ein tiefliegendes Ergebnis der Quantenfeldtheorie, der \emph{Spin-Statistik-Zusammenhang}. Beweisen lässt sich Postulat 6 nicht; es muss angenommen werden. Damit sind die sechs Postulate beisammen, und wir können im Prinzip jedes System der Quantenmechanik modellieren.

\section*{Statistische Physik}
,,Im Prinzip`` ist dabei die entscheidende Einschränkung. Wir betrachten nun makroskopische Objekte, die aus üblicherweise $10^{19}$-$10^{23}$ Teilchen pro $\text{cm}^3$ bestehen. Daher ist es praktisch nicht möglich, den \emph{mikroskopischen Zustand} -- den Quantenmechanischen Zustand des Gesamtsystems -- aller Teilchen zu messen oder sogar anzugeben. 

\begin{greybox}[Warum ein Kubikzentimeter schon zu viel ist]
Modellieren wir jedes Teilchen so einfach wie irgend möglich, als Qubit mit $\mathcal{H}\equiv \mathbb{C}^2$. Weil sich die Dimensionen im Tensorprodukt multiplizieren, hat der Hilbertraum eines Balls aus $N\approx 10^{23}$ Teilchen dann die Dimension $2^N = 2^{10^{23}}$; so viele komplexe Zahlen wären schon für einen \emph{reinen} Zustand hinzuschreiben. Auch ein Quantenrechner hilft nicht: er bräuchte ebenso viele Qubits, wie das System Teilchen hat.
\end{greybox}

Eine vollständige Beschreibung des Systemzustands ist zwar nicht möglich, aber das hindert uns nicht daran, interessante physikalische Vorhersagen zu treffen. Das liegt daran, dass wir in der Praxis nur makroskopische Größen wie Druck, Volumen usw. messen können. Deshalb suchen wir nun nach einer grobkörnigen Beschreibung des Systems, die nur die makroskopisch relevanten Größen beschreibt. Die Angabe von einigen wenigen makroskopischen Größen wird als \emph{Makrozustand} bezeichnet. Vorsicht: ein Makrozustand ist \emph{kein} Quantenzustand, sondern nur eine Liste von Zahlen -- er wird weder in Ket-Notation noch mit einem Dichteoperator geschrieben. Die Wörter klingen ähnlich, die Objekte sind völlig verschieden.

Der tatsächlich realisierte Mikrozustand ist unbekannt, aber es ist gegenenfalls möglich die \emph{Wahrscheinlichkeiten} der Mikrozustände anzugeben, die mit einem Makrozustand verträglich sind. Die Menge aller Mikrozustände, die einen gegebenen Makrozustand repräsentieren, heißt zusammen mit der Angabe der Wahrscheinlichkeiten: \emph{statistisches Ensemble} oder \emph{statistische Gesamtheit}.

\begin{beispiel}
  Angenommen, 10 Münzen liegen in einer Reihe auf dem Tisch. Der Mikrozustand ist durch z.B.\ $KKZZZZKZKZ$ gegeben. D.h.\ der Mikrozustand beschreibt den vollständigen Zustand. Eine mögliche makroskopische Größe ist z.B.\ die Anzahl der Münzen, die Kopf zeigen. In diesem Beispiel lautet dieser Zahl $4$. Der Makrozustand in diesem Fall wird durch dieser Zahl beschrieben. Es gibt mehrere Mikrozustände die zu diesem Makrozustand gehören, z.B.\ $KKKKZZZZZZ$ und $KZKZKZKZZZ$ usw.
\end{beispiel}

Statistische Ensemblen werden in der Quantenmechanik wie folgt beschrieben. Ein statistisches Ensemble ist durch angabe von eine Liste von Zustände $\rho_j$ und dazugehörigen Wahrscheinlichkeiten $p_j$ gegeben: $\{(p_j,\rho_j)\,|\, j=1,2, \ldots, N\}$. Das \emph{Präparierverfahren} $\rho$ das dieses Ensemble darstellt wird mit einem Dichteoperator beschrieben:
\begin{equation}
  \rho \equiv \sum_{j=1}^N p_j \rho_j.
\end{equation}

\begin{beispiel}
  Ein Qubit ($\mathcal{H}\equiv \mathbb{C}^2$) wird in dem folgenden statistischen Ensemble präpariert: Der Zustand $\rho_1 = |0\rangle\langle 0|$ taucht mit Wahrscheinlichkeit $p_1=1/2$ auf und $\rho_2 = |+\rangle\langle +|$ mit $p_2 = 1/2$. Der dazugehörige Zustand wird durch
  \begin{equation}
    \rho = \frac12 |0\rangle\langle 0| + \frac12|+\rangle\langle+| \equiv \begin{pmatrix}
      \frac34 & \frac14 \\ \frac14 & \frac14
    \end{pmatrix}
  \end{equation}
  gegeben.
\end{beispiel}

Zu einem Makrozustand gehören in aller Regel unendlich viele verträgliche Ensembles -- statt $p_1 = p_2 = \frac12$ wären ebenso gut $\frac34$ und $\frac14$ möglich. Welches das richtige ist, sagt uns die Quantenmechanik nicht; dafür brauchen wir ein Prinzip.

Eine makroskopische Größe wird als Erwartungswert einer Observable erhalten: Man identifiziert eine Observable $\widehat{A}$, d.h., einen hermiteschen Operator, der eine physikalische Größe darstellt und berechnet
\begin{equation}
  \langle \widehat{A}\rangle \equiv \tr(\widehat{A} \rho),
\end{equation} 
wobei $\rho$ das statistische Ensemble darstellt.

Ziel der statistischen Physik ist es, Systeme mit vielen Teilchen statistisch zu modellieren, wobei nur die Informationen aus einer Reihe makroskopischer Größen (z.B.\ Energie, Magnetisierung) zur Verfügung stehen. Es gibt eine Vielzahl von Herausforderungen, die es zu bewältigen gilt. Die größte Herausforderung besteht darin, die Wahrscheinlichkeiten zu bestimmen, mit denen Mikrozustände bei einem bestimmten Makrozustand auftreten. Dabei wird die Grundannahme der statistischen Mechanik eingeführt:
\begin{center}
\parbox[c][5cm][c]{14cm}{\begin{tcolorbox}[colback=banana,arc=0cm,frame hidden,enhanced]{In einem isolierten System im thermischen Gleichgewicht sind alle zugänglichen Mikrozustände gleich wahrscheinlich.}\end{tcolorbox}}
\end{center}

Man kann diese Annahme (noch) nicht beweisen. Auf der mikroskopischen Ebene erwarten wir, dass jeder Prozess, der das System von Zustand $X$ in den Zustand $Y$ bringt, reversibel sein sollte, sodass das System genauso leicht von Zustand $Y$ in den Zustand $X$ gelangen kann. In diesem Fall sollte das System keinen Zustand einem anderen bevorzugen. Trotzdem ist es nicht offensichtlich, dass alle der als „zugänglich“ angenommenen Mikrozustände tatsächlich in einer vernünftigen Zeitspanne erreicht werden können. Wir werden in der Tat sehen, dass in einem großen System die Zahl der „zugänglichen“ Mikrozustände so groß ist, dass in einer typischen Lebensdauer nur ein winziger Bruchteil von ihnen tatsächlich vorkommen kann. Was wir also annehmen, ist, dass die Mikrozustände, die tatsächlich über einen hinreichend langen, aber nicht unvorstellbar großen Zeitraum vorkommen, einen repräsentativen Querschnitt darstellen. Wir nehmen an, dass die Übergänge „zufällig“ sind, in dem Sinne, dass es kein Muster gibt, über das wir uns Gedanken machen müssten.

Klassisch ist ,,gleich wahrscheinlich`` dabei harmlos: man zählt die zugänglichen Mikrozustände ab und teilt durch ihre Anzahl. Wegen der Superpositionen gibt es in der Quantenmechanik aber unendlich viele Zustände -- welche sollen wir zählen? Die Antwort verlangt eine neue Größe, die \emph{Entropie}, und mit ihr beginnt die nächste Vorlesung.
%!TEX root = TheoretischePhysikC.tex

\section{Vorlesung: Entropie}
Wie wird der Begriff ,,gleich wahrscheinlich verteilte Mikrozustände`` in der Quantenmechanik definiert? Um diese Frage zu beantworten müssen wir das Konzept von \emph{Entropie} einführen. Diese Größe wird als Maß für die \emph{Informationsinhalt} eines Zustandes betrachtet. Oder als ein Maß an \emph{fehlender Information}. Je mehr Informationsgehalt ein Zustand besitzt, desto weniger wissen wir, welcher Mikrozustand auftritt.

Drei Lesarten laufen nebeneinander her: fehlende Information, Unordnung, Informationsgehalt -- dieselbe Zahl, verschieden gelesen. Nötig wird sie, weil im Hilbertraum das klassische Abzählen der Mikrozustände versagt: wegen der Superpositionen gibt es unendlich viele.

\begin{intuition}
	Ein Speicher fasse zehn Bit. Kennen wir seinen Inhalt nicht, so fehlen uns zehn Bit; stehen alle Bits auf Null, fehlt uns keine Information mehr -- und es ist auch kein Speicherplatz mehr frei.
\end{intuition}

Sei $\mathcal{H} \equiv \mathbb{C}^d$ ein Hilbertraum und $\rho$ ein Dichteoperator, der das statistische Ensemble $\{(p_j,\rho_j)\,|\, j=1,2, \ldots, N\}$ darstellt:
\begin{equation}
  \rho \equiv \sum_{j=1}^N p_j \rho_j.
\end{equation}

Die von Neumann-Entropie von $\rho$ wird wie folgt definiert
\begin{equation}
  \eqbox{S(\rho) = -k\tr(\rho\log(\rho)) = -k\langle \log(\rho)\rangle,}    
\end{equation} 
wobei $k=1,38064852 \times 10^{-23} \, \text{m}^2 \,\text{kg} \,\text{s}^{-2} \,\text{K}^{-1}$ die Boltzmann-Konstante ist und $\log(\rho)$ wird mit Hilfe vom Funktionalkalkül berechnet. Der Logarithmus ist hier, wie überall in der Physik, der natürliche; und $\log(\rho)$ muss selbst ein Operator sein, sonst gäbe es das Produkt $\rho\log(\rho)$ gar nicht.

Weil jeder Dichteoperator $\rho$ hermitesch ist, können wir $\rho$ bezüglich einer Eigenwert-Eigenvektor-Zerlegung darstellen:
\begin{equation}
  \rho = \sum_{j=1}^d r_j |r_j\rangle\langle r_j|,
\end{equation}
wobei $r_j\ge 0$ und $\langle r_j|r_k\rangle = \delta_{jk}$. Der Logarithmus von $\rho$ wird als
\begin{equation}
  \log(\rho) = \sum_{j=1}^d \log(r_j) |r_j\rangle\langle r_j|
\end{equation}
berechnet. Ist ein Eigenwert $r_j = 0$, so ist $\log(r_j) = -\infty$; darauf kommen wir im ersten Beispiel zurück.

\begin{merksatz}[Vorsicht mit dem Logarithmus einer Matrix]
	$\log(\rho)$ ist \emph{nicht} der eintragsweise Logarithmus der Einträge von $\rho$. Erst diagonalisieren, dann den Logarithmus auf die Eigenwerte anwenden. Eintragsweise rechnen darf man nur, wenn $\rho$ bereits diagonal ist.
\end{merksatz}

Ein Beispiel: für
\begin{equation}
	\rho = \begin{pmatrix} \frac{1}{2} & \frac{1}{4}\\[2pt] \frac{1}{4} & \frac{1}{2}\end{pmatrix}
	\quad \text{ist} \quad
	\log(\rho) = \frac{1}{2}\begin{pmatrix} \log\frac{3}{16} & \log 3\\[2pt] \log 3 & \log\frac{3}{16}\end{pmatrix}
	\neq \begin{pmatrix} \log\frac{1}{2} & \log\frac{1}{4}\\[2pt] \log\frac{1}{4} & \log\frac{1}{2}\end{pmatrix},
\end{equation}
denn die Eigenwerte sind $3/4$ und $1/4$.

Daher wird die von Neumann-Entropie von $\rho$ wie folgt berechnet:
\begin{equation}
  S(\rho) = -k\sum_{j=1}^d \log(r_j) \tr(|r_j\rangle\langle r_j|\rho) = -k\sum_{j=1}^d r_j\log(r_j).
\end{equation}
Das ist der Befreiungsschlag: die Entropie hängt \emph{nur} von den Eigenwerten ab, die Eigenvektoren tauchen nicht mehr auf. Eine unitäre Drehung ändert das Spektrum nicht, also auch nicht die Entropie:
\begin{equation}
	\eqbox{S(\rho) = S(U\rho U^{\dagger}) \quad \text{für alle unitären } U.}
\end{equation}

\begin{übung}
	Beweisen Sie die unitäre Invarianz $S(\rho) = S(U\rho U^{\dagger})$.
\end{übung}

\begin{beispiel}
  Ein reiner Zustand $|\psi\rangle$ entspricht einem statischen Ensemble, in dem $|\psi\rangle$ mit Wahrscheinlichkeit $1$ auftritt. Der dazugehörige Dichteoperator ist durch
  \begin{equation}
    \rho = |\psi\rangle\langle \psi|
  \end{equation}
  gegeben. Dieser Operator befindet sich bereits in diagonaler Form wobei $r_1 = 1$ (und $r_2 = r_3 = \cdots = r_d =0$). Daher lautet die von Neumann-Entropie
  \begin{equation}
    S(\rho) = -k\sum_{j=1}^d r_j\log(r_j) = -k r_1\log(r_1) = 0.
  \end{equation}
  Hierbei nimmt die Funktion $f(x) = -x\log(x)$ den Wert $0$ an der Stelle $x=0$ an. Genauer ist $0\log(0)$ als Grenzwert $\lim_{x\to 0} x\log(x)$ zu lesen; die Kurve $-x\log(x)$ läuft bei $x=0$ sichtbar auf Null zu. In diesem Beispiel beträgt die von Neumann-Entropie Null. Die Interpretation ist hier, dass das statistische Ensemble keine Informationen enthält. Man weiß immer, in welchem Zustand sich das System befindet. Es fehlt keine Information über den Zustand des Systems. Ein reiner Zustand ist eben kein Ensemble.
\end{beispiel}

\begin{beispiel}
  Angenommen, $\mathcal{H}$ ist der Hilbertraum $\mathbb{C}^d$. Nehmen wir weiter an, dass $\{|j\rangle \}_{j=1}^d$ eine Orthonormalbasis für $\mathcal{H}$ ist. Der vollständig gemischte Zustand $\rho$ ist gegeben durch
  \begin{equation}
    \rho \equiv \sum_{j=1}^d \frac{1}{d} |j\rangle\langle j|.
  \end{equation}
  Dieser Zustand entspricht einem statischen Ensemble, in dem $|j\rangle$ mit gleich verteilter Wahrscheinlichkeit $1/d$ auftritt. Dieser Operator befindet sich auch bereits in diagonaler Form wobei $r_1 = 1/d = r_2 = \cdots = r_d$. Die von Neumann-Entropie wird wie folgt berechnet:
  \begin{equation}
    S(\rho) = -k\sum_{j=1}^d r_j\log(r_j) = -k \times d \times  (1/d) \times \log(1/d) = k \log(d).
  \end{equation}
  Die von Neumann-Entropie erreicht in diesem Fall den maximal möglichen Wert. Die Interpretation ist hier, dass das statistische Ensemble die maximale Information enthält. Man weiß nicht, in welchem Zustand sich das System befindet. Es fehlen alle Information über den Zustand des Systems.
\end{beispiel}

% fig_vl19_entropie.tex — Die von-Neumann-Entropie als Summe der Terme
% -r log r: der einzelne Term, und die Entropie eines Qubits.
%
% WHY THIS FIGURE.  The lecture states three things it never shows: that
% -x log x runs into zero at x = 0 (the text says the curve does so
% "sichtbar", so the curve had better be there — it is the one picture
% Osborne draws at the board, v23 [64:14]), that S = 0 for a pure state
% and S = k log d at the other extreme, and that only the spectrum
% enters.  What prose cannot hand over is the SHAPE between those two
% extremes: that S rises with infinite slope out of the pure state — a
% one-percent admixture already buys eight percent of the maximum — and
% then flattens off, so that near the maximally mixed state the entropy
% hardly notices how the weight is distributed.  That shape, and the
% fact that the whole reachable range of a qubit is the strip below
% k log 2, is the content of the plate.
%
% GEOMETRY FROM THE FORMULA.  Nothing is drawn by eye.  Both curves are
% the lecture's own expression S(rho) = -k sum_j r_j log(r_j) evaluated
% on a spectrum, in units of k:
%
%   (a)  eta(r) = -r log r          one eigenvalue's contribution,
%   (b)  S/k    = eta(p) + eta(1-p) a qubit with spectrum {p, 1-p}.
%
% Every marked value is exact, not eyeballed:
%   * eta has its maximum at r = 1/e, of height 1/e = 0.36788 — the
%     y tick and the grey guide are placed there;
%   * eta(1/2) = (1/2) log 2 = 0.34657, and twice that is the peak of
%     panel (b): the red dot in (a) is literally half the red dot in
%     (b).  That is why both are red — same physical situation, two
%     eigenvalues of equal weight;
%   * S/k at p = 1/2 is log 2 = 0.69315 = log d for d = 2, the ceiling
%     of the lecture's second example;
%   * S/k at p = 3/4 is -(3/4)log(3/4) - (1/4)log(1/4) = 0.56234, the
%     value of the worked example in the text, whose two matrices
%     (1/4)diag(3,1) and (1/4)[[2,1],[1,2]] share the spectrum {3/4,1/4}
%     and therefore land on the SAME dot.  Drawing them arriving at one
%     point is the unitary invariance S(rho) = S(U rho U^dagger) shown
%     rather than restated;
%   * the dotted rule between p = 1/4 and p = 3/4 is the same statement
%     one level down: a spectrum is a set, so p and 1-p are one state
%     and must have one entropy.
%
% DESIGN DECISIONS
%   * colour ramps with the entropy value, and carries nothing else:
%     open circle (spacecadet on white) where the term vanishes —
%     r = 0 never occurs, r = 1 occurs for certain, neither hides any
%     information; birrenorange for the intermediate example; birrenred
%     for the maximum.  The curves themselves are the house data colour
%     cgblue in both panels because they are the same function;
%   * the band above log 2 is hatched, not filled: it is unreachable for
%     d = 2, and hatching survives greyscale printing.  The grey rule
%     bounding it is a reference line, so it is neutral, like the guide
%     to the maximum of eta;
%   * no y labels: the panel titles already name the plotted quantity,
%     and the width saved goes into the annotations, which sit on their
%     objects rather than in a legend;
%   * the words under the p axis ("rein", "maximal gemischt") name the
%     states, not the numbers — the reading of the figure is physical;
%   * annotation positions were set against the 12pt amsart build, not
%     against scripts/tikz-preview.sh: the preview is a 10pt standalone,
%     so every \scriptsize label is noticeably smaller there and
%     clearances that look ample in the PNG can close up on the page.
%     The eta(1/2) label therefore sits above the curve, centred on
%     r = 1/2 and tied to its dot by a hairline leader (to the right of
%     the dot the label grew into panel (b); a rule carried down to the
%     axis instead crossed the words of the limit annotation), and the
%     annotation block of (b) is centred on p = 1/2, the axis of symmetry
%     of the curve, which is where the free space is widest;
%   * grey does two jobs and they are kept apart: dashed rules are
%     levels (1/e, log 2), solid hairlines are leaders from a label to
%     the thing it names.

\begin{center}
\begin{tikzpicture}

% ─────────────────────────────────────────────────────────────
%  Exact constants (see header)
% ─────────────────────────────────────────────────────────────
\def\einve{0.3678794}       % 1/e             — maximum of eta
\def\halflogtwo{0.3465736}  % (1/2)log 2      — eta(1/2)
\def\logtwo{0.6931472}      % log 2           — S/k of the mixed qubit
\def\Sdreiviertel{0.5623351}% S({3/4,1/4})/k  — the worked example

% ─────────────────────────────────────────────────────────────
%  (a)  The single term eta(r) = -r log r
% ─────────────────────────────────────────────────────────────
\begin{axis}[
    name=term,
    tpcplot,
    width=5.8cm, height=5.4cm,
    grid=none, clip=false,
    xmin=-0.05, xmax=1.12, ymin=-0.035, ymax=0.50,
    xtick={0,0.5,1}, xticklabels={$0$,$\tfrac12$,$1$},
    ytick={0,\einve}, yticklabels={$0$,$1/e$},
    xlabel={$r$},
    title={(a) \tpcifenglish{one term}{ein Term}: $\eta(r)=-r\log r$},
    domain=0.0002:1, samples=300,
]
  % One reference line only: the height of the maximum.  The equal-weight
  % point gets a leader down from its label instead of a rule through the
  % whole panel — a full-height rule at r = 1/2 runs straight through the
  % words of the limit annotation below, and dashed grey is reserved here
  % for level lines, solid hairlines for leaders.
  \draw[levellinegrey] (axis cs:-0.05,\einve) -- (axis cs:\einve,\einve);

  \addplot[wavefn, cgblue] {-x*ln(x)};

  % Both zeros of eta: nothing is missing there
  \node[circle, draw=spacecadet, fill=white, line width=0.7pt,
        inner sep=1.1pt] at (axis cs:0,0) {};
  \node[circle, draw=spacecadet, fill=white, line width=0.7pt,
        inner sep=1.1pt] at (axis cs:1,0) {};

  % Equal weights: half of the qubit maximum in (b)
  \node[circle, fill=birrenred, inner sep=1.55pt]
    at (axis cs:0.5,\halflogtwo) {};
  \node[lbl, font=\scriptsize, anchor=south]
    at (axis cs:0.5,0.405) {$\eta(\tfrac12)=\tfrac12\log 2$};
  \draw[birrengray, line width=0.4pt]
    (axis cs:0.5,0.399) -- (axis cs:0.5,0.357);

  % The limit that defines 0 log 0
  \node[lbl, font=\scriptsize, anchor=west]
    at (axis cs:0.085,0.055) {$\lim_{r\to0} r\log r = 0$};
  \draw[birrengray, line width=0.4pt]
    (axis cs:0.078,0.047) -- (axis cs:0.022,0.013);

  % The state names hang a fixed 6.2 mm below the axis line, here and in
  % (b).  In axis units the same drop would land at two different heights
  % on the page: the panels cover 0.535 and 0.945 units in the same
  % 5.4 cm, so one unit is nearly twice as long in (a) as in (b).
  \node[lbl, text=birrengray, font=\scriptsize, anchor=north, yshift=-6.2mm]
    at (axis cs:0,-0.035) {\tpcifenglish{never}{nie}};
  \node[lbl, text=birrengray, font=\scriptsize, anchor=north, yshift=-6.2mm]
    at (axis cs:1,-0.035) {\tpcifenglish{certain}{sicher}};
\end{axis}

% ─────────────────────────────────────────────────────────────
%  (b)  A qubit: S/k = eta(p) + eta(1-p)
% ─────────────────────────────────────────────────────────────
\begin{axis}[
    name=qubit,
    at={(term.south east)}, anchor=south west, xshift=12mm,
    tpcplot,
    width=6.8cm, height=5.4cm,
    grid=none, clip=false,
    xmin=-0.04, xmax=1.07, ymin=-0.045, ymax=0.90,
    xtick={0,0.25,0.5,0.75,1},
    xticklabels={$0$,$\tfrac14$,$\tfrac12$,$\tfrac34$,$1$},
    ytick={0,\logtwo}, yticklabels={$0$,$\log 2$},
    xlabel={$p$},
    title={(b) \tpcifenglish{qubit}{Qubit}: $S/k=\eta(p)+\eta(1-p)$},
    domain=0.0004:0.9996, samples=320,
]
  % Unreachable for d = 2: hatched, bounded by a neutral reference rule
  \fill[forbidden] (axis cs:-0.04,\logtwo) rectangle (axis cs:1.07,0.90);
  \draw[levellinegrey] (axis cs:-0.04,\logtwo) -- (axis cs:1.07,\logtwo);
  % the backing is `tpcplot`'s legend style, but opaque: at the legend's
  % 0.88 the hatching still shows through the words
  \node[lbl, text=birrengray, font=\scriptsize, inner sep=2.5pt,
        fill=white, rounded corners=1.5pt]
    at (axis cs:0.5,0.795)
    {\tpcifenglish{unreachable}{unerreichbar}: $S_{\max}=k\log d$};

  \addplot[wavefn, cgblue] {-x*ln(x)-(1-x)*ln(1-x)};

  % One spectrum, two places on the p axis
  \draw[classicallim] (axis cs:0.25,\Sdreiviertel)
                   -- (axis cs:0.75,\Sdreiviertel);
  \node[circle, fill=birrenorange, inner sep=1.55pt]
    at (axis cs:0.25,\Sdreiviertel) {};
  \node[circle, fill=birrenorange, inner sep=1.55pt]
    at (axis cs:0.75,\Sdreiviertel) {};

  % Maximum and the two pure states
  \node[circle, fill=birrenred, inner sep=1.55pt]
    at (axis cs:0.5,\logtwo) {};
  \node[circle, draw=spacecadet, fill=white, line width=0.7pt,
        inner sep=1.1pt] at (axis cs:0,0) {};
  \node[circle, draw=spacecadet, fill=white, line width=0.7pt,
        inner sep=1.1pt] at (axis cs:1,0) {};

  % The two matrices of the worked example arrive at the same dot
  \node[lbl, font=\footnotesize, anchor=north] (mats) at (axis cs:0.5,0.44)
    {$\tfrac14\!\left(\begin{smallmatrix}3&0\\0&1\end{smallmatrix}\right)$
     \tpcifenglish{and}{und}
     $\tfrac14\!\left(\begin{smallmatrix}2&1\\1&2\end{smallmatrix}\right)$};
  \draw[-{Stealth[length=4pt]}, birrengray, line width=0.5pt,
        shorten >=2.5pt] (mats.north east) -- (axis cs:0.75,\Sdreiviertel);
  \node[lbl, text=birrengray, font=\scriptsize, anchor=north, align=center]
    at (axis cs:0.5,0.315)
    {\tpcifenglish{same spectrum,}{gleiches Spektrum,}\\
     \tpcifenglish{same entropy}{gleiche Entropie}};

  \node[lbl, text=birrengray, font=\scriptsize, anchor=north, yshift=-6.2mm]
    at (axis cs:0,-0.045) {\tpcifenglish{pure}{rein}};
  \node[lbl, text=birrengray, font=\scriptsize, anchor=north, yshift=-6.2mm]
    at (axis cs:1,-0.045) {\tpcifenglish{pure}{rein}};
  \node[lbl, text=birrengray, font=\scriptsize, anchor=north, yshift=-6.2mm]
    at (axis cs:0.5,-0.045)
    {\tpcifenglish{max.\ mixed}{maximal gemischt}};
\end{axis}

\end{tikzpicture}
\end{center}


Die Abbildung zeigt links den einzelnen Term $\eta(r)=-r\log r$: er verschwindet
an beiden Enden -- bei $r=0$ tritt der Zustand nie auf, bei $r=1$ mit Sicherheit.
Rechts steht die Summe zweier solcher Terme, die Entropie eines Qubits. Achten Sie
auf die steile Flanke: schon ein Prozent Beimischung zum reinen Zustand kostet acht
Prozent des Maximums $k\log 2$; die beiden Matrizen mit demselben Spektrum landen
dagegen auf einem einzigen Punkt.

Wegen des Vorfaktors $k$ ist die Entropie eines einzelnen Qubits winzig; groß wird sie erst bei sehr vielen Teilchen.

\begin{übung}\label{ex:entropie-additiv}
	Es sei $\rho = \rho_1\otimes\rho_2\otimes\cdots\otimes\rho_N$ der Zustand von $N$ Teilchen. Zeigen Sie die Additivität
	\begin{equation}
		S(\rho) = \sum_{j=1}^N S(\rho_j).
	\end{equation}
	Für gleiche Faktoren ist das $N$-fache der Entropie eines einzelnen Teilchens -- die Entropie ist \emph{extensiv}.
\end{übung}
%!TEX root = TheoretischePhysikC.tex

\section{Vorlesung: Mikrokanonisches Ensemble}
Systeme in der Natur bestehen aus $10^{16}$ bis $10^{23}$ Teilchen; ihr Hilbertraum wird so groß, dass eine exakte Behandlung so gut wie unmöglich wird -- selbst im Prinzip. Der Ausweg ist die statistische Physik, und sie kostet Information.

Angenommen, ein abgeschlossenes Quantensystem (Hilbertraum: $\mathcal{H}\equiv \mathbb{C}^d$, Hamiltonoperator: $\widehat{H}$) wird mit fester Teilchenzahl, festem Volumen und Energie nur im Intervall $[E, E+\Delta]$ mit kleinem $\Delta$ präpariert. Wir besitzen zu wenig Information: die Teilchenzahl usw. Diese Informationen reicht nicht, um den Zustand des Systems vollständig zu beschreiben. Wie können wir den Zustand des Systems modellieren? Laut zentraler Annahme der statistischen Physik sollen alle Zustände mit einer Energie in $[E, E+\Delta]$ mit gleichem statistischen Gewicht auftreten. Das ist ein Prinzip, nicht die Wahrheit: ohne diese Annahme kämen wir an die Wahrscheinlichkeiten gar nicht heran.

Die Energie-Eigenzustände $|E_j\rangle$ und dazugehörige Energie-Eigenwerte $E_j$ des Systems liefern uns eine Diagonalisierung von $H$:
\begin{equation}
  \widehat{H} = \sum_{j=1}^d E_j |E_j\rangle \langle E_j|.
\end{equation}
Damit ist keine Einschränkung verbunden: jeder Hamiltonoperator ist hermitesch und deshalb diagonalisierbar. Endlichdimensional zu rechnen erleichtert nur die Darstellung.
Mithilfe von $\{E_j, |E_j\rangle \}_{j=1}^d$ identifizierung wir das folgende statistische Ensemble als Modell für den Zustand des Systems. Der Eigenzustand $|E_j\rangle$ taucht mit Wahrscheinlichkeit $0$ auf, wenn $E_j\not\in [E,E+\Delta]$. Im Fall, dass $E_j\in [E,E+\Delta]$ taucht $|E_j\rangle$ mit Wahrscheinlichkeit $1/N_E$ auf, wobei $N_E$ die Anzahl der Zustände in $[E, E+\Delta]$ ist. Wir stellen daher die Hypothese auf, dass der folgende Quantenzustand das System am besten modelliert (angesichts der uns vorliegenden Informationen):
\begin{equation}
  \rho = \sum_{E_j\in [E,E+\Delta]} \frac{1}{N_E} |E_j\rangle \langle E_j|.
\end{equation} 
Dieses statistische Ensemble wird als \emph{mikrokanonisches Ensemble} bezeichnet. Das mikrokanonische Ensemble wird als Modell für den Zustand eines Systems verwendet, wenn wir die Teilchenzahl, das Volumen und die Energie des Systems genau kennen.

\begin{beispiel}
  Ein Dreiniveausystem, dessen Hamiltonoperator bereits diagonal vorliegt:
  \begin{equation}
    \widehat{H} = \begin{pmatrix} 0 & 0 & 0\\ 0 & 1 & 0\\ 0 & 0 & 1{,}1\end{pmatrix},
  \end{equation}
  mit den Standardbasisvektoren als Energieeigenzuständen, $E_1 = 0$, $E_2 = 1$ und $E_3 = 1{,}1$. Gemessen sei die Energie im Bereich $[0{,}9;\, 1{,}2]$, also $E = 0{,}9$ und $\Delta = 0{,}3$. Darin liegen $E_2$ und $E_3$, nicht aber $E_1$, also ist $N_E = 2$ und
  \begin{equation}
    \rho = \frac{1}{2}|E_2\rangle\langle E_2| + \frac{1}{2}|E_3\rangle\langle E_3|
         = \begin{pmatrix} 0 & 0 & 0\\ 0 & \frac{1}{2} & 0\\ 0 & 0 & \frac{1}{2}\end{pmatrix}.
  \end{equation}
  Beachten Sie, was die Messung leistet: der Grundzustand tritt mit Wahrscheinlichkeit genau null auf, dort ist das System nie.
\end{beispiel}

Die Entropie des mikrokanonischen Ensembles wird wie folgt berechnet:
\begin{equation}
  S(\rho) = -k\sum_{E_j\in [E,E+\Delta]} \frac{1}{N_E} \log\left(\frac{1}{N_E}\right) = k \log(N_E).
\end{equation}
Kurz ist die Rechnung, weil $\rho$ hier schon diagonal ist; das Minuszeichen hebt sich gegen $\log(1/N_E) = -\log(N_E)$ auf. Im Beispiel oben ist $S(\rho) = k\log(2)$.

% fig_vl20_mikro.tex — Die Energieschale: das Abzählen der Mikrozustände
% in [E, E+\Delta], und was aus dem Zählen wird, wenn N wächst.
%
% WHY THIS FIGURE.  The lecture writes down
%       rho = sum_{E_j in [E,E+Delta]} (1/N_E) |E_j><E_j|,   S = k log(N_E)
% and makes N_E concrete only on a three-level toy, where N_E = 2.  Two
% things stay invisible there, and both are shapes rather than sentences:
%   (i)  what "counting the states in the shell" looks like when the
%        states are degenerate — the count is a binomial coefficient, and
%        the multiplicities are wildly unequal from level to level;
%   (ii) that this toy count IS the thermodynamic entropy: k log(N_E)
%        divided by N converges to a smooth concave curve of the energy
%        per particle, with maximum log 2, whose slope is what the next
%        lecture calls 1/T.
% Panel (a) can be counted with the finger; panel (b) says that the same
% count, done for 10^23 spins, is the entropy of thermodynamics.
%
% GEOMETRY FROM THE FORMULA.  Nothing is drawn by eye.  The model is N
% non-interacting spins in a field, single-spin energies 0 and epsilon,
%
%       H |s> = k epsilon |s>,   k = number of up spins in s,
%
% so the spectrum is the equally spaced ladder E_k = k epsilon, and the
% degeneracy of level k is the number of strings with k up spins,
%
%       g_k = C(N,k).
%
% Panel (a), N = 6:  g = (1,6,15,20,15,6,1), sum = 2^6 = 64 dots, one dot
% per microstate, placed at height E_k = k epsilon and spread horizontally
% at unit spacing about x = 0 — so the width of a row IS its degeneracy
% and the rows may be compared by eye.  The shell is E = 1.5 epsilon,
% Delta = 2 epsilon, i.e. [1.5, 3.5] epsilon, which admits exactly the
% levels k = 2 and k = 3:
%
%       N_E = C(6,2) + C(6,3) = 15 + 20 = 35,   p_j = 1/35,
%       S   = k log(35).
%
% Deliberately two levels, not one: N_E counts eigenvalues in an interval,
% and the interval need not respect the level structure — that is exactly
% what the lecture's own example does (E_2 and E_3 both fall in [0.9,1.2]).
%
% Panel (b), the same count for a shell holding one level, per spin:
%
%       S/(kN) = (1/N) log C(N,k),   e = E/(N epsilon) = k/N,
%
% plotted from exact binomial coefficients for N = 6 and N = 24, against
% its N -> infinity limit, the binary entropy
%
%       s(e) = -e log(e) - (1-e) log(1-e),
%
% whose maximum s(1/2) = log 2 is the ceiling S <= k log(d) of the
% previous lecture, here d = 2^N.  The finite-N clouds sit below the
% curve and creep up towards it: the correction is -log(N)/(2N), which is
% why N = 6 is already the right shape and 10^23 is indistinguishable
% from the limit.
%
% DESIGN DECISIONS
%   * colour is probability, and nothing else: cgblue = the states that
%     actually occur (p_j = 1/N_E), grey open circles = the states with
%     p_j = 0.  The same filled/open distinction is repeated in (b), where
%     the two filled points are precisely the two levels of the shell of
%     (a) — that is the link between the panels, and it needs no words;
%   * energies outside the shell are hatched, not filled, exactly as the
%     classically forbidden region of fig_vl11 is: unreachable is
%     unreachable, and hatching survives greyscale;
%   * the energy levels themselves are neutral grey dashed rules; the two
%     shell edges are solid spacecadet, because they are the one piece of
%     input information the ensemble is built from;
%   * in (b) the finite sizes are dot clouds (the energy really is
%     discrete there, e = k/N) and only the thermodynamic limit is drawn
%     as a continuous curve.  The three sizes N = 6, 24, infinity are the
%     ordered ramp cgblue -> birrenorange -> birrenred of `tpclevels`;
%     munsell is skipped for N = 24 for the reason fig_vl11 gives — beside
%     cgblue it is the same teal, so it could not carry the code.  N = 6
%     keeps the cgblue of panel (a), because it IS the system of panel (a);
%   * (b) plots shells that hold a single level, so there N_E = C(N,k),
%     whereas the shell of (a) straddles two.  The corner of (b) says so:
%     otherwise the reader would expect log(35) where log(20) is drawn;
%   * the degeneracy column on the right of (a) is ruled off and summed,
%     so that sum_k g_k = 2^N = dim(H) is on the plate: the shell keeps 35
%     of 64, and the reader can check both numbers by counting dots;
%   * every curve is labelled where it is, no legend box.

\begin{center}
\begin{tikzpicture}

% ─────────────────────────────────────────────────────────────
%  (a)  All 2^6 microstates, sorted by energy; the shell picks two rows
% ─────────────────────────────────────────────────────────────
\begin{axis}[
    name=schale,
    tpclevelplot,
    width=7.4cm, height=4.8cm, scale only axis,
    xmin=-15.4, xmax=13.8,
    ymin=-2.5,  ymax=7.3,
    title={(a) \tpcifenglish{energy shell of $N=6$ spins}%
                             {Energieschale von $N=6$ Spins}},
]

  % Energies outside [E, E+Delta]: p_j = 0, hatched like a forbidden region
  \fill[forbidden] (axis cs:-11,-0.55) rectangle (axis cs:11,1.5);
  \fill[forbidden] (axis cs:-11,3.5)   rectangle (axis cs:11,6.55);

  % The ladder E_k = k epsilon: neutral reference rules
  \draw[levellinegrey] (axis cs:-10.6,0) -- (axis cs:10.6,0);
  \draw[levellinegrey] (axis cs:-10.6,1) -- (axis cs:10.6,1);
  \draw[levellinegrey] (axis cs:-10.6,2) -- (axis cs:10.6,2);
  \draw[levellinegrey] (axis cs:-10.6,3) -- (axis cs:10.6,3);
  \draw[levellinegrey] (axis cs:-10.6,4) -- (axis cs:10.6,4);
  \draw[levellinegrey] (axis cs:-10.6,5) -- (axis cs:10.6,5);
  \draw[levellinegrey] (axis cs:-10.6,6) -- (axis cs:10.6,6);

  % Which way is up: energy increases along the page
  \draw[axisstyle] (axis cs:-14.4,4.25) -- (axis cs:-14.4,6.5);
  \node[lbl, text=birrengray, font=\scriptsize, rotate=90]
    at (axis cs:-13.5,5.35) {\tpcifenglish{energy}{Energie}};

  % The two edges of the shell — the whole input to the ensemble
  \draw[potential] (axis cs:-11,1.5) -- (axis cs:11,1.5);
  \draw[potential] (axis cs:-11,3.5) -- (axis cs:11,3.5);
  \node[lbl, anchor=east] at (axis cs:-11.3,1.5) {$E$};
  \node[lbl, anchor=east] at (axis cs:-11.3,3.5) {$E+\Delta$};

  % Microstates outside the shell: p_j = 0
  \addplot[only marks, mark=o, mark size=1.25pt,
           mark options={draw=birrengray, fill=white, line width=0.4pt}]
    coordinates {
      (0.0,0) (-2.5,1) (-1.5,1) (-0.5,1) (0.5,1) (1.5,1) (2.5,1)
      (-7.0,4) (-6.0,4) (-5.0,4) (-4.0,4) (-3.0,4) (-2.0,4) (-1.0,4)
      (0.0,4) (1.0,4) (2.0,4) (3.0,4) (4.0,4) (5.0,4) (6.0,4) (7.0,4)
      (-2.5,5) (-1.5,5) (-0.5,5) (0.5,5) (1.5,5) (2.5,5) (0.0,6)
    };

  % Microstates inside the shell: p_j = 1/N_E
  \addplot[only marks, mark=*, mark size=1.35pt,
           mark options={draw=cgblue, fill=cgblue}]
    coordinates {
      (-7.0,2) (-6.0,2) (-5.0,2) (-4.0,2) (-3.0,2) (-2.0,2) (-1.0,2)
      (0.0,2) (1.0,2) (2.0,2) (3.0,2) (4.0,2) (5.0,2) (6.0,2) (7.0,2)
      (-9.5,3) (-8.5,3) (-7.5,3) (-6.5,3) (-5.5,3) (-4.5,3) (-3.5,3)
      (-2.5,3) (-1.5,3) (-0.5,3) (0.5,3) (1.5,3) (2.5,3) (3.5,3)
      (4.5,3) (5.5,3) (6.5,3) (7.5,3) (8.5,3) (9.5,3)
    };

  % Probabilities, written where they hold
  \node[lbl, text=cgblue] at (axis cs:0,2.5) {$p_j=1/35$};
  \node[lbl, text=birrengray, fill=white, inner sep=1.5pt]
    at (axis cs:0,5.5) {$p_j=0$};
  \node[lbl, text=birrengray, fill=white, inner sep=1.5pt]
    at (axis cs:0,0.5) {$p_j=0$};

  % Degeneracies g_k = C(6,k), read off as the width of each row
  \node[lbl, text=birrengray, font=\scriptsize, anchor=west]
    at (axis cs:11.25,6.9) {$g_k=\tbinom{6}{k}$};
  \node[lbl, text=birrengray, anchor=west] at (axis cs:11.3,0) {$1$};
  \node[lbl, text=birrengray, anchor=west] at (axis cs:11.3,1) {$6$};
  \node[lbl, text=cgblue,     anchor=west] at (axis cs:11.3,2) {$15$};
  \node[lbl, text=cgblue,     anchor=west] at (axis cs:11.3,3) {$20$};
  \node[lbl, text=birrengray, anchor=west] at (axis cs:11.3,4) {$15$};
  \node[lbl, text=birrengray, anchor=west] at (axis cs:11.3,5) {$6$};
  \node[lbl, text=birrengray, anchor=west] at (axis cs:11.3,6) {$1$};

  % The column of degeneracies adds up to the dimension of the Hilbert space
  \draw[birrengray!70, line width=0.4pt]
    (axis cs:11.2,-0.5) -- (axis cs:12.9,-0.5);
  \node[lbl, text=birrengray, font=\scriptsize, anchor=north west]
    at (axis cs:11.2,-0.55) {$2^6=64$};

  % What the model is, and what one dot means
  \node[lbl, text=birrengray, font=\scriptsize, anchor=west]
    at (axis cs:-15.4,7.0)
    {$E_k=k\varepsilon$; \tpcifenglish{one dot $=$ one microstate}%
                                     {ein Punkt $=$ ein Mikrozustand}};

  % The count itself
  \node[lbl, anchor=north] at (axis cs:-0.7,-0.95)
    {$N_E=\tbinom{6}{2}+\tbinom{6}{3}=15+20=35$};
  \node[lbl, anchor=north] at (axis cs:-0.7,-1.80)
    {$S=k\log(N_E)=k\log(35)$};
\end{axis}

% ─────────────────────────────────────────────────────────────
%  (b)  The same count per spin, and its thermodynamic limit
% ─────────────────────────────────────────────────────────────
\begin{axis}[
    at={(schale.south east)}, anchor=south west, xshift=20mm,
    tpcplot,
    width=4.2cm, height=4.8cm, scale only axis,
    xmin=-0.03, xmax=1.03, ymin=0, ymax=0.79,
    xtick={0,0.5,1}, xticklabels={$0$,$\tfrac12$,$1$},
    ytick={0,0.2,0.4,0.6},
    xlabel={$e=E/(N\varepsilon)$},
    ylabel={$S/(kN)$},
    every axis x label/.style={at={(0.5,-0.115)}, anchor=north, font=\small},
    every axis y label/.style={at={(axis description cs:-0.17,0.5)},
                               rotate=90, anchor=south, font=\small},
    title={(b) \tpcifenglish{entropy per spin}{Entropie pro Spin}},
]
  % Ceiling S <= k log(d), d = 2^N: one log(2) per spin
  \addplot[levellinegrey, domain=0:1, samples=2] {0.6931};
  \node[lbl, text=birrengray, font=\scriptsize, anchor=south west]
    at (axis cs:0.025,0.70) {$\log 2$};

  % N -> infinity: the binary entropy
  \addplot[birrenred, line width=1.1pt, domain=0.0008:0.9992, samples=241]
    {-x*ln(x)-(1-x)*ln(1-x)};
  \node[lbl, text=birrenred, anchor=west] at (axis cs:0.70,0.655)
    {$N\to\infty$};

  % N = 24, exact binomials
  \addplot[only marks, mark=*, mark size=0.9pt,
           mark options={draw=birrenorange, fill=birrenorange}]
    coordinates {
      (0.0000,0.0000) (0.0417,0.1324) (0.0833,0.2342) (0.1250,0.3172)
      (0.1667,0.3863) (0.2083,0.4441) (0.2500,0.4921) (0.2917,0.5314)
      (0.3333,0.5628) (0.3750,0.5868) (0.4167,0.6037) (0.4583,0.6138)
      (0.5000,0.6171) (0.5417,0.6138) (0.5833,0.6037) (0.6250,0.5868)
      (0.6667,0.5628) (0.7083,0.5314) (0.7500,0.4921) (0.7917,0.4441)
      (0.8333,0.3863) (0.8750,0.3172) (0.9167,0.2342) (0.9583,0.1324)
      (1.0000,0.0000)
    };
  \node[lbl, text=birrenorange, font=\scriptsize, anchor=north]
    at (axis cs:0.5,0.585) {$N=24$};

  % N = 6 — the system of panel (a).  Filled: the two levels of the shell.
  \addplot[only marks, mark=o, mark size=1.4pt,
           mark options={draw=cgblue, fill=white, line width=0.5pt}]
    coordinates {
      (0.0000,0.0000) (0.1667,0.2986) (0.6667,0.4513) (0.8333,0.2986)
      (1.0000,0.0000)
    };
  \addplot[only marks, mark=*, mark size=1.4pt,
           mark options={draw=cgblue, fill=cgblue}]
    coordinates {(0.3333,0.4513) (0.5000,0.4993)};
  \node[lbl, text=cgblue, font=\scriptsize, anchor=north]
    at (axis cs:0.5,0.425) {$N=6$};

  % What is plotted
  \node[lbl, text=birrengray, font=\scriptsize, anchor=south]
    at (axis cs:0.5,0.10) {\tpcifenglish{shell $=$ one level}%
                                        {Schale $=$ ein Niveau}};
  \node[lbl, text=birrengray, font=\scriptsize, anchor=south]
    at (axis cs:0.5,0.025) {$N_E=\tbinom{N}{k}$, $e=k/N$};
\end{axis}

\end{tikzpicture}
\end{center}


Wie groß $N_E$ tatsächlich wird, sieht man erst bei entarteten Niveaus. Tafel (a) zählt alle $2^6$ Mikrozustände von sechs Spins ab; die Schale $[E,E+\Delta]$ lässt genau $35$ übrig, jeden mit demselben Gewicht $1/35$, alle übrigen mit Gewicht null. Tafel (b) zeigt dieselbe Abzählung pro Spin -- gefüllt sind die beiden Niveaus aus (a) -- und wohin sie für große $N$ läuft.

Diesen Zustand \emph{maximiert} die Entropie unter den Nebenbedingungen, dass die Energie in $[E,E+\Delta]$ liegt (und die Teilchenzahl und Volumen des Systems fest bleibt). Das ist zunächst nur eine Beobachtung; man kann sie aber umdrehen und zum Prinzip erheben.

Die Methode, bei der man die Hypothese so aufstellt, dass der Systemzustand die Entropie \emph{maximiert} wird, wurde 1957 von Edwin Thompson Jaynes in Anlehnung an Methoden der statistischen Mechanik und der Shannon'schen Informationstheorie eingeführt. Grundlage ist, in Abwesenheit von Information die Entropie der A-priori-Wahrscheinlichkeit zu maximieren, da jede andere Zuweisung willkürliche Einschränkungen der betrachteten Situation träfe. Die Maximum-Entropie-Methode legt sich so wenig wie möglich fest. Jaynes zufolge ist dies aber nur der letzte Schritt, um nach Einfüllen aller vorhandenen Information etwaige noch bestehende Lücken zu schließen.

In der statistischen Physik bedeutet dies: 
\begin{center}
\parbox[c][5cm][c]{14cm}{\begin{tcolorbox}[colback=banana,arc=0cm,frame hidden,enhanced]{Unter allen Zuständen eines physikalischen Systems, die kompatibel mit dem vorhandenen Wissen über das System sind, ist jener zu wählen, welcher die Entropie maximiert.}\end{tcolorbox}}
\end{center}

Warum lohnt dieser Tausch? Die Grundannahme spricht von \emph{gleich wahrscheinlichen} Mikrozuständen. Klassisch bilden diese eine wohlbestimmte Menge; quantenmechanisch sind ganze kontinuierliche Familien von Superpositionen zugelassen, und dann ist unklar, was überhaupt abgezählt werden soll. Das Maximum-Entropie-Prinzip kennt dieses Problem nicht, denn es braucht nur eine Zielfunktion und Nebenbedingungen:
\begin{equation}
  \max_{\rho}\, S(\rho) \quad \text{unter den Nebenbedingungen, die unser Wissen quantifizieren,}
\end{equation}
hier also, dass die Energie in $[E,E+\Delta]$ liegt. So gehen wir von nun an vor: erst die Nebenbedingungen, dann die Maximierung.

%!TEX root = TheoretischePhysikC.tex

\section{Vorlesung: Temperatur}
Man betrachte zwei Systeme $1$ und $2$ in Kontakt, so dass Energie, aber keine Teilchen ausgetauscht werden können. Das Gesamtsystems ist isoliert. Die Gesamtenergie beträgt $E$. Die Wahrscheinlichkeit, System $1$ bei Energie $E_1$ zu finden ist durch
\begin{equation}\label{eq:probenergy}
	w(E_1) \propto N_{E_1}^{(1)}N_{E-E_1}^{(2)}
\end{equation}
gegeben. Hierbei sollen $N_{E}^{(1)}$ und $N_{E}^{(2)}$ als Funktionen von $E_1$ betrachtet werden. 

Der Grund liegt im Gesamtsystem: dessen Hamiltonoperator ist nach (\ref{eq:H-tensor}) die Summe $\widehat{H} = \widehat{H}_1\otimes \mathbb{I} + \mathbb{I}\otimes \widehat{H}_2$, Energien addieren sich:
\begin{equation}
	\widehat{H}\left(|E_j\rangle\otimes |E_k\rangle\right) = (E_j+E_k)\,|E_j\rangle\otimes |E_k\rangle.
\end{equation}
Gebunden ist nur die Summe $E_j+E_k$, gesucht ist ihre Aufteilung. Zu festem $E_1$ ist jedes Teilsystem mikrokanonisch, das Gesamtensemble das Produkt beider,
\begin{equation}
	\rho = \rho_1\otimes \rho_2, \qquad p_{jk} = \frac{1}{N_{E_1}^{(1)}}\cdot\frac{1}{N_{E-E_1}^{(2)}},
\end{equation}
und Abzählen liefert $w(E_1)$.

Die Maximum-Wahrscheinlichkeit wird mit Hilfe der Ableitung gefunden werden:
\begin{equation}
	\frac{d}{dE_1} w(E_1) = 0.
\end{equation}
Nach Einsetzen von (\ref{eq:probenergy}) erhalten wir
\begin{equation}
	N_{E-E_1}^{(2)}\left(\frac{d}{dE_1}N_{E_1}^{(1)}\right) - N_{E_1}^{(1)}\left(\frac{d}{d(E-E_1)}N_{E-E_1}^{(2)}\right) = 0.
\end{equation}
Umgeformt:
\begin{equation}
	\frac{1}{N_{E_1}^{(1)}}\frac{d}{dE_1}N_{E_1}^{(1)} =  \frac{1}{N_{E-E_1}^{(2)}}\frac{d}{d(E-E_1)}N_{E-E_1}^{(2)}.
\end{equation}
Nicht zufällig stehen hier logarithmische Ableitungen: unter Tensorprodukten ist die Entropie additiv,
\begin{equation}
	S(\rho_1\otimes \rho_2) = S(\rho_1)+S(\rho_2).
\end{equation}
Mit 
\begin{equation}
	S_1(E_1) = k\log(N_{E_1}^{(1)})\quad \text{und} \quad S_2(E-E_1) = k\log(N_{E-E_1}^{(2)})
\end{equation}
wird diese Relation wie folgt aussehen
\begin{equation}
	\frac{dS_1(E_1)}{dE_1} = \frac{dS_2(E_2)}{dE_2}\bigg|_{E_2 = E-E_1}.
\end{equation}
Daher wird die \emph{Temperatur} $T$ durch 
\begin{equation}
	\frac{1}{T} = \frac{\partial S(E)}{\partial E}
\end{equation}
definiert.\footnote{Benannt wird eigentlich die \emph{inverse} Temperatur; so ist es in der Physik gelaufen.} Gemäß dieser Definition wird festgestellt, dass thermisches Gleichgewicht zwischen den Systemen $1$ und $2$ bedeutet, dass die Temperaturen gleich sind.

% fig_vl21_kontakt.tex — Zwei Systeme in thermischem Kontakt: wo die
% Gesamtentropie ihr Maximum hat, und wie scharf dieses Maximum wird.
%
% WHY THIS FIGURE.  The lecture derives dS_1/dE_1 = dS_2/dE_2 by
% differentiating w(E_1) and then *defines* 1/T by it.  Two things stay
% invisible in the prose.  First, the condition is a statement about
% SLOPES: the sum is maximal where the rising curve climbs exactly as
% fast as the falling one drops — and that place is in general NOT at
% E_1 = E/2.  Second, the whole argument only means something because
% the maximum of w is sharp; how sharp is a scaling statement
% (width ~ 1/sqrt(N)) that a sentence cannot show.
%
% GEOMETRY FROM THE FORMULA.  Nothing is sketched.  The lecture leaves
% N_E abstract, so the plate fixes the simplest counting with a closed
% form — N oscillators sharing E/eps quanta (Einstein solid), whose
% state count in the Stirling form is
%
%       S(N,E)/k = N f(E/(N eps)),   f(x) = (1+x)log(1+x) - x log x.
%
% System 1 gets N_1 = 1, system 2 gets N_2 = 3 (three times as large),
% and the total energy is E = 2 eps.  Then, with u = E_1/E,
%
%       S_1/k = f(2u),        S_2/k = 3 f(2(1-u)/3).
%
% Because f'(x) = log(1 + 1/x), the stationarity condition
% dS_1/dE_1 = dS_2/dE_2 reduces to E_1/N_1 = E_2/N_2, so the maximum
% sits exactly at
%
%       E_1^* = E N_1/(N_1+N_2) = E/4,
%
% and the common slope there is f'(1/2) = log 3 (in units 1/eps); in the
% plotted variable u the two tangents therefore have the exactly mirrored
% slopes +2log3 and -2log3.  Every coordinate in panel (a) — the peak
% position, the tangency points, the tangent slopes, the value of
% S_1+S_2 at E/2 — is that arithmetic, not a drawing.
%
% Panel (b) is panel (a) exponentiated.  Scaling both systems by the
% same factor lambda (N_i -> lambda N_i, E -> lambda E) leaves u^* and
% the shape function g(u) = S_tot(u)/k untouched, so
%
%       w(E_1) ∝ N^{(1)} N^{(2)} = e^{S_tot/k} = e^{lambda g(u)},
%
% which is *exactly* the w(E_1) of eq. (eq:probenergy), normalised to its
% own maximum — panel (a) is the lambda = 1 member of that family.  The
% width shrinks like 1/sqrt(lambda): with g''(u^*) = -64/9 the Gaussian
% width is 0.375/sqrt(lambda).
%
% NUMERICS.  Writing that exponent as lambda*(S_1+S_2-Smax) is fatal at
% lambda = 100: pgfmath's log carries an absolute error of order 1e-4,
% the subtraction of two numbers near 3.82 keeps it, and multiplying by
% 100 turns it into a visibly ragged crest that falls ~2% short of 1.
% So the difference is formed analytically first.  With d = 2u - 1/2
% (zero at the maximum) the linear terms cancel exactly and
%
%   g(u) - g(u^*) = (3/2+d)log(1+d/(3/2)) - (1/2+d)log(1+d/(1/2))
%                 + (9/2-d)log(1-d/(9/2)) - (3/2-d)log(1-d/(3/2)),
%
% every logarithm now of the form log(1+small).  This vanishes at d = 0
% by construction, so all three curves reach 1 exactly, and its second
% derivative there is -16/9 in d, i.e. g''(u^*) = -64/9 as above.
%
% DECISIONS.
%   * colour carries which system a curve belongs to: system 1 blue,
%     system 2 orange, the sum red.  A tangent is dashed in the colour
%     of its own curve; construction lines that belong to no system
%     (the verticals at E_1^* and E/2) are neutral grey;
%   * panel (b) is an ordered ramp of the *same* red as the sum in (a),
%     because w = exp(S_tot/k) is that curve and nothing else;
%   * E_1 < 0 and E_2 < 0 are unreachable and hatched, not filled — the
%     constraint E_1 + E_2 = E is what makes this a one-parameter
%     problem at all;
%   * the y axis of (a) carries no ticks: the counting behind the shape
%     is a stand-in, only the slopes and the position of the maximum
%     are the physics.  It is labelled S, not S/k, so that a tangent's
%     slope is the lecture's own 1/T = dS/dE rather than 1/(kT).  The
%     numbers behind the curve are S/k; with no ticks and k > 0 that is
%     the same curve up to a scale, and the scale is what the axis name
%     fixes;
%   * E/2 is marked in grey solely so that one can see the maximum is
%     not there — temperatures equalise, energies do not;
%   * every label sits on its own object; the panels were laid out by
%     first computing the gaps between the curves, so no annotation
%     crosses a curve.

\begin{center}
\begin{tikzpicture}

% ─────────────────────────────────────────────────────────────
%  Shared abbreviations.  f(x) = (1+x)ln(1+x) - x ln x, evaluated at
%  x = 2u (system 1) and x = 2(1-u)/3 (system 2).
% ─────────────────────────────────────────────────────────────
\def\Sone{((1+2*x)*ln(1+2*x)-2*x*ln(2*x))}
\def\Stwo{(3*((1+(2-2*x)/3)*ln(1+(2-2*x)/3)-((2-2*x)/3)*ln((2-2*x)/3)))}
\def\ustar{0.25}          % E_1^*/E = N_1/(N_1+N_2)
\def\Sonestar{0.9547713}  % f(1/2)
\def\Stwostar{2.8643138}  % 3 f(1/2)
\def\Smax{3.8190851}      % 4 f(1/2)
\def\slp{2.1972246}       % d(S_1/k)/du = 2 log 3
\def\Shalf{3.6355738}     % S_1+S_2 at E_1 = E/2

% g(u) - g(u^*), formed analytically so that no large exponent ever
% amplifies a cancellation (see NUMERICS above).  d = 2u - 1/2.
\def\dd{(2*x-0.5)}
\def\Gdiff{( (1.5+\dd)*ln(1+\dd/1.5) - (0.5+\dd)*ln(1+\dd/0.5)
             +(4.5-\dd)*ln(1-\dd/4.5) - (1.5-\dd)*ln(1-\dd/1.5) )}

% ─────────────────────────────────────────────────────────────
%  (a)  The two entropies and their sum against E_1
% ─────────────────────────────────────────────────────────────
\begin{axis}[
    name=ent,
    tpcpotential,
    width=9.0cm, height=6.6cm,
    enlargelimits=false,
    xmin=-0.20, xmax=1.18, ymin=-0.45, ymax=4.80,
    xtick={0.25,0.5,1}, xticklabels={$E_1^{*}$,$E/2$,$E$},
    ytick=\empty,
    xlabel={$E_1$}, ylabel={$S$},
    clip=false,
    domain=0.0015:0.9985, samples=241,
    title={(a) \tpcifenglish{The two entropies and their sum}%
                            {Die beiden Entropien und ihre Summe}},
]

  % Unreachable: E_1 < 0 on the left, E_2 = E - E_1 < 0 on the right
  \fill[forbidden] (axis cs:-0.20,0) rectangle (axis cs:0,4.60);
  \fill[forbidden] (axis cs:1,0)     rectangle (axis cs:1.12,4.60);
  \node[lbl, text=birrengray, font=\scriptsize, rotate=90]
    at (axis cs:-0.10,2.3) {$E_1<0$};
  \node[lbl, text=birrengray, font=\scriptsize, rotate=90]
    at (axis cs:1.06,2.3)  {$E_2<0$};

  % Reference verticals: equal energies (dotted) and the maximum (dashed)
  \draw[classicallim]  (axis cs:0.5,0)     -- (axis cs:0.5,\Shalf);
  \draw[levellinegrey] (axis cs:\ustar,0)  -- (axis cs:\ustar,\Smax);

  % The three curves
  \addplot[wavefn, cgblue]              {\Sone};
  \addplot[wavefn, birrenorange]        {\Stwo};
  \addplot[line width=1.3pt, birrenred] {\Sone+\Stwo};

  % Tangents at E_1^*: slopes +2log3, -2log3 and 0 — mirror images
  \addplot[cgblue, line width=0.7pt, densely dashed, domain=0.14:0.56,
           samples=2] {\Sonestar+\slp*(x-\ustar)};
  \addplot[birrenorange, line width=0.7pt, densely dashed,
           domain=0.14:0.56, samples=2] {\Stwostar-\slp*(x-\ustar)};
  \addplot[birrenred, line width=0.7pt, densely dashed, domain=0.14:0.56,
           samples=2] {\Smax};

  % Points of tangency, and where the naive guess E_1 = E/2 would sit
  \node[statepoint, inner sep=1.3pt] at (axis cs:\ustar,\Sonestar) {};
  \node[statepoint, fill=birrenorange, inner sep=1.3pt]
    at (axis cs:\ustar,\Stwostar) {};
  \node[statepoint, fill=birrenred, inner sep=1.3pt]
    at (axis cs:\ustar,\Smax) {};
  \node[statepoint, fill=birrengray, inner sep=1.1pt]
    at (axis cs:0.5,\Shalf) {};

  % Labels, each in the gap belonging to its own curve
  \node[lbl, text=birrengray, font=\scriptsize, anchor=north east]
    at (axis cs:0.99,4.55) {$N_2=3\,N_1$};
  \node[lbl, anchor=south] at (axis cs:\ustar,3.92) {$T_1=T_2$};
  \node[lbl, text=birrenred, anchor=south west]
    at (axis cs:0.64,3.50) {$S_1+S_2$};
  \node[lbl, text=birrenorange, anchor=south west]
    at (axis cs:0.52,2.42) {$S_2(E-E_1)$};
  % The slope labels start just right of E_1^*, where the slopes are
  % taken.  Their heights are set by the two things they must not touch:
  % the label of S_2 clears the falling S_2 curve at its right end, the
  % label of S_1 clears the rising +1/T_1 tangent at the same place.
  \node[lbl, text=birrenorange, font=\scriptsize, anchor=north west]
    at (axis cs:0.28,2.36) {$-1/T_2$};
  \node[lbl, text=cgblue, font=\scriptsize, anchor=south west]
    at (axis cs:0.28,1.42) {$+1/T_1$};
  \node[lbl, text=cgblue, anchor=north west]
    at (axis cs:0.52,1.15) {$S_1(E_1)$};

  \node[lbl, anchor=north east] at (axis cs:-0.02,-0.05) {$0$};
\end{axis}

% ─────────────────────────────────────────────────────────────
%  (b)  The same total entropy, exponentiated: w(E_1) and its width
% ─────────────────────────────────────────────────────────────
\begin{axis}[
    tpcplot,
    at={(ent.south east)}, anchor=south west, xshift=13mm,
    width=5.4cm, height=6.6cm,
    xmin=0, xmax=1, ymin=0, ymax=1.4,
    xtick={0,0.25,1}, xticklabels={$0$,$E_1^{*}$,$E$},
    ytick={0,0.5,1},
    xlabel={$E_1$}, ylabel={$w/w_{\max}$},
    every axis y label/.style={at={(axis description cs:-0.19,0.5)},
                               rotate=90, anchor=south, font=\small},
    domain=0.0015:0.9985, samples=1201,
    legend pos=north east,
    legend style={font=\scriptsize, draw=light-gray},
    title={(b) \tpcifenglish{How sharp the maximum is}%
                            {Wie scharf das Maximum wird}},
]
  \draw[levellinegrey] (axis cs:\ustar,0) -- (axis cs:\ustar,1.02);

  \addplot[birrenred!50, line width=1.0pt, line join=round]
    {exp(max(1*\Gdiff,-16))};
  \addlegendentry{$N$}
  \addplot[birrenred!72, line width=1.0pt, line join=round]
    {exp(max(10*\Gdiff,-16))};
  \addlegendentry{$10\,N$}
  \addplot[birrenred, line width=1.1pt, line join=round]
    {exp(max(100*\Gdiff,-16))};
  \addlegendentry{$100\,N$}

  \node[lbl, text=birrengray, font=\scriptsize, anchor=north east]
    at (axis cs:0.99,1.00) {\tpcifenglish{width}{Breite}
                            $\propto 1/\sqrt{N}$};
\end{axis}

\end{tikzpicture}
\end{center}


Die Abbildung zeigt das Maximum als Steigungsgleichheit: $S_1$ steigt so schnell, wie $S_2$ fällt. Es liegt nicht bei $E_1=E/2$ und wird mit wachsendem System extrem scharf.

\begin{merksatz}
	Im Gleichgewicht gleichen sich die Temperaturen an, nicht die Energien -- begründet, nicht angenommen.
\end{merksatz}

%!TEX root = TheoretischePhysikC.tex


\section{Vorlesung: Kanonisches Ensemble}
Man betrachte ein System im kontakt mit der Umgebung (,,Wärmebad``), so dass Energie mit der Umgebung ausgetauscht werden kann. Angenommen, das System wird mit Hilbertraum $\mathcal{H}$ and Hamiltonoperator $\widehat{H}$ modelliert. Der Austausch von Energie kann mathematisch modelliert in dem die Energie des Systems einen festen Mittelwert $E = \langle \widehat{H}\rangle$ besitzen.

Diese Situation ist der Normalfall: Das elektromagnetische Feld in diesem Hörsaal ist in einem thermischen Zustand, und jedes System im Raum wechselwirkt mit ihm. Den wahren Zustand könnten wir nicht einmal im Prinzip bestimmen; was wir aufstellen, ist eine \emph{Hypothese} mit denselben Erwartungswerten.

Gesucht wird ein Zustand $\rho$, der die Entropie unter der Bedingung maximiert, dass der Mittelwert der Energie konstant gehalten wird. Dieses Problem wird wie folgt mathematisch nach dem Prinzip der maximalen Entropie beschrieben:
\begin{equation}
	\begin{split}
		\max_{\rho} &\,S(\rho) \\
		\text{so dass} &\, E=\tr(\rho \widehat{H}).
	\end{split}
\end{equation}

Zwei Nebenbedingungen fehlen hier noch: Damit $\rho$ überhaupt einen Zustand beschreibt, muss $\tr(\rho)=1$ gelten und $\rho \ge 0$ sein. Vollständig lautet die Aufgabe also
\begin{equation}
	\max_{\rho} S(\rho) \quad\text{so dass}\quad E=\tr(\rho \widehat{H}),\quad \tr(\rho)=1,\quad \rho\ge 0.
\end{equation}
Gesucht ist eine Matrix mit $d^2$ Freiheitsgraden, und $S$ ist nichtlinear in diesen Einträgen; dass es eine geschlossene Lösung gibt, ist Glück.

Solche Maximierungsprobleme können mit Hilfe von Lagrange-Multiplikatoren gelöst werden: 
\begin{equation}
	\max_{\rho} f(\rho) = -k\tr(\rho\log(\rho)) + \beta \tr(\rho \widehat{H})
\end{equation}
Zu jeder Nebenbedingung gehört ein Multiplikator; oben steht nur der zur Energie. Die Normierung trägt den Term $\gamma\tr(\rho)$, der gleich als $\log(\mathcal{Z})\mathbb{I}$ wiederkehrt; die Positivität prüfen wir am Ende nach.

\begin{greybox}[Lagrange-Multiplikatoren]
	Gesucht sei das Maximum von $f(x,y)$ unter der Nebenbedingung $g(x,y)=0$. Man führt eine Konstante $\lambda$ ein und maximiert stattdessen die nebenbedingungsfreie Funktion $F(x,y,\lambda) = f(x,y) + \lambda g(x,y)$; die Lösung hängt noch von $\lambda$ ab, das man zum Schluss aus der Nebenbedingung bestimmt. Zur Anschauung denke man sich $\lambda$ sehr groß und im Vorzeichen so, dass eine Verletzung bestraft wird: Das Maximum kann dann nur dort liegen, wo die Bedingung gilt.
\end{greybox}

Die Lösung wird als Funktion von $\beta$ dargestellt: $\rho \equiv \rho(\beta)$. Danach wird $\beta$ gesucht, damit die Bedingung $E=\tr(\rho \widehat{H})$ erfüllt wird. 

Um das Maximierungsproblem zu lösen, muss die Variation von $f(\rho)$ bezüglich aller möglichen Variationen $\delta \rho$ von $\rho$ stationär sein:
\begin{equation}
	0= f(\rho + \delta\rho) -f(\rho)
\end{equation}
Man kann zeigen, dass die Funktion ihren Maximalwert nur dann erhält, wenn
\begin{equation}
	\log(\rho) + \beta \widehat{H} + \log(\mathcal{Z})\mathbb{I} = 0,
\end{equation}
wobei $\mathcal{Z}$ eine Konstante ist.
Die Lösung dieser Gleichung wird daher durch eine Exponentialfunktion gegeben:
\begin{equation}
	\rho \propto e^{-\beta \widehat{H}}.
\end{equation}
Erlaubt ist dieser Schritt, weil $\log(\mathcal{Z})$ mit der Identität multipliziert wurde: $\mathbb{I}$ kommutiert mit allem, nur deshalb zerfällt $e^{-\beta \widehat{H}-\log(\mathcal{Z})\mathbb{I}}$ in ein Produkt.
Die Proportionalitätskonstante ist genau durch $\mathcal{Z}$ gegeben und dadurch erhalten wir
\begin{equation}
	\rho = \frac{e^{-\beta \widehat{H}}}{\mathcal{Z}},
\end{equation}
wobei die Bedingung $\tr(\rho) = 1$ bedeutet, dass die \emph{Zustandssumme} $\mathcal{Z}$ durch
\begin{equation}
	\mathcal{Z} = \tr\left(e^{-\beta \widehat{H}}\right) = \sum_{j=1}^d e^{-\beta E_j},
\end{equation}
gegeben wird, wobei $E_j$ die Energieeigenwerte von $\widehat{H}$ sind. Die Positivität ist damit ebenfalls erledigt, denn $e^{-\beta \widehat{H}}$ hat nur positive Eigenwerte. Der Mittelwert $\langle \widehat{H}\rangle$ wird wie folgt berechnet
\begin{equation}
	\langle \widehat{H}\rangle = \sum_{j=1}^d E_j \frac{e^{-\beta E_j}}{\mathcal{Z}} = -\frac{1}{\mathcal{Z}}\frac{\partial \mathcal{Z}}{\partial \beta}.
\end{equation}
Das Minuszeichen der letzten Gleichheit stammt vom Ableiten: Jeder Term $e^{-\beta E_j}$ zieht ein $-E_j$ herunter. Kompakter geschrieben ist also $\langle \widehat{H}\rangle = -\frac{\partial}{\partial \beta}\log(\mathcal{Z})$.

Die Entropie von $\rho=\frac{e^{-\beta \widehat{H}}}{\mathcal{Z}}$ wird wie folgt berechnet
\begin{equation}
	\eqbox{S(\rho) = k\beta\langle\widehat{H}\rangle + k\log(\mathcal{Z}).}
\end{equation}
Definiere Temperatur analog zum mikrokanonischen Ensemble:
\begin{equation}
	\frac{1}{T} = \frac{\partial S}{\partial \langle\widehat{H}\rangle} = k\beta \Rightarrow \beta = \frac{1}{kT}.
\end{equation}
Daher
\begin{equation}
	\eqbox{\rho = \frac{e^{-\frac{\widehat{H}}{kT}}}{\mathcal{Z}}.}
\end{equation}
Die Faktoren $e^{-E_j/kT}$ heißen \emph{Boltzmann-Faktoren}.

Praktisch jedes System in thermischem Kontakt mit einem Wärmebad wird mit diesem \emph{kanonischen Zustand} modelliert, klassisch wie quantenmechanisch. Ab jetzt gilt zudem als Definition: \emph{Thermisches Gleichgewicht} heißt, dass der Zustand durch diese Gleichung gegeben ist.

\section*{Hauptsätze der Thermodynamik}

\begin{center}
\parbox[c][5cm][c]{14cm}{\begin{tcolorbox}[colback=banana,arc=0cm,frame hidden,enhanced]{\textbf{Hauptsatz 0}. Ist System $A$ im thermischen Gleichgewicht mit System $B$ und System $B$ im thermischen Gleichgewicht mit System $C$, so ist System $A$ im thermischen Gleichgewicht mit System $C$.}\end{tcolorbox}}
\end{center}
\begin{center}
\parbox[c][5cm][c]{14cm}{\begin{tcolorbox}[colback=banana,arc=0cm,frame hidden,enhanced]{\textbf{Hauptsatz 1}. Es gibt eine Größe $U$: ,,Innere Energie``. Die Änderung $\Delta U = Q+W$ der inneren Energie entspricht der zugeführten Wärme $Q$ plus der verrichteten Arbeit $W$. (Energieerhaltung.)}\end{tcolorbox}}
\end{center}
Hierbei ist \emph{Wärme} $Q$ als jeder spontane Energiefluss von einem Objekt zu einem anderen definiert, der von einer Temperaturdifferenz zwischen den beiden Objekten herrührt. \emph{Arbeit} $W$ ist in der Thermodynamik als jeder andere Transfer von Energie in ein System oder aus einem System heraus definiert. 
\begin{center}
\parbox[c][5cm][c]{14cm}{\begin{tcolorbox}[colback=banana,arc=0cm,frame hidden,enhanced]{\textbf{Hauptsatz 2}. Es gibt eine Größe ,,Entropie`` und bei adiabatischen Prozessen gilt $\Delta S \ge 0$. (Adiabatisch bedeutet: ohne Austausch von Wärme mit der Umgebung.)}\end{tcolorbox}}
\end{center}
\begin{merksatz}[Die Perspektive dieses Kurses]
	Die Thermodynamik ist nicht fundamental. Sie ist eine Folgerung der Quantenmechanik -- eine grobkörnige Beschreibung der Naturgesetze für große Systeme.
\end{merksatz}
\begin{center}
\parbox[c][5cm][c]{14cm}{\begin{tcolorbox}[colback=banana,arc=0cm,frame hidden,enhanced]{\textbf{Hauptsatz 3}. Man kann den absoluten Temperaturnullpunkt $T=0$ nicht vollständig erreichen.}\end{tcolorbox}}
\end{center}

Ganz richtig ist dieser Hauptsatz allerdings nicht, denn es gibt Ausnahmen. Er ist äquivalent zu einer Aussage über die Entropie: Normalerweise verschwindet sie für $T\to 0$. Bei manchen Systemen -- Glas, gewissen Eiskristallen, Spingläsern -- bleibt am Nullpunkt Entropie übrig; für sie gilt der dritte Hauptsatz nicht.
 




\end{document}
